\section{Maximal contact}
\label{sec:maximal-contact}

This section prepares the induction on dimension in the resolution of marked ideals
(Section~\ref{sec:hironaka}): Taylor expansions with holomorphic remainder, marked ideals of
maximal order and their tangent directions (hypersurfaces of maximal contact), and the behaviour of
derivative ideals under blowups. The source is Włodarczyk's paper, Sections 2.6, 2.7, 2.10 and 2.11,
in the formalism of marked ideals of Bierstone and Milman. The route differs from the source as follows.
\begin{enumerate}
\item The category is that of complex manifolds. Where Włodarczyk uses a formal automorphism of the
completed local ring (the glueing lemma), the analytic route uses convergent coordinate changes;
the expansions it needs are the Taylor expansions with \emph{holomorphic} remainder of
Lemmas~\ref{mc:taylor-tail} and~\ref{mc:normal-vanishing}.
\item The statements on marked ideals are written for a marked ideal \((\mathfrak a,\mu)\) on a complex
manifold \(X\) (that is, \(N=X\) in Definition~\ref{hir:marked-ideal}), with the derivative ideals of
Definition~\ref{hir:derivative-ideal}. For a marked ideal \((M,N,E,\mathfrak a,d)\) with \(N\subsetneq M\) they
are applied to the manifold \(N\) and the ideal \(\mathfrak a|_N\); the derivative ideal is then the one of
\(N\) (derivatives tangent to \(N\)), not the restriction of the derivative ideal of \(M\): for
\(N=\{u=0\}\subset\mathbf C^2\) and \(\mathfrak a=(u)\) one has \(\mathfrak a|_N=0\) but \(1=\partial_uu\in\mathcal D(\mathfrak a)\).
Its ambient representative is the ideal of Definition~\ref{mc:derivative-along}.
\end{enumerate}

Włodarczyk's analytic version of the algorithm (\texttt{wlodarczyk09-analytic-resolution}, Sections~5.5--5.6, printed pp.~50--53; its statement
numbers are its own) states these results for germs \(M_Z\) of manifolds at compact sets. Its definition of \(\mathcal H\) and Lemma~5.5.1\,(1) are
Definition~\ref{mc:homogenized-ideal} and Lemma~\ref{mc:homogenized-basic}; Lemma~5.5.2 is Lemma~\ref{mc:homogenized-equivalent} (part (1)) and
Lemma~\ref{mc:homogenized-tower} (part (2)); the Glueing Lemma~5.5.3 is Theorem~\ref{mc:glueing-lemma} (parts (1), (3), (4); part (2) is empty for
\(E=\emptyset\)) together with Theorems~\ref{mc:glueing-tower-pullback} and~\ref{mc:glueing-lift} (part (5)); Definition~5.6.1 and Lemma~5.6.2 are
Definition~\ref{mc:coefficient-ideal} and Lemmas~\ref{mc:coefficient-equivalent} and~\ref{mc:coefficient-wequiv}; Lemma~5.6.3 is Lemma~\ref{mc:restrict-order};
Proposition~5.6.4 is Lemmas~\ref{mc:coefficient-restriction-support}, \ref{mc:coefficient-restriction-tower} and~\ref{mc:coefficient-restriction-full}. The
differences: the glueing automorphism \(\phi_{uv}\) of a germ \(M_S\) becomes a biholomorphism between open neighbourhoods of the point, and the
noetherian step of its proof becomes coherence (Lemma~\ref{mc:coherent-germs}); the restricted coefficient ideal keeps all of \(E\) (the source keeps the
traces on \(S\) of the members transversal to \(S\)) and has the mark \(\mu!\) instead of \(\operatorname{lcm}(1,\dots,\mu)\); the hypotheses of
Lemma~5.6.3 and Proposition~5.6.4 that the support contain no (germ of) \(S\) and that \(S\) have normal crossings with \(E\) are not needed; and maximal
order is used only to provide tangent directions (Lemma~\ref{mc:tangent-direction-exists}).

\subsection{Taylor expansions with holomorphic remainder}

For a complex normed space \(E\), a function \(F\) holomorphic on an open subset of \(E\) and \(v\in E\), write
\(\partial_vF(z)=DF(z)\,v\), and \(\partial_{v_1}\cdots\partial_{v_m}F\) for iterated directional derivatives; they are
holomorphic on the same open set.

\begin{lemma}[Taylor expansion along a line]
\label{mc:taylor-tail}
\uses{dem:germ-noetherian}
\source{wlodarczyk05-hironaka}
\lean{Hironaka.HolomorphicTaylorTail.dirDeriv,Hironaka.HolomorphicTaylorTail.iterDeriv,Hironaka.HolomorphicTaylorTail.iterDeriv_append,Hironaka.HolomorphicTaylorTail.differentiableOn_iterDeriv,Hironaka.HolomorphicTaylorTail.iteratedDeriv_line,Hironaka.HolomorphicTaylorTail.iterDeriv_comp_clm,Hironaka.HolomorphicTaylorTail.lastCoeff_eq_iteratedDeriv,Hironaka.HolomorphicTaylorTail.exists_taylor_tail}
\leanok
Let \(F\) be holomorphic near \(q\in E\), \(v\in E\) and \(\mu\in\mathbb N\). There is a function \(G\), holomorphic near
\((q,0)\in E\times\mathbf C\), such that near \((q,0)\)
\[
 F(z+tv)=\sum_{j<\mu}\frac{1}{j!}\,\partial_v^jF(z)\,t^j+t^\mu\,G(z,t).
\]
\end{lemma}
\begin{proof}
\leanok
Choose a ball \(B(q,r)\) on which \(F\) is holomorphic, \(R=r/(2(\|v\|+1))\), \(\delta=r/2\) and \(\rho=R/2\), and put
\(h(w,t)=F(q+w+tv)\), holomorphic on \(\{\|w\|<\delta\}\times\{|t|<R\}\). The parametric Taylor division of the
proof of Proposition~\ref{dem:germ-noetherian} gives \(h=Q(h)\,t^\mu+\sum_{j<\mu}c_j(h)(w)\,t^j\) with \(Q(h)\) jointly
holomorphic and \(c_j(h)(w)=(2\pi i)^{-1}\oint_{|\zeta|=\rho}h(w,\zeta)\zeta^{-j-1}\,d\zeta\). By the Cauchy formula for
derivatives, \(c_j(h)(w)=\frac1{j!}\frac{d^j}{dt^j}h(w,t)|_{t=0}\). Along the line \(s\mapsto z+sv\), by induction on \(j\)
and the chain rule, \(\frac{d^j}{ds^j}F(z+sv)=\partial_v^jF(z+sv)\). With \(w=z-q\) this is the claim, \(G(z,t)=Q(h)(z-q,t)\).
\end{proof}

Let \(n\ge0\), \(0\le k\le n\), and let \(P_k=\{u\in\mathbf C^n: u_i=0\text{ for }i\ge k\}\) be the coordinate plane \(\mathbf C^k\times0\);
the coordinates \(u_i\), \(i\ge k\), are its \emph{normal coordinates}.

\begin{definition}[Powers of the ideal of normal coordinates]
\label{mc:normal-power}
\lean{Hironaka.HolomorphicTaylorTail.InNormalPow,Hironaka.HolomorphicTaylorTail.InNormalPow.eventually_differentiableAt,Hironaka.HolomorphicTaylorTail.InNormalPow.congr,Hironaka.HolomorphicTaylorTail.inNormalPow_zero_fun,Hironaka.HolomorphicTaylorTail.InNormalPow.add,Hironaka.HolomorphicTaylorTail.InNormalPow.const_mul_holo,Hironaka.HolomorphicTaylorTail.InNormalPow.coord_mul,Hironaka.HolomorphicTaylorTail.InNormalPow.coord_pow_mul,Hironaka.HolomorphicTaylorTail.InNormalPow.mono_k,Hironaka.HolomorphicTaylorTail.inNormalPow_finset_sum}
\leanok
Let \(q\in\mathbf C^n\). A function \(F\) lies \emph{near \(q\) in the \(\mu\)-th power of the ideal of the normal coordinates}
of \(P_k\) if, for \(\mu=0\), \(F\) is holomorphic near \(q\), and, for \(\mu\ge1\), \(F=\sum_{i\ge k}u_iH_i\) near \(q\) with each
\(H_i\) in the \((\mu-1)\)-th power near \(q\). These functions are closed under sums, products with functions holomorphic
near \(q\) and local equality, contain \(0\), and \(u_i^jF\) lies in the \((m+j)\)-th power when \(F\) lies in the \(m\)-th power
and \(i\ge k\); the \(\mu\)-th power for \(P_{k+1}\) is contained in the \(\mu\)-th power for \(P_k\).
\end{definition}
These closure properties are immediate inductions on \(\mu\). In particular such an \(F\) is a finite sum of
products \(u_{i_1}\cdots u_{i_\mu}H\) with \(i_1,\dots,i_\mu\ge k\) and \(H\) holomorphic near \(q\), i.e.\ \(F\in\mathcal I_{P_k}^\mu\) near \(q\).

\begin{lemma}[Taylor expansion in normal coordinates, vanishing form]
\label{mc:normal-vanishing}
\uses{mc:taylor-tail,mc:normal-power}
\source{wlodarczyk05-hironaka}
\lean{Hironaka.HolomorphicTaylorTail.killCoord,Hironaka.HolomorphicTaylorTail.killCoord_apply_apply,Hironaka.HolomorphicTaylorTail.killCoord_add,Hironaka.HolomorphicTaylorTail.killCoord_of_apply_eq_zero,Hironaka.HolomorphicTaylorTail.inNormalPow_of_iterDeriv_eq_zero_aux,Hironaka.HolomorphicTaylorTail.inNormalPow_of_iterDeriv_eq_zero}
\leanok
Let \(q\in P_k\), \(\mu\in\mathbb N\), and let \(F\) be holomorphic near \(q\). Assume that every iterated directional derivative
\(\partial_{v_1}\cdots\partial_{v_m}F\) of order \(m<\mu\) vanishes at every point of \(P_k\) near \(q\). Then near \(q\) the function \(F\)
lies in the \(\mu\)-th power of the ideal of the normal coordinates of \(P_k\).
\end{lemma}
\begin{proof}
\leanok
Induction on the number \(n-k\) of normal coordinates, for all \(\mu\) and \(F\). If \(k\ge n\), then \(P_k=\mathbf C^n\); for
\(\mu=0\) there is nothing to prove, and for \(\mu\ge1\) the hypothesis with \(m=0\) says \(F=0\) near \(q\).

Let \(k<n\), \(i=k\), \(e=e_i\), and let \(\pi(u)=u-u_ie\) kill the \(i\)-th coordinate; \(\pi\) is linear, \(\pi(q)=q\) and
\(\pi(u)+u_ie=u\). Lemma~\ref{mc:taylor-tail} at \(q\) in the direction \(e\), evaluated at \((z,t)=(\pi(u),u_i)\), gives near \(q\)
\[
 F(u)=\sum_{j<\mu}\frac1{j!}\,u_i^j\,g_j(u)+u_i^\mu\,G(\pi(u),u_i),\qquad g_j=(\partial_e^jF)\circ\pi .
\]
The last term lies in the \(\mu\)-th power. For \(j<\mu\), \(g_j\) is holomorphic near \(q\), and for \(y\in P_{k+1}\) near \(q\) and
directions \(v_1,\dots,v_m\), the chain rule for the linear map \(\pi\) gives
\(\partial_{v_1}\cdots\partial_{v_m}g_j(y)=\partial_{\pi v_1}\cdots\partial_{\pi v_m}\partial_e^jF(\pi(y))\), an iterated derivative of \(F\) of
order \(m+j\) at \(\pi(y)\in P_k\) (its \(i\)-th coordinate is \(0\), and its coordinates of index \(>k\) are those of \(y\)). For
\(m<\mu-j\) it vanishes, so by induction \(g_j\) lies in the \((\mu-j)\)-th power for \(P_{k+1}\), hence for \(P_k\), and
\(u_i^jg_j\) lies in the \(\mu\)-th power. Sums of such functions do too.
\end{proof}
Together with the forward half of Lemma~\ref{hir:order-taylor} (an iterated derivative of order \(m\) of a germ of
\(\mathfrak m_y^\mu\) lies in \(\mathfrak m_y^{\mu-m}\)), applied in flattening coordinates of a closed submanifold \(C\), this gives
the inclusion \(\mathfrak a\subset\mathcal I_C^\mu\) of Lemma~\ref{hir:hadamard-bridge}.

\subsection{Marked ideals of maximal order and tangent directions}

In this subsection \((\mathfrak a,\mu)\) is a marked ideal on a complex manifold \(X\) with \(\mathfrak a\) coherent and \(\mu\ge1\),
\(\mathcal D^i\) are the derivative ideals of Definition~\ref{hir:derivative-ideal}, and
\(\operatorname{cosupp}(\mathfrak a,\mu)=\{x:\operatorname{ord}_x\mathfrak a\ge\mu\}\) (Definition~\ref{hir:cosupport}).

\begin{definition}[Maximal order]
\label{mc:maximal-order}
\uses{hir:derivative-ideal,hir:cosupport}
\source{wlodarczyk05-hironaka}
\lean{Hironaka.MaximalContact.IsMaximalOrder}
\leanok
\((\mathfrak a,\mu)\) is \emph{of maximal order} if \(\mathcal D^\mu(\mathfrak a)=\mathcal O_X\), i.e.\ \(1\in\mathcal D^\mu(\mathfrak a)_x\) for every
\(x\in X\).
\end{definition}
By Lemma~\ref{hir:order-taylor} this says \(\operatorname{ord}_x\mathfrak a\le\mu\) for every \(x\), Włodarczyk's definition.

\begin{lemma}[Derivatives of a marked ideal of maximal order]
\label{mc:maximal-order-derivative}
\uses{mc:maximal-order}
\source{wlodarczyk05-hironaka}
\lean{Hironaka.MaximalContact.iterDerivIdeal_iterDerivIdeal,Hironaka.MaximalContact.IsMaximalOrder.iterDerivIdeal}
\leanok
If \((\mathfrak a,\mu)\) is of maximal order and \(0\le i\le\mu\), then \((\mathcal D^i(\mathfrak a),\mu-i)\) is of maximal order.
\end{lemma}
\begin{proof}
\leanok
\(\mathcal D^{\mu-i}(\mathcal D^i(\mathfrak a))=\mathcal D^\mu(\mathfrak a)=\mathcal O_X\), by the definition \(\mathcal D^{j+1}=\mathcal D\circ\mathcal D^j\) and induction on
\(\mu-i\).
\end{proof}

\begin{definition}[Tangent direction]
\label{mc:tangent-direction}
\uses{hir:derivative-ideal}
\source{wlodarczyk05-hironaka}
\lean{Hironaka.MaximalContact.IsTangentDirection}
\leanok
Let \(U\subset X\) be open. A \emph{tangent direction} of \((\mathfrak a,\mu)\) on \(U\) is a function \(u\) whose germ at every
point of \(U\) lies in \(\mathcal D^{\mu-1}(\mathfrak a)\) and which has multiplicity one: \(du(x)\ne0\) (equivalently, some partial
derivative \(\partial^e_vu(x)\) in some chart \(e\) at \(x\) is nonzero) for every \(x\in U\) with \(u(x)=0\).
\end{definition}
By the implicit function theorem the zero set \(V(u)\) of a tangent direction is a closed submanifold of \(U\) of codimension
one (or empty), a \emph{hypersurface of maximal contact}.

\begin{lemma}[The cosupport lies on every tangent direction]
\label{mc:tangent-direction-support}
\uses{mc:tangent-direction,hir:giraud-support}
\source{wlodarczyk05-hironaka}
\lean{Hironaka.MaximalContact.cosupport_add_subset_cosupport_iterDerivIdeal,Hironaka.MaximalContact.apply_eq_zero_of_mem_cosupport,Hironaka.MaximalContact.IsTangentDirection.cosupport_inter_subset}
\leanok
If \(u\) is a tangent direction of \((\mathfrak a,\mu)\) on \(U\), then \(\operatorname{cosupp}(\mathfrak a,\mu)\cap U\subset V(u)\).
\end{lemma}
\begin{proof}
\leanok
By Lemma~\ref{hir:giraud-support} (the inclusion \(\operatorname{cosupp}(\mathfrak a,j+1)\subset\operatorname{cosupp}(\mathcal D(\mathfrak a),j)\), applied
\(\mu-1\) times), \(x\in\operatorname{cosupp}(\mathfrak a,\mu)\) implies \(x\in\operatorname{cosupp}(\mathcal D^{\mu-1}(\mathfrak a),1)\), i.e.\ every germ of
\(\mathcal D^{\mu-1}(\mathfrak a)_x\) vanishes at \(x\); in particular \(u(x)=0\).
\end{proof}

\begin{lemma}[Existence of tangent directions]
\label{mc:tangent-direction-exists}
\uses{mc:maximal-order,mc:tangent-direction,mc:tangent-direction-support,hir:derivative-ideal-generators}
\source{wlodarczyk05-hironaka}
\lean{Hironaka.MaximalContact.exists_isTangentDirection}
\leanok
If \((\mathfrak a,\mu)\) is of maximal order and \(x\in\operatorname{cosupp}(\mathfrak a,\mu)\), then \((\mathfrak a,\mu)\) has a tangent direction on
some open neighbourhood of \(x\).
\end{lemma}
\begin{proof}
\leanok
\(\mathcal D^{\mu-1}(\mathfrak a)\) is coherent (Lemma~\ref{hir:derivative-ideal-generators}); let \(f_1,\dots,f_m\) generate it on an open
neighbourhood \(U_0\) of \(x\) inside the domain of a chart \(e\). By the same lemma, \(\mathcal D(\mathcal D^{\mu-1}(\mathfrak a))=\mathcal D^\mu(\mathfrak a)\) is
generated near \(x\) by the \(f_i\) and the \(\partial^e_{e_j}f_i\); since it contains \(1\), not all of these vanish at \(x\). As in the proof
of Lemma~\ref{mc:tangent-direction-support}, every germ of \(\mathcal D^{\mu-1}(\mathfrak a)_x\) vanishes at \(x\), so \(f_i(x)=0\) for all \(i\), and
\(\partial^e_{e_j}f_i(x)\ne0\) for some \(i,j\). Then \(df_i(x)\ne0\), hence \(df_i\ne0\) on an open neighbourhood \(U\subset U_0\) of \(x\), and
\(u=f_i|_U\) is a tangent direction on \(U\).
\end{proof}

\subsection{Derivative ideals under a blowup}

In this subsection \(\pi\colon Y\to X\) is the blowup of the complex manifold \(X\) along \(C\)
(Definition~\ref{ana:blowups}), \(D\) its exceptional divisor, so that \(\mathcal O_Y(-qD)=\pi^*(\mathcal I_C^q)\)
(Lemma~\ref{bcalc:exceptional-divisor}), and \(\sigma^c(\mathfrak a,\mu)=(\pi^*\mathfrak a:\mathcal O_Y(-\mu D))\) is the
controlled transform (Definition~\ref{bcalc:controlled-transform-def}). Marked ideals are taken with \(N\) the whole
manifold. For a chart \(e\) and \(v\in\mathbf C^n\), \(\partial^e_v\) is the partial derivative of
Definition~\ref{hir:derivative-ideal}; \(e_z\) is a chart at \(z\) and \(e_1,\dots,e_n\) the standard basis.

\begin{lemma}[Derivatives lower powers of an ideal by one]
\label{mc:derivative-power}
\uses{hir:derivative-ideal,hir:derivative-order,hir:order-basic}
\source{wlodarczyk05-hironaka}
\lean{Hironaka.BlowupDerivativeIdeal.memPow_of_eq',Hironaka.BlowupDerivativeIdeal.memPow_add',Hironaka.BlowupDerivativeIdeal.memPow_sum',Hironaka.BlowupDerivativeIdeal.memPow_congr',Hironaka.BlowupDerivativeIdeal.mem_idealPow_one_of_mem,Hironaka.BlowupDerivativeIdeal.chartDeriv_mem_idealPow',Hironaka.BlowupDerivativeIdeal.derivIdeal_le_idealPow,Hironaka.BlowupDerivativeIdeal.iterDerivIdeal_le_idealPow}
\leanok
Let \(\mathfrak b\) be an ideal of holomorphic germs, \(e\) a chart at \(x\) and \(v\in\mathbf C^n\). If \(g\in(\mathfrak b^{q+1})_x\), then
\(\partial^e_vg\in(\mathfrak b^q)_x\). Consequently \(\mathfrak a\subset\mathfrak b^{q+i}\) implies \(\mathcal D^i(\mathfrak a)\subset\mathfrak b^q\).
\end{lemma}
\begin{proof}
\leanok
The first claim is proved as in Lemma~\ref{hir:derivative-order}, with \(\mathfrak m_x\) replaced by \(\mathfrak b\): for \(q=0\),
\(\partial^e_vg\) is holomorphic; a germ of \(\mathfrak b^{q+2}\) is \(\sum_kc_ka_kb_k\) with \(a_k\in\mathfrak b^{q+1}\), \(b_k\in\mathfrak b\), and in
\(\partial(c_ka_kb_k)=(\partial c_k)a_kb_k+c_k(\partial a_k)b_k+c_ka_k\partial b_k\) the first and last terms lie in \(\mathfrak b^{q+1}\) and the
middle one does by induction. If \(\mathfrak a\subset\mathfrak b^{q+1}\), the generators of \(\mathcal D(\mathfrak a)_x\) (germs of
\(\mathfrak a_x\subset\mathfrak b^{q+1}_x\subset\mathfrak b^q_x\) and their partial derivatives) lie in \(\mathfrak b^q_x\), which is closed
under holomorphic combinations. Induction on \(i\), with \(\mathcal D^{i+1}=\mathcal D\circ\mathcal D^i\), gives the last claim.
\end{proof}

\begin{lemma}[Powers of a locally principal ideal]
\label{mc:principal-power}
\uses{ana:ideal-sheaves}
\lean{Hironaka.BlowupDerivativeIdeal.holomorphicAt_pow,Hironaka.BlowupDerivativeIdeal.idealPow_mem_iff_of_generator}
\leanok
If an ideal \(\mathfrak l\) is generated by a holomorphic function \(y\) at every point of an open set \(U\), then \(\mathfrak l^q\) is
generated by \(y^q\) at every point of \(U\).
\end{lemma}
\begin{proof}
\leanok
Induction on \(q\); \(\mathfrak l^0=\mathcal O\) is generated by \(1\). A generator \(ab\) of \(\mathfrak l^{q+1}\) with \(a=cy^q\), \(b=dy\) near
the point is \((cd)y^{q+1}\); conversely \(y^{q+1}=y^q\cdot y\) with \(y^q\in\mathfrak l^q\) and \(y\in\mathfrak l\).
\end{proof}

\begin{lemma}[Derivatives in local coordinates]
\label{mc:coordinate-derivatives}
\uses{ana:local-holomorphic-inverse,hir:derivative-ideal-generators}
\lean{Hironaka.BlowupDerivativeIdeal.exists_coordDeriv_coeff,Hironaka.BlowupDerivativeIdeal.chartDeriv_eventuallyEq_sum_flat,Hironaka.BlowupDerivativeIdeal.sum_triple}
\leanok
\begin{enumerate}
\item Let \(W\) be a holomorphic map from a neighbourhood of \(z\in Y\) to \(\mathbf C^n\) such that \(W\circ e_z^{-1}\) has an invertible
derivative at \(e_z(z)\). There are functions \(c_{lm}\), holomorphic at \(z\), such that for every \(G\) holomorphic near \(W(z)\), near \(z\)
\[
 (\partial_lG)\circ W=\sum_mc_{lm}\,\partial^{e_z}_{e_m}(G\circ W).
\]
\item Let \(e\) be a chart at \(x\in X\), \(\Phi\) a local biholomorphism at \(e(x)\), \(f\) holomorphic near \(x\), \(F=f\circ e^{-1}\circ\Phi^{-1}\),
\(e'\) a chart at \(x\) and \(v\in\mathbf C^n\). There are functions \(\alpha_i\), holomorphic at \(x\), with
\(\partial^{e'}_vf=\sum_i\alpha_i\,(\partial_iF)\circ\Phi\circ e\) near \(x\).
\end{enumerate}
\end{lemma}
\begin{proof}
\leanok
(1) By the inverse function theorem (Lemma~\ref{ana:local-holomorphic-inverse}) \(\Theta=W\circ e_z^{-1}\) is a biholomorphism of a
neighbourhood of \(e_z(z)\) onto an open set. Put \(u_l(p)=D\Theta^{-1}(W(p))\,e_l\) and \(c_{lm}(p)=u_l(p)_m\), holomorphic at \(z\). For \(p\) near
\(z\) and \(\zeta=e_z(p)\), \(W(p)=\Theta(\zeta)\), \((G\circ W)\circ e_z^{-1}=G\circ\Theta\) near \(\zeta\), and \(D\Theta(\zeta)\,u_l(p)=e_l\) because
\(\Theta\circ\Theta^{-1}\) is the identity near \(\Theta(\zeta)\). Hence
\(\sum_mc_{lm}(p)\,\partial^{e_z}_{e_m}(G\circ W)(p)=D(G\circ\Theta)(\zeta)\,u_l(p)=DG(W(p))\,e_l\).
(2) By Lemma~\ref{hir:derivative-ideal-generators}, \(\partial^{e'}_vf=\sum_l\gamma_l\,\partial^e_{e_l}f\) with \(\gamma_l\) holomorphic at \(x\). Near
\(e(q)\), \(f\circ e^{-1}=F\circ\Phi\), so \(\partial^e_{e_l}f(q)=DF(\Phi(e(q)))\,D\Phi(e(q))\,e_l=\sum_iD\Phi(e(q))_{il}\,(\partial_iF)(\Phi(e(q)))\); take
\(\alpha_i=\sum_l\gamma_l\,D\Phi(e(\cdot))_{il}\).
\end{proof}

For \(0\le k\le j<n\) let \(\sigma_j\colon\mathbf C^n\to\mathbf C^n\) be the monomial map of Lemma~\ref{bcalc:monomial-lift}:
\(\sigma_j(w)_i=w_iw_j\) for \(i\ge k\), \(i\ne j\), and \(\sigma_j(w)_i=w_i\) otherwise.

\begin{lemma}[The derivation \(w_j\,\sigma_j^*(\partial_i)\)]
\label{mc:monomial-derivation}
\uses{bcalc:monomial-lift}
\source{wlodarczyk05-hironaka}
\lean{Hironaka.BlowupDerivativeIdeal.lineDir,Hironaka.BlowupDerivativeIdeal.lineDir_apply,Hironaka.BlowupDerivativeIdeal.monomial_add_smul_single,Hironaka.BlowupDerivativeIdeal.fderiv_comp_monomial,Hironaka.BlowupDerivativeIdeal.monomialCoeff,Hironaka.BlowupDerivativeIdeal.differentiable_monomialCoeff,Hironaka.BlowupDerivativeIdeal.monomialCoeff_self,Hironaka.BlowupDerivativeIdeal.sum_monomialCoeff_mul_lineDir,Hironaka.BlowupDerivativeIdeal.mul_fderiv_monomial}
\leanok
Let \(F\) be holomorphic near \(\sigma_j(w)\). For every \(i\),
\[
 w_j\,(\partial_iF)(\sigma_j(w))=\sum_lP_{il}(w)\,\partial_l(F\circ\sigma_j)(w),
\]
where the polynomials \(P_{il}\) are: \(P_{il}=w_j\delta_{il}\) for \(i<k\); \(P_{il}=\delta_{il}\) for \(k\le i\ne j\); and \(P_{jj}=w_j\),
\(P_{jl}=-w_l\) for \(k\le l\ne j\), \(P_{jl}=0\) for \(l<k\). In particular \(P_{ij}=\delta_{ij}w_j\).
\end{lemma}
\begin{proof}
\leanok
On coordinate lines \(\sigma_j\) is affine: \(\sigma_j(w+te_l)=\sigma_j(w)+t\,b_l(w)\) with \(b_l=w_je_l\) for \(k\le l\ne j\), \(b_l=e_l\) for \(l<k\), and
\(b_j=e_j+\sum_{k\le i\ne j}w_ie_i\). Differentiating \(t\mapsto F(\sigma_j(w+te_l))\) at \(0\) gives \(\partial_l(F\circ\sigma_j)(w)=DF(\sigma_j(w))\,b_l(w)\). It
remains to check \(w_je_i=\sum_lP_{il}b_l\): for \(i\ne j\) only \(l=i\) contributes and \(P_{ii}b_i=w_je_i\); for \(i=j\),
\(w_jb_j-\sum_{k\le l\ne j}w_l\,w_je_l=w_je_j\). Apply \(DF(\sigma_j(w))\).
\end{proof}
This is the computation of Włodarczyk's proof of Lemma ``le:inclusions'': \(y\,\sigma^*(\partial/\partial u_i)\), with \(y=w_j\) the equation of the
exceptional divisor, is a derivation with holomorphic coefficients.

\begin{lemma}[The derivations \(y\,\pi^*(\partial)\) near a point of the blowup]
\label{mc:exceptional-derivation}
\uses{mc:coordinate-derivatives,mc:monomial-derivation,bcalc:monomial-lift,bcalc:exceptional-divisor,bex:hadamard}
\source{wlodarczyk05-hironaka,bierstone-milman08-functoriality}
\lean{Hironaka.BlowupDerivativeIdeal.idOn,Hironaka.BlowupDerivativeIdeal.IsDerivationGenerator,Hironaka.BlowupDerivativeIdeal.derivationData_of_coordinates,Hironaka.BlowupDerivativeIdeal.isDerivationGenerator_of_notMem,Hironaka.BlowupDerivativeIdeal.exists_isDerivationGenerator_of_mem,Hironaka.BlowupDerivativeIdeal.exists_isDerivationGenerator}
\leanok
Let \(z\in Y\) and \(x=\pi(z)\). There are an open set \(U\ni z\) and a function \(y\) holomorphic on \(U\) such that
\begin{enumerate}
\item \(y\) generates \(\pi^*\mathcal I_C=\mathcal O_Y(-D)\) at every point of \(U\), and vanishes identically on no nonempty open subset of \(U\);
\item for every \(f\) holomorphic near \(x\), every chart \(e'\) at \(x\) and every \(v\), there are functions \(b_m\) and \(\rho\), holomorphic at \(z\),
such that near \(z\), for the derivation \(D'=\sum_mb_m\,\partial^{e_z}_{e_m}\),
\[
 y\cdot(\partial^{e'}_vf)\circ\pi=D'(f\circ\pi)\qquad\text{and}\qquad D'(y)=\rho\,y .
\]
\end{enumerate}
\end{lemma}
\begin{proof}
\leanok
If \(x\notin C\), take \(U=\pi^{-1}(X\setminus C)\) and \(y=1\): \(1\in\mathcal I_C\) off the closed set \(C\), so \(\pi^*\mathcal I_C\) is the unit ideal on \(U\).
The derivative of \(\pi\) at \(z\) is bijective (Lemma~\ref{bcalc:exceptional-divisor}(2)), so \(W=e_x\circ\pi\) satisfies the hypothesis of
Lemma~\ref{mc:coordinate-derivatives}(1). With \(\Phi\) the identity in Lemma~\ref{mc:coordinate-derivatives}(2), \(F=f\circ e_x^{-1}\) and \(F\circ W=f\circ\pi\),
\[
 (\partial^{e'}_vf)\circ\pi=\sum_i\alpha_i(\pi)\,(\partial_iF)\circ W=\sum_m\Bigl(\sum_i\alpha_i(\pi)\,c_{im}\Bigr)\partial^{e_z}_{e_m}(f\circ\pi),
\]
and \(D'(1)=0\).

If \(x\in C\), let \(\Phi\circ e_x\) flatten \(C\) onto \(\mathbf C^k\times0\) near \(x\) and let \(j\ge k\), \(O\ni z\) and \(W\colon O\to\mathbf C^n\) be a monomial
chart and the local inverse of its lift given by Lemma~\ref{bcalc:monomial-lift}: \(\sigma_j\circ W=\Phi\circ e_x\circ\pi\) on \(O\), and
\(W\circ e_z^{-1}\) has an injective, hence invertible, derivative at \(e_z(z)\). Put \(U=O\) and \(y=W_j\).
\emph{Generation}: \(W_j=(\Phi\circ e_x)_j\circ\pi\) and \((\Phi\circ e_x)_j\in\mathcal I_C\), so \(W_j\in\pi^*\mathcal I_C\). Conversely, a germ \(a\) of \(\mathcal I_C\)
at \(\pi(p)\) is \(\sum_{i\ge k}a_i\,(\Phi\circ e_x)_i\) (Lemma~\ref{bex:hadamard}), and \((\Phi\circ e_x)_i\circ\pi=\sigma_j(W)_i\) is \(W_iW_j\) for \(i\ne j\) and \(W_j\) for
\(i=j\), so \(a\circ\pi\in(W_j)\). \emph{Nonvanishing}: where \(W_j=0\), \(\sigma_j(W)\in\mathbf C^k\times0\), i.e.\ \(\pi\in C\); and \(\pi^{-1}(C)\) has empty
interior (Lemma~\ref{bcalc:exceptional-divisor}(1)). \emph{Derivation}: with \(F=f\circ e_x^{-1}\circ\Phi^{-1}\) and \(G=F\circ\sigma_j\), so that
\(G\circ W=f\circ\pi\) near \(z\), Lemmas~\ref{mc:coordinate-derivatives} and~\ref{mc:monomial-derivation} give
\[
 W_j\,(\partial^{e'}_vf)\circ\pi=\sum_i\alpha_i(\pi)\,W_j\,(\partial_iF)(\sigma_j(W))=\sum_i\alpha_i(\pi)\sum_lP_{il}(W)\,(\partial_lG)\circ W=\sum_mb_m\,\partial^{e_z}_{e_m}(f\circ\pi)
\]
with \(b_m=\sum_i\alpha_i(\pi)\sum_lP_{il}(W)\,c_{lm}\). Applying the same identities to \(G=w_j\) (so \(G\circ W=W_j\) and \(\partial_lG=\delta_{lj}\)),
\(D'(W_j)=\sum_i\alpha_i(\pi)\sum_lP_{il}(W)\,\delta_{lj}=\sum_i\alpha_i(\pi)P_{ij}(W)=\alpha_j(\pi)\,W_j\).
\end{proof}

\begin{lemma}[Giraud, Villamayor: derivatives and controlled transforms]
\label{mc:inclusions}
\uses{mc:exceptional-derivation,mc:principal-power,mc:derivative-power,bcalc:controlled-transform-def,bcalc:exceptional-divisor,hir:derivative-ideal,hir:hadamard-bridge}
\source{wlodarczyk05-hironaka,bierstone-milman08-functoriality}
\lean{Hironaka.BlowupDerivativeIdeal.chartDeriv_pow,Hironaka.BlowupDerivativeIdeal.eventuallyEq_of_pow_mul_eventuallyEq,Hironaka.BlowupDerivativeIdeal.comp_mem_pullback,Hironaka.BlowupDerivativeIdeal.controlledTransform_derivIdeal_le,Hironaka.BlowupDerivativeIdeal.controlledTransform_iterDerivIdeal_le,Hironaka.BlowupDerivativeIdeal.controlledTransform_iterDerivIdeal_le_of_subset_cosupport}
\leanok
Let \(\mathfrak a\) be a coherent ideal on \(X\) and \(\mu,r\in\mathbb N\) with \(\mathfrak a\subset\mathcal I_C^{\mu+r}\) (by Lemma~\ref{hir:hadamard-bridge} this holds
when \(C\subset\operatorname{cosupp}(\mathfrak a,\mu+r)\)). Then
\[
 \sigma^c(\mathcal D^r(\mathfrak a),\mu)\subset\mathcal D^r(\sigma^c(\mathfrak a,\mu+r)).
\]
\end{lemma}
\begin{proof}
\leanok
First let \(r=1\). Fix \(z\in Y\), \(x=\pi(z)\), and \(U\), \(y\) as in Lemma~\ref{mc:exceptional-derivation}; by Lemma~\ref{mc:principal-power},
\(\mathcal O_Y(-qD)\) is generated by \(y^q\) near \(z\). Write \(\mathfrak c=\sigma^c(\mathfrak a,\mu+1)\).

(i) If \(s\in\mathfrak a_x\), then \(s\circ\pi\in\pi^*\mathfrak a\subset\pi^*(\mathcal I_C^{\mu+1})=\mathcal O_Y(-(\mu+1)D)\), so \(s\circ\pi=y^{\mu+1}H\) near \(z\) with \(H\)
holomorphic, and \(H\in\mathfrak c_z\): every germ of \(\mathcal O_Y(-(\mu+1)D)_z\) is \(d\,y^{\mu+1}\), and \(d\,y^{\mu+1}H=d\,(s\circ\pi)\in\pi^*\mathfrak a\).

(ii) For every generator \(t\) of \(\mathcal D(\mathfrak a)_x\) there is \(K\in\mathcal D(\mathfrak c)_z\) with \(y\,(t\circ\pi)=y^{\mu+1}K\) near \(z\). If
\(t\in\mathfrak a_x\), take \(K=yH\) with \(H\) from (i). If \(t=\partial^{e'}_vf\) with \(f\in\mathfrak a_x\), write \(f\circ\pi=y^{\mu+1}H\) as in (i) and take \(D'\),
\(\rho\) from Lemma~\ref{mc:exceptional-derivation}; by the product rule
\[
 y\,(t\circ\pi)=D'(y^{\mu+1}H)=(\mu+1)\,y^{\mu}D'(y)\,H+y^{\mu+1}D'(H)=y^{\mu+1}\bigl((\mu+1)\rho H+D'(H)\bigr),
\]
and \(K=(\mu+1)\rho H+\sum_mb_m\,\partial^{e_z}_{e_m}H\in\mathcal D(\mathfrak c)_z\) since \(H\in\mathfrak c_z\). By linearity (ii) holds for every
\(t\in\mathcal D(\mathfrak a)_x\).

(iii) Let \(g\in\sigma^c(\mathcal D(\mathfrak a),\mu)_z\). Then \(y^\mu g\in\pi^*(\mathcal D(\mathfrak a))_z\), i.e.\ \(y^\mu g=\sum_ic_i\,(t_i\circ\pi)\) near \(z\) with
\(t_i\in\mathcal D(\mathfrak a)_x\), and by (ii) \(y^{\mu+1}g=\sum_ic_i\,y^{\mu+1}K_i\). As \(y\) vanishes identically on no nonempty open subset of \(U\),
\(g-\sum_ic_iK_i\), a continuous function vanishing on the dense set \(\{y\ne0\}\) near \(z\), vanishes near \(z\); so \(g\in\mathcal D(\mathfrak c)_z\).

For general \(r\) induct on \(r\): \(\mathcal D^r(\mathfrak a)\) is coherent (Lemma~\ref{hir:derivative-ideal-generators}) and contained in
\(\mathcal I_C^{\mu+1}\) (Lemma~\ref{mc:derivative-power}), so by the case \(r=1\) and induction (with \(\mu+1\) in place of \(\mu\)),
\(\sigma^c(\mathcal D^{r+1}(\mathfrak a),\mu)\subset\mathcal D(\sigma^c(\mathcal D^r(\mathfrak a),\mu+1))\subset\mathcal D(\mathcal D^r(\sigma^c(\mathfrak a,\mu+1+r)))\).
\end{proof}

\begin{lemma}[Villamayor: maximal order survives admissible blowups]
\label{mc:maximal-order-transform}
\uses{mc:inclusions,mc:maximal-order,hir:hadamard-bridge}
\source{wlodarczyk05-hironaka}
\lean{Hironaka.MaximalContact.IsMaximalOrder.controlledTransform,Hironaka.MaximalContact.IsMaximalOrder.controlledTransform_of_subset_cosupport}
\leanok
Let \((\mathfrak a,\mu)\) be a coherent marked ideal of maximal order with \(\mathfrak a\subset\mathcal I_C^\mu\) (for instance
\(C\subset\operatorname{cosupp}(\mathfrak a,\mu)\)). Then \((\sigma^c(\mathfrak a,\mu),\mu)\) is of maximal order.
\end{lemma}
\begin{proof}
\leanok
By Lemma~\ref{mc:inclusions} with \(r=\mu\) and mark \(0\), \(\sigma^c(\mathcal D^\mu(\mathfrak a),0)\subset\mathcal D^\mu(\sigma^c(\mathfrak a,\mu))\). Since
\(1\in\mathcal D^\mu(\mathfrak a)\) everywhere, \(1=1\circ\pi\in\pi^*\mathcal D^\mu(\mathfrak a)=\sigma^c(\mathcal D^\mu(\mathfrak a),0)\).
\end{proof}
\begin{lemma}[Giraud: the controlled transform of a tangent direction]
\label{mc:giraud-transform}
\uses{mc:inclusions,mc:derivative-power,mc:tangent-direction,bcalc:function-transform,hir:hadamard-bridge}
\source{wlodarczyk05-hironaka}
\lean{Hironaka.MaximalContactBlowup.mem_iterDerivIdeal_controlledTransform,Hironaka.MaximalContactBlowup.exists_controlledTransform_mem_iterDerivIdeal,Hironaka.MaximalContactBlowup.exists_controlledTransform_mem_iterDerivIdeal_of_subset_cosupport}
\leanok
Let \(\mathfrak a\) be coherent with \(\mathfrak a\subset\mathcal I_C^{\mu+1}\) (for instance \(C\subset\operatorname{cosupp}(\mathfrak a,\mu+1)\)), \(h\) a
local equation of the exceptional divisor near \(z\in Y\), and \(u\) a germ of \(T(\mathfrak a)=\mathcal D^\mu(\mathfrak a)\) at \(\pi(z)\). Then
\(u\circ\pi=h\,u'\) near \(z\) with \(u'\) holomorphic at \(z\), and every such \(u'\) is a germ of \(\mathcal D^\mu(\sigma^c(\mathfrak a,\mu+1))\) at \(z\).
\end{lemma}
\begin{proof}
\leanok
By Lemma~\ref{mc:derivative-power}, \(\mathcal D^\mu(\mathfrak a)\subset\mathcal I_C\), so \(u'\) exists (Lemma~\ref{bcalc:function-transform}) and
\(u'\in\sigma^c(\mathcal D^\mu(\mathfrak a),1)\). Lemma~\ref{mc:inclusions} with \(r=\mu\) and mark \(1\) gives
\(\sigma^c(\mathcal D^\mu(\mathfrak a),1)\subset\mathcal D^\mu(\sigma^c(\mathfrak a,\mu+1))\).
\end{proof}
This is part (1) of Włodarczyk's Lemma ``le:Gi''; part (3), \(V(u')\) is the strict transform of \(V(u)\), is
Lemma~\ref{bcalc:giraud-strict}, and part (2), multiplicity one of \(u'\), follows from it because a generator of the ideal of a
smooth hypersurface has nonvanishing differential; for a tangent direction defined only on an open set, parts (2)--(3) are
Lemma~\ref{bcalc:giraud-local}.

\begin{definition}[Derivatives of a marked ideal]
\label{mc:derivative-marked}
\uses{hir:marked-ideal,hir:derivative-ideal,hir:cosupport}
\source{wlodarczyk05-hironaka}
\lean{Hironaka.DerivativeIdealTower.derivMarked,Hironaka.DerivativeIdealTower.derivMarked_N,Hironaka.DerivativeIdealTower.derivMarked_E,Hironaka.DerivativeIdealTower.derivMarked_ideal,Hironaka.DerivativeIdealTower.derivMarked_d,Hironaka.DerivativeIdealTower.cosupport_subset_of_le,Hironaka.DerivativeIdealTower.controlledTransform_le_of_le,Hironaka.DerivativeIdealTower.cosupp_subset_cosupp_derivMarked}
\leanok
For a marked ideal \(\underline{\mathcal I}=(M,N,E,\mathfrak a,d)\) and \(j\in\mathbb N\), \(\mathcal D^j(\underline{\mathcal I})=(M,N,E,\mathcal D^j(\mathfrak a),d-j)\).
If \(\mathfrak b\subset\mathfrak c\), then \(\operatorname{cosupp}(\mathfrak c,d)\subset\operatorname{cosupp}(\mathfrak b,d)\) and
\(\sigma^{\mathrm c}(\mathfrak b,\mu)\subset\sigma^{\mathrm c}(\mathfrak c,\mu)\). For \(N=M\) and \(j\le d\),
\(\operatorname{cosupp}\underline{\mathcal I}\subset\operatorname{cosupp}\mathcal D^j(\underline{\mathcal I})\).
\end{definition}
The two monotonicity statements are immediate from the definitions (a germ of \(\mathfrak b\) is a germ of \(\mathfrak c\); \(\sigma^*\mathfrak b\subset\sigma^*\mathfrak c\)), and the
last one is Lemma~\ref{hir:giraud-support} applied \(j\) times, \(\operatorname{cosupp}(\mathfrak a,(d-j)+j)\subset\operatorname{cosupp}(\mathcal D^j\mathfrak a,d-j)\).

\begin{lemma}[Derivative ideals along multiple test blow-ups]
\label{mc:inclusions-tower}
\uses{mc:derivative-marked,mc:inclusions,hir:geo-transform,hir:geo-transform-ambient,hir:bridge-consumers,hir:giraud-support,hir:derivative-ideal-generators}
\source{wlodarczyk05-hironaka}
\lean{Hironaka.DerivativeIdealTower.inclusions_tower}
\leanok
Let \(\underline{\mathcal I}=(M,M,E,\mathfrak a,d)\) with \(\mathfrak a\) coherent, \(j\le d\), and let \(T\) be a multiple test blow-up of \(\underline{\mathcal I}\)
(Definition~\ref{hir:geo-transform}). Then \(T\) is a multiple test blow-up of \(\mathcal D^j(\underline{\mathcal I})\), and
\[
 [\mathcal D^j(\underline{\mathcal I})]_T\subset\mathcal D^j(\underline{\mathcal I}_T),
\]
i.e.\ the transformed ideal of \(\mathcal D^j(\underline{\mathcal I})\) is contained in \(\mathcal D^j\) of the transformed ideal of \(\underline{\mathcal I}\); both transforms
have \(N\) the whole space, the same \(E\), and coherent ideals.
\end{lemma}
\begin{proof}
\leanok
Induction on \(T\); for the empty tower both sides are \(\mathcal D^j(\mathfrak a)\), coherent by Lemma~\ref{hir:derivative-ideal-generators}. Let \(T\) be followed by
the blowup \(\sigma\) along an admissible centre \(C\) of \(\underline{\mathcal I}_T=(M_T,M_T,E_T,\mathfrak a_T,d)\), and write
\(\mathfrak b_T\subset\mathcal D^j(\mathfrak a_T)\) for the transformed ideal of \(\mathcal D^j(\underline{\mathcal I})\), with mark \(d-j\) and the same \(E_T\).
\(C\) is admissible for \((M_T,M_T,E_T,\mathfrak b_T,d-j)\): it is a closed submanifold with normal crossings with \(E_T\), and
\(C\subset\operatorname{cosupp}(\mathfrak a_T,d)\subset\operatorname{cosupp}(\mathcal D^j\mathfrak a_T,d-j)\subset\operatorname{cosupp}(\mathfrak b_T,d-j)\)
(Definition~\ref{mc:derivative-marked}). By Lemma~\ref{hir:bridge-consumers}, \(\mathfrak a_T\subset\mathcal I_C^{d}\) and \(\mathfrak b_T\subset\mathcal I_C^{d-j}\), so both
transformed ideals are the coherent controlled transforms (Lemma~\ref{hir:geo-transform-ambient}), and by monotonicity and
Lemma~\ref{mc:inclusions} (with mark \(d-j\) and \(r=j\))
\[
 \sigma^{\mathrm c}(\mathfrak b_T,d-j)\subset\sigma^{\mathrm c}(\mathcal D^j\mathfrak a_T,d-j)\subset\mathcal D^j(\sigma^{\mathrm c}(\mathfrak a_T,d)).
\]
\(N\) stays the whole space and both \(E\) are transformed alike.
\end{proof}
This is Włodarczyk's Lemma ``le:inclusions2''. His Lemma ``le:Giraud'' (the cosupports of the transforms lie on the strict transforms of a
hypersurface of maximal contact) is Theorem~\ref{mc:giraud-tower} below: it follows from Lemma~\ref{mc:giraud-transform} and parts (2)--(3) of
``le:Gi'' on open sets (Lemma~\ref{bcalc:giraud-local}) by induction on the tower.

\subsection{Homogenized ideals and the glueing inclusion}

In this subsection \((\mathfrak a,\mu)\) is a coherent marked ideal on a complex manifold \(M\) (with \(N=M\)), \(\mu\ge1\), and
\(T(\mathfrak a)=\mathcal D^{\mu-1}(\mathfrak a)\).

\begin{definition}[Homogenized ideal]
\label{mc:homogenized-ideal}
\uses{hir:derivative-ideal}
\source{wlodarczyk05-hironaka,wlodarczyk09-analytic-resolution:page-0020}
\lean{Hironaka.HomogenizedIdeal.tangentIdeal,Hironaka.HomogenizedIdeal.HomGen,Hironaka.HomogenizedIdeal.homogenizedIdeal}
\leanok
The \emph{homogenized ideal} of \((\mathfrak a,\mu)\) is
\[
 \mathcal H(\mathfrak a)=\mathfrak a+\mathcal D(\mathfrak a)\cdot T(\mathfrak a)+\dots+\mathcal D^i(\mathfrak a)\cdot T(\mathfrak a)^i+\dots+\mathcal D^{\mu-1}(\mathfrak a)\cdot T(\mathfrak a)^{\mu-1},
\]
i.e.\ \(\mathcal H(\mathfrak a)_x\) is generated by the germs of \(\mathcal D^i(\mathfrak a)\cdot T(\mathfrak a)^i\) at \(x\), \(i<\mu\).
\end{definition}

\begin{lemma}[Elementary properties of the homogenized ideal]
\label{mc:homogenized-basic}
\uses{mc:homogenized-ideal}
\source{wlodarczyk05-hironaka,wlodarczyk09-analytic-resolution:page-0020}
\lean{Hironaka.HomogenizedIdeal.mem_of_isInSpanAt,Hironaka.HomogenizedIdeal.mem_congr,Hironaka.HomogenizedIdeal.mul_mem_mul,Hironaka.HomogenizedIdeal.pow_mem_idealPow,Hironaka.HomogenizedIdeal.iterDerivIdeal_succ_mem_congr,Hironaka.HomogenizedIdeal.le_homogenizedIdeal,Hironaka.HomogenizedIdeal.homogenizedIdeal_one}
\leanok
\(\mathfrak a\subset\mathcal H(\mathfrak a)\), and \(\mathcal H(\mathfrak a)=\mathfrak a\) if \(\mu=1\).
\end{lemma}
\begin{proof}
\leanok
The summand \(i=0\) is \(\mathfrak a\cdot T(\mathfrak a)^0=\mathfrak a\cdot\mathcal O=\mathfrak a\). For \(\mu=1\) it is the only summand, and the coherent ideal
\(\mathfrak a\) is closed under holomorphic combinations of its germs.
\end{proof}

\begin{lemma}[Iterated coordinate derivatives]
\label{mc:coordinate-iterated}
\uses{mc:coordinate-derivatives,hir:derivative-ideal}
\lean{Hironaka.HomogenizedIdeal.iterDeriv_comp_mem_iterDerivIdeal}
\leanok
Let \(W\) be local holomorphic coordinates at \(x\) (as in Lemma~\ref{mc:coordinate-derivatives}(1)) and \(F\) holomorphic near \(W(x)\) with
\(F\circ W\in\mathfrak a_x\). Then \((\partial_l^jF)\circ W\in\mathcal D^j(\mathfrak a)_x\) for every \(l\) and \(j\).
\end{lemma}
\begin{proof}
\leanok
Induction on \(j\). By Lemma~\ref{mc:coordinate-derivatives}(1) applied to \(G=\partial_l^jF\),
\((\partial_l^{j+1}F)\circ W=\sum_mc_{lm}\,\partial^{e}_{e_m}\bigl((\partial_l^jF)\circ W\bigr)\) near \(x\), a holomorphic combination of partial derivatives of
a germ of \(\mathcal D^j(\mathfrak a)_x\), hence a germ of \(\mathcal D^{j+1}(\mathfrak a)_x\).
\end{proof}

\begin{lemma}[The glueing inclusion]
\label{mc:glueing-inclusion}
\uses{mc:homogenized-ideal,mc:homogenized-basic,mc:coordinate-iterated,mc:taylor-tail}
\source{wlodarczyk05-hironaka,wlodarczyk09-analytic-resolution:page-0021}
\lean{Hironaka.HomogenizedIdeal.comp_mem_homogenizedIdeal}
\leanok
Let \(W\) be local holomorphic coordinates at \(x\) with a holomorphic local inverse \(\Psi\) (\(\Psi\circ W=e\) near \(x\), \(e\) the chart at
\(x\)), let \(h\in T(\mathfrak a)_x\) with \(h(x)=0\), let \(1\le i\le n\), and let
\[
 \varphi(p)=e^{-1}\Psi\bigl(W(p)+h(p)\,e_i\bigr)
\]
be the coordinate change near \(x\) that replaces the coordinate \(W_i\) by \(W_i+h\). Then \(f\circ\varphi\in\mathcal H(\mathfrak a)_x\) for every
\(f\in\mathfrak a_x\).
\end{lemma}
\begin{proof}
\leanok
Put \(F=f\circ e^{-1}\circ\Psi\), holomorphic near \(W(x)\) with \(F\circ W=f\) near \(x\). By Lemma~\ref{mc:taylor-tail} there is \(G\), holomorphic near
\((W(x),0)\), with \(F(w+te_i)=\sum_{j<\mu}\frac1{j!}\,\partial_i^jF(w)\,t^j+t^\mu G(w,t)\) near \((W(x),0)\). Since \((W(p),h(p))\to(W(x),0)\), near \(x\)
\[
 f\circ\varphi=\sum_{j<\mu}\frac1{j!}\,\bigl((\partial_i^jF)\circ W\bigr)\,h^j+G(W,h)\cdot h\cdot h^{\mu-1}.
\]
By Lemma~\ref{mc:coordinate-iterated}, \((\partial_i^jF)\circ W\in\mathcal D^j(\mathfrak a)_x\), and \(h^j\in T(\mathfrak a)^j_x\); so the \(j\)-th term is a multiple of a germ of
\(\mathcal D^j(\mathfrak a)\cdot T(\mathfrak a)^j\), and the last one of a germ of \(\mathcal D^{\mu-1}(\mathfrak a)\cdot T(\mathfrak a)^{\mu-1}\) (as \(h\in T(\mathfrak a)=\mathcal D^{\mu-1}(\mathfrak a)\)).
\(G(W,h)\) is holomorphic at \(x\).
\end{proof}
\begin{lemma}[Products with powers of the tangent ideal]
\label{mc:homogenized-product}
\uses{mc:homogenized-ideal,mc:maximal-order-derivative,hir:order-basic}
\source{wlodarczyk05-hironaka}
\lean{Hironaka.HomogenizedIdeal.mul_mem_homogenizedIdeal}
\leanok
For \(i<\mu\), \(\mathcal H(\mathcal D^i(\mathfrak a),\mu-i)\cdot T(\mathfrak a)^i\subset\mathcal H(\mathfrak a,\mu)\).
\end{lemma}
\begin{proof}
\leanok
The tangent ideal of \((\mathcal D^i(\mathfrak a),\mu-i)\) is \(\mathcal D^{\mu-i-1}(\mathcal D^i(\mathfrak a))=\mathcal D^{\mu-1}(\mathfrak a)=T(\mathfrak a)\)
(Lemma~\ref{mc:maximal-order-derivative}). So a generator of \(\mathcal H(\mathcal D^i(\mathfrak a),\mu-i)\cdot T(\mathfrak a)^i\) is a combination of products
\(a\,b'\,b\) with \(a\in\mathcal D^{i+j}(\mathfrak a)\), \(b'\in T(\mathfrak a)^j\), \(b\in T(\mathfrak a)^i\) and \(j<\mu-i\); and \(b'b\in T(\mathfrak a)^{i+j}\)
(Lemma~\ref{hir:order-basic}(1)), so \(ab'b\) is a germ of \(\mathcal D^{i+j}(\mathfrak a)\cdot T(\mathfrak a)^{i+j}\) with \(i+j<\mu\).
\end{proof}

\begin{definition}[Coordinate change]
\label{mc:coordinate-change}
\uses{mc:coordinate-derivatives}
\source{wlodarczyk05-hironaka,wlodarczyk09-analytic-resolution:page-0021}
\lean{Hironaka.HomogenizedIdeal.CoordinateChange,Hironaka.HomogenizedIdeal.CoordinateChange.coordMap,Hironaka.HomogenizedIdeal.CoordinateChange.continuousAt_W,Hironaka.HomogenizedIdeal.CoordinateChange.coordMap_self,Hironaka.HomogenizedIdeal.CoordinateChange.exists_chart_form,Hironaka.HomogenizedIdeal.CoordinateChange.continuousAt_coordMap,Hironaka.HomogenizedIdeal.CoordinateChange.holomorphicAt_comp,Hironaka.HomogenizedIdeal.CoordinateChange.comp_eventuallyEq,Hironaka.HomogenizedIdeal.CoordinateChange.comp_isInSpanAt,Hironaka.HomogenizedIdeal.CoordinateChange.comp_mem_homogenizedIdeal'}
\leanok
A \emph{coordinate change at \(x\)} consists of local holomorphic coordinates \(W\) at \(x\) with a holomorphic local inverse \(\Psi\)
(\(\Psi\circ W=e\) near \(x\)), a function \(h\) holomorphic at \(x\) with \(h(x)=0\), and an index \(i\); its map is
\(\varphi(p)=e^{-1}\Psi(W(p)+h(p)e_i)\). Then \(\varphi(x)=x\), \(\varphi\) is holomorphic near \(x\), and composition with \(\varphi\) carries germs at \(x\) to
germs at \(x\), holomorphic germs to holomorphic germs and holomorphic combinations to holomorphic combinations.
\end{definition}

\begin{proposition}[The coordinate change preserves the homogenized ideal]
\label{mc:homogenized-invariance}
\uses{mc:coordinate-change,mc:glueing-inclusion,mc:homogenized-basic,mc:homogenized-product,hir:derivative-ideal-generators}
\source{wlodarczyk05-hironaka,wlodarczyk09-analytic-resolution:page-0021}
\lean{Hironaka.HomogenizedIdeal.CoordinateChange.comp_mem_idealPow_tangentIdeal,Hironaka.HomogenizedIdeal.CoordinateChange.comp_mem_homogenizedIdeal_of_mem}
\leanok
Let \(\varphi\) be a coordinate change at \(x\) with \(h\in T(\mathfrak a)_x\). Then \(f\circ\varphi\in\mathcal H(\mathfrak a)_x\) for every \(f\in\mathcal H(\mathfrak a)_x\), and
\(b\circ\varphi\in(T(\mathfrak a)^j)_x\) for every \(b\in(T(\mathfrak a)^j)_x\).
\end{proposition}
\begin{proof}
\leanok
For \(T(\mathfrak a)\): the marked ideal \((T(\mathfrak a),1)\) is coherent (Lemma~\ref{hir:derivative-ideal-generators}), its tangent ideal is \(T(\mathfrak a)\ni h\), and
\(\mathcal H(T(\mathfrak a),1)=T(\mathfrak a)\) (Lemma~\ref{mc:homogenized-basic}); so Lemma~\ref{mc:glueing-inclusion} gives \(t\circ\varphi\in T(\mathfrak a)_x\) for
\(t\in T(\mathfrak a)_x\), and products and combinations give the claim for \(T(\mathfrak a)^j\). For \(\mathcal H(\mathfrak a)\) it suffices to treat a product
\(ab\) with \(a\in\mathcal D^i(\mathfrak a)_x\), \(b\in(T(\mathfrak a)^i)_x\), \(i<\mu\). The marked ideal \((\mathcal D^i(\mathfrak a),\mu-i)\) is coherent with tangent ideal
\(T(\mathfrak a)\ni h\), so Lemma~\ref{mc:glueing-inclusion} gives \(a\circ\varphi\in\mathcal H(\mathcal D^i(\mathfrak a),\mu-i)_x\); with \(b\circ\varphi\in(T(\mathfrak a)^i)_x\),
Lemma~\ref{mc:homogenized-product} gives \((ab)\circ\varphi\in\mathcal H(\mathfrak a)_x\).
\end{proof}
In Włodarczyk's Lemma ``le:homo0'' the coordinates are \(W=(u,u_2,\dots,u_n)\) for a tangent direction \(u\), and \(\varphi\) is the coordinate change with
\(u\circ\varphi=v\) for a second tangent direction \(v=u+h\), \(h\in T(\mathfrak a)\); the formal automorphism of the source becomes the convergent \(\varphi\), and
Proposition~\ref{mc:homogenized-invariance} is the inclusion \(\varphi^*\mathcal H(\mathfrak a)\subset\mathcal H(\mathfrak a)\) of its part (1). Equality, from the
inverse change, is part of Theorem~\ref{mc:glueing-lemma}.

\subsection{The glueing lemma}

In this subsection \((\mathfrak a,\mu)\) is a coherent marked ideal on \(M\) (with \(N=M\)), \(\mu\ge1\) and \(T(\mathfrak a)=\mathcal D^{\mu-1}(\mathfrak a)\). For a
map \(\varphi\colon M\to M\) and an ideal \(\mathfrak b\), \(\varphi^*\mathfrak b\) is the ideal generated at \(y\) by the germs \(g\circ\varphi\) with
\(g\in\mathfrak b_{\varphi(y)}\).

\begin{lemma}[Coherent ideals near a point]
\label{mc:coherent-germs}
\uses{ana:ideal-sheaves}
\lean{Hironaka.HomogenizedGlueing.holomorphicAt_of_mem,Hironaka.HomogenizedGlueing.eventually_mem,Hironaka.HomogenizedGlueing.lin_mem,Hironaka.HomogenizedGlueing.sub_mem,Hironaka.HomogenizedGlueing.isInSpanAt_comp,Hironaka.HomogenizedGlueing.eventually_pullback_le}
\leanok
Let \(\mathfrak b\) be a coherent ideal and \(x\in M\).
\begin{enumerate}
\item If \(g\in\mathfrak b_x\), then \(g\in\mathfrak b_y\) for every \(y\) near \(x\). The germs of \(\mathfrak b_x\) are closed under complex linear combinations.
\item Let \(\varphi\) be holomorphic near \(x\) with \(\varphi(x)=x\), and assume \(g\circ\varphi\in\mathfrak b_x\) for every \(g\in\mathfrak b_x\). Then
\((\varphi^*\mathfrak b)_y\subset\mathfrak b_y\) for every \(y\) near \(x\).
\end{enumerate}
\end{lemma}
\begin{proof}
\leanok
Let \(s_1,\dots,s_N\), holomorphic on an open set \(U\ni x\), generate \(\mathfrak b\) at every point of \(U\).
(1) Near \(x\), \(g=\sum_kc_ks_k\) with \(c_k\) holomorphic near \(x\); the same identity holds near every \(y\) close to \(x\), so \(g\in\mathfrak b_y\).
A linear combination of germs of \(\mathfrak b_x\) is a holomorphic combination of the \(s_k\).
(2) Each \(s_k\) is a germ of \(\mathfrak b_x\), so \(s_k\circ\varphi\in\mathfrak b_x\) and, by (1), \(s_k\circ\varphi\in\mathfrak b_y\) for \(y\) near \(x\); for such
\(y\) also \(\varphi(y)\in U\) and \(\varphi\) is holomorphic near \(y\). A generator \(g\circ\varphi\) of \((\varphi^*\mathfrak b)_y\), \(g\in\mathfrak b_{\varphi(y)}\),
satisfies \(g=\sum_kc_ks_k\) near \(\varphi(y)\), hence \(g\circ\varphi=\sum_k(c_k\circ\varphi)(s_k\circ\varphi)\) near \(y\), a holomorphic combination of
germs of \(\mathfrak b_y\).
\end{proof}

\begin{lemma}[The homogenized ideal is coherent]
\label{mc:homogenized-coherent}
\uses{mc:homogenized-ideal,hir:derivative-ideal-generators,ana:ideal-sheaves}
\lean{Hironaka.HomogenizedGlueing.finSumIdeal,Hironaka.HomogenizedGlueing.finSumIdeal_succ,Hironaka.HomogenizedGlueing.isCoherent_finSumIdeal,Hironaka.HomogenizedGlueing.isCoherent_homogenizedIdeal}
\leanok
If \(\mathfrak a\) is coherent, then \(\mathcal H(\mathfrak a)\) is coherent.
\end{lemma}
\begin{proof}
\leanok
If \(\mathfrak b\) and \(\mathfrak c\) are generated at every point of open sets \(U\) and \(U'\) by \(s_1,\dots,s_N\) and \(t_1,\dots,t_P\), then at every point
of \(U\cap U'\) the sum \(\mathfrak b+\mathfrak c\) (generated stalkwise by the germs of \(\mathfrak b\) and of \(\mathfrak c\)) is generated by the \(s_k\) and the
\(t_l\), and the product \(\mathfrak b\mathfrak c\) by the \(s_kt_l\); so finite sums, products and powers of coherent ideals are coherent. The ideals
\(\mathcal D^i(\mathfrak a)\) are coherent (Lemma~\ref{hir:derivative-ideal-generators}), hence so are the summands \(\mathcal D^i(\mathfrak a)\cdot T(\mathfrak a)^i\) and,
by induction on the number of summands, \(\mathcal H(\mathfrak a)\).
\end{proof}

\begin{lemma}[Coordinates from two tangent directions]
\label{mc:coordinate-pair}
\uses{mc:coordinate-change,ana:local-holomorphic-inverse,hir:derivative-ideal-generators}
\source{wlodarczyk05-hironaka,wlodarczyk09-analytic-resolution:page-0021}
\lean{Hironaka.HomogenizedGlueing.exists_apply_ne_zero_and,Hironaka.HomogenizedGlueing.killMap,Hironaka.HomogenizedGlueing.killMap_apply,Hironaka.HomogenizedGlueing.killMap_apply_self,Hironaka.HomogenizedGlueing.eq_smul_of_killMap_eq_zero,Hironaka.HomogenizedGlueing.pairCoord,Hironaka.HomogenizedGlueing.pairCoord_apply_self,Hironaka.HomogenizedGlueing.pairCoord_sub,Hironaka.HomogenizedGlueing.fderiv_ne_zero_of_chartDeriv,Hironaka.HomogenizedGlueing.pairCoord_coordinates,Hironaka.HomogenizedGlueing.exists_localInverse,Hironaka.HomogenizedGlueing.holomorphicAt_sub,Hironaka.HomogenizedGlueing.CoordinatePair,Hironaka.HomogenizedGlueing.exists_coordinatePair,Hironaka.HomogenizedGlueing.eventually_mdifferentiableAt_coordMap,Hironaka.HomogenizedGlueing.eventually_coordMap_eq_self}
\leanok
Let \(u\), \(v\) be holomorphic at \(x\) with \(u(x)=v(x)\), \(du(x)\ne0\) and \(dv(x)\ne0\). There are local holomorphic coordinates
\(W=(u,u_2,\dots,u_n)\) and \(W'=(v,u_2,\dots,u_n)\) at \(x\) (up to the order of the coordinates) such that the coordinate changes
(Definition~\ref{mc:coordinate-change}) \(\varphi\), given by \(W\), \(h=v-u\) and the index of \(u\), and \(\psi\), given by \(W'\), \(h'=u-v\) and the same
index, satisfy near \(x\): \(\varphi\) and \(\psi\) are holomorphic, \(\psi\circ\varphi=\mathrm{id}\), \(\varphi\circ\psi=\mathrm{id}\), \(u\circ\varphi=v\) and
\(v\circ\psi=u\). Moreover a coordinate change fixes every point near \(x\) at which its function \(h\) vanishes.
\end{lemma}
\begin{proof}
\leanok
Let \(e\) be the chart at \(x\), \(\ell=d(u\circ e^{-1})(e(x))\) and \(\ell'=d(v\circ e^{-1})(e(x))\); both are nonzero, because a partial derivative in
another chart is a holomorphic combination of the partial derivatives in \(e\) (Lemma~\ref{hir:derivative-ideal-generators}). Choose
\(w\in\mathbf C^n\) with \(\ell(w)\ne0\ne\ell'(w)\): take \(w_1\), \(w_2\) with \(\ell(w_1)\ne0\) and \(\ell'(w_2)\ne0\); if \(\ell'(w_1)=0\) and
\(\ell(w_2)=0\), then \(w=w_1+w_2\) works, otherwise \(w_1\) or \(w_2\) does. Let \(k\) be an index with \(w_k\ne0\) and
\(P(\zeta)=\zeta-(\zeta_k/w_k)\,w\), a linear map with \(P(\zeta)_k=0\) and kernel \(\mathbf Cw\). Put \(W=P\circ e+u\,e_k\) and
\(W'=P\circ e+v\,e_k=W+(v-u)\,e_k\); their coordinates other than the \(k\)-th are the common functions \((P\circ e)_j\). The derivative of
\(W\circ e^{-1}\) at \(e(x)\) is \(\xi\mapsto P\xi+\ell(\xi)\,e_k\). If it vanishes on \(\xi\), its \(k\)-th coordinate gives \(\ell(\xi)=0\), then
\(P\xi=0\), so \(\xi=cw\) and \(0=\ell(\xi)=c\,\ell(w)\) gives \(\xi=0\). So \(W\), and likewise \(W'\), are local coordinates at \(x\), and by the
inverse function theorem (Lemma~\ref{ana:local-holomorphic-inverse}) \(W\circ e^{-1}\) has a holomorphic inverse \(\Psi\) near \(W(x)\):
\(\Psi\circ W=e\) near \(x\) and \(W\circ e^{-1}\circ\Psi=\mathrm{id}\) near \(W(x)\); similarly \(\Psi'\) for \(W'\). Then
\(\varphi=e^{-1}\Psi(W+(v-u)e_k)=e^{-1}\Psi W'\) and \(\psi=e^{-1}\Psi'W\). Since \(W'(x)=W(x)\), near \(x\) we have \(W\circ\varphi=W'\), so
\(u\circ\varphi=(W\circ\varphi)_k=W'_k=v\) and \(\psi\circ\varphi=e^{-1}\Psi'(W\circ\varphi)=e^{-1}\Psi'W'=e^{-1}e=\mathrm{id}\); symmetrically
\(v\circ\psi=u\) and \(\varphi\circ\psi=\mathrm{id}\). In the chart \(e\), a coordinate change is \(e^{-1}\circ H\circ e\) with \(H\) holomorphic near \(e(x)\) and
\(H(e(x))=e(x)\) (Definition~\ref{mc:coordinate-change}), hence holomorphic near \(x\). Where its \(h\) vanishes it is \(e^{-1}\Psi W=e^{-1}e=\mathrm{id}\).
\end{proof}

\begin{theorem}[Glueing Lemma]
\label{mc:glueing-lemma}
\uses{mc:coordinate-pair,mc:homogenized-invariance,mc:homogenized-coherent,mc:coherent-germs,mc:tangent-direction,mc:tangent-direction-support,hir:derivative-ideal-generators}
\source{wlodarczyk05-hironaka,wlodarczyk09-analytic-resolution:page-0021}
\lean{Hironaka.HomogenizedGlueing.CoordinatePair.mem_tangentIdeal_h,Hironaka.HomogenizedGlueing.CoordinatePair.mem_tangentIdeal_h',Hironaka.HomogenizedGlueing.CoordinatePair.comp_mem_homogenizedIdeal_iff,Hironaka.HomogenizedGlueing.CoordinatePair.eventually_coordMap_eq_self_of_mem_cosupport,Hironaka.HomogenizedGlueing.CoordinatePair.eventually_pullback_homogenizedIdeal,Hironaka.HomogenizedGlueing.glueing_lemma}
\leanok
Let \(u\) and \(v\) be tangent directions of \((\mathfrak a,\mu)\) on open sets containing \(x\in\operatorname{cosupp}(\mathfrak a,\mu)\), and let
\(\varphi\), \(\psi\) be the coordinate changes of Lemma~\ref{mc:coordinate-pair}. Then, near \(x\):
\begin{enumerate}
\item \(\varphi^*\mathcal H(\mathfrak a)=\mathcal H(\mathfrak a)\), i.e.\ the two ideals have the same germs at every point near \(x\); at \(x\),
\(f\circ\varphi\in\mathcal H(\mathfrak a)_x\) if and only if \(f\in\mathcal H(\mathfrak a)_x\);
\item \(u\circ\varphi=v\);
\item \(\varphi\) and \(\psi\) fix every point of \(\operatorname{cosupp}(\mathfrak a,\mu)\).
\end{enumerate}
\end{theorem}
\begin{proof}
\leanok
The germs of \(u\) and \(v\) at \(x\) lie in \(T(\mathfrak a)_x\), so they vanish at \(x\) (Lemma~\ref{mc:tangent-direction-support}) and, being tangent
directions, have nonzero differentials there; Lemma~\ref{mc:coordinate-pair} applies. As \(T(\mathfrak a)\) is coherent
(Lemma~\ref{hir:derivative-ideal-generators}), \(h=v-u\) and \(-h\) are germs of \(T(\mathfrak a)\) at \(x\) (Lemma~\ref{mc:coherent-germs}(1)).
(2) is part of Lemma~\ref{mc:coordinate-pair}.
(1) At \(x\): Proposition~\ref{mc:homogenized-invariance} for \(\varphi\) gives \(f\circ\varphi\in\mathcal H(\mathfrak a)_x\) for \(f\in\mathcal H(\mathfrak a)_x\);
conversely, if \(f\circ\varphi\in\mathcal H(\mathfrak a)_x\), the same proposition for \(\psi\) gives \(f=(f\circ\varphi)\circ\psi\in\mathcal H(\mathfrak a)_x\).
Near \(x\): \(\mathcal H(\mathfrak a)\) is coherent (Lemma~\ref{mc:homogenized-coherent}), so Lemma~\ref{mc:coherent-germs}(2), applied to \(\varphi\) and to
\(\psi\), gives \(\varphi^*\mathcal H(\mathfrak a)\subset\mathcal H(\mathfrak a)\) and \(\psi^*\mathcal H(\mathfrak a)\subset\mathcal H(\mathfrak a)\) at every point near \(x\).
For \(y\) near \(x\) and \(g\in\mathcal H(\mathfrak a)_y\), the germ of \(g\circ\psi\) at \(\varphi(y)\) is a generator of
\((\psi^*\mathcal H(\mathfrak a))_{\varphi(y)}\subset\mathcal H(\mathfrak a)_{\varphi(y)}\) (as \(\psi(\varphi(y))=y\)), and \(g=(g\circ\psi)\circ\varphi\) near \(y\);
so \(g\in(\varphi^*\mathcal H(\mathfrak a))_y\).
(3) By Lemma~\ref{mc:coherent-germs}(1), the germs of \(u\) and \(v\) at every \(y\) near \(x\) lie in \(T(\mathfrak a)_y\), so \(u\) and \(v\) vanish on
\(\operatorname{cosupp}(\mathfrak a,\mu)\) near \(x\) (Lemma~\ref{mc:tangent-direction-support}); there \(h=0\), and \(\varphi\) and \(\psi\) fix the point
(Lemma~\ref{mc:coordinate-pair}).
\end{proof}
This is Włodarczyk's Glueing Lemma ``le:homo'', parts (1)--(3), for \(E=\emptyset\), the only case the source uses: its étale neighbourhoods
\(\phi_u,\phi_v\colon\overline X\to X\) become an open neighbourhood of \(x\) with \(\phi_v=\mathrm{id}\) and \(\phi_u=\varphi\). The hypothesis of maximal
order in the source serves only to provide tangent directions (Lemma~\ref{mc:tangent-direction-exists}). Part (3) is stated for
\(\operatorname{cosupp}(\mathfrak a,\mu)\), which contains \(\operatorname{cosupp}(\mathcal H(\mathfrak a),\mu)\) since \(\mathfrak a\subset\mathcal H(\mathfrak a)\).
Part (4) of the source, the equivariance of the multiple test blow-ups of \((\mathcal H(\mathfrak a),\mu)\) under \(\varphi\), is not part of this theorem;
its pullback form is Theorem~\ref{mc:glueing-tower-pullback} below.

\subsection{The glueing lemma along multiple test blow-ups}

In this subsection \((\mathfrak a,\mu)\) is a coherent marked ideal on \(M\) (with \(N=M\)), \(\mu\ge1\), and
\(\underline{\mathcal H}=(M,M,\emptyset,\mathcal H(\mathfrak a),\mu)\) is the homogenized marked ideal. For an open \(U\subset M\),
\(\underline{\mathcal H}|_U\) is its restriction (Lemma~\ref{hir:open-restriction-marked}).

\begin{lemma}[Biholomorphisms on explicit open sets]
\label{mc:local-inverse-biholomorphism}
\uses{hir:open-biholomorphism}
\lean{Hironaka.HomogenizedGlueingTower.LocalInversePair,Hironaka.HomogenizedGlueingTower.LocalInversePair.continuousOn_left,Hironaka.HomogenizedGlueingTower.LocalInversePair.continuousOn_right,Hironaka.HomogenizedGlueingTower.LocalInversePair.dom,Hironaka.HomogenizedGlueingTower.LocalInversePair.cod,Hironaka.HomogenizedGlueingTower.LocalInversePair.mem_dom_iff,Hironaka.HomogenizedGlueingTower.LocalInversePair.mem_cod_iff,Hironaka.HomogenizedGlueingTower.LocalInversePair.apply_mem_cod,Hironaka.HomogenizedGlueingTower.LocalInversePair.apply_mem_dom,Hironaka.HomogenizedGlueingTower.LocalInversePair.contMDiff_dom,Hironaka.HomogenizedGlueingTower.LocalInversePair.contMDiff_cod,Hironaka.HomogenizedGlueingTower.LocalInversePair.toDiffeomorph,Hironaka.HomogenizedGlueingTower.LocalInversePair.toDiffeomorph_apply,Hironaka.HomogenizedGlueingTower.LocalInversePair.toDiffeomorph_symm_apply}
\leanok
Let \(W\subset M\) be open and let \(\varphi,\psi\colon M\to M\) be holomorphic at every point of \(W\) with \(\psi\circ\varphi=\mathrm{id}\) and
\(\varphi\circ\psi=\mathrm{id}\) on \(W\). Then \(U=W\cap\varphi^{-1}(W)\) and \(U'=W\cap\psi^{-1}(U)\) are open, and \(\varphi\) restricts to a biholomorphism
\(\Phi\colon U\to U'\) with inverse \(\psi\).
\end{lemma}
\begin{proof}
\leanok
\(\varphi\) and \(\psi\) are continuous on the open set \(W\), so \(U\) and \(U'\) are open. For \(p\in U\), \(\varphi(p)\in W\) and \(\psi(\varphi(p))=p\in U\), so
\(\varphi(U)\subset U'\); \(\psi(U')\subset U\) by definition, and the two compositions are the identity because \(U,U'\subset W\). Both restrictions are
holomorphic, hence smooth (Lemma~\ref{hir:open-biholomorphism}), as maps into the open submanifolds \(U'\) and \(U\).
\end{proof}
Unlike Lemma~\ref{hir:open-biholomorphism}, the open sets are explicit, so that properties of \(\varphi\) on \(W\) hold on \(U\).

\begin{lemma}[The glueing biholomorphism]
\label{mc:glueing-biholomorphism}
\uses{mc:glueing-lemma,mc:local-inverse-biholomorphism,mc:homogenized-coherent,mc:homogenized-basic,hir:open-restriction-marked,hir:marked-ideal}
\source{wlodarczyk05-hironaka,wlodarczyk09-analytic-resolution:page-0021}
\lean{Hironaka.HomogenizedGlueingTower.homogenizedMarked,Hironaka.HomogenizedGlueingTower.cosupp_homogenizedMarked_subset,Hironaka.HomogenizedGlueingTower.GlueingData,Hironaka.HomogenizedGlueingTower.exists_glueingData,Hironaka.HomogenizedGlueingTower.GlueingData.dom,Hironaka.HomogenizedGlueingTower.GlueingData.cod,Hironaka.HomogenizedGlueingTower.GlueingData.equiv,Hironaka.HomogenizedGlueingTower.GlueingData.equiv_apply,Hironaka.HomogenizedGlueingTower.GlueingData.mem_dom_and_mem_cod,Hironaka.HomogenizedGlueingTower.GlueingData.apply_eq_of_mem_cosupport,Hironaka.HomogenizedGlueingTower.GlueingData.preimage_inter_eq,Hironaka.HomogenizedGlueingTower.GlueingData.preimage_zeroSet,Hironaka.HomogenizedGlueingTower.GlueingData.pullback_comp_val,Hironaka.HomogenizedGlueingTower.GlueingData.pullback_restrict,Hironaka.HomogenizedGlueingTower.GlueingData.preimage_inter_eq_of_isAdmissibleCentre}
\leanok
Let \(u\) and \(v\) be tangent directions of \((\mathfrak a,\mu)\) on open sets containing \(x\in\operatorname{cosupp}(\mathfrak a,\mu)\). There are an open
\(W\ni x\) and maps \(\varphi,\psi\) as in Lemma~\ref{mc:local-inverse-biholomorphism} such that, with \(\Phi\colon U\to U'\) the resulting
biholomorphism, \(x\in U\cap U'\), \(\Phi(x)=x\), and
\begin{enumerate}
\item \(\Phi^*(\underline{\mathcal H}|_{U'})=\underline{\mathcal H}|_U\);
\item \(\Phi^{-1}(V(u)\cap U')=V(v)\cap U\);
\item \(\Phi\) fixes every point of \(\operatorname{cosupp}(\mathfrak a,\mu)\cap U\), and \(\Phi^{-1}(C\cap U')=C\cap U\) for every
\(C\subset\operatorname{cosupp}(\mathfrak a,\mu)\), in particular for every admissible centre \(C\) of \(\underline{\mathcal H}\).
\end{enumerate}
\end{lemma}
\begin{proof}
\leanok
Theorem~\ref{mc:glueing-lemma} gives \(\varphi\), \(\psi\) and a neighbourhood of \(x\) on which \(\varphi,\psi\) are holomorphic and mutually inverse,
\(\varphi^*\mathcal H(\mathfrak a)=\mathcal H(\mathfrak a)\) stalkwise, \(u\circ\varphi=v\), and \(\varphi\) fixes \(\operatorname{cosupp}(\mathfrak a,\mu)\); let
\(W\ni x\) be an open subset of it. As \(\varphi(x)=x\), \(x\in U\); then \(\psi(x)=\psi(\varphi(x))=x\in U\), so \(x\in U'\).
(1) Both marked ideals have \(N\) the whole space, \(E=\emptyset\) and mark \(\mu\). The ideal of \(\Phi^*(\underline{\mathcal H}|_{U'})\) is the pullback of
\(\mathcal H(\mathfrak a)\) along \(\varphi|_U\) (pullback is functorial for coherent ideals; \(\mathcal H(\mathfrak a)\) is coherent by
Lemma~\ref{mc:homogenized-coherent}). At \(p\in U\), a generator \(b\circ\varphi\), \(b\in\mathcal H(\mathfrak a)_{\varphi(p)}\), is a germ of
\((\varphi^*\mathcal H(\mathfrak a))_p=\mathcal H(\mathfrak a)_p\), so its restriction to \(U\) is a generator of \((\mathcal H(\mathfrak a)|_U)_p\); conversely a germ
\(g\in\mathcal H(\mathfrak a)_p=(\varphi^*\mathcal H(\mathfrak a))_p\) is a holomorphic combination of such \(b\circ\varphi\), and so is its restriction to \(U\).
(2) \(u(\varphi(p))=v(p)\) for \(p\in W\).
(3) Points of \(\operatorname{cosupp}(\mathfrak a,\mu)\cap W\) are fixed. If \(p\in U\) and \(\varphi(p)\in C\), then \(\varphi(p)\in W\) is fixed:
\(\varphi(\varphi(p))=\varphi(p)\); applying \(\psi\) (\(\varphi(p),p\in W\)) gives \(\varphi(p)=\psi(\varphi(\varphi(p)))=\psi(\varphi(p))=p\), so \(p\in C\).
Conversely \(p\in C\cap U\) is fixed, so \(\varphi(p)=p\in C\). Finally \(\mathfrak a\subset\mathcal H(\mathfrak a)\) (Lemma~\ref{mc:homogenized-basic}), so an
admissible centre of \(\underline{\mathcal H}\) lies in \(\operatorname{cosupp}(\mathcal H(\mathfrak a),\mu)\subset\operatorname{cosupp}(\mathfrak a,\mu)\).
\end{proof}

\begin{definition}[Strict transforms along a tower]
\label{mc:strict-along}
\uses{hir:blowup-tower,hir:geo-transform,hir:tower-transport,bcalc:tower-restrict,hir:open-restriction-marked}
\lean{Hironaka.HomogenizedGlueingTower.strictAlong,Hironaka.HomogenizedGlueingTower.strictAlong_refl,Hironaka.HomogenizedGlueingTower.strictAlong_step,Hironaka.HomogenizedGlueingTower.N_geoTransformAlong,Hironaka.HomogenizedGlueingTower.strictAlong_restrict,Hironaka.HomogenizedGlueingTower.strictAlong_transport}
\leanok
For a tower \(T\) over \(X\) and \(A\subset X\), the \emph{strict transform} \(A_T\) is \(A\) for the empty tower and
\(\overline{\sigma^{-1}(A_{T'}\setminus C)}\) for a blowup \(\sigma\) along \(C\) on top of \(T'\). It is the \(N\) of the transform along \(T\) of any marked
ideal with \(N=A\); for \(U\subset X\) open, \((A\cap U)_{T|_U}=A_T\cap T^{-1}(U)\); and for a biholomorphism \(\Phi\colon X'\to X\),
\((\Phi^{-1}A)_{\Phi^*T}=\beta^{-1}(A_T)\), where \(\beta\) is the last-stage biholomorphism of the transported tower (Definition~\ref{hir:tower-transport}).
\end{definition}
\begin{proof}
\leanok
Induction on the tower. The \(N\) of a transform is the strict transform of \(N\) (Definition~\ref{hir:geo-transform}). Restriction: closures commute with
preimages under an open inclusion (Lemma~\ref{hir:open-restriction-marked}). Transport: at a step the new blowup is \(\beta^{-1}\circ\sigma\) along
\(\beta^{-1}C\), and \((\beta^{-1}\circ\sigma)^{-1}(\beta^{-1}A'\setminus\beta^{-1}C)=\sigma^{-1}(A'\setminus C)\).
\end{proof}

\begin{theorem}[Glueing Lemma along multiple test blow-ups, pullback form]
\label{mc:glueing-tower-pullback}
\uses{mc:glueing-biholomorphism,mc:strict-along,hir:tower-transport-open,hir:open-restriction-tower,hir:tower-transport-resolution,hir:open-restriction-marked}
\source{wlodarczyk05-hironaka,wlodarczyk09-analytic-resolution:page-0021}
\lean{Hironaka.HomogenizedGlueingTower.isClosedSubmanifold_univ,Hironaka.HomogenizedGlueingTower.isCoherent_restrict,Hironaka.HomogenizedGlueingTower.GlueingData.pullbackTower,Hironaka.HomogenizedGlueingTower.GlueingData.isGeoTestTower_restrict_dom,Hironaka.HomogenizedGlueingTower.GlueingData.isGeoTestTower_pullbackTower,Hironaka.HomogenizedGlueingTower.GlueingData.geoTransformAlong_pullbackTower_step,Hironaka.HomogenizedGlueingTower.GlueingData.cosupp_pullbackTower,Hironaka.HomogenizedGlueingTower.GlueingData.isGeoResolution_pullbackTower,Hironaka.HomogenizedGlueingTower.GlueingData.strictAlong_pullbackTower_zeroSet,Hironaka.HomogenizedGlueingTower.glueing_lemma_tower}
\leanok
With the notation of Lemma~\ref{mc:glueing-biholomorphism}, let \(T\) be a multiple test blow-up of \(\underline{\mathcal H}\), and let
\(\varphi^*T=\Phi^*(T|_{U'})\) be its \(\varphi\)-pullback, a tower over \(U\) (Lemma~\ref{hir:tower-transport-open}). Then:
\begin{enumerate}
\item \(\varphi^*T\) and \(T|_U\) are multiple test blow-ups of \(\underline{\mathcal H}|_U\), and \(\varphi^*T\) is a resolution if \(T\) is one;
\item from the first blowup on, the transforms of \(\underline{\mathcal H}|_U\) along \(\varphi^*T\) are the restrictions of the transforms of
\(\underline{\mathcal H}\) along \(T\) to the preimage of \(U'\), and at every stage the cosupport is the preimage, under the last-stage biholomorphism, of the
trace of the cosupport along \(T\);
\item the strict transform of \(V(v)\cap U\) along \(\varphi^*T\) is the preimage, under the last-stage biholomorphism, of
\(V(u)_T\cap T^{-1}(U')\).
\end{enumerate}
\end{theorem}
\begin{proof}
\leanok
\(\underline{\mathcal H}\) has \(N=M\), a closed submanifold, and coherent ideal. (1) and (2): Lemma~\ref{hir:open-restriction-tower} for \(T|_U\);
Lemma~\ref{hir:tower-transport-open} for \(\varphi^*T\), a multiple test blow-up (a resolution) of \(\Phi^*(\underline{\mathcal H}|_{U'})\), which is
\(\underline{\mathcal H}|_U\) by Lemma~\ref{mc:glueing-biholomorphism}(1); the cosupports by Lemma~\ref{hir:tower-transport-resolution} and
Lemma~\ref{hir:open-restriction-marked}. (3) By Lemma~\ref{mc:glueing-biholomorphism}(2), \(V(v)\cap U=\Phi^{-1}(V(u)\cap U')\); apply
Definition~\ref{mc:strict-along} (transport, then restriction to \(U'\)).
\end{proof}
This is the part of Włodarczyk's \emph{le:homo} (4) that follows from the functoriality of multiple test blow-ups under étale maps (his
Proposition \emph{pr:funct}): with \(\phi_v\) the inclusion of \(U\) and \(\phi_u=\varphi\), the towers \(\phi_v^*(X_i)=T|_U\) and
\(\phi_u^*(X_i)=\varphi^*T\) are multiple test blow-ups of the same marked ideal \(\underline{\mathcal H}|_U\), with the same first centre
(Lemma~\ref{mc:glueing-biholomorphism}(3)). The source's (4)(a)--(c) further identify the two towers stage by stage: \(\varphi\) lifts to
biholomorphisms \(\varphi_i\) between the open preimages of \(U\) and \(U'\) in \(X_i\), commuting with the blowups and fixing
\(\operatorname{cosupp}(\underline{\mathcal H}_i)\) pointwise; this is not part of the theorem.

\subsection{Hypersurfaces of maximal contact along multiple test blow-ups}

\begin{lemma}[Strict transforms are local over the base]
\label{mc:strict-transform-local}
\uses{bcalc:strict-submanifold}
\lean{Hironaka.MaximalContactTower.geoStrictTransform_inter_preimage}
\leanok
Let \(\pi\colon Y\to X\) be continuous, \(O\subset X\) open, and \(A,B,C\subset X\) with \(A\cap O=B\cap O\). Then
\(\overline{\pi^{-1}(A\setminus C)}\cap\pi^{-1}(O)=\overline{\pi^{-1}(B\setminus C)}\cap\pi^{-1}(O)\).
\end{lemma}
\begin{proof}
\leanok
By symmetry it suffices to show \(\subset\). As \(V=\pi^{-1}(O)\) is open, \(V\cap\overline S\subset\overline{V\cap S}\) for every \(S\), and
\(V\cap\pi^{-1}(A\setminus C)=\pi^{-1}((A\cap O)\setminus C)=\pi^{-1}((B\cap O)\setminus C)\subset\pi^{-1}(B\setminus C)\).
\end{proof}

\begin{theorem}[Giraud: hypersurfaces of maximal contact along multiple test blow-ups]
\label{mc:giraud-tower}
\uses{mc:giraud-transform,bcalc:giraud-local,bcalc:function-transform,mc:tangent-direction,mc:tangent-direction-support,mc:strict-along,mc:strict-transform-local,mc:inclusions-tower,hir:geo-transform-ambient,hir:bridge-consumers}
\source{wlodarczyk05-hironaka}
\lean{Hironaka.MaximalContactTower.exists_tangentDirection_along,Hironaka.MaximalContactTower.cosupp_subset_strictAlong,Hironaka.MaximalContactTower.centre_subset_strictAlong}
\leanok
Let \(\underline{\mathcal I}=(X,X,E,\mathfrak a,\mu)\) with \(\mathfrak a\) coherent and \(\mu\ge1\), and let \(u\) be a tangent direction of
\((\mathfrak a,\mu)\) on an open set \(U\). Let \(T\) be a multiple test blow-up of \(\underline{\mathcal I}\), with transform
\(\underline{\mathcal I}_T=(X_T,X_T,E_T,\mathfrak a_T,\mu)\), and let \(V(u)_T\) be the strict transform of \(V(u)\cap U\) along \(T\)
(Definition~\ref{mc:strict-along}). Then:
\begin{enumerate}
\item every point \(z\) of \(X_T\) over \(U\) has an open neighbourhood \(O\) over \(U\) with a tangent direction \(w\) of \((\mathfrak a_T,\mu)\) on
\(O\) such that \(V(w)\cap O=V(u)_T\cap O\);
\item \(\operatorname{cosupp}(\underline{\mathcal I}_T)\) lies in \(V(u)_T\) over \(U\); in particular so does every admissible centre of
\(\underline{\mathcal I}_T\).
\end{enumerate}
\end{theorem}
\begin{proof}
\leanok
(1) Induction on \(T\). For the empty tower take \(O=U\), \(w=u\). Let \(T'\) be \(T\) followed by the blowup \(\pi\) along an admissible centre
\(C\) of \(\underline{\mathcal I}_T\), and \(z'\) a point over \(U\). By induction \(z=\pi(z')\) has a neighbourhood \(O\) with a tangent direction
\(w\) of \((\mathfrak a_T,\mu)\) and \(V(w)\cap O=V(u)_T\cap O\). Along \(T\) the ideals stay coherent and \(N\) stays the whole space
(Lemma~\ref{mc:inclusions-tower}), so \(\mathfrak a_{T'}=\sigma^c(\mathfrak a_T,\mu)\) (Lemma~\ref{hir:geo-transform-ambient}), and
\(\mathfrak a_T\subset\mathcal I_C^{\mu}\) as \(C\subset\operatorname{cosupp}(\mathfrak a_T,\mu)\) (Lemma~\ref{hir:bridge-consumers}). Since \(w\) is a germ of
\(T(\mathfrak a_T)\), it vanishes on \(C\cap O\) (Lemma~\ref{mc:tangent-direction-support}). Let \(h\) be a local equation of the exceptional
divisor near \(z'\) (Lemma~\ref{bcalc:function-transform}); by Lemma~\ref{mc:giraud-transform}, \(w\circ\pi=h\,w'\) near \(z'\) with \(w'\)
holomorphic. Let \(O'\) be an open neighbourhood of \(z'\) inside \(\pi^{-1}(O)\) and the domain of \(h\), at every point of which \(w'\) is
holomorphic and \(w\circ\pi=h\,w'\) holds nearby. At every \(y\in O'\): \(w'\) is a germ of \(T(\sigma^c(\mathfrak a_T,\mu))=T(\mathfrak a_{T'})\)
(Lemma~\ref{mc:giraud-transform} at \(y\)); \(w'(y)=0\) if and only if \(y\) lies on the strict transform of \(V(w)\cap O\), and then \(w'\) has
multiplicity one at \(y\) (Lemma~\ref{bcalc:giraud-local}). Finally the strict transforms of \(V(w)\cap O\) and of \(V(u)_T\) agree over \(O\)
(Lemma~\ref{mc:strict-transform-local}), and the latter is \(V(u)_{T'}\). So \(O'\) and \(w'\) satisfy (1) at \(z'\).
(2) A point \(z\) of the cosupport over \(U\) lies in the \(O\) of (1); the germ of \(w\) at \(z\) lies in \(T(\mathfrak a_T)\), so \(w(z)=0\)
(Lemma~\ref{mc:tangent-direction-support}) and \(z\in V(u)_T\). Admissible centres lie in the cosupport.
\end{proof}
This is Włodarczyk's Lemma ``le:Giraud'', with the multiple test blow-ups of \((\mathfrak a,\mu)\) on the whole space read over \(U\) (for those of
\((\mathfrak a|_U,\mu)\), apply the theorem to \(X=U\)). Maximal order is not used; in the source it only provides tangent directions
(Lemma~\ref{mc:tangent-direction-exists}).

\subsection{Lifting biholomorphisms along a blowup}

In this subsection \(\pi\colon Y'\to Y\) is the blowup of a complex manifold along \(C\).

\begin{lemma}[Transforms on the blown-up side]
\label{mc:transform-source-naturality}
\uses{hir:geo-transform,ana:ideal-sheaves,hir:open-restriction-ideal}
\lean{Hironaka.BlowupBiholomorphismLift.holomorphicAt_comp_diffeo,Hironaka.BlowupBiholomorphismLift.holomorphicAt_comp_diffeo_symm,Hironaka.BlowupBiholomorphismLift.mem_pullback_diffeo_iff,Hironaka.BlowupBiholomorphismLift.comp_mem_pullback_diffeo,Hironaka.BlowupBiholomorphismLift.mem_restrictOn_pullback_diffeo_iff,Hironaka.BlowupBiholomorphismLift.colonAlong_pullback_diffeo,Hironaka.BlowupBiholomorphismLift.geoStrictTransform_comp_diffeo,Hironaka.BlowupBiholomorphismLift.geoTransform_comp_diffeo}
\leanok
Let \(\chi\colon Y_1\to Y_2\) be a biholomorphism and \(\pi\colon Y_2\to X\) holomorphic. For a coherent ideal \(\mathfrak c\) on \(Y_2\),
\(f\in(\chi^*\mathfrak c)_p\) if and only if \(f\circ\chi^{-1}\in\mathfrak c_{\chi(p)}\), and the quotient along a set commutes with \(\chi^*\). Consequently,
for a marked ideal \(\underline{\mathcal I}\) on \(X\) with coherent ideal, a set \(C\) with coherent ideal \(\mathcal I_C\) and coherent transform along
\(\pi\), the transform of \(\underline{\mathcal I}\) along \(\pi\circ\chi\) is \(\chi^*\) of its transform along \(\pi\).
\end{lemma}
\begin{proof}
\leanok
If \(f\) is a holomorphic combination of germs \(b\circ\chi\), \(b\in\mathfrak c_{\chi(p)}\), then \(f\circ\chi^{-1}\) is a combination of the \(b\), so lies in
\(\mathfrak c\) (coherent); conversely \(f=(f\circ\chi^{-1})\circ\chi\). Restricted combinations along \(\chi^{-1}W\) and along \(W\) correspond in the same way,
and the colon conditions correspond through \(g\mapsto g\circ\chi\). For the transform: strict transforms satisfy
\(\overline{(\pi\chi)^{-1}(A\setminus C)}=\chi^{-1}\overline{\pi^{-1}(A\setminus C)}\) since \(\chi\) is a homeomorphism; pullbacks compose, and
\(\chi^*\) commutes with powers and with the quotient along the strict transform of \(N\).
\end{proof}

\begin{lemma}[Lifting a biholomorphism of open sets]
\label{mc:lift-biholomorphism}
\uses{ana:blowups,bcalc:restrict-open,hir:biholomorphism-blowup}
\source{wlodarczyk05-hironaka,wlodarczyk09-analytic-resolution:page-0022}
\lean{Hironaka.BlowupBiholomorphismLift.preOpen,Hironaka.BlowupBiholomorphismLift.restrictMap,Hironaka.BlowupBiholomorphismLift.isBlowupAlong_restrictMap,Hironaka.BlowupBiholomorphismLift.exists_lift}
\leanok
Let \(A,B\subset Y\) be open and \(\Phi\colon A\to B\) a biholomorphism with \(\Phi^{-1}(C\cap B)=C\cap A\). There is a biholomorphism
\(\Phi_1\colon\pi^{-1}(A)\to\pi^{-1}(B)\) with \(\pi\circ\Phi_1=\Phi\circ\pi\).
\end{lemma}
\begin{proof}
\leanok
The restrictions \(\pi_A\colon\pi^{-1}(A)\to A\) and \(\pi_B\colon\pi^{-1}(B)\to B\) are blowups along \(C\cap A\) and \(C\cap B\)
(Lemma~\ref{bcalc:restrict-open}), and \(\Phi^{-1}\circ\pi_B\) is a blowup of \(A\) along \(\Phi^{-1}(C\cap B)=C\cap A\)
(Lemma~\ref{hir:biholomorphism-blowup}). Two blowups of \(A\) along the same centre factor through each other by the universal property
(Definition~\ref{ana:blowups}), and the factorizations are mutually inverse biholomorphisms over \(A\) (uniqueness in the universal property).
Let \(\Phi_1\) be the one with \(\Phi^{-1}\circ\pi_B\circ\Phi_1=\pi_A\).
\end{proof}
This is the step ``both morphisms lift by the universal property of a blow-up'' in Włodarczyk's proof of ``le:homo'' (4).

\subsection{Maps acting trivially modulo an ideal}

For an ideal \(\mathfrak J\) on a complex manifold \(Y\), \(V(\mathfrak J)\) is the set of points at which all germs of \(\mathfrak J\) vanish. A map
\(f\colon Y\to Y\) \emph{acts trivially modulo \(\mathfrak J\) at \(y\)} if \(f(y)=y\), \(f\) is holomorphic near \(y\), and \(g\circ f-g\in\mathfrak J_y\)
for every germ \(g\) holomorphic at \(y\).

\begin{lemma}[Trivial action: basic facts]
\label{mc:trivial-action}
\uses{mc:coherent-germs,bcalc:controlled-transform,ana:ideal-sheaves}
\lean{Hironaka.GlueingTrivialAction.zeroSet,Hironaka.GlueingTrivialAction.TrivialAt,Hironaka.GlueingTrivialAction.le_vanishingIdeal_of_subset_zeroSet,Hironaka.GlueingTrivialAction.mul_mem,Hironaka.GlueingTrivialAction.pullback_le_controlledTransform,Hironaka.GlueingTrivialAction.comp_mem_pullback,Hironaka.GlueingTrivialAction.controlledTransform_mem_congr,Hironaka.GlueingTrivialAction.mem_zeroSet_of_mem_zeroSet_controlledTransform,Hironaka.GlueingTrivialAction.lowProj,Hironaka.GlueingTrivialAction.lowProj_apply,Hironaka.GlueingTrivialAction.mixCoord,Hironaka.GlueingTrivialAction.mixCoord_apply,Hironaka.GlueingTrivialAction.mixCoord_zero,Hironaka.GlueingTrivialAction.mixCoord_n,Hironaka.GlueingTrivialAction.mixCoord_succ,Hironaka.GlueingTrivialAction.mdifferentiableAt_mixCoord,Hironaka.GlueingTrivialAction.mixCoord_self,Hironaka.GlueingTrivialAction.sub_mem_of_coord_sub_mem,Hironaka.GlueingTrivialAction.trivialAt_of_coords,Hironaka.GlueingTrivialAction.eqOn_of_eqOn_diff,Hironaka.GlueingTrivialAction.mdifferentiableAt_of_coords}
\leanok
Let \(\mathfrak J\) be coherent.
\begin{enumerate}
\item If \(C\subset V(\mathfrak J)\), then \(\mathfrak J\subset\mathcal I_C\). For a blowup \(\pi\) along \(C\), \(\pi^*\mathfrak J\subset\sigma^c(\mathfrak J,\mu)\), so
\(\pi\) maps \(V(\sigma^c(\mathfrak J,\mu))\) into \(V(\mathfrak J)\).
\item (Hadamard) If \(a,b\colon Y\to\mathbf C^n\) are holomorphic near \(z\), \(a(z)=b(z)\), and every \(a_i-b_i\) is a germ of \(\mathfrak J\) at
\(z\), then so is \(G\circ a-G\circ b\) for every \(G\) holomorphic near \(b(z)\).
\item If \(W\) is holomorphic near \(z\) with a holomorphic \(\ell\) near \(W(z)\) such that \(\ell\circ W=\mathrm{id}\) near \(z\), and \(f\) fixes
\(z\), is holomorphic near \(z\), and every coordinate of \(W\circ f-W\) is a germ of \(\mathfrak J\) at \(z\), then \(f\) acts trivially modulo
\(\mathfrak J\) at \(z\).
\item Continuous maps into a Hausdorff space that agree on \(V\setminus S\), \(V\) open and \(S\) with empty interior, agree on \(V\).
\end{enumerate}
\end{lemma}
\begin{proof}
\leanok
(1) A germ of \(\mathfrak J\) at \(y\) is a germ of \(\mathfrak J\) at nearby points (Lemma~\ref{mc:coherent-germs}), hence vanishes on \(C\) near \(y\).
\(\pi^*\mathfrak J\) is coherent, so \(g\,f\in\pi^*\mathfrak J\) for \(f\in\pi^*\mathfrak J\) and holomorphic \(g\); and \(g\in\mathfrak J_{\pi(z)}\) gives
\(g\circ\pi\in\sigma^c(\mathfrak J,\mu)_z\), so \(g(\pi(z))=0\) if \(z\in V(\sigma^c(\mathfrak J,\mu))\).
(2) Let \(c_k\) agree with \(a\) in the coordinates \(<k\) and with \(b\) in the others, so \(c_0=b\), \(c_n=a\) and
\(c_{k+1}=c_k+(a_k-b_k)e_k\). By Taylor's formula with holomorphic remainder along \(e_k\), \(G(w+te_k)=G(w)+t\,T_k(w,t)\) near \((b(z),0)\);
hence \(G(c_{k+1})-G(c_k)=(a_k-b_k)\,T_k(c_k,a_k-b_k)\in\mathfrak J_z\), and \(G\circ a-G\circ b\) is the sum of these.
(3) For \(g\) holomorphic at \(z\), \(G=g\circ\ell\) is holomorphic near \(W(z)\) and \(g=G\circ W\), \(g\circ f=G\circ W\circ f\) near \(z\) (as \(f(z)=z\));
apply (2) to \(a=W\circ f\), \(b=W\).
(4) Every point of \(V\) is in the closure of \(V\setminus S\), since an open set meeting \(V\) is not contained in \(S\).
\end{proof}

\begin{lemma}[A local equation in a monomial chart]
\label{mc:monomial-local-equation}
\uses{bcalc:function-transform,mc:monomial-derivation}
\lean{Hironaka.GlueingTrivialAction.isLocalEquation_monomial,Hironaka.GlueingTrivialAction.extendLifted,Hironaka.GlueingTrivialAction.extendLifted_of_mem,Hironaka.GlueingTrivialAction.mdifferentiableAt_extendLifted,Hironaka.GlueingTrivialAction.eq_of_monomial_eq}
\leanok
Let \(C\) be flattened near \(y_0\) by coordinates \(F\) onto \(\mathbf C^k\times0\), and let \(W\) be a local inverse of the \(j\)-th monomial chart
\(\sigma_j\) of the blowup (\(j\ge k\)) on an open \(O\), so that \(\sigma_j\circ W=F\circ\pi\) on \(O\). Then \(W_j\) is a local equation of the
exceptional divisor on \(O\). The monomial lift extends to a holomorphic map on the model space near its domain, and \(\sigma_j\) is injective
on points with nonzero \(j\)-th coordinate.
\end{lemma}
\begin{proof}
\leanok
As in the proof of Lemma~\ref{mc:monomial-derivation}: by Hadamard's lemma \(\mathcal I_C\) is generated near \(y_0\) by the \(F_i\), \(i\ge k\), and
\(F_i\circ\pi=W_iW_j\) (\(i\neq j\)) or \(W_j\) (\(i=j\)); so \(\pi^*\mathcal I_C=(W_j)\) on \(O\), and \(W_j\) vanishes identically on no open subset because
\(\pi^{-1}(C)\) has empty interior. Injectivity: the \(j\)-th coordinate of \(\sigma_j(w)\) is \(w_j\), and \(w_iw_j\) determines \(w_i\) when
\(w_j\neq0\).
\end{proof}

\begin{theorem}[Trivial action lifts along a blowup]
\label{mc:trivial-action-lift}
\uses{mc:trivial-action,mc:monomial-local-equation,bcalc:function-transform,bcalc:controlled-transform}
\source{wlodarczyk05-hironaka,wlodarczyk09-analytic-resolution:page-0022}
\lean{Hironaka.GlueingTrivialAction.trivialAt_lift_of_notMem,Hironaka.GlueingTrivialAction.trivialAt_lift_of_mem,Hironaka.GlueingTrivialAction.trivialAt_lift}
\leanok
Let \(\pi\colon Y\to X\) be the blowup along \(C\), \(\mathfrak J\) coherent with \(C\subset V(\mathfrak J)\), and \(\mathfrak J_1=\sigma^c(\mathfrak J,1)\). Let
\(f\colon X\to X\) act trivially modulo \(\mathfrak J\) at \(\pi(z)\) and map points off \(C\) near \(\pi(z)\) off \(C\), and let \(f_1\) be holomorphic near
\(z\) with \(\pi\circ f_1=f\circ\pi\) near \(z\). If \(z\in V(\mathfrak J_1)\), then \(f_1\) acts trivially modulo \(\mathfrak J_1\) at \(z\); in particular
\(f_1(z)=z\).
\end{theorem}
\begin{proof}
\leanok
\(\pi(z)\in V(\mathfrak J)\) by Lemma~\ref{mc:trivial-action}(1).
\emph{Off the centre}, \(\pi(z)\notin C\): \(\pi\) has a holomorphic section \(s\) over \(X\setminus C\) with \(s\circ\pi=\mathrm{id}\) there. As
\(\pi(f_1(z))=f(\pi(z))=\pi(z)\notin C\), \(f_1(z)=s(\pi(z))=z\). For \(g\) holomorphic at \(z\), \(G=g\circ s\) is holomorphic at \(\pi(z)\),
\(g=G\circ\pi\) and \(g\circ f_1=G\circ f\circ\pi\) near \(z\), so \(g\circ f_1-g=(G\circ f-G)\circ\pi\in\pi^*\mathfrak J\subset\mathfrak J_1\).
\emph{Over the centre}, \(\pi(z)\in C\): take flattening coordinates \(F\) at \(\pi(z)\), a monomial chart \(j\) with local inverse \(W\) at \(z\), and the local
equation \(W_j\) (Lemma~\ref{mc:monomial-local-equation}). Each \(F_i\circ f-F_i\) lies in \(\mathfrak J\subset\mathcal I_C\), so
\((F_i\circ f-F_i)\circ\pi=W_j\,\delta_i\) with \(\delta_i\in(\mathfrak J_1)_z\) (Lemma~\ref{bcalc:function-transform}); \(\delta_i(z)=0\) as \(z\in V(\mathfrak J_1)\).
Put \(G_i=(W_i+\delta_i)/(1+\delta_j)\) for \(i\ge k\), \(i\neq j\), and \(G_i=W_i+W_j\delta_i\) otherwise. Then \(G(z)=W(z)\) and
\(\sigma_j\circ G=F\circ f\circ\pi\) near \(z\), so \(\ell_j\circ G\) (\(\ell_j\) the monomial lift) is a holomorphic map with
\(\pi\circ\ell_j\circ G=f\circ\pi\). Off \(\pi^{-1}(C)\), \(\pi\) is injective and both \(f_1\) and \(\ell_j\circ G\) take values off \(\pi^{-1}(C)\) with the same
image under \(\pi\), so they agree; \(\pi^{-1}(C)\) has empty interior, hence \(f_1=\ell_j\circ G\) near \(z\) (Lemma~\ref{mc:trivial-action}(4)) and
\(f_1(z)=\ell_j(W(z))=z\). Likewise \(\sigma_j(W\circ f_1)=F\circ f\circ\pi=\sigma_j(G)\), and off \(\pi^{-1}(C)\) the \(j\)-th coordinate
\(G_j=W_j(1+\delta_j)\) is nonzero, so \(W\circ f_1=G\) off \(\pi^{-1}(C)\), hence near \(z\). The coordinates of \(W\circ f_1-W\) are
\(W_j\delta_i\) and \((\delta_i-W_i\delta_j)/(1+\delta_j)\), germs of \(\mathfrak J_1\); Lemma~\ref{mc:trivial-action}(3) concludes.
\end{proof}
This is the inductive step in Włodarczyk's proof of ``le:homo'' (4), in the analytic setting: there \(\phi_u^*(f)-\phi_v^*(f)\in(h)\subset T(\mathcal I)\),
and after a blowup the same holds for the lifts modulo the transformed \(T(\mathcal I)_1\), so the lifts agree on \(\operatorname{supp}(T(\mathcal I)_1,1)\).

\subsection{The glueing lemma along multiple test blow-ups: the lifts}

In this subsection \((\mathfrak a,\mu)\) is a coherent marked ideal on \(M\), \(\mu\ge1\), \(\underline{\mathcal H}=(M,M,\emptyset,\mathcal H(\mathfrak a),\mu)\),
and \(\mathfrak J=\mathcal D^{\mu-1}(\mathcal H(\mathfrak a))\); \(T(\mathfrak a)=\mathcal D^{\mu-1}(\mathfrak a)\subset\mathfrak J\) since \(\mathfrak a\subset\mathcal H(\mathfrak a)\).

\begin{lemma}[Coordinate changes act trivially]
\label{mc:coordinate-change-trivial}
\uses{mc:coordinate-change,mc:trivial-action,mc:taylor-tail,mc:coherent-germs}
\source{wlodarczyk05-hironaka,wlodarczyk09-analytic-resolution:page-0021,wlodarczyk09-analytic-resolution:page-0022}
\lean{Hironaka.HomogenizedGlueingLift.eventually_mdifferentiableAt_W,Hironaka.HomogenizedGlueingLift.eventually_trivialAt_coordMap}
\leanok
Let \(\varphi\) be a coordinate change at \(x\) (Definition~\ref{mc:coordinate-change}) with function \(h\), and \(\mathfrak J'\) a coherent ideal with
\(h\in\mathfrak J'_x\). Then \(\varphi\) acts trivially modulo \(\mathfrak J'\) at every point of \(V(\mathfrak J')\) near \(x\).
\end{lemma}
\begin{proof}
\leanok
Near \(x\), \(h\in\mathfrak J'_y\) (Lemma~\ref{mc:coherent-germs}), so \(h(y)=0\) for \(y\in V(\mathfrak J')\), and \(\varphi\) fixes \(y\). For \(g\) holomorphic at
\(y\), put \(F=g\circ e^{-1}\circ\Psi\), holomorphic near \(W(y)\); then \(g=F\circ W\) and \(g\circ\varphi=F(W+h\,e_i)\) near \(y\). Taylor's formula
with holomorphic remainder along \(e_i\) (Lemma~\ref{mc:taylor-tail}) gives \(F(w+te_i)=F(w)+t\,R(w,t)\) near \((W(y),0)\), so
\(g\circ\varphi-g=h\cdot R(W,h)\in\mathfrak J'_y\).
\end{proof}

\begin{lemma}[Glueing data acting trivially]
\label{mc:glueing-trivial}
\uses{mc:glueing-biholomorphism,mc:coordinate-change-trivial,mc:coordinate-pair,mc:homogenized-basic,mc:trivial-action,mc:tangent-direction-support}
\source{wlodarczyk05-hironaka,wlodarczyk09-analytic-resolution:page-0022}
\lean{Hironaka.HomogenizedGlueingLift.extendMap,Hironaka.HomogenizedGlueingLift.extendMap_of_mem,Hironaka.HomogenizedGlueingLift.extendMap_val,Hironaka.HomogenizedGlueingLift.contMDiffAt_extendMap,Hironaka.HomogenizedGlueingLift.eventually_mdifferentiableAt_extendMap,Hironaka.HomogenizedGlueingLift.trivialAt_congr,Hironaka.HomogenizedGlueingLift.trivialAt_mono,Hironaka.HomogenizedGlueingLift.zeroSet_mono,Hironaka.HomogenizedGlueingLift.cosupport_one_subset_zeroSet,Hironaka.HomogenizedGlueingLift.glueIdeal,Hironaka.HomogenizedGlueingLift.isCoherent_glueIdeal,Hironaka.HomogenizedGlueingLift.tangentIdeal_le_glueIdeal,Hironaka.HomogenizedGlueingLift.exists_glueingData_trivial}
\leanok
In Lemma~\ref{mc:glueing-biholomorphism} the open set \(W\) can be chosen so that, moreover, \(\varphi\) and \(\psi\) act trivially modulo \(\mathfrak J\) at
every point of \(V(\mathfrak J)\cap W\). Trivial action depends only on the germ of the map and passes to larger ideals; points of
\(\operatorname{cosupp}(\mathfrak J',1)\) lie in \(V(\mathfrak J')\).
\end{lemma}
\begin{proof}
\leanok
\(\varphi\) and \(\psi\) are the coordinate changes of Lemma~\ref{mc:coordinate-pair}, with \(h=v-u\) and \(u-v\), germs of \(T(\mathfrak a)\subset\mathfrak J\)
at \(x\); apply Lemma~\ref{mc:coordinate-change-trivial} and shrink \(W\). The other statements are immediate (a germ of \(\mathfrak J'\) at a point of
\(\operatorname{cosupp}(\mathfrak J',1)\) has order \(\ge1\), Lemma~\ref{mc:tangent-direction-support}).
\end{proof}

\begin{lemma}[One step of the lift]
\label{mc:lift-step}
\uses{mc:lift-biholomorphism,mc:transform-source-naturality,hir:biholomorphism-transform,hir:open-restriction-marked,mc:strict-along}
\lean{Hironaka.HomogenizedGlueingLift.restrictMap_comp_lift,Hironaka.HomogenizedGlueingLift.preimage_lift_geoStrictTransform,Hironaka.HomogenizedGlueingLift.pullback_lift_geoTransform}
\leanok
Let \(\pi\colon Z\to Y\) be the blowup along a closed submanifold \(C\), \(\Phi\colon A\to B\) a biholomorphism of open sets with
\(\Phi^{-1}(C\cap B)=C\cap A\), and \(\Phi_1\) its lift (Lemma~\ref{mc:lift-biholomorphism}). If \(\Phi^{-1}(S_u\cap B)=S_v\cap A\), then \(\Phi_1\)
maps the strict transform of \(S_v\) over \(A\) onto that of \(S_u\) over \(B\). If \(\Phi^*(\underline{\mathcal I}|_B)=\underline{\mathcal I}|_A\) for a marked ideal
\(\underline{\mathcal I}\) with coherent ideal and coherent transform \(\underline{\mathcal I}'\) along \(\pi\), then
\(\Phi_1^*(\underline{\mathcal I}'|_{\pi^{-1}B})=\underline{\mathcal I}'|_{\pi^{-1}A}\).
\end{lemma}
\begin{proof}
\leanok
Restriction to the open sets commutes with strict transforms and transforms (Lemma~\ref{hir:open-restriction-marked}); along
\(\pi_B\circ\Phi_1=\Phi\circ\pi_A\), Lemma~\ref{mc:transform-source-naturality} moves \(\Phi_1^*\) inside the transform, and the naturality of the
transform under the biholomorphism \(\Phi\) of the base (Lemma~\ref{hir:biholomorphism-transform}) with \(\Phi^{-1}(C\cap B)=C\cap A\) and
\(\Phi^*(\underline{\mathcal I}|_B)=\underline{\mathcal I}|_A\) gives the transform of \(\underline{\mathcal I}|_A\) along \(\pi_A\). For strict transforms,
\((\Phi\circ\pi_A)^{-1}(S_u\setminus C)=\pi_A^{-1}(S_v\setminus C)\) over \(A\).
\end{proof}

\begin{theorem}[Glueing Lemma: the lifts along multiple test blow-ups]
\label{mc:glueing-lift}
\uses{mc:glueing-trivial,mc:trivial-action-lift,mc:lift-step,mc:lift-biholomorphism,mc:inclusions-tower,mc:glueing-biholomorphism,mc:strict-along,hir:geo-transform-ambient,mc:tangent-direction-support}
\source{wlodarczyk05-hironaka,wlodarczyk09-analytic-resolution:page-0021,wlodarczyk09-analytic-resolution:page-0022}
\lean{Hironaka.HomogenizedGlueingLift.tanMarked,Hironaka.HomogenizedGlueingLift.IsGlueingLift,Hironaka.HomogenizedGlueingLift.IsGlueingLift.preimage_inter_eq,Hironaka.HomogenizedGlueingLift.exists_isGlueingLift,Hironaka.HomogenizedGlueingLift.IsGlueingLift.apply_eq_of_mem_cosupp,Hironaka.HomogenizedGlueingLift.glueing_lemma_lift}
\leanok
Let \(u\), \(v\) be tangent directions of \((\mathfrak a,\mu)\) on open sets containing \(x\in\operatorname{cosupp}(\mathfrak a,\mu)\), and let
\(\Phi\colon U\to U'\) be the glueing biholomorphism of Lemma~\ref{mc:glueing-trivial}. For every multiple test blow-up \(T\) of \(\underline{\mathcal H}\)
with last stage \(X_T\) there is a biholomorphism \(\Phi_T\) from the preimage \(U_T\) of \(U\) onto the preimage \(U'_T\) of \(U'\) in \(X_T\) such that
\begin{enumerate}
\item \(\Phi_T\) lifts \(\varphi\): \(T\circ\Phi_T=\varphi\circ T\);
\item \(\Phi_T^*(\underline{\mathcal H}_T|_{U'_T})=\underline{\mathcal H}_T|_{U_T}\);
\item \(\Phi_T^{-1}(V(u)_T\cap U'_T)=V(v)_T\cap U_T\) for the strict transforms;
\item \(\Phi_T\) and \(\Phi_T^{-1}\) act trivially modulo the transformed ideal \(\mathfrak J_T\) at its zeros; in particular \(\Phi_T\) fixes
\(\operatorname{cosupp}(\underline{\mathcal H}_T)\cap U_T\) pointwise.
\end{enumerate}
\end{theorem}
\begin{proof}
\leanok
Induction on \(T\); for the empty tower take \(\Phi\) (Lemmas~\ref{mc:glueing-biholomorphism} and~\ref{mc:glueing-trivial}). Let \(T'\) be \(T\) followed by
the blowup \(\pi\) along an admissible centre \(C\) of \(\underline{\mathcal H}_T\). By Lemma~\ref{mc:inclusions-tower}, \(T'\) is a multiple test
blow-up of \((\mathfrak J,1)\), so \(C\subset\operatorname{cosupp}(\mathfrak J_T,1)\subset V(\mathfrak J_T)\), and \(\mathfrak J_{T'}=\sigma^c(\mathfrak J_T,1)\)
(Lemma~\ref{hir:geo-transform-ambient}). Since \(\Phi_T\) and \(\Phi_T^{-1}\) fix the points of \(V(\mathfrak J_T)\), \(\Phi_T^{-1}(C\cap U'_T)=C\cap U_T\);
let \(\Phi_{T'}\) be the lift (Lemma~\ref{mc:lift-biholomorphism}). (1) is clear, (2) and (3) are Lemma~\ref{mc:lift-step}, and (4) is
Theorem~\ref{mc:trivial-action-lift} applied to \(\Phi_T\) and to \(\Phi_T^{-1}\) (which map points off \(C\) off \(C\)). Finally
\(\operatorname{cosupp}(\underline{\mathcal H}_T)=\operatorname{cosupp}(\mathcal H_T,\mu)\subset\operatorname{cosupp}(\mathcal D^{\mu-1}\mathcal H_T,1)\subset V(\mathfrak J_T)\)
(Lemma~\ref{mc:tangent-direction-support} and \(\mathfrak J_T\subset\mathcal D^{\mu-1}\mathcal H_T\), Lemma~\ref{mc:inclusions-tower}).
\end{proof}
This is Włodarczyk's ``le:homo'' (4)(a)--(c) in the analytic setting: with \(\phi_v\) the inclusion of \(U\) and \(\phi_u=\varphi\), the lifts are
\(\phi_{vi}\) the inclusion of \(U_T\) and \(\phi_{ui}=\Phi_T\); (4)(a) is (1)--(2), (4)(b) is (3), (4)(c) is (4). Together with
Theorem~\ref{mc:glueing-tower-pullback}, the \(\varphi\)-pullback of \(T\) and its restriction to \(U\) are identified stage by stage.

\subsection{Sums of marked ideals and the coefficient ideal}

In this subsection all marked ideals are \((M,M,E,\mathfrak a,d)\) with \(N=M\); for ideals \(\mathfrak J_0,\dots,\mathfrak J_{m-1}\),
\(\sum_{i<m}\mathfrak J_i\) is the ideal generated stalkwise by their germs.

\begin{lemma}[Finite sums of ideals]
\label{mc:finite-sum}
\uses{mc:homogenized-coherent,mc:coherent-germs,hir:controlled-operations,hir:cosupport-basic,bcalc:controlled-transform-def}
\lean{Hironaka.CoefficientIdeal.finSumIdeal_one,Hironaka.CoefficientIdeal.finSumIdeal_congr,Hironaka.CoefficientIdeal.finSumIdeal_le,Hironaka.CoefficientIdeal.isCoherent_finSumIdeal_of_lt,Hironaka.CoefficientIdeal.mem_cosupport_finSumIdeal,Hironaka.CoefficientIdeal.controlledTransform_finSumIdeal}
\leanok
Let \(\mathfrak J_0,\dots,\mathfrak J_{m-1}\) be coherent. Then \(\sum_{i<m}\mathfrak J_i\) is coherent; it is contained in every coherent ideal containing all
\(\mathfrak J_i\); \(\operatorname{cosupp}(\sum_{i<m}\mathfrak J_i,\mu)=\bigcap_{i<m}\operatorname{cosupp}(\mathfrak J_i,\mu)\); and if \(m\ge1\) and
\(\mathfrak J_i\subset\mathcal I_C^\mu\) for all \(i\), then \(\sigma^{\mathrm c}(\sum_{i<m}\mathfrak J_i,\mu)=\sum_{i<m}\sigma^{\mathrm c}(\mathfrak J_i,\mu)\).
\end{lemma}
\begin{proof}
\leanok
Coherence: Lemma~\ref{mc:homogenized-coherent}. A germ of the sum is a holomorphic combination of germs of the \(\mathfrak J_i\); this gives the
inclusion and, since a combination of germs vanishing to order \(\mu\) vanishes to order \(\mu\), the cosupport (the reverse inclusion because each
\(\mathfrak J_i\) is contained in the sum). The controlled transform: induction on \(m\), starting from \(\sum_{i<1}\mathfrak J_i=\mathfrak J_0\), with
\(\sum_{i<m+1}\mathfrak J_i=\sum_{i<m}\mathfrak J_i+\mathfrak J_m\) and \(\sigma^{\mathrm c}(\mathfrak a+\mathfrak b,\mu)=\sigma^{\mathrm c}(\mathfrak a,\mu)+\sigma^{\mathrm c}(\mathfrak b,\mu)\)
for coherent \(\mathfrak a,\mathfrak b\subset\mathcal I_C^\mu\) (Lemma~\ref{hir:controlled-operations}); the partial sums are contained in the coherent ideal \(\mathcal I_C^\mu\).
\end{proof}

\begin{definition}[Sum of finitely many marked ideals]
\label{mc:marked-finite-sum}
\uses{hir:marked-sum-product,mc:finite-sum,hir:marked-ideal-operations,hir:geo-transform-ambient,hir:bridge-consumers}
\source{wlodarczyk05-hironaka}
\lean{Hironaka.CoefficientIdeal.sumExp,Hironaka.CoefficientIdeal.mul_sumExp,Hironaka.CoefficientIdeal.one_le_sumExp,Hironaka.CoefficientIdeal.msum,Hironaka.CoefficientIdeal.IsAmbientFamily,Hironaka.CoefficientIdeal.IsAmbientFamily.isCoherent_msum,Hironaka.CoefficientIdeal.cosupp_msum,Hironaka.CoefficientIdeal.isAdmissibleCentre_msum_iff,Hironaka.CoefficientIdeal.geoTransform_msum,Hironaka.CoefficientIdeal.IsAmbientFamily.geoTransform}
\leanok
For marked ideals \(\underline{\mathcal I}_i=(M,N_i,E_i,\mathfrak a_i,\mu_i)\), \(i<m\),
\[
 \underline{\mathcal I}_0+\dots+\underline{\mathcal I}_{m-1}=\Bigl(M,N_0,E_0,\sum_{i<m}\mathfrak a_i^{\prod_{j\ne i}\mu_j},\prod_{j<m}\mu_j\Bigr).
\]
If \(m\ge1\), \(N_i=M\), \(E_i=E_0\), \(\mathfrak a_i\) is coherent and \(\mu_i\ge1\) for all \(i\), then the sum has coherent ideal and cosupport
\(\bigcap_i\operatorname{cosupp}\underline{\mathcal I}_i\); a centre is admissible for it if and only if it is admissible for every \(\underline{\mathcal I}_i\), and then the
transform of the sum is the sum of the transforms (Definition~\ref{hir:geo-transform}), whose summands again satisfy these hypotheses.
\end{definition}
\begin{proof}
\leanok
Write \(\mu=\prod_j\mu_j\) and \(e_i=\prod_{j\ne i}\mu_j\ge1\), so \(\mu_ie_i=\mu\). By Lemma~\ref{mc:finite-sum} and
\(\operatorname{cosupp}(\mathfrak a^{e},\nu e)=\operatorname{cosupp}(\mathfrak a,\nu)\) for \(e\ge1\) (Lemma~\ref{hir:marked-ideal-operations}),
\(\operatorname{cosupp}(\sum_i\mathfrak a_i^{e_i},\mu)=\bigcap_i\operatorname{cosupp}(\mathfrak a_i^{e_i},\mu_ie_i)=\bigcap_i\operatorname{cosupp}(\mathfrak a_i,\mu_i)\). The
admissibility statement follows, as all \(E_i\) agree. For an admissible centre \(C\), \(\mathfrak a_i\subset\mathcal I_C^{\mu_i}\) (Lemma~\ref{hir:bridge-consumers}), so
\(\mathfrak a_i^{e_i}\subset\mathcal I_C^{\mu}\), and by Lemmas~\ref{hir:geo-transform-ambient}, \ref{mc:finite-sum} and~\ref{hir:controlled-operations}
\(\sigma^{\mathrm c}(\sum_i\mathfrak a_i^{e_i},\mu)=\sum_i\sigma^{\mathrm c}(\mathfrak a_i^{e_i},\mu_ie_i)=\sum_i\sigma^{\mathrm c}(\mathfrak a_i,\mu_i)^{e_i}\); the strict transform of \(M\)
is the whole blowup, the \(E\) are transformed alike, the marks do not change, and the controlled transforms are coherent.
\end{proof}

\begin{lemma}[Multiple test blow-ups of a sum]
\label{mc:marked-finite-sum-tower}
\uses{mc:marked-finite-sum,hir:geo-transform}
\source{wlodarczyk05-hironaka}
\lean{Hironaka.CoefficientIdeal.geoTransformAlong_msum}
\leanok
Under the hypotheses of Definition~\ref{mc:marked-finite-sum}, the multiple test blow-ups of \(\underline{\mathcal I}_0+\dots+\underline{\mathcal I}_{m-1}\) are exactly
the simultaneous multiple test blow-ups of the \(\underline{\mathcal I}_i\), and along them the transform of the sum is the sum of the transforms.
\end{lemma}
\begin{proof}
\leanok
Induction on the tower, using at each step that the transforms satisfy the hypotheses again and Definition~\ref{mc:marked-finite-sum}.
\end{proof}
This is part (1) of Włodarczyk's Lemma ``le:operations'' for \(m\) summands; for \(m=2\) it is Lemma~\ref{hir:marked-operations-transform}(1).

\begin{definition}[Coefficient ideal]
\label{mc:coefficient-ideal}
\uses{mc:marked-finite-sum,mc:derivative-marked}
\source{wlodarczyk05-hironaka,wlodarczyk09-analytic-resolution:page-0022}
\lean{Hironaka.CoefficientIdeal.derivMarked_zero,Hironaka.CoefficientIdeal.coeffMarked,Hironaka.CoefficientIdeal.isAmbientFamily_derivMarked}
\leanok
The \emph{coefficient ideal} of \(\underline{\mathcal I}=(M,N,E,\mathfrak a,\mu)\) is
\[
 \mathcal C(\underline{\mathcal I})=\underline{\mathcal I}+\mathcal D(\underline{\mathcal I})+\dots+\mathcal D^{\mu-1}(\underline{\mathcal I}),
\]
i.e.\ \((\mathfrak a,\mu)+(\mathcal D\mathfrak a,\mu-1)+\dots+(\mathcal D^{\mu-1}\mathfrak a,1)\) (Definitions~\ref{mc:derivative-marked} and~\ref{mc:marked-finite-sum}).
For \(N=M\) and \(\mathfrak a\) coherent, the summands satisfy the hypotheses of Definition~\ref{mc:marked-finite-sum}.
\end{definition}

\begin{lemma}[The coefficient ideal is equivalent to the marked ideal]
\label{mc:coefficient-equivalent}
\uses{mc:coefficient-ideal,mc:marked-finite-sum-tower,mc:inclusions-tower,mc:derivative-marked}
\source{wlodarczyk05-hironaka,wlodarczyk09-analytic-resolution:page-0022}
\lean{Hironaka.CoefficientIdeal.coeff_tower}
\leanok
Let \(\underline{\mathcal I}=(M,M,E,\mathfrak a,\mu)\) with \(\mathfrak a\) coherent and \(\mu\ge1\). A tower is a multiple test blow-up of
\(\mathcal C(\underline{\mathcal I})\) if and only if it is one of \(\underline{\mathcal I}\), and along it
\(\operatorname{cosupp}\mathcal C(\underline{\mathcal I})_T=\operatorname{cosupp}\underline{\mathcal I}_T\).
\end{lemma}
\begin{proof}
\leanok
By Lemma~\ref{mc:marked-finite-sum-tower} the multiple test blow-ups of \(\mathcal C(\underline{\mathcal I})\) are the simultaneous ones of the
\(\mathcal D^i(\underline{\mathcal I})\), \(i<\mu\); the summand \(i=0\) is \(\underline{\mathcal I}\), and by Lemma~\ref{mc:inclusions-tower} every multiple test blow-up
of \(\underline{\mathcal I}\) is one of every \(\mathcal D^i(\underline{\mathcal I})\). Along such a tower,
\(\operatorname{cosupp}\mathcal C(\underline{\mathcal I})_T=\bigcap_{i<\mu}\operatorname{cosupp}[\mathcal D^i(\underline{\mathcal I})]_T\); the term \(i=0\) is
\(\operatorname{cosupp}\underline{\mathcal I}_T\), and for every \(i\), \(\operatorname{cosupp}\underline{\mathcal I}_T\subset\operatorname{cosupp}\mathcal D^i(\underline{\mathcal I}_T)\subset
\operatorname{cosupp}[\mathcal D^i(\underline{\mathcal I})]_T\) by Definition~\ref{mc:derivative-marked} and \([\mathcal D^i(\underline{\mathcal I})]_T\subset\mathcal D^i(\underline{\mathcal I}_T)\)
(Lemma~\ref{mc:inclusions-tower}).
\end{proof}
This is Włodarczyk's Lemma ``le:coeff'', \(\mathcal C(\underline{\mathcal I})\simeq\underline{\mathcal I}\); the source assumes maximal order, which the proof does not use.

\subsection{Restriction of the coefficient ideal to a submanifold}

In this subsection \(\underline{\mathcal B}=(M,M,E,\mathfrak b,d)\) has \(N=M\) and \(\mathfrak b\) coherent, and \(\underline{\mathcal I}=(M,M,E,\mathfrak a,\mu)\) has
\(N=M\), \(\mathfrak a\) coherent and \(\mu\ge1\). Sources: Włodarczyk's Lemmas ``le:cf'' and ``le:S''; in the analytic paper
(\texttt{wlodarczyk09-analytic-resolution}) Lemma~5.6.2, Lemma~5.6.3 and Proposition~5.6.4.

\begin{definition}[Restriction of a marked ideal to a set]
\label{mc:restrict-marked}
\uses{hir:marked-ideal,hir:order,mc:strict-along,mc:coefficient-ideal}
\source{wlodarczyk05-hironaka}
\lean{Hironaka.CoefficientIdealRestriction.withN,Hironaka.CoefficientIdealRestriction.withN_N,Hironaka.CoefficientIdealRestriction.withN_E,Hironaka.CoefficientIdealRestriction.withN_ideal,Hironaka.CoefficientIdealRestriction.withN_d,Hironaka.CoefficientIdealRestriction.cosupp_withN,Hironaka.CoefficientIdealRestriction.CentresIn,Hironaka.CoefficientIdealRestriction.centresIn_refl,Hironaka.CoefficientIdealRestriction.centresIn_step_iff,Hironaka.CoefficientIdealRestriction.N_geoTransformAlong_withN,Hironaka.CoefficientIdealRestriction.N_coeffMarked,Hironaka.CoefficientIdealRestriction.isCoherent_coeffMarked}
\leanok
For a set \(S\subset M\), the \emph{restriction} \(\underline{\mathcal B}|_S=(M,S,E,\mathfrak b,d)\) has the same \(E\), ideal and mark and \(N=S\); its cosupport is
\(\{x\in S\mid \operatorname{ord}_x(\mathfrak b|_S)\ge d\}\) (Definition~\ref{hir:order}). A tower \(T\) has its \emph{centres in the strict transforms of \(S\)} if
each centre \(C_i\) lies in the strict transform \(S_i\) of \(S\) (Definition~\ref{mc:strict-along}). The coefficient ideal \(\mathcal C(\underline{\mathcal I})\)
has \(N=M\) and a coherent ideal.
\end{definition}

\begin{lemma}[The order does not drop under restriction]
\label{mc:restrict-order}
\uses{mc:restrict-marked,hir:order}
\source{wlodarczyk05-hironaka,wlodarczyk09-analytic-resolution:page-0022,wlodarczyk09-analytic-resolution:page-0023}
\lean{Hironaka.CoefficientIdealRestriction.cosupport_univ_inter_subset,Hironaka.CoefficientIdealRestriction.cosupp_inter_subset_cosupp_withN}
\leanok
\(\operatorname{cosupp}\underline{\mathcal B}\cap S\subset\operatorname{cosupp}(\underline{\mathcal B}|_S)\) for every set \(S\).
\end{lemma}
\begin{proof}
\leanok
A germ of \(\mathfrak m_x^d\) agrees on \(S\) with itself.
\end{proof}
This is the first statement of Włodarczyk's Lemma ``le:cf'' (Lemma~5.6.3 of the analytic paper).

\begin{lemma}[Restriction commutes with the transform]
\label{mc:restrict-step}
\uses{mc:restrict-marked,mc:restrict-order,hir:geo-transform,hir:restricted-equivalence,hir:restricted-equivalence-basic,hir:marked-operations-along-step,hir:geo-transform-wellformed,ana:restricted-calculus,hir:hadamard-bridge-along}
\source{wlodarczyk05-hironaka}
\lean{Hironaka.CoefficientIdealRestriction.geoTransform_withN_restrictEquiv}
\leanok
Let \(S\subset M\) be a closed submanifold and \(\sigma\colon M'\to M\) the blowup along a centre \(C\) admissible for \(\underline{\mathcal B}\) and for
\(\underline{\mathcal B}|_S\), with strict transform \(S'\) of \(S\). Then the transform of \(\underline{\mathcal B}|_S\) is restricted-equivalent
(Definition~\ref{hir:restricted-equivalence}) to the restriction to \(S'\) of the transform of \(\underline{\mathcal B}\):
\(\sigma^{\mathrm c}((\mathfrak b,d)|_S)=(\sigma^{\mathrm c}(\mathfrak b,d))|_{S'}\).
\end{lemma}
\begin{proof}
\leanok
Let \(\mathfrak l=(\sigma^*\mathcal I_C)^d\), \(\mathfrak s'\) and \(\mathfrak c'\) the two transformed ideals; both are coherent. Applied to
\(\underline{\mathcal B}|_S\), the defining equation of the transform along \(S'\) (Lemma~\ref{hir:marked-operations-along-step}) gives
\(\operatorname{Res}_{S'}\sigma^*\mathfrak b=\operatorname{Res}_{S'}(\mathfrak l\,\mathfrak s')\); applied to \(\underline{\mathcal B}\) (with \(N=M\), whose strict
transform is \(M'\)) and restricted to \(S'\) it gives \(\operatorname{Res}_{S'}\sigma^*\mathfrak b=\operatorname{Res}_{S'}(\mathfrak l\,\mathfrak c')\). So
\(\operatorname{Res}_{S'}(\mathfrak s'\mathfrak l)=\operatorname{Res}_{S'}(\mathfrak c'\mathfrak l)\), and \(\mathfrak l\) is invertible along \(S'\)
(Lemma~\ref{hir:geo-transform-wellformed}), so it cancels (Lemma~\ref{ana:restricted-calculus}).
\end{proof}
This is the second statement of Włodarczyk's Lemma ``le:cf''. The source computes it in coordinates \((x,y)\) adapted to \(S\supset C\); the proof here
needs no coordinates beyond those already used for the defining equation of the transform.

\begin{lemma}[Restriction along multiple test blow-ups]
\label{mc:restrict-tower}
\uses{mc:restrict-step,mc:restrict-order,mc:strict-along,hir:test-tower,hir:restricted-equivalence-basic,bcalc:strict-ideal}
\source{wlodarczyk05-hironaka}
\lean{Hironaka.CoefficientIdealRestriction.centresIn_of_isGeoTestTower_withN,Hironaka.CoefficientIdealRestriction.restrict_tower}
\leanok
Let \(S\subset M\) be a closed submanifold. A multiple test blow-up of \(\underline{\mathcal B}|_S\) has its centres in the strict transforms of \(S\).
Every multiple test blow-up \(T\) of \(\underline{\mathcal B}\) with centres in the strict transforms of \(S\) is a multiple test blow-up of
\(\underline{\mathcal B}|_S\), and along it \([\underline{\mathcal B}|_S]_T\) is restricted-equivalent to \((\underline{\mathcal B}_T)|_{S_T}\); moreover
\(\underline{\mathcal B}_T\) has \(N\) the whole space, both transformed ideals are coherent and \(S_T\) is a closed submanifold.
\end{lemma}
\begin{proof}
\leanok
An admissible centre lies in the cosupport, hence in \(N\), which is the strict transform of \(S\). For the second statement induct on \(T\). At a step
with centre \(C\subset S_T\) admissible for \(\underline{\mathcal B}_T\), \(C\) is admissible for \((\underline{\mathcal B}_T)|_{S_T}\) by
Lemma~\ref{mc:restrict-order}, hence for the restricted-equivalent \([\underline{\mathcal B}|_S]_T\); restricted-equivalent marked ideals have the same
transform (Lemma~\ref{hir:restricted-equivalence-basic}), and Lemma~\ref{mc:restrict-step} applies. The strict transform of a closed submanifold
containing the centre is a closed submanifold (Lemma~\ref{bcalc:strict-ideal}).
\end{proof}
This is the last statement of Włodarczyk's Lemma ``le:cf'': \([(\mathcal I,\mu)|_S]_i=(\mathcal I_i,\mu)|_{S_i}\).

\begin{lemma}[The coefficient ideal is equivalent to the marked ideal: the equivalence relation]
\label{mc:coefficient-wequiv}
\uses{mc:coefficient-equivalent,hir:marked-equivalence}
\source{wlodarczyk05-hironaka,wlodarczyk09-analytic-resolution:page-0022}
\lean{Hironaka.CoefficientIdealRestriction.coeff_wEquiv}
\leanok
\(\mathcal C(\underline{\mathcal I})\simeq\underline{\mathcal I}\) in the sense of Definition~\ref{hir:marked-equivalence}.
\end{lemma}
\begin{proof}
\leanok
Both have the same \(E\), and Lemma~\ref{mc:coefficient-equivalent} gives the rest.
\end{proof}
This is Lemma~5.6.2 of the analytic paper.

\begin{lemma}[The support of the restricted coefficient ideal]
\label{mc:coefficient-restriction-support}
\uses{mc:restrict-marked,mc:restrict-order,mc:coefficient-equivalent,mc:finite-sum,hir:order-power,hir:giraud-support}
\source{wlodarczyk05-hironaka,wlodarczyk09-analytic-resolution:page-0023}
\lean{Hironaka.CoefficientIdealRestriction.cosupport_coeff_subset,Hironaka.CoefficientIdealRestriction.cosupp_coeff_withN}
\leanok
For every set \(S\subset M\), \(\operatorname{cosupp}\underline{\mathcal I}\cap S=\operatorname{cosupp}(\mathcal C(\underline{\mathcal I})|_S)\).
\end{lemma}
\begin{proof}
\leanok
The inclusion \(\subset\) is Lemma~\ref{mc:coefficient-equivalent} (for the empty tower) followed by Lemma~\ref{mc:restrict-order}. Conversely, let
\(x\in S\) with \(\operatorname{ord}_x(\mathcal C(\mathfrak a)|_S)\ge\mu_{\mathcal C}\). The ideal of \(\mathcal C(\underline{\mathcal I})\) contains
\(g^e\) for every \(g\in\mathcal D^{\mu-1}\mathfrak a\), where \(e\) is the exponent of the last summand \((\mathcal D^{\mu-1}\mathfrak a,1)\), and
\(\mu_{\mathcal C}=1\cdot e\ge1\). So \(g^e\), hence \(g\), vanishes at \(x\): \(x\in V(\mathcal D^{\mu-1}\mathfrak a)=\operatorname{cosupp}(\mathfrak a,\mu)\)
(Lemma~\ref{hir:giraud-support}).
\end{proof}
This is the first statement of Włodarczyk's Lemma ``le:S'' (Proposition~5.6.4 of the analytic paper), with \(\mu_{\mathcal C}=\mu!\) in place of
\(\operatorname{lcm}(1,\dots,\mu)\), which changes nothing. The source proves \(\supset\) by expanding the germs of \(\mathfrak a\) in coordinates normal to \(S\);
the last summand of \(\mathcal C\) makes this unnecessary, and \(S\) need not be a submanifold.

\begin{lemma}[The restricted coefficient ideal along multiple test blow-ups]
\label{mc:coefficient-restriction-tower}
\uses{mc:restrict-tower,mc:coefficient-equivalent,mc:restrict-order,mc:restrict-marked}
\source{wlodarczyk05-hironaka,wlodarczyk09-analytic-resolution:page-0023}
\lean{Hironaka.CoefficientIdealRestriction.coeff_restrict_tower}
\leanok
Let \(S\subset M\) be a closed submanifold and \(T\) a multiple test blow-up of \(\underline{\mathcal I}\) with centres in the strict transforms of \(S\).
Then \(T\) is a multiple test blow-up of \(\mathcal C(\underline{\mathcal I})|_S\), along it \([\mathcal C(\underline{\mathcal I})|_S]_T\) is
restricted-equivalent to \(\mathcal C(\underline{\mathcal I})_T|_{S_T}\), and
\(\operatorname{cosupp}\underline{\mathcal I}_T\cap S_T\subset\operatorname{cosupp}[\mathcal C(\underline{\mathcal I})|_S]_T\).
\end{lemma}
\begin{proof}
\leanok
By Lemma~\ref{mc:coefficient-equivalent}, \(T\) is a multiple test blow-up of \(\mathcal C(\underline{\mathcal I})\) and
\(\operatorname{cosupp}\mathcal C(\underline{\mathcal I})_T=\operatorname{cosupp}\underline{\mathcal I}_T\). Lemma~\ref{mc:restrict-tower} for
\(\underline{\mathcal B}=\mathcal C(\underline{\mathcal I})\) gives the first two statements, and the inclusion follows from Lemma~\ref{mc:restrict-order}.
\end{proof}
This is part (1) and the inclusion \(\subset\) of part (2) of Włodarczyk's Lemma ``le:S'' (Proposition~5.6.4 of the analytic paper). The reverse
inclusion \(\operatorname{cosupp}[\mathcal C(\underline{\mathcal I})|_S]_T\subset\operatorname{cosupp}\underline{\mathcal I}_T\) for a nonempty tower, and
with it part (3), is Lemma~\ref{mc:coefficient-restriction-full} below (the source's tracking of the Taylor coefficients in coordinates normal to the
\(S_i\) through the monomial charts).

\subsection{The restricted coefficient ideal along multiple test blow-ups: the full statement}

This subsection proves the remaining inclusion of Włodarczyk's Lemma ``le:S'' (Proposition~5.6.4 of the analytic paper), following the source:
the Taylor coefficients of the transforms in the normal coordinates of the strict transforms \(S_i\) of \(S\) are controlled transforms of
derivatives. Normal derivatives play the role of the coefficients \(c_\alpha=(\alpha!)^{-1}\partial^\alpha_xf|_S\).

\begin{lemma}[Order from normal derivatives, in \(\mathbb C^n\)]
\label{mc:normal-taylor-model}
\uses{mc:taylor-tail,mc:normal-vanishing}
\source{wlodarczyk05-hironaka}
\lean{Hironaka.NormalTaylorOrder.bv,Hironaka.NormalTaylorOrder.NormalIdx,Hironaka.NormalTaylorOrder.NormalIdx.mono,Hironaka.NormalTaylorOrder.NormalIdx.append,Hironaka.NormalTaylorOrder.normalIdx_replicate,Hironaka.NormalTaylorOrder.InNormalPow.mul_of_one,Hironaka.NormalTaylorOrder.InNormalPow.comp_killCoord,Hironaka.NormalTaylorOrder.killCoord_bv,Hironaka.NormalTaylorOrder.inNormalPow_of_normalDeriv_aux,Hironaka.NormalTaylorOrder.inNormalPow_of_normalDeriv,Hironaka.NormalTaylorOrder.inNormalPow_one_of_apply_eq_zero}
\leanok
Let \(F\) be holomorphic near \(0\in\mathbb C^n\). If for every list \(L\) of normal directions \(e_l\), \(l\ge k\), with \(|L|<\mu\), the derivative \(\partial^LF\)
agrees on the plane \(\mathbb C^k\times0\) near \(0\) with a germ of \(\mathfrak m_0^{\mu-|L|}\), then \(F\in\mathfrak m_0^\mu\).
\end{lemma}
\begin{proof}
\leanok
Induction on the number of normal coordinates. Expand in the first one, \(u_k\), with the holomorphic Taylor tail (Lemma~\ref{mc:taylor-tail}):
\(F(u)=\sum_{j<\mu}(j!)^{-1}(\partial^j_{e_k}F)(\mathrm{pr}\,u)\,u_k^j+u_k^\mu G\), where \(\mathrm{pr}\) kills \(u_k\). The coefficient
\((\partial^j_{e_k}F)\circ\mathrm{pr}\) satisfies the hypothesis for the plane \(\mathbb C^{k+1}\times0\) with \(\mu-j\) (derivatives commute with \(\mathrm{pr}\),
which maps \(\mathbb C^{k+1}\times0\) into \(\mathbb C^k\times0\) and preserves \(\mathfrak m_0^m\)), so it lies in \(\mathfrak m_0^{\mu-j}\) by induction.
\end{proof}

\begin{lemma}[Order from normal derivatives along a submanifold]
\label{mc:normal-taylor-order}
\uses{mc:normal-taylor-model,hir:order,hir:normal-power-pullback}
\source{wlodarczyk05-hironaka,wlodarczyk09-analytic-resolution:page-0023}
\lean{Hironaka.NormalTaylorOrder.inNormalPow_comp_of_mem_idealPow,Hironaka.NormalTaylorOrder.mem_idealPow_of_normalDeriv}
\leanok
Let \(\psi\) be holomorphic coordinates centred at \(p\in S\) in which \(S\) is \(\mathbb C^k\times0\) near \(p\), and \(f\) holomorphic at \(p\). If each normal
derivative \(\partial^L(f\circ\psi^{-1})\circ\psi\) (\(|L|<\mu\)) vanishes to order \(\mu-|L|\) along \(S\) at \(p\), then \(f\in\mathfrak m_p^\mu\).
\end{lemma}
\begin{proof}
\leanok
A germ of \(\mathfrak m_p^m\) read in the coordinates lies in \(\mathfrak m_0^m\) (induction on \(m\)); apply Lemma~\ref{mc:normal-taylor-model} and pull back
(Lemma~\ref{hir:normal-power-pullback}).
\end{proof}

\begin{definition}[Flat charts and normal derivatives]
\label{mc:flat-chart}
\uses{mc:normal-taylor-order,fpgeom:complex-analytic-subvariety}
\source{wlodarczyk05-hironaka}
\lean{Hironaka.NormalDerivativeChart.FlatChart,Hironaka.NormalDerivativeChart.mem_congr_of_isCoherent,Hironaka.NormalDerivativeChart.FlatChart.ψ,Hironaka.NormalDerivativeChart.FlatChart.inv,Hironaka.NormalDerivativeChart.FlatChart.nd,Hironaka.NormalDerivativeChart.FlatChart.ofRegularAtDim,Hironaka.NormalDerivativeChart.FlatChart.mem_flatDomain,Hironaka.NormalDerivativeChart.FlatChart.ψ_self,Hironaka.NormalDerivativeChart.FlatChart.zero_mem_target,Hironaka.NormalDerivativeChart.FlatChart.inv_zero,Hironaka.NormalDerivativeChart.FlatChart.eventually_mem_flatDomain,Hironaka.NormalDerivativeChart.FlatChart.mdifferentiableAt_ψ,Hironaka.NormalDerivativeChart.FlatChart.eventually_mdifferentiableAt_ψ,Hironaka.NormalDerivativeChart.FlatChart.mdifferentiableAt_inv,Hironaka.NormalDerivativeChart.FlatChart.eventually_mdifferentiableAt_inv,Hironaka.NormalDerivativeChart.FlatChart.eventually_inv_ψ,Hironaka.NormalDerivativeChart.FlatChart.eventually_ψ_inv,Hironaka.NormalDerivativeChart.FlatChart.eventually_flat,Hironaka.NormalDerivativeChart.FlatChart.tendsto_inv,Hironaka.NormalDerivativeChart.FlatChart.tendsto_ψ,Hironaka.NormalDerivativeChart.FlatChart.eventually_differentiableAt_comp_inv,Hironaka.NormalDerivativeChart.FlatChart.nd_nil,Hironaka.NormalDerivativeChart.FlatChart.mem_idealPow_of_vanishesToOrderOn,Hironaka.NormalDerivativeChart.FlatChart.injective_fderiv_Φ}
\leanok
A \emph{flat chart} for \(S\) at \(p\) is the preferred chart followed by a local biholomorphism \(\Phi\) with \(\Phi(e\,p)=0\) in which \(S\) is the plane
\(\mathbb C^k\times0\) near \(p\) (the data of a regular point). The \emph{normal derivatives} of \(f\) are \(\partial^L(f\circ\psi^{-1})\circ\psi\),
\(\psi=\Phi\circ e\). Lemma~\ref{mc:normal-taylor-order} applies to flat charts.
\end{definition}

\begin{lemma}[Normal derivatives of an ideal]
\label{mc:normal-derivative-ideal}
\uses{mc:flat-chart,hir:derivative-ideal,mc:coordinate-derivatives}
\source{wlodarczyk05-hironaka}
\lean{Hironaka.NormalDerivativeChart.FlatChart.nd_mem_iterDerivIdeal}
\leanok
For a coherent \(\mathfrak b\) and \(g\in\mathfrak b_p\), every derivative of \(g\) in a flat chart along a list \(L\) lies in \(\mathcal D^{|L|}(\mathfrak b)_p\).
\end{lemma}
\begin{proof}
\leanok
A coordinate derivative is a holomorphic combination of chart derivatives (Lemma~\ref{mc:coordinate-derivatives}); induct on \(L\).
\end{proof}

\begin{lemma}[Normal derivatives through monomial charts and tangential changes]
\label{mc:normal-derivative-monomial}
\uses{mc:flat-chart}
\source{wlodarczyk05-hironaka}
\lean{Hironaka.NormalDerivativeMonomial.fderiv_apply_eq_of_line,Hironaka.NormalDerivativeMonomial.monomial_add_smul_single,Hironaka.NormalDerivativeMonomial.iterDeriv_comp_monomial,Hironaka.NormalDerivativeMonomial.iterDeriv_mul_of_const,Hironaka.NormalDerivativeMonomial.iterDeriv_pow_mul,Hironaka.NormalDerivativeMonomial.iterDeriv_congr_of_eventuallyEq,Hironaka.NormalDerivativeMonomial.iterDeriv_comp_translate,Hironaka.NormalDerivativeMonomial.iterDeriv_comp_add}
\leanok
In \(\mathbb C^n\), for the monomial map \(\sigma_j\) (\(\sigma_j(w)_i=w_iw_j\) for \(i\ge k\), \(i\ne j\)) and directions \(e_l\), \(l\ne j\), \(l\ge k\):
\(\sigma_j(w+te_l)=\sigma_j(w)+tw_je_l\), hence \(\partial^L(G\circ\sigma_j)=w_j^{|L|}(\partial^LG)\circ\sigma_j\); a factor that is constant along the directions of
\(L\) (such as \(w_j^m\)) commutes with \(\partial^L\); derivatives only depend on germs and commute with translations and with maps \(\Psi\) with
\(\Psi(v+te_l)=\Psi(v)+te_l\).
\end{lemma}
\begin{proof}
\leanok
Derivatives along lines: if \(f(w+tv)=g(a+tb)\) for all \(t\), then \(Df(w)v=Dg(a)b\). Induct on \(L\).
\end{proof}

\begin{lemma}[Nested flat charts in product form]
\label{mc:nested-product-chart}
\uses{mc:flat-chart,mc:normal-derivative-monomial,bcalc:nested-flattening}
\source{wlodarczyk05-hironaka}
\lean{Hironaka.NormalDerivativeMonomial.planeProj,Hironaka.NormalDerivativeMonomial.planeProj_apply,Hironaka.NormalDerivativeMonomial.planeProj_mem,Hironaka.NormalDerivativeMonomial.planeProj_of_mem,Hironaka.NormalDerivativeMonomial.planeProj_planeProj,Hironaka.NormalDerivativeMonomial.planeProj_bv,Hironaka.NormalDerivativeMonomial.mem_plane_iff_sub_planeProj,Hironaka.NormalDerivativeMonomial.prodTo,Hironaka.NormalDerivativeMonomial.prodInv,Hironaka.NormalDerivativeMonomial.planeProj_prodTo,Hironaka.NormalDerivativeMonomial.planeProj_prodInv,Hironaka.NormalDerivativeMonomial.prodTo_sub_planeProj,Hironaka.NormalDerivativeMonomial.prodInv_prodTo,Hironaka.NormalDerivativeMonomial.prodEquiv,Hironaka.NormalDerivativeMonomial.prodInv_add_smul_bv,Hironaka.NormalDerivativeMonomial.FlatChart.nd_eventuallyEq_of_prod,Hironaka.NormalDerivativeMonomial.FlatChart.exists_nested}
\leanok
Let \(c\) be a flat chart for \(S\) (dimension \(k_N\)) at \(p\) and \(C\subset S\) regular of dimension \(k_C\le k_N\) at \(p\). There is a flat chart \(c'\) for \(S\) at
\(p\) in which also \(C\) is \(\mathbb C^{k_C}\times0\), with the same normal derivatives as \(c\).
\end{lemma}
\begin{proof}
\leanok
Flatten \(C\) inside the plane by Lemma~\ref{bcalc:nested-flattening}, getting \(\Phi_m\) that preserves \(\mathbb C^{k_N}\times0\); replace it by the product-form
\(T(u)=\Phi_m(Pu)+(u-Pu)\), \(P\) the projection onto the plane, with explicit inverse \(T^{-1}(v)=\Phi_m^{-1}(Pv)+(v-Pv)\). Then \(T^{-1}\) translates along the
normal directions, so \(T\circ\Phi\) has the same normal derivatives (Lemma~\ref{mc:normal-derivative-monomial}). This is the source's ``taking if necessary
linear combinations of \(y\)'': only tangential coordinates change.
\end{proof}

\begin{lemma}[The coefficient invariant]
\label{mc:normal-coefficient-invariant}
\uses{mc:flat-chart,mc:normal-derivative-ideal,mc:normal-taylor-order,mc:restrict-marked,mc:derivative-marked,ana:local-holomorphic-inverse}
\source{wlodarczyk05-hironaka}
\lean{Hironaka.CoefficientRestrictionTower.vanishesToOrderOn_of_isInSpanOnAt,Hironaka.CoefficientRestrictionTower.GoodAt,Hironaka.CoefficientRestrictionTower.GoodAt.orderGE,Hironaka.CoefficientRestrictionTower.goodAt_refl,Hironaka.CoefficientRestrictionTower.isInSpanAt_comp_map,Hironaka.CoefficientRestrictionTower.FlatChart.exists_of_map,Hironaka.CoefficientRestrictionTower.FlatChart.eventually_map_inv,Hironaka.CoefficientRestrictionTower.FlatChart.holomorphicAt_nd}
\leanok
Let \(\underline{\mathcal I}=(X,X,E,\mathfrak a,\mu)\), \(\mathfrak a\) coherent, \(S\) a closed submanifold and \(T\) a tower. A point \(p\in S_T\) is \emph{good} if there are a flat chart
for \(S_T\) at \(p\) and generators of \((\mathfrak a_T)_p\) whose normal derivatives \(\partial^L\), \(|L|<\mu\), restrict on \(S_T\) near \(p\) into
\([\mathcal D^{|L|}\mathfrak a]_T\). At a good point where every \([\mathcal D^i\mathfrak a]_T\) has order \(\ge\mu-i\) along \(S_T\), every germ of \(\mathfrak a_T\) lies in
\(\mathfrak m_p^\mu\). Every point of \(S\) is good for the empty tower.
\end{lemma}
\begin{proof}
\leanok
The first statement is Lemma~\ref{mc:normal-taylor-order} for the generators. For the empty tower take all germs of \(\mathfrak a\) as generators and
Lemma~\ref{mc:normal-derivative-ideal}.
\end{proof}

\begin{lemma}[The coefficient invariant along multiple test blow-ups]
\label{mc:normal-coefficient-step}
\uses{mc:normal-coefficient-invariant,mc:nested-product-chart,mc:normal-derivative-monomial,mc:monomial-local-equation,mc:restrict-tower,mc:inclusions-tower,bcalc:strict-submanifold,bcalc:monomial-lift,bcalc:isomorphism-off-centre,bcalc:function-transform,hir:geo-transform-ambient,hir:bridge-consumers,hir:geo-transform-wellformed,ana:restricted-colon}
\source{wlodarczyk05-hironaka}
\lean{Hironaka.CoefficientRestrictionTower.goodAt_step_of_notMem,Hironaka.CoefficientRestrictionTower.nd_monomial_identity,Hironaka.CoefficientRestrictionTower.goodAt_step_of_mem,Hironaka.CoefficientRestrictionTower.goodAt_tower}
\leanok
Along a multiple test blow-up of \(\underline{\mathcal I}\) with centres in the strict transforms of \(S\), every point of \(S_T\) is good.
\end{lemma}
\begin{proof}
\leanok
Induct on the tower; let \(\pi\) blow up \(C\subset S_T\) and \(z\in S_{T'}\). If \(\pi z\notin C\), \(\pi\) is a local biholomorphism near \(z\): pull back the chart and the
generators. If \(\pi z\in C\), refine the chart at \(\pi z\) to a nested product-form chart (Lemma~\ref{mc:nested-product-chart}) and take the monomial chart \(W\) at
\(z\) (Lemma~\ref{bcalc:strict-submanifold}): \(\sigma_j\circ W\) are the old coordinates, \(S_{T'}=W^{-1}(\mathbb C^k\times0)\), and \(W_j\) is a local equation of the
exceptional divisor (Lemma~\ref{mc:monomial-local-equation}). The new chart is \(W-W(z)\); the new generators are \(h'\) with \(g\circ\pi=W_j^\mu h'\) (they generate the
controlled transform after cancelling \(W_j^\mu\)). Lemma~\ref{mc:normal-derivative-monomial} gives
\(W_j^\mu\,\partial^Lh'=W_j^{|L|}(\partial^Lg)\circ\pi\); on \(S_T\) the right side is \(W_j^{|L|}\sum a_i(s_i\circ\pi)\) with \(s_i\in[\mathcal D^{|L|}\mathfrak a]_T\subset\mathcal I_C^{\mu-|L|}\),
so \(s_i\circ\pi=W_j^{\mu-|L|}s_i'\) with \(s_i'\in[\mathcal D^{|L|}\mathfrak a]_{T'}\); cancel \(W_j^\mu\) along \(S_{T'}\), where it vanishes on no relatively open set
(\(S_{T'}\) is the closure of \(\pi^{-1}(S_T\setminus C)\)).
\end{proof}

\begin{lemma}[The restricted coefficient ideal along multiple test blow-ups: full statement]
\label{mc:coefficient-restriction-full}
\uses{mc:normal-coefficient-step,mc:coefficient-restriction-tower,mc:coefficient-restriction-support,mc:marked-finite-sum-tower,mc:finite-sum,hir:order-power-along,hir:restricted-equivalence-basic}
\source{wlodarczyk05-hironaka,wlodarczyk09-analytic-resolution:page-0023}
\lean{Hironaka.CoefficientRestrictionTower.cosupp_coeff_withN_tower_subset,Hironaka.CoefficientRestrictionTower.cosupp_coeff_withN_tower,Hironaka.CoefficientRestrictionTower.coeff_restrict_tower_iff}
\leanok
Let \(\underline{\mathcal I}=(X,X,E,\mathfrak a,\mu)\) with \(\mathfrak a\) coherent and \(\mu\ge1\), and \(S\) a closed submanifold. A tower is a multiple test blow-up of
\(\mathcal C(\underline{\mathcal I})|_S\) if and only if it is a multiple test blow-up of \(\underline{\mathcal I}\) with centres in the strict transforms of \(S\), and along it
\(\operatorname{cosupp}[\mathcal C(\underline{\mathcal I})|_S]_T=\operatorname{cosupp}\underline{\mathcal I}_T\cap S_T\).
\end{lemma}
\begin{proof}
\leanok
The inclusion \(\supset\) is Lemma~\ref{mc:coefficient-restriction-tower}. Conversely, at \(x\) in the left side, \([\mathcal C(\underline{\mathcal I})|_S]_T\) is restricted-equivalent
to \(\mathcal C(\underline{\mathcal I})_T|_{S_T}=\sum_i[\mathcal D^i\mathfrak a]_T^{e_i}\) (Lemma~\ref{mc:marked-finite-sum-tower}); along \(S_T\) the order of the sum is the minimum of the
scaled orders (Lemma~\ref{hir:order-power-along}), so \([\mathcal D^i\mathfrak a]_T\) has order \(\ge\mu-i\) along \(S_T\) at \(x\), and the good point \(x\)
(Lemma~\ref{mc:normal-coefficient-step}) lies in \(\operatorname{cosupp}\underline{\mathcal I}_T\). The test blow-ups correspond by induction on the tower: a centre inside \(S_T\) is admissible
for one side iff for the other, since the cosupports correspond and \(E\) is the same.
\end{proof}
This is Włodarczyk's Lemma ``le:S'' (2)--(3) for the restriction keeping \(E\) (no member of \(E\) is dropped).

\begin{lemma}[Reduction to a hypersurface of maximal contact]
\label{mc:contact-equivalence}
\uses{mc:coefficient-restriction-full,mc:giraud-tower,mc:tangent-direction,bcalc:hypersurface,hir:marked-equivalence}
\source{wlodarczyk05-hironaka}
\lean{Hironaka.ContactEquivalence.isClosedSubmanifold_zeroSet_of_isTangentDirection,Hironaka.ContactEquivalence.centresIn_of_isTangentDirection,Hironaka.ContactEquivalence.contact_wEquiv}
\leanok
Let \(\underline{\mathcal I}=(X,X,E,\mathfrak a,\mu)\) with \(\mathfrak a\) coherent and \(\mu\ge1\), and \(u\) a tangent direction of \((\mathfrak a,\mu)\) on all of \(X\). Then
\(V(u)\) is a closed submanifold, every multiple test blow-up of \(\underline{\mathcal I}\) has its centres in the strict transforms of \(V(u)\), and
\(\underline{\mathcal I}\simeq\mathcal C(\underline{\mathcal I})|_{V(u)}\) (Definition~\ref{hir:marked-equivalence}).
\end{lemma}
\begin{proof}
\leanok
\(V(u)\) is a hypersurface (Lemma~\ref{bcalc:hypersurface}). By Lemma~\ref{mc:giraud-tower} the admissible centres and the cosupports of the transforms lie in the
strict transforms of \(V(u)\); Lemma~\ref{mc:coefficient-restriction-full} identifies the multiple test blow-ups and gives
\(\operatorname{cosupp}[\mathcal C(\underline{\mathcal I})|_{V(u)}]_T=\operatorname{cosupp}\underline{\mathcal I}_T\cap V(u)_T=\operatorname{cosupp}\underline{\mathcal I}_T\).
\end{proof}
This is Włodarczyk's Lemma ``le:codim2'' with \(U=X\) and without its hypothesis on the codimension of the support, which the proof does not use. The version
for a tangent direction on an open \(U\) is Lemma~\ref{mc:contact-equivalence-open}.

\subsection{The homogenized ideal is equivalent to the marked ideal}

In this subsection \(\underline{\mathcal I}=(M,M,E,\mathfrak a,\mu)\) has \(N=M\), \(\mathfrak a\) coherent and \(\mu\ge1\), and
\(\mathcal H(\underline{\mathcal I})=(M,M,E,\mathcal H(\mathfrak a),\mu)\). The subsection proves the two equivalences of
Lemma~\ref{hir:homogenized-coefficient}; the commutation of \(\mathcal H\) and \(\mathcal C\) with restrictions to open sets and biholomorphisms is
Lemma~\ref{mc:homogenized-etale}.

\begin{definition}[Sums \(\sum_i A_iB^i\)]
\label{mc:hom-sum}
\uses{mc:homogenized-ideal,mc:finite-sum,mc:homogenized-coherent}
\source{wlodarczyk05-hironaka}
\lean{Hironaka.HomogenizedEquivalence.homSum,Hironaka.HomogenizedEquivalence.homogenizedIdeal_eq_homSum,Hironaka.HomogenizedEquivalence.homSum_congr,Hironaka.HomogenizedEquivalence.isCoherent_homSum,Hironaka.HomogenizedEquivalence.le_homSum}
\leanok
For ideals \(A_0,\dots,A_{\mu-1},B\) put \(\sum_{i<\mu}A_iB^i\). By definition \(\mathcal H(\mathfrak a)=\sum_{i<\mu}\mathcal D^i(\mathfrak a)\,T(\mathfrak a)^i\) with
\(T(\mathfrak a)=\mathcal D^{\mu-1}(\mathfrak a)\). The sum is coherent when the \(A_i\) and \(B\) are, and it contains \(A_0\) when \(\mu\ge1\).
\end{definition}

\begin{lemma}[Order of \(\sum_i A_iB^i\)]
\label{mc:hom-sum-order}
\uses{mc:hom-sum,hir:order,hir:marked-ideal-operations}
\source{wlodarczyk05-hironaka}
\lean{Hironaka.HomogenizedEquivalence.orderGE_mul,Hironaka.HomogenizedEquivalence.orderGE_idealPow,Hironaka.HomogenizedEquivalence.orderGE_homSum}
\leanok
Along any set \(N\): \(\operatorname{ord}_x(\mathfrak a\mathfrak b|_N)\ge p+q\) when \(\operatorname{ord}_x(\mathfrak a|_N)\ge p\) and
\(\operatorname{ord}_x(\mathfrak b|_N)\ge q\); \(\operatorname{ord}_x(\mathfrak b^i|_N)\ge i\) when \(\operatorname{ord}_x(\mathfrak b|_N)\ge1\); hence
\(\operatorname{ord}_x(\sum_iA_iB^i|_N)\ge\mu\) when \(\operatorname{ord}_x(A_i|_N)\ge\mu-i\) for \(i<\mu\) and \(\operatorname{ord}_x(B|_N)\ge1\).
\end{lemma}
\begin{proof}
\leanok
Orders add under products of germs, and a holomorphic combination of germs of order \(\ge q\) has order \(\ge q\).
\end{proof}

\begin{lemma}[Controlled transform of \(\sum_i A_iB^i\)]
\label{mc:hom-sum-transform}
\uses{mc:hom-sum,mc:finite-sum,hir:controlled-operations}
\source{wlodarczyk05-hironaka}
\lean{Hironaka.HomogenizedEquivalence.controlledTransform_homSum}
\leanok
Let \(\sigma\) be the blowup along \(C\), \(\mu\ge1\), and \(A_i\subset\mathcal I_C^{\mu-i}\) (\(i<\mu\)), \(B\subset\mathcal I_C\) coherent. Then
\(\sigma^{\mathrm c}(\sum_iA_iB^i,\mu)=\sum_i\sigma^{\mathrm c}(A_i,\mu-i)\,\sigma^{\mathrm c}(B,1)^i\).
\end{lemma}
\begin{proof}
\leanok
Each summand lies in \(\mathcal I_C^{\mu}\); the controlled transform with a common mark commutes with finite sums (Lemma~\ref{mc:finite-sum}), and
\(\sigma^{\mathrm c}(A_iB^i,(\mu-i)+i)=\sigma^{\mathrm c}(A_i,\mu-i)\,\sigma^{\mathrm c}(B^i,i)\), \(\sigma^{\mathrm c}(B^i,i)=\sigma^{\mathrm c}(B,1)^i\)
(Lemma~\ref{hir:controlled-operations}).
\end{proof}

\begin{lemma}[The homogenized ideal along multiple test blow-ups]
\label{mc:homogenized-tower}
\uses{mc:hom-sum,mc:hom-sum-order,mc:hom-sum-transform,mc:inclusions-tower,mc:derivative-marked,hir:geo-transform-ambient,hir:bridge-consumers,hir:marked-equivalence-basic}
\source{wlodarczyk05-hironaka,wlodarczyk09-analytic-resolution:page-0020}
\lean{Hironaka.HomogenizedEquivalence.homMarked,Hironaka.HomogenizedEquivalence.mem_cosupport_derivMarked_tower,Hironaka.HomogenizedEquivalence.cosupp_subset_of_eq_homSum,Hironaka.HomogenizedEquivalence.le_of_eq_homSum,Hironaka.HomogenizedEquivalence.hom_tower,Hironaka.HomogenizedEquivalence.isGeoTestTower_of_hom}
\leanok
A tower is a multiple test blow-up of \(\mathcal H(\underline{\mathcal I})\) if and only if it is one of \(\underline{\mathcal I}\). Along such a tower \(T\),
\(\mathcal H(\underline{\mathcal I})_T\) has \(N\) the whole space and
\[
 \mathcal H(\mathfrak a)_T=\sum_{i<\mu}[\mathcal D^i\mathfrak a]_T\,[T(\mathfrak a)]_T^{\,i},
\]
where \([\cdot]_T\) is the transform of \((\mathcal D^i\mathfrak a,\mu-i)\) resp.\ \((T(\mathfrak a),1)\); in particular
\(\mathfrak a_T\subset\mathcal H(\mathfrak a)_T\) and \(\operatorname{cosupp}\underline{\mathcal I}_T\subset\operatorname{cosupp}\mathcal H(\underline{\mathcal I})_T\).
\end{lemma}
\begin{proof}
\leanok
Induct on \(T\). By Lemma~\ref{mc:inclusions-tower} a multiple test blow-up of \(\underline{\mathcal I}\) is one of every \(\mathcal D^i(\underline{\mathcal I})\), and
\(\operatorname{cosupp}\underline{\mathcal I}_T\subset\operatorname{cosupp}[\mathcal D^i\mathfrak a]_T\) (mark \(\mu-i\)); so by the formula and Lemma~\ref{mc:hom-sum-order}
\(\operatorname{cosupp}\underline{\mathcal I}_T\subset\operatorname{cosupp}\mathcal H(\underline{\mathcal I})_T\), and an admissible centre for
\(\underline{\mathcal I}_T\) is admissible for \(\mathcal H(\underline{\mathcal I})_T\) (the two have the same \(E\)). It is admissible for every
\([\mathcal D^i\mathfrak a]_T\), so these lie in the powers of \(\mathcal I_C\) given by their marks (Lemma~\ref{hir:bridge-consumers}), the transforms are controlled
transforms (Lemma~\ref{hir:geo-transform-ambient}), and Lemma~\ref{mc:hom-sum-transform} gives the formula after the blowup. Conversely the term \(i=0\) is
\(\mathfrak a_T\), so \(\mathfrak a_T\subset\mathcal H(\mathfrak a)_T\), \(\operatorname{cosupp}\mathcal H(\underline{\mathcal I})_T\subset\operatorname{cosupp}\underline{\mathcal I}_T\),
and a multiple test blow-up of \(\mathcal H(\underline{\mathcal I})\) is one of \(\underline{\mathcal I}\) (induction on \(T\)).
\end{proof}

\begin{lemma}[\(\mathcal H(\underline{\mathcal I})\simeq\underline{\mathcal I}\) and \(\mathcal C(\mathcal H(\underline{\mathcal I}))\simeq\underline{\mathcal I}\)]
\label{mc:homogenized-equivalent}
\uses{mc:homogenized-tower,mc:coefficient-wequiv,hir:marked-equivalence,hir:marked-equivalence-basic,mc:homogenized-coherent}
\source{wlodarczyk05-hironaka,wlodarczyk09-analytic-resolution:page-0020,wlodarczyk09-analytic-resolution:page-0022}
\lean{Hironaka.HomogenizedEquivalence.hom_wEquiv,Hironaka.HomogenizedEquivalence.coeff_hom_wEquiv}
\leanok
\(\mathcal H(\underline{\mathcal I})\simeq\underline{\mathcal I}\) and \(\mathcal C(\mathcal H(\underline{\mathcal I}))\simeq\underline{\mathcal I}\) (Definition~\ref{hir:marked-equivalence}).
\end{lemma}
\begin{proof}
\leanok
The same \(E\), the same multiple test blow-ups, and the two inclusions of cosupports of Lemma~\ref{mc:homogenized-tower}. The second statement is
Lemma~\ref{mc:coefficient-wequiv} for \(\mathcal H(\underline{\mathcal I})\) (coherent by Lemma~\ref{mc:homogenized-coherent}) and transitivity.
\end{proof}
This is Włodarczyk's lemma after ``le:hom'' (\(\mathcal H(\mathcal I,\mu)\simeq(\mathcal I,\mu)\)).

\begin{lemma}[Maximal order of \(\mathcal H\) and \(\mathcal C\)]
\label{mc:homogenized-maximal-order}
\uses{mc:maximal-order,mc:maximal-order-derivative,mc:homogenized-basic,mc:coefficient-ideal,hir:companion-maximal-order,hir:order-power,hir:giraud-support}
\source{wlodarczyk05-hironaka}
\lean{Hironaka.HomogenizedEquivalence.isMaximalOrder_homogenizedIdeal,Hironaka.HomogenizedEquivalence.isMaximalOrder_coeffMarked,Hironaka.HomogenizedEquivalence.isMaximalOrder_coeff_hom}
\leanok
If \((\mathfrak a,\mu)\) is of maximal order then so are \((\mathcal H(\mathfrak a),\mu)\) and the coefficient ideal \(\mathcal C(\underline{\mathcal I})\) (with its mark
\(\mu_{\mathcal C}\)); hence so is \(\mathcal C(\mathcal H(\underline{\mathcal I}))\).
\end{lemma}
\begin{proof}
\leanok
\(\mathfrak a\subset\mathcal H(\mathfrak a)\) gives \(\mathcal D^\mu\mathcal H(\mathfrak a)\supset\mathcal D^\mu\mathfrak a=\mathcal O\). For \(\mathcal C\): \((\mathcal D^{\mu-1}\mathfrak a,1)\)
is of maximal order, so at each \(x\) some \(g\in\mathcal D^{\mu-1}\mathfrak a\) has \(g\notin\mathfrak m_x^2\) (Giraud--Villamayor), and \(g^e\notin\mathfrak m_x^{e+1}\)
lies in the ideal of \(\mathcal C\), whose mark is \(\mu_{\mathcal C}=e\); so \(\operatorname{ord}_x\le\mu_{\mathcal C}\) everywhere, which is maximal order
(Lemma~\ref{hir:companion-maximal-order}).
\end{proof}

\begin{lemma}[Pullbacks of finite sums and of \(\sum_iA_iB^i\)]
\label{mc:pullback-finite-sum}
\uses{mc:finite-sum,mc:hom-sum,hir:controlled-operations}
\source{wlodarczyk05-hironaka}
\lean{Hironaka.HomogenizedEtale.pullback_finSumIdeal,Hironaka.HomogenizedEtale.pullback_homSum,Hironaka.HomogenizedEtale.pullback_homogenizedIdeal_of_iterDeriv,Hironaka.HomogenizedEtale.pullback_coeffMarked_of_iterDeriv,Hironaka.HomogenizedEtale.pullback_homMarked_of_iterDeriv}
\leanok
For a holomorphic map \(f\) and coherent ideals, \(f^*\sum_{i<m}\mathfrak J_i=\sum_{i<m}f^*\mathfrak J_i\) (\(m\ge1\)) and \(f^*\sum_iA_iB^i=\sum_i(f^*A_i)(f^*B)^i\). If
\(f^*\mathcal D^j(\mathfrak a)=\mathcal D^j(f^*\mathfrak a)\) for all \(j\), then \(f^*\mathcal H(\mathfrak a)=\mathcal H(f^*\mathfrak a)\), \(f^*\mathcal H(\underline{\mathcal I})=\mathcal H(f^*\underline{\mathcal I})\) and
\(f^*\mathcal C(\underline{\mathcal I})=\mathcal C(f^*\underline{\mathcal I})\).
\end{lemma}
\begin{proof}
\leanok
Induction on \(m\) with \(f^*(\mathfrak a+\mathfrak b)=f^*\mathfrak a+f^*\mathfrak b\); pullbacks commute with products and powers (Lemma~\ref{hir:controlled-operations}). The rest is
the definition of \(\mathcal H\) and \(\mathcal C\) as such sums.
\end{proof}

\begin{lemma}[Derivative ideals commute with restriction to open sets]
\label{mc:derivative-open-restriction}
\uses{hir:derivative-ideal,hir:derivative-invariance,hir:open-restriction-ideal}
\source{wlodarczyk05-hironaka}
\lean{Hironaka.HomogenizedEtale.chartDeriv_comp_val,Hironaka.HomogenizedEtale.pullback_val_derivIdeal,Hironaka.HomogenizedEtale.pullback_val_iterDerivIdeal}
\leanok
For an open \(U\subset M\) with inclusion \(\iota\), \(\iota^*\mathcal D^j(\mathfrak a)=\mathcal D^j(\iota^*\mathfrak a)\).
\end{lemma}
\begin{proof}
\leanok
\(\supset\) holds for every holomorphic map (Lemma~\ref{hir:derivative-invariance}). Conversely a partial derivative in a chart \(e\) of \(M\) is, by the chain rule, a
holomorphic combination of partial derivatives in the preferred chart at \(\iota y\), whose restriction to \(U\) is the preferred chart of \(U\) at \(y\); so its
restriction is a combination of partial derivatives of \(\iota^*\)-germs. Induct on \(j\).
\end{proof}

\begin{lemma}[The homogenized and coefficient ideals are étale-local]
\label{mc:homogenized-etale}
\uses{mc:pullback-finite-sum,mc:derivative-open-restriction,hir:derivative-invariance,mc:homogenized-tower,mc:coefficient-ideal}
\source{wlodarczyk05-hironaka}
\lean{Hironaka.HomogenizedEtale.pullback_homogenizedIdeal_of_inverse,Hironaka.HomogenizedEtale.pullback_coeffMarked_of_inverse,Hironaka.HomogenizedEtale.pullback_val_homogenizedIdeal,Hironaka.HomogenizedEtale.pullback_val_coeffMarked,Hironaka.HomogenizedEtale.pullback_val_homMarked}
\leanok
\(\mathcal H\) and \(\mathcal C\) commute with restriction to open sets and with biholomorphisms: \(\iota^*\mathcal H(\mathfrak a)=\mathcal H(\iota^*\mathfrak a)\),
\(\iota^*\mathcal C(\underline{\mathcal I})=\mathcal C(\iota^*\underline{\mathcal I})\), and the same for a biholomorphism \(\varphi\) in place of \(\iota\).
\end{lemma}
\begin{proof}
\leanok
Lemma~\ref{mc:pullback-finite-sum} with Lemma~\ref{mc:derivative-open-restriction} for \(\iota\), resp.\ Lemma~\ref{hir:derivative-invariance} for \(\varphi\).
\end{proof}
These are Włodarczyk's Lemmas ``le:etale2'' and ``le:etale3'' for the étale maps used in the analytic setting (open inclusions and biholomorphisms).

\begin{lemma}[Tangent directions restrict to open sets]
\label{mc:tangent-direction-open}
\uses{mc:tangent-direction,mc:derivative-open-restriction,hir:derivative-invariance}
\source{wlodarczyk05-hironaka}
\lean{Hironaka.ContactEquivalenceOpen.exists_chartDeriv_ne_zero_complexChartAt,Hironaka.ContactEquivalenceOpen.mem_pullback_val_of_mem,Hironaka.ContactEquivalenceOpen.isTangentDirection_pullback_val}
\leanok
Let \(u\) be a tangent direction of \((\mathfrak a,\mu)\) on an open set \(U\), \(\mathfrak a\) consisting of holomorphic germs, and \(\iota\colon U\to M\) the
inclusion. Then \(\iota^*u=u|_U\) is a tangent direction of \((\iota^*\mathfrak a,\mu)\) on all of \(U\). If some partial derivative of a germ \(u\)
holomorphic at \(x\) does not vanish at \(x\) in some chart, some coordinate partial derivative in the preferred chart at \(x\) does not.
\end{lemma}
\begin{proof}
\leanok
\(\iota^*u\in\iota^*\mathcal D^{\mu-1}(\mathfrak a)=\mathcal D^{\mu-1}(\iota^*\mathfrak a)\) (Lemma~\ref{mc:derivative-open-restriction}). By the chain rule a
partial derivative in any chart is a holomorphic combination of the coordinate partial derivatives in the preferred chart at \(x\)
(Lemma~\ref{hir:derivative-invariance}), so one of these does not vanish at a zero \(x\in U\) of \(u\); and the preferred chart of \(U\) at \(x\) is the
restriction of the preferred chart of \(M\) at \(x\).
\end{proof}

\begin{lemma}[Reduction to a hypersurface of maximal contact on an open set]
\label{mc:contact-equivalence-open}
\uses{mc:contact-equivalence,mc:tangent-direction-open,mc:homogenized-etale,hir:open-restriction-marked,hir:marked-equivalence}
\source{wlodarczyk05-hironaka}
\lean{Hironaka.ContactEquivalenceOpen.contact_wEquiv_open}
\leanok
Let \(\underline{\mathcal I}=(X,X,E,\mathfrak a,\mu)\) with \(\mathfrak a\) coherent and \(\mu\ge1\), and \(u\) a tangent direction of \((\mathfrak a,\mu)\) on an
open set \(U\). Then \(\underline{\mathcal I}|_U\simeq(U,V(u)\cap U,E|_U,\mathcal C(\mathfrak a)|_U,\mu_{\mathcal C})\), i.e.\
\(\underline{\mathcal I}|_U\simeq\mathcal C(\underline{\mathcal I})|_U\) with the submanifold \(V(u)\cap U\) (Definition~\ref{hir:marked-equivalence} on
\(U\)).
\end{lemma}
\begin{proof}
\leanok
By Lemma~\ref{mc:tangent-direction-open}, \(u|_U\) is a tangent direction of \(\underline{\mathcal I}|_U=(U,U,E|_U,\iota^*\mathfrak a,\mu)\) on all of \(U\);
Lemma~\ref{mc:contact-equivalence} on \(U\) gives \(\underline{\mathcal I}|_U\simeq\mathcal C(\underline{\mathcal I}|_U)|_{V(u)\cap U}\), and
\(\mathcal C(\underline{\mathcal I}|_U)=\mathcal C(\underline{\mathcal I})|_U\) (Lemma~\ref{mc:homogenized-etale}).
\end{proof}
This is Włodarczyk's Lemma ``le:codim2'' (with any \(E\), and without its hypothesis on the codimension of the support).

\subsection{Derivative ideals along a submanifold}

\begin{definition}[Derivative ideal along a submanifold]
\label{mc:derivative-along}
\uses{hir:derivative-ideal,hir:marked-ideal}
\source{bierstone-milman08-functoriality,wlodarczyk05-hironaka}
\lean{Hironaka.DerivativeIdealAlong.fieldDeriv,Hironaka.DerivativeIdealAlong.IsTangentField,Hironaka.DerivativeIdealAlong.DerivGenAlong,Hironaka.DerivativeIdealAlong.derivIdealAlong,Hironaka.DerivativeIdealAlong.iterDerivIdealAlong,Hironaka.DerivativeIdealAlong.holomorphicAt_fieldDeriv}
\leanok
Let \(N\subset M\) and \(x\in M\). A holomorphic vector field \(\xi=\sum_m\xi_m\,\partial^e_{e_m}\) near \(x\) (in a chart \(e\) at \(x\)) is \emph{tangent to
\(N\) at \(x\)} if \(\xi(g)\in(\mathcal I_N)_x\) for every \(g\in(\mathcal I_N)_x\). The \emph{derivative ideal of \(\mathfrak a\) along \(N\)},
\(\mathcal D_N(\mathfrak a)\), is generated at \(x\) by the germs of \(\mathfrak a_x\), the germs of \((\mathcal I_N)_x\), and the derivatives \(\xi(f)\) of
\(f\in\mathfrak a_x\) along vector fields \(\xi\) tangent to \(N\) at \(x\); \(\mathcal D_N^0(\mathfrak a)=\mathfrak a\) and
\(\mathcal D_N^{j+1}(\mathfrak a)=\mathcal D_N(\mathcal D_N^j(\mathfrak a))\).
\end{definition}
For a closed submanifold \(N\) the derivations of \(\mathcal O_N\) are induced by the ambient derivations preserving \(\mathcal I_N\), so the restriction
of \(\mathcal D_N(\mathfrak a)\) to \(N\) is the derivative ideal of \(\mathfrak a|_N\) in the sense of Bierstone and Milman.

\begin{lemma}[Elementary properties of \(\mathcal D_N\)]
\label{mc:derivative-along-basic}
\uses{mc:derivative-along,hir:derivative-ideal}
\lean{Hironaka.DerivativeIdealAlong.le_derivIdealAlong,Hironaka.DerivativeIdealAlong.vanishingIdeal_le_derivIdealAlong,Hironaka.DerivativeIdealAlong.derivIdealAlong_mono,Hironaka.DerivativeIdealAlong.isTangentField_univ,Hironaka.DerivativeIdealAlong.derivIdealAlong_univ,Hironaka.DerivativeIdealAlong.iterDerivIdealAlong_univ,Hironaka.DerivativeIdealAlong.vanishingIdeal_mem_of_isInSpanAt,Hironaka.DerivativeIdealAlong.derivIdealAlong_vanishingIdeal}
\leanok
\(\mathfrak a+\mathcal I_N\subset\mathcal D_N(\mathfrak a)\), and \(\mathcal D_N\) is monotone. For \(N=M\), \(\mathcal D_M^j(\mathfrak a)=\mathcal D^j(\mathfrak a)\). Moreover
\(\mathcal D_N(\mathcal I_N)=\mathcal I_N\).
\end{lemma}
\begin{proof}
\leanok
The first claims hold by construction. For \(N=M\), \(\mathcal I_M\) consists of the germs vanishing near the point, every vector field is tangent to
\(M\), \(\xi(f)=\sum_m\xi_m\partial^e_{e_m}f\) is a combination of partial derivatives, and conversely \(\partial^e_vf=\sum_mv_m\partial^e_{e_m}f\) is the derivative
along the field with constant coefficients \(v_m\); induct on \(j\). Finally the generators of \(\mathcal D_N(\mathcal I_N)\) are germs of \(\mathcal I_N\) (by tangency for
the derivatives), and \(\mathcal I_N\) is closed under holomorphic combinations.
\end{proof}

\begin{lemma}[Product rule and chain rule along a vector field]
\label{mc:field-chain-rule}
\uses{mc:derivative-along,hir:derivative-ideal-generators,hir:order-along-flat}
\lean{Hironaka.DerivativeIdealAlongCoherent.fieldDeriv_eventuallyEq_sum_mul,Hironaka.DerivativeIdealAlongCoherent.flatPartial,Hironaka.DerivativeIdealAlongCoherent.holomorphicAt_flatPartial,Hironaka.DerivativeIdealAlongCoherent.flatPartial_eq,Hironaka.DerivativeIdealAlongCoherent.chartDeriv_eq_sum_flat,Hironaka.DerivativeIdealAlongCoherent.fieldDeriv_eventuallyEq_sum_flatPartial}
\leanok
Let \(\xi\) be a holomorphic vector field near \(y\), written in a chart \(e\) at \(y\).
\begin{enumerate}
\item If \(f=\sum_ia_is_i\) near \(y\) with \(a_i,s_i\) holomorphic near \(y\), then \(\xi(f)=\sum_i\big(\xi(a_i)\,s_i+a_i\,\xi(s_i)\big)\) near \(y\).
\item Let \(\psi=\Phi\circ e_x\) be flattening coordinates of a submanifold at a point \(x\) as in Lemma~\ref{hir:order-along-flat}, with
domain \(D_\Phi\ni y\), and write \(\partial/\partial\psi_j\) for the derivative along the coordinate field \(\xi_{e_j}\) of that lemma, so that
\((\partial h/\partial\psi_j)(p)=D(h\circ\psi^{-1})(\psi(p))\,e_j\). Then for \(h\) holomorphic near \(y\),
\[\xi(h)=\sum_{j=1}^n\xi(\psi_j)\,\frac{\partial h}{\partial\psi_j}\qquad\text{near }y.\]
\end{enumerate}
\end{lemma}
\begin{proof}
\leanok
(1) Apply the product rule for the partial derivatives \(\partial^e_{e_m}\) (Lemma~\ref{hir:derivative-ideal-generators}) to each term of
\(\xi=\sum_m\xi_m\partial^e_{e_m}\) and exchange the two finite sums.
(2) Let \(p\) be near \(y\), in the domain of \(e\) and in \(D_\Phi\). Near \(e(p)\), \(h\circ e^{-1}=(h\circ\psi^{-1})\circ G\) with
\(G=\psi\circ e^{-1}=\Phi\circ(e_x\circ e^{-1})\), which is holomorphic at \(e(p)\) (a transition map of the atlas followed by \(\Phi\)). By the chain rule
\(\partial^e_vh(p)=D(h\circ\psi^{-1})(\psi(p))\,DG(e(p))\,v\), and the \(j\)-th coordinate of \(DG(e(p))\,v\) is \(\partial^e_v\psi_j(p)\); expanding
\(DG(e(p))\,v=\sum_j(\partial^e_v\psi_j(p))\,e_j\) gives \(\partial^e_vh=\sum_j\partial^e_v\psi_j\,\partial h/\partial\psi_j\). Multiply by \(\xi_m\) with \(v=e_m\) and sum
over \(m\).
\end{proof}

\begin{proposition}[Generators of \(\mathcal D_N\) in one flat chart]
\label{mc:derivative-along-generators}
\uses{mc:field-chain-rule,mc:derivative-along,mc:derivative-along-basic,hir:order-along-flat,bex:hadamard}
\lean{Hironaka.DerivativeIdealAlongCoherent.derivFamilyAlong,Hironaka.DerivativeIdealAlongCoherent.isInSpanAt_of_mem_vanishingIdeal,Hironaka.DerivativeIdealAlongCoherent.single_mem_complexCoordinatePlane,Hironaka.DerivativeIdealAlongCoherent.derivIdealAlong_mem_iff_of_generators}
\leanok
Let \(\psi=\Phi\circ e_x\) be flattening coordinates of \(N\) at \(x\) as in Lemma~\ref{hir:order-along-flat}: on the domain \(D_\Phi\), \(p\in N\) if and only if
\(\psi(p)\in\mathbf C^k\times0\). Let \(s_1,\dots,s_r\), holomorphic on an open set \(U\subset D_\Phi\), generate \(\mathfrak a\) at every point of \(U\).
Then at every point \(y\in U\), on \(N\) or not, \(\mathcal D_N(\mathfrak a)_y\) is generated by
\[s_1,\dots,s_r,\qquad \psi_j\ (k\le j<n),\qquad \frac{\partial s_i}{\partial\psi_j}\ (1\le i\le r,\ 0\le j<k)\]
(coordinates indexed from \(0\)): the generators of \(\mathfrak a\), the normal coordinates, and the tangential partial derivatives of the generators.
\end{proposition}
\begin{proof}
\leanok
The listed functions lie in \(\mathcal D_N(\mathfrak a)_y\): the \(s_i\) lie in \(\mathfrak a_y\); \(\psi_j\) (\(j\ge k\)) vanishes on \(N\cap D_\Phi\), so lies in
\((\mathcal I_N)_y\); and for \(j<k\), \(e_j\in\mathbf C^k\times0\), so the coordinate field \(\xi_{e_j}\) is tangent to \(N\) at \(y\)
(Lemma~\ref{hir:order-along-flat}\,(2)) and \(\partial s_i/\partial\psi_j=\xi_{e_j}(s_i)\) is a generator of Definition~\ref{mc:derivative-along}.

Conversely we show that each generator of \(\mathcal D_N(\mathfrak a)_y\) lies in the ideal \(J_y\) spanned by the listed functions. Germs of \(\mathfrak a_y\) are
combinations of the \(s_i\). A germ of \((\mathcal I_N)_y\) is \(\sum_{j\ge k}b_j\psi_j\) with \(b_j\) holomorphic, by Hadamard's lemma in flattening coordinates
(Lemma~\ref{bex:hadamard}\,(2), valid at every point of \(D_\Phi\)). Let finally \(\xi\) be tangent to \(N\) at \(y\) and \(f\in\mathfrak a_y\), \(f=\sum_la_ls_{\sigma(l)}\)
near \(y\). By Lemma~\ref{mc:field-chain-rule}\,(1), \(\xi(f)=\sum_l\xi(a_l)s_{\sigma(l)}+a_l\xi(s_{\sigma(l)})\), so it suffices that \(\xi(s_i)\in J_y\). By
Lemma~\ref{mc:field-chain-rule}\,(2), \(\xi(s_i)=\sum_j\xi(\psi_j)\,\partial s_i/\partial\psi_j\). For \(j<k\) the factor \(\partial s_i/\partial\psi_j\) is listed. For
\(j\ge k\), \(\psi_j\in(\mathcal I_N)_y\), hence \(\xi(\psi_j)\in(\mathcal I_N)_y\) by tangency, and \(\xi(\psi_j)\) is a combination of the normal coordinates by
Hadamard's lemma again.
\end{proof}
At a point \(y\notin N\) of \(D_\Phi\) some normal coordinate \(\psi_j(y)\) is nonzero, so both sides are the unit ideal; the statement does not need to treat this case
separately.

\begin{corollary}[Coherence of \(\mathcal D_N\)]
\label{mc:derivative-along-coherent}
\uses{mc:derivative-along-generators,mc:derivative-along-basic,ana:blowups}
\lean{Hironaka.DerivativeIdealAlongCoherent.derivIdealAlong_mem_of_notMem,Hironaka.DerivativeIdealAlongCoherent.isCoherent_derivIdealAlong,Hironaka.DerivativeIdealAlongCoherent.isCoherent_iterDerivIdealAlong}
\leanok
Let \(\mathfrak a\) be coherent and \(N\subset M\) a closed submanifold. Then \(\mathcal D_N^j(\mathfrak a)\) is coherent for every \(j\ge0\). Off \(N\),
\(\mathcal D_N(\mathfrak a)\) is the unit ideal.
\end{corollary}
\begin{proof}
\leanok
At \(x\in N\), take flattening coordinates of \(N\) at \(x\) and generators of \(\mathfrak a\) on an open \(U\ni x\); on \(U\cap D_\Phi\) the finitely many functions of
Proposition~\ref{mc:derivative-along-generators} are holomorphic and generate \(\mathcal D_N(\mathfrak a)\). At \(x\notin N\), on the open set \(M\setminus N\) every
holomorphic germ lies in \(\mathcal I_N\subset\mathcal D_N(\mathfrak a)\), so the constant \(1\) generates. Induct on \(j\).
\end{proof}

\subsection{Tangent directions along a submanifold}

In this subsection \(N\subset M\) is a set, \(\mathfrak a\) an ideal of holomorphic germs on \(M\), \(\mu\ge1\), and \(\mathcal D_N\) the derivative ideal along \(N\)
(Definition~\ref{mc:derivative-along}). These are the single-stage statements of item (3) of Lemma~\ref{hir:along-calculus}.

\begin{definition}[Tangent direction along a submanifold]
\label{mc:tangent-direction-along}
\uses{mc:derivative-along,mc:tangent-direction}
\source{wlodarczyk05-hironaka}
\lean{Hironaka.TangentDirectionAlong.IsTangentDirectionAlong}
\leanok
A \emph{tangent direction of \((\mathfrak a|_N,\mu)\) along \(N\) on an open set \(U\)} is a function \(u\) whose germ at every point of \(U\) lies in
\(\mathcal D_N^{\mu-1}(\mathfrak a)\), such that at each zero \(y\in U\) of \(u\) some holomorphic vector field \(\xi\) tangent to \(N\) at \(y\) has
\(\xi(u)(y)\ne0\); that is, the restriction \(u|_N\) has a nonvanishing differential at its zeros.
\end{definition}

\begin{lemma}[Tangent directions along \(N\): basic facts]
\label{mc:tangent-direction-along-basic}
\uses{mc:tangent-direction-along,mc:derivative-along-basic,hir:giraud-along,hir:order-basic,mc:tangent-direction,mc:tangent-direction-support}
\source{wlodarczyk05-hironaka}
\lean{Hironaka.TangentDirectionAlong.cosupport_add_subset_cosupport_iterDerivIdealAlong,Hironaka.TangentDirectionAlong.apply_eq_zero_of_mem_cosupport,Hironaka.TangentDirectionAlong.IsTangentDirectionAlong.cosupport_inter_subset,Hironaka.TangentDirectionAlong.fieldDeriv_const,Hironaka.TangentDirectionAlong.isTangentDirectionAlong_univ_iff}
\leanok
\begin{enumerate}
\item \(\operatorname{cosupp}(\mathfrak a|_N,\mu+j)\subset\operatorname{cosupp}((\mathcal D_N^j\mathfrak a)|_N,\mu)\) for all \(j\); every germ of \(\mathcal D_N^{\mu-1}(\mathfrak a)\) at a point of
\(\operatorname{cosupp}(\mathfrak a|_N,\mu)\) vanishes there; hence \(\operatorname{cosupp}(\mathfrak a|_N,\mu)\cap U\subset V(u)\) for a tangent direction \(u\) along \(N\) on \(U\).
\item For \(N=M\), the tangent directions along \(N\) are the tangent directions of Definition~\ref{mc:tangent-direction}.
\end{enumerate}
\end{lemma}
\begin{proof}
\leanok
(1) Iterate the first inclusion of Giraud--Villamayor along \(N\) (Lemma~\ref{hir:giraud-along}); a germ of order \(\ge1\) along \(N\) at a point of \(N\) vanishes
there (Lemma~\ref{hir:order-basic}). (2) \(\mathcal D_M^{\mu-1}=\mathcal D^{\mu-1}\) (Lemma~\ref{mc:derivative-along-basic}); every vector field is tangent to \(M\), and
\(\xi(u)=\sum_m\xi_m\partial^e_{e_m}u\ne0\) forces some \(\partial^e_{e_m}u\ne0\), while \(\partial^e_vu\) is the derivative along the field with constant coefficients \(v_m\).
\end{proof}

\begin{proposition}[Existence of tangent directions along \(N\)]
\label{mc:tangent-direction-along-exists}
\uses{mc:tangent-direction-along,mc:tangent-direction-along-basic,mc:derivative-along-generators,mc:derivative-along-coherent,hir:giraud-along,hir:order-along-flat}
\source{wlodarczyk05-hironaka}
\lean{Hironaka.TangentDirectionAlong.exists_isTangentDirectionAlong}
\leanok
Let \(\mathfrak a\) be coherent, \(N\) a closed submanifold, \(\mu\ge1\), and \(x\in\operatorname{cosupp}(\mathfrak a|_N,\mu)\) with \(\operatorname{ord}_x(\mathfrak a|_N)\le\mu\). Then
there is a tangent direction of \((\mathfrak a|_N,\mu)\) along \(N\) on an open neighbourhood of \(x\).
\end{proposition}
\begin{proof}
\leanok
Since \(x\notin\operatorname{cosupp}(\mathfrak a|_N,\mu+1)=N\cap V(\mathcal D_N^\mu\mathfrak a)\) (Lemma~\ref{hir:giraud-along}), some germ \(g\) of
\(\mathcal D_N^\mu(\mathfrak a)=\mathcal D_N(\mathcal D_N^{\mu-1}\mathfrak a)\) has \(g(x)\ne0\). Take flattening coordinates \(\psi\) of \(N\) at \(x\) and generators \(s_i\) of the
coherent ideal \(\mathcal D_N^{\mu-1}(\mathfrak a)\) (Corollary~\ref{mc:derivative-along-coherent}) near \(x\). By Proposition~\ref{mc:derivative-along-generators}, \(g\) is a
combination of the \(s_i\), the normal coordinates \(\psi_j\) (\(j\ge k\)) and the tangential partials \(\partial s_i/\partial\psi_j\) (\(j<k\)), so one of these does not vanish
at \(x\). The \(s_i\) vanish at \(x\) (Lemma~\ref{mc:tangent-direction-along-basic}), and so do the \(\psi_j\), \(j\ge k\), as \(x\in N\). Hence
\(\partial s_i/\partial\psi_j(x)\ne0\) for some \(i\) and some \(j<k\); then \(u=s_i\) is a tangent direction on a neighbourhood \(U\) of \(x\) on which
\(\partial s_i/\partial\psi_j\ne0\): \(u\) is a germ of \(\mathcal D_N^{\mu-1}(\mathfrak a)\) at every point of \(U\), and the coordinate field \(\xi_{e_j}\) is tangent to \(N\)
(Lemma~\ref{hir:order-along-flat}) with \(\xi_{e_j}(u)=\partial u/\partial\psi_j\).
\end{proof}
Only the maximal order at the point \(x\) is used: \(\operatorname{ord}_x(\mathfrak a|_N)=\mu\).

\begin{lemma}[Hypersurfaces of maximal contact inside \(N\)]
\label{mc:contact-hypersurface-along}
\uses{mc:tangent-direction-along,mc:field-chain-rule,hir:order-along-flat,bex:hadamard,bcalc:snc-chart,bcalc:strict-ideal,bcalc:monomial-lift,bcalc:snc-simultaneous,ana:blowups,ana:local-holomorphic-inverse}
\source{wlodarczyk05-hironaka}
\lean{Hironaka.TangentDirectionAlong.exists_flatPartial_ne_zero,Hironaka.TangentDirectionAlong.contactCoords,Hironaka.TangentDirectionAlong.contactCoords_apply,Hironaka.TangentDirectionAlong.contactCoords_mem_plane_iff,Hironaka.TangentDirectionAlong.mdifferentiableAt_contactCoords,Hironaka.TangentDirectionAlong.injective_fderiv_contactCoords,Hironaka.TangentDirectionAlong.regularAtDim_inter_zeroSet,Hironaka.TangentDirectionAlong.exists_vanishingIdeal_inter_zeroSet_mem_iff,Hironaka.TangentDirectionAlong.IsTangentDirectionAlong.regularAtDim_contact,Hironaka.TangentDirectionAlong.isClosed_preimage_val_inter_zeroSet}
\leanok
Let \(\psi=\Phi\circ e_x\) be flattening coordinates of \(N\) at \(x\in N\), with \(N=\psi^{-1}(\mathbf C^k\times0)\) on their domain, let \(u\) be holomorphic on an open
\(U\ni x\) with \(u(x)=0\), and let \(\xi(u)(x)\ne0\) for a vector field \(\xi\) tangent to \(N\) at \(x\). Then
\begin{enumerate}
\item \(\partial u/\partial\psi_j(x)\ne0\) for some \(j<k\);
\item \(N\cap V(u)\cap U\) is regular of dimension \(k-1\) at \(x\), and near \(x\) its ideal is generated by \(u\) and the normal coordinates \(\psi_i\), \(i\ge k\):
\(\mathcal I_{N\cap V(u)}=\mathcal I_N+(u)\).
\end{enumerate}
Consequently, for a tangent direction \(u\) of \((\mathfrak a|_N,\mu)\) along a closed submanifold \(N\) on \(U\), \(N\cap V(u)\cap U\) is closed in \(U\) and regular, of
dimension \(\dim_yN-1\), at each of its points \(y\).
\end{lemma}
\begin{proof}
\leanok
(1) By Lemma~\ref{mc:field-chain-rule}, \(\xi(u)(x)=\sum_j\xi(\psi_j)(x)\,\partial u/\partial\psi_j(x)\); for \(j\ge k\), \(\psi_j\in\mathcal I_N\), so \(\xi(\psi_j)\in\mathcal I_N\) by
tangency and \(\xi(\psi_j)(x)=0\).
(2) Let \(l=k-1\) and \(\tau\) the transposition of \(j\) and \(l\). Put \(W_r=\psi_{\tau(r)}\) for \(r\ne l\) and \(W_l=u\). On the domain of \(\psi\), \(W\in\mathbf C^{k-1}\times0\)
if and only if \(u=0\) and \(\psi_r=0\) for \(r\ge k\) (\(\tau\) fixes these \(r\)), i.e.\ if and only if the point lies in \(N\cap V(u)\). In the coordinates \(\psi\), \(DW\) is
the selection matrix of \(\tau\) with the row \(l\) replaced by \(du\); a vector \(w\) in its kernel has \(w_m=0\) for \(m\ne j\), and then the row \(l\) gives
\(w_j\,\partial u/\partial\psi_j(x)=0\). So \(DW(x)\) is injective, and \(W-W(x)\), read in the chart at \(x\), is a local holomorphic coordinate system centred at \(x\)
(Lemma~\ref{bcalc:snc-chart}) in which \(N\cap V(u)\cap U\) is \(\mathbf C^{k-1}\times0\); since \(W(x)\in\mathbf C^{k-1}\times0\), \(W-W(x)\) lies in the plane exactly when \(W\) does.
Hadamard's lemma in these coordinates (Lemma~\ref{bcalc:strict-ideal}) generates the ideal by \(W_r\), \(r\ge k-1\), which are \(u\) and the \(\psi_r\), \(r\ge k\).
The consequence: at a point \(y\) of \(N\cap V(u)\cap U\) apply (2) with flattening coordinates of \(N\) at \(y\) and the field given by the tangent direction; closedness in
\(U\) holds because \(N\) is closed and \(u\) is continuous on \(U\).
\end{proof}
This prepares the restriction to a hypersurface of maximal contact inside \(N\): in the induction on dimension the marked ideal is restricted from \(N\) to
\(N\cap V(u)\), whose ideal is \(\mathcal I_N+(u)\). The behaviour along multiple test blow-ups (the cosupports stay in the strict transforms of \(N\cap V(u)\)) needs
the transforms along \(N\), item (2) of Lemma~\ref{hir:along-calculus}, and is not part of these statements.

\subsection{The derivative ideal along a submanifold is étale-local}

A holomorphic map \(\varphi\colon M'\to M\) between complex manifolds of the same dimension is a \emph{local biholomorphism} if near every \(y\in M'\) it has a
holomorphic local inverse: a map \(\psi\), holomorphic near \(\varphi(y)\), with \(\psi\circ\varphi=\mathrm{id}\) near \(y\) and \(\varphi\circ\psi=\mathrm{id}\) near
\(\varphi(y)\). Biholomorphisms and inclusions of open sets are local biholomorphisms.

\begin{lemma}[Transfer of tangent vector fields]
\label{mc:field-transfer}
\uses{mc:derivative-along,mc:field-chain-rule,hir:derivative-chain-rule,ana:ideal-sheaves}
\source{wlodarczyk05-hironaka}
\lean{Hironaka.DerivativeIdealAlongEtale.fieldDeriv_congr,Hironaka.DerivativeIdealAlongEtale.vanishingIdeal_mem_congr,Hironaka.DerivativeIdealAlongEtale.isTangentField_congr_set,Hironaka.DerivativeIdealAlongEtale.eventually_differentiableAt_chartMap_of_eventually,Hironaka.DerivativeIdealAlongEtale.holomorphicAt_chainCoeff_of_eventually,Hironaka.DerivativeIdealAlongEtale.chartDeriv_comp_eventuallyEq_sum_of_eventually,Hironaka.DerivativeIdealAlongEtale.transferField,Hironaka.DerivativeIdealAlongEtale.IsLocalInverseAt,Hironaka.DerivativeIdealAlongEtale.IsLocalInverseAt.symm,Hironaka.DerivativeIdealAlongEtale.IsLocalInverseAt.tendsto,Hironaka.DerivativeIdealAlongEtale.fieldDeriv_transfer,Hironaka.DerivativeIdealAlongEtale.holomorphicAt_transferField,Hironaka.DerivativeIdealAlongEtale.comp_mem_vanishingIdeal,Hironaka.DerivativeIdealAlongEtale.comp_mem_vanishingIdeal_preimage,Hironaka.DerivativeIdealAlongEtale.isTangentField_transfer}
\leanok
Let \(\psi\) be a holomorphic local inverse of \(\varphi\) at \(y\), \(e\) a chart at \(\varphi(y)\), \(e'\) a chart at \(y\), and \(\xi\) a holomorphic vector field near
\(\varphi(y)\) written in \(e\). Let \(\xi'\) be the field near \(y\), written in \(e'\), with coefficients \(\xi'_j=\sum_m(\xi_m\circ\varphi)\,(c_{mj}\circ\varphi)\), where
\(c_{mj}\) is the \(j\)-th coordinate of \(D(e'\circ\psi\circ e^{-1})\,e_m\). Then near \(y\), \(\xi'(g)=\xi(g\circ\psi)\circ\varphi\) for every \(g\) holomorphic at \(y\);
\(g\mapsto g\circ\psi\) maps germs vanishing on \(\varphi^{-1}N\) at \(y\) to germs vanishing on \(N\) at \(\varphi(y)\), and \(b\mapsto b\circ\varphi\) maps germs vanishing on \(N\)
at \(\varphi(y)\) to germs vanishing on \(\varphi^{-1}N\) at \(y\); hence, if \(\xi\) is tangent to \(N\) at \(\varphi(y)\), then \(\xi'\) is tangent to \(\varphi^{-1}N\) at \(y\).
Tangency at a point only depends on the set near that point.
\end{lemma}
\begin{proof}
\leanok
By the chain rule in charts for \(\psi\) (Lemma~\ref{hir:derivative-chain-rule}, whose proof only uses that the map is holomorphic near the point),
\(\partial^e_{e_m}(g\circ\psi)=\sum_jc_{mj}\,(\partial^{e'}_{e_j}g)\circ\psi\) near \(\varphi(y)\); compose with \(\varphi\), use \(\psi\circ\varphi=\mathrm{id}\) near \(y\), multiply by
\(\xi_m\circ\varphi\) and sum. If \(g\) vanishes on \(\varphi^{-1}N\) near \(y\) and \(q\in N\) is near \(\varphi(y)\), then \(\varphi(\psi(q))=q\in N\) and \(\psi(q)\) is near \(y\),
so \(g(\psi(q))=0\). For tangency: if \(g\) vanishes on \(\varphi^{-1}N\) near \(y\), then \(\xi(g\circ\psi)\) vanishes on \(N\) near \(\varphi(y)\), and so
\(\xi'(g)=\xi(g\circ\psi)\circ\varphi\) vanishes on \(\varphi^{-1}N\) near \(y\). The last statement holds because the ideal of a set at a point only depends on
the set near the point.
\end{proof}

\begin{proposition}[\(\mathcal D_N\) commutes with local biholomorphisms]
\label{mc:derivative-along-etale}
\uses{mc:field-transfer,mc:derivative-along,mc:derivative-along-basic,mc:field-chain-rule,hir:giraud-along,ana:ideal-sheaves}
\source{wlodarczyk05-hironaka}
\lean{Hironaka.DerivativeIdealAlongEtale.comp_mem_pullback,Hironaka.DerivativeIdealAlongEtale.pullback_derivIdealAlong_le,Hironaka.DerivativeIdealAlongEtale.derivIdealAlong_pullback_le,Hironaka.DerivativeIdealAlongEtale.pullback_derivIdealAlong,Hironaka.DerivativeIdealAlongEtale.pullback_iterDerivIdealAlong,Hironaka.DerivativeIdealAlongEtale.isLocalInverseAt_of_inverse,Hironaka.DerivativeIdealAlongEtale.pullback_iterDerivIdealAlong_of_inverse,Hironaka.DerivativeIdealAlongEtale.openInv,Hironaka.DerivativeIdealAlongEtale.isLocalInverseAt_val,Hironaka.DerivativeIdealAlongEtale.pullback_val_iterDerivIdealAlong}
\leanok
Let \(\varphi\colon M'\to M\) be a local biholomorphism, \(\mathfrak a\) an ideal of holomorphic germs on \(M\) and \(N\subset M\). Then
\(\varphi^*\mathcal D_N^j(\mathfrak a)=\mathcal D^j_{\varphi^{-1}N}(\varphi^*\mathfrak a)\) for every \(j\). In particular this holds for biholomorphisms and, with
\(\iota\colon U\to M\) the inclusion of an open set, \(\iota^*\mathcal D_N^j(\mathfrak a)=\mathcal D^j_{N\cap U}(\mathfrak a|_U)\).
\end{proposition}
\begin{proof}
\leanok
It suffices to treat \(j=1\) (induction; the germs of \(\mathcal D_N^j\) are holomorphic, Lemma~\ref{hir:giraud-along}). Fix \(y\) and a local inverse \(\psi\) at \(y\).
\(\subset\): a germ of \(\varphi^*\mathcal D_N(\mathfrak a)\) at \(y\) is a combination of \(b\circ\varphi\), \(b\) a combination of generators of \(\mathcal D_N(\mathfrak a)\) at
\(\varphi(y)\); the pullbacks of germs of \(\mathfrak a\) and of \(\mathcal I_N\) are germs of \(\varphi^*\mathfrak a\) and of \(\mathcal I_{\varphi^{-1}N}\), and for \(\xi\) tangent to \(N\) and
\(f\in\mathfrak a\), \(\xi(f)\circ\varphi=\xi(f\circ\varphi\circ\psi)\circ\varphi=\xi'(f\circ\varphi)\) with \(\xi'\) the transferred field, tangent to \(\varphi^{-1}N\)
(Lemma~\ref{mc:field-transfer}). \(\supset\): germs of \(\varphi^*\mathfrak a\) lie in \(\varphi^*\mathcal D_N(\mathfrak a)\) since \(\mathfrak a\subset\mathcal D_N(\mathfrak a)\)
(Lemma~\ref{mc:derivative-along-basic}); a germ \(t\) vanishing on \(\varphi^{-1}N\) is \((t\circ\psi)\circ\varphi\) with \(t\circ\psi\in\mathcal I_N\subset\mathcal D_N(\mathfrak a)\); and
for \(\xi'\) tangent to \(\varphi^{-1}N\) at \(y\) and \(f'=\sum_kc_k\,(b_k\circ\varphi)\), \(b_k\in\mathfrak a_{\varphi(y)}\), the product rule (Lemma~\ref{mc:field-chain-rule}) gives
\(\xi'(f')=\sum_k\xi'(c_k)\,(b_k\circ\varphi)+c_k\,\xi'(b_k\circ\varphi)\), and \(\xi'(b_k\circ\varphi)=\xi(b_k)\circ\varphi\) for the field \(\xi\) transferred back along \(\psi\)
(Lemma~\ref{mc:field-transfer} for the local inverse \(\varphi\) of \(\psi\) at \(\varphi(y)\)), which is tangent to \(\psi^{-1}\varphi^{-1}N\), a set equal to \(N\) near
\(\varphi(y)\). A biholomorphism has a global inverse; for the inclusion of an open set \(U\) take \(\psi=\) the identity on \(U\) (and any constant elsewhere), which
is holomorphic near every point of \(U\).
\end{proof}
This is Włodarczyk's Lemmas le:etale2'' and le:etale3'' for the derivative ideals along a submanifold, with local biholomorphisms in place of étale maps.

\subsection{Homogenized and coefficient ideals along a submanifold}

In this subsection \(N\subset M\) is a set, \(\mathfrak a\) an ideal of holomorphic germs on \(M\), \(\mu\ge1\), and \(\mathcal D_N\) the derivative ideal along \(N\)
(Definition~\ref{mc:derivative-along}). These are the single-stage statements of item (4) of Lemma~\ref{hir:along-calculus} for the homogenized and the
coefficient ideal; the constructions are those of Definitions~\ref{mc:homogenized-ideal} and~\ref{mc:coefficient-ideal} with \(\mathcal D_N\) in place of \(\mathcal D\).

\begin{definition}[Homogenized and coefficient ideals along \(N\)]
\label{mc:homogenized-along}
\uses{mc:derivative-along,mc:derivative-along-basic,mc:hom-sum,mc:marked-finite-sum,mc:derivative-marked,mc:homogenized-ideal,mc:coefficient-ideal,hir:marked-ideal}
\source{wlodarczyk05-hironaka,wlodarczyk09-analytic-resolution:page-0020,wlodarczyk09-analytic-resolution:page-0022,bierstone-milman08-functoriality}
\lean{Hironaka.HomogenizedIdealAlong.tangentIdealAlong,Hironaka.HomogenizedIdealAlong.homogenizedIdealAlong,Hironaka.HomogenizedIdealAlong.derivMarkedAlong,Hironaka.HomogenizedIdealAlong.derivMarkedAlong_N,Hironaka.HomogenizedIdealAlong.derivMarkedAlong_E,Hironaka.HomogenizedIdealAlong.derivMarkedAlong_ideal,Hironaka.HomogenizedIdealAlong.derivMarkedAlong_d,Hironaka.HomogenizedIdealAlong.coeffMarkedAlong,Hironaka.HomogenizedIdealAlong.homMarkedAlong,Hironaka.HomogenizedIdealAlong.homogenizedIdealAlong_univ,Hironaka.HomogenizedIdealAlong.derivMarkedAlong_eq_derivMarked,Hironaka.HomogenizedIdealAlong.coeffMarkedAlong_eq_coeffMarked,Hironaka.HomogenizedIdealAlong.homMarkedAlong_eq_homMarked,Hironaka.HomogenizedIdealAlong.coeffMarkedAlong_d}
\leanok
Put \(T_N(\mathfrak a)=\mathcal D_N^{\mu-1}(\mathfrak a)\) and
\[
 \mathcal H_N(\mathfrak a)=\sum_{i<\mu}\mathcal D_N^i(\mathfrak a)\cdot T_N(\mathfrak a)^i .
\]
For a marked ideal \(\underline{\mathcal I}=(M,N,E,\mathfrak a,d)\) put \(\mathcal D_N^j(\underline{\mathcal I})=(M,N,E,\mathcal D_N^j\mathfrak a,d-j)\), the
\emph{coefficient ideal along \(N\)}
\[
 \mathcal C_N(\underline{\mathcal I})=\underline{\mathcal I}+\mathcal D_N(\underline{\mathcal I})+\dots+\mathcal D_N^{d-1}(\underline{\mathcal I})
\]
(the sum of Definition~\ref{mc:marked-finite-sum}; its ideal is \(\sum_{i<d}(\mathcal D_N^i\mathfrak a)^{e_i}\) with \(e_i=\prod_{j<d,\,j\ne i}(d-j)\), its mark
\(\mu_{\mathcal C}=\prod_{j<d}(d-j)=e_{d-1}\)), and \(\mathcal H_N(\underline{\mathcal I})=(M,N,E,\mathcal H_N(\mathfrak a),d)\). For \(N=M\) these are \(\mathcal H(\mathfrak a)\),
\(\mathcal D^j(\underline{\mathcal I})\), \(\mathcal C(\underline{\mathcal I})\) and \(\mathcal H(\underline{\mathcal I})\).
\end{definition}
The last statement holds because \(\mathcal D_M^j=\mathcal D^j\) (Lemma~\ref{mc:derivative-along-basic}); the mark of \(\mathcal C_N\) is \(e_{d-1}\) since
\((d-(d-1))\,e_{d-1}=\prod_{j<d}(d-j)\).

\begin{lemma}[Coherence and cosupports of \(\mathcal H_N\) and \(\mathcal C_N\)]
\label{mc:homogenized-along-basic}
\uses{mc:homogenized-along,mc:derivative-along-coherent,mc:tangent-direction-along-basic,mc:hom-sum,mc:hom-sum-order,mc:finite-sum,hir:order-power-along,mc:marked-finite-sum}
\source{wlodarczyk05-hironaka,wlodarczyk09-analytic-resolution:page-0020}
\lean{Hironaka.HomogenizedIdealAlong.isCoherent_tangentIdealAlong,Hironaka.HomogenizedIdealAlong.isCoherent_homogenizedIdealAlong,Hironaka.HomogenizedIdealAlong.isCoherent_coeffMarkedAlong,Hironaka.HomogenizedIdealAlong.isCoherent_homMarkedAlong,Hironaka.HomogenizedIdealAlong.le_homogenizedIdealAlong,Hironaka.HomogenizedIdealAlong.orderGE_homogenizedIdealAlong,Hironaka.HomogenizedIdealAlong.cosupport_homogenizedIdealAlong,Hironaka.HomogenizedIdealAlong.cosupp_homMarkedAlong,Hironaka.HomogenizedIdealAlong.cosupport_coeffSummand,Hironaka.HomogenizedIdealAlong.cosupp_coeffMarkedAlong}
\leanok
\begin{enumerate}
\item If \(\mathfrak a\) is coherent and \(N\) is a closed submanifold, \(\mathcal H_N(\mathfrak a)\) and the ideal of \(\mathcal C_N(\underline{\mathcal I})\) are coherent.
\item \(\mathfrak a\subset\mathcal H_N(\mathfrak a)\) and \(\operatorname{cosupp}(\mathcal H_N(\mathfrak a)|_N,\mu)=\operatorname{cosupp}(\mathfrak a|_N,\mu)\), for every set \(N\).
\item If \(N\) is regular at each of its points and \(d\ge1\), the \(i\)-th summand of \(\mathcal C_N(\underline{\mathcal I})\) has
\(\operatorname{cosupp}((\mathcal D_N^i\mathfrak a)^{e_i}|_N,\mu_{\mathcal C})=\operatorname{cosupp}((\mathcal D_N^i\mathfrak a)|_N,d-i)\), and
\(\operatorname{cosupp}\mathcal C_N(\underline{\mathcal I})=\operatorname{cosupp}\underline{\mathcal I}\).
\end{enumerate}
\end{lemma}
\begin{proof}
\leanok
(1) The \(\mathcal D_N^i(\mathfrak a)\) are coherent (Corollary~\ref{mc:derivative-along-coherent}), and finite sums, products and powers of coherent ideals are coherent
(Definition~\ref{mc:hom-sum}, Lemma~\ref{mc:finite-sum}).
(2) The summand \(i=0\) is \(\mathfrak a\cdot\mathcal O=\mathfrak a\), so \(\operatorname{cosupp}(\mathcal H_N(\mathfrak a)|_N,\mu)\subset\operatorname{cosupp}(\mathfrak a|_N,\mu)\). Conversely, at
\(x\in\operatorname{cosupp}(\mathfrak a|_N,\mu)=\operatorname{cosupp}(\mathfrak a|_N,(\mu-i)+i)\), \(\operatorname{ord}_x((\mathcal D_N^i\mathfrak a)|_N)\ge\mu-i\) for \(i<\mu\)
(Lemma~\ref{mc:tangent-direction-along-basic}(1)), in particular \(\operatorname{ord}_x(T_N(\mathfrak a)|_N)\ge1\), and Lemma~\ref{mc:hom-sum-order} gives
\(\operatorname{ord}_x(\mathcal H_N(\mathfrak a)|_N)\ge\mu\).
(3) \((d-i)\,e_i=\mu_{\mathcal C}\), so the first statement is \(\operatorname{cosupp}(\mathfrak b^{\nu}|_N,(d-i)\nu)=\operatorname{cosupp}(\mathfrak b|_N,d-i)\) with \(\nu=e_i\ge1\)
(Lemma~\ref{hir:order-power-along}). The cosupport of a finite sum with a common mark is the intersection of the cosupports (Lemma~\ref{mc:finite-sum}, whose proof
is valid along any \(N\)), so \(\operatorname{cosupp}\mathcal C_N(\underline{\mathcal I})=\bigcap_{i<d}\operatorname{cosupp}((\mathcal D_N^i\mathfrak a)|_N,d-i)\). The term \(i=0\) is
\(\operatorname{cosupp}\underline{\mathcal I}\), and it is contained in every other term by Lemma~\ref{mc:tangent-direction-along-basic}(1).
\end{proof}

\begin{definition}[Maximal order along \(N\)]
\label{mc:maximal-order-along}
\uses{mc:maximal-order,hir:cosupport,hir:giraud-along,hir:giraud-support,hir:companion-maximal-order,hir:cosupport-basic}
\source{wlodarczyk05-hironaka,bierstone-milman08-functoriality}
\lean{Hironaka.HomogenizedIdealAlong.IsMaximalOrderAlong,Hironaka.HomogenizedIdealAlong.isMaximalOrderAlong_univ_iff,Hironaka.HomogenizedIdealAlong.isMaximalOrderAlong_iff_inter_zeroLocus}
\leanok
\((\mathfrak a|_N,\mu)\) is \emph{of maximal order along \(N\)} if \(\operatorname{ord}_x(\mathfrak a|_N)\le\mu\) for every \(x\), i.e.\
\(\operatorname{cosupp}(\mathfrak a|_N,\mu+1)=\emptyset\). For \(N=M\) and \(\mathfrak a\) coherent this is maximal order (Definition~\ref{mc:maximal-order}). If \(N\) is
regular at each of its points, it holds if and only if \(\mathcal D_N^\mu(\mathfrak a)\) vanishes at no point of \(N\), i.e.\ \((\mathcal D_N^\mu\mathfrak a)|_N=\mathcal O_N\).
\end{definition}
For \(N=M\): if \(\mathcal D^\mu(\mathfrak a)=\mathcal O\), a point of \(\operatorname{cosupp}(\mathfrak a,1+\mu)\subset\operatorname{cosupp}(\mathcal D^\mu\mathfrak a,1)\)
(Lemma~\ref{hir:giraud-support}) would be a zero of the constant \(1\); the converse is Lemma~\ref{hir:companion-maximal-order}. On a submanifold,
\(\operatorname{cosupp}(\mathfrak a|_N,\mu+1)=N\cap V(\mathcal D_N^\mu\mathfrak a)\) (Lemma~\ref{hir:giraud-along}).

\begin{lemma}[Orders of powers along \(N\), strict form]
\label{mc:order-power-strict}
\uses{hir:local-retraction,hir:order-product,hir:companion-maximal-order,hir:point-power-pullback,hir:order-power-along}
\source{wlodarczyk05-hironaka}
\lean{Hironaka.HomogenizedIdealAlong.not_vanishesToOrderOn_pow}
\leanok
Let \(N\) be regular at \(x\in N\), \(g\) holomorphic near \(x\), \(o,e\in\mathbb N\). If \(\operatorname{ord}_x(g|_N)\le o\), then
\(\operatorname{ord}_x(g^e|_N)\le eo\).
\end{lemma}
\begin{proof}
\leanok
Let \(r\) be a local retraction onto \(N\) (Lemma~\ref{hir:local-retraction}) and \(\hat g=g\circ r\), which agrees with \(g\) on \(N\) near \(x\). If \(g^e=h\) on \(N\) near
\(x\) with \(h\in\mathfrak m_x^{eo+1}\), then \(\hat g^e=h\circ r\in\mathfrak m_x^{eo+1}\) near \(x\) (Lemma~\ref{hir:point-power-pullback}, \(r(x)=x\)). Since orders add under
products (Lemma~\ref{hir:order-product}, in the form \(g\notin\mathfrak m_x^{o+1}\Rightarrow g^e\notin\mathfrak m_x^{eo+1}\) of Lemma~\ref{hir:companion-maximal-order}),
\(\hat g\in\mathfrak m_x^{o+1}\), and then \(g|_N\) vanishes to order \(o+1\) at \(x\), with witness \(\hat g\).
\end{proof}

\begin{proposition}[Maximal order of \(\mathcal H_N\) and \(\mathcal C_N\)]
\label{mc:homogenized-along-maximal-order}
\uses{mc:homogenized-along,mc:homogenized-along-basic,mc:maximal-order-along,mc:order-power-strict,hir:giraud-along,mc:homogenized-maximal-order}
\source{wlodarczyk05-hironaka}
\lean{Hironaka.HomogenizedIdealAlong.not_orderGE_homogenizedIdealAlong,Hironaka.HomogenizedIdealAlong.IsMaximalOrderAlong.homogenizedIdealAlong,Hironaka.HomogenizedIdealAlong.not_orderGE_coeffMarkedAlong,Hironaka.HomogenizedIdealAlong.IsMaximalOrderAlong.coeffMarkedAlong,Hironaka.HomogenizedIdealAlong.maximalOrderAt_hom_coeff}
\leanok
Let \(\underline{\mathcal I}=(M,N,E,\mathfrak a,d)\) with \(d\ge1\).
\begin{enumerate}
\item If \(\operatorname{ord}_x(\mathfrak a|_N)\le d\), then \(\operatorname{ord}_x(\mathcal H_N(\mathfrak a)|_N)\le d\). Hence \(\mathcal H_N(\underline{\mathcal I})\) is of maximal order along
\(N\) when \(\underline{\mathcal I}\) is.
\item Let \(N\) be regular at each of its points. If \(x\in N\) and \(\operatorname{ord}_x(\mathfrak a|_N)\le d\), then \(\operatorname{ord}_x(\mathcal C_N(\underline{\mathcal I})|_N)\le\mu_{\mathcal C}\).
Hence \(\mathcal C_N(\underline{\mathcal I})\) is of maximal order along \(N\) when \(\underline{\mathcal I}\) is.
\item In particular, at a point \(x\in\operatorname{cosupp}\underline{\mathcal I}\) with \(\operatorname{ord}_x(\mathfrak a|_N)=d\), \(x\) lies in the cosupports of \(\mathcal H_N(\underline{\mathcal I})\)
and \(\mathcal C_N(\underline{\mathcal I})\), and the orders there equal the marks.
\end{enumerate}
\end{proposition}
\begin{proof}
\leanok
(1) \(\mathfrak a\subset\mathcal H_N(\mathfrak a)\) (Lemma~\ref{mc:homogenized-along-basic}), so a lower bound for the order of \(\mathcal H_N(\mathfrak a)|_N\) bounds that of
\(\mathfrak a|_N\).
(2) By Giraud--Villamayor along \(N\) (Lemma~\ref{hir:giraud-along}), \(\operatorname{cosupp}(\mathfrak a|_N,d+1)=\operatorname{cosupp}((\mathcal D_N^{d-1}\mathfrak a)|_N,2)\), so
some \(g\in\mathcal D_N^{d-1}(\mathfrak a)_x\) has \(\operatorname{ord}_x(g|_N)\le1\). Then \(g^{\mu_{\mathcal C}}\) is a germ of the last summand
\((\mathcal D_N^{d-1}\mathfrak a)^{e_{d-1}}\), \(e_{d-1}=\mu_{\mathcal C}\), of the ideal of \(\mathcal C_N(\underline{\mathcal I})\), and
\(\operatorname{ord}_x(g^{\mu_{\mathcal C}}|_N)\le\mu_{\mathcal C}\) by Lemma~\ref{mc:order-power-strict}. Off \(N\) the cosupports are empty.
(3) Lemma~\ref{mc:homogenized-along-basic}(2)--(3) and (1)--(2).
\end{proof}
For \(N=M\) this is Lemma~\ref{mc:homogenized-maximal-order}.

\begin{lemma}[\(\mathcal H_N\) and \(\mathcal C_N\) are étale-local]
\label{mc:homogenized-along-etale}
\uses{mc:homogenized-along,mc:derivative-along-etale,mc:derivative-along-coherent,mc:pullback-finite-sum,hir:controlled-operations}
\source{wlodarczyk05-hironaka}
\lean{Hironaka.HomogenizedIdealAlong.pullback_tangentIdealAlong,Hironaka.HomogenizedIdealAlong.pullback_homogenizedIdealAlong,Hironaka.HomogenizedIdealAlong.pullback_coeffMarkedAlong,Hironaka.HomogenizedIdealAlong.pullback_homMarkedAlong,Hironaka.HomogenizedIdealAlong.isLocalInverseAt_of_contMDiff,Hironaka.HomogenizedIdealAlong.pullback_homogenizedIdealAlong_of_inverse,Hironaka.HomogenizedIdealAlong.pullback_coeffMarkedAlong_of_inverse,Hironaka.HomogenizedIdealAlong.pullback_homMarkedAlong_of_inverse,Hironaka.HomogenizedIdealAlong.pullback_val_homogenizedIdealAlong,Hironaka.HomogenizedIdealAlong.pullback_val_coeffMarkedAlong,Hironaka.HomogenizedIdealAlong.pullback_val_homMarkedAlong}
\leanok
Let \(\varphi\colon M'\to M\) be a holomorphic local biholomorphism, \(\mathfrak a\) coherent and \(N\) a closed submanifold. Then
\(\varphi^*T_N(\mathfrak a)=T_{\varphi^{-1}N}(\varphi^*\mathfrak a)\), \(\varphi^*\mathcal H_N(\mathfrak a)=\mathcal H_{\varphi^{-1}N}(\varphi^*\mathfrak a)\),
\(\varphi^*\mathcal C_N(\underline{\mathcal I})=\mathcal C_{\varphi^{-1}N}(\varphi^*\underline{\mathcal I})\) and
\(\varphi^*\mathcal H_N(\underline{\mathcal I})=\mathcal H_{\varphi^{-1}N}(\varphi^*\underline{\mathcal I})\) (\(d\ge1\)). In particular this holds for biholomorphisms and for
inclusions of open sets.
\end{lemma}
\begin{proof}
\leanok
\(\varphi^*\mathcal D_N^j(\mathfrak a)=\mathcal D^j_{\varphi^{-1}N}(\varphi^*\mathfrak a)\) (Proposition~\ref{mc:derivative-along-etale}); pullbacks commute with finite sums, products
and powers of coherent ideals (Lemma~\ref{mc:pullback-finite-sum}), and the \(\mathcal D_N^j(\mathfrak a)\) are coherent (Corollary~\ref{mc:derivative-along-coherent}). The
submanifold of \(\varphi^*\underline{\mathcal I}\) is \(\varphi^{-1}N\) and the marks are unchanged. A biholomorphism and the inclusion of an open set are holomorphic local
biholomorphisms (Proposition~\ref{mc:derivative-along-etale}).
\end{proof}
These are Włodarczyk's Lemmas ``le:etale2'' and ``le:etale3'' for \(\mathcal H\) and \(\mathcal C\) along a submanifold.

\subsection{The glueing lemma along a submanifold}

In this subsection \(\mathfrak a\) is a coherent ideal, \(N\subset M\) a closed submanifold, \(\mu\ge1\), \(T_N(\mathfrak a)=\mathcal D_N^{\mu-1}(\mathfrak a)\) and
\(\mathcal H_N(\mathfrak a)\) the homogenized ideal along \(N\) (Definition~\ref{mc:homogenized-along}). A \emph{flat chart} of \(N\) at \(x\in N\) is a local
biholomorphism \(\Theta\) of \(\mathbf C^n\) defined near \(e(x)\), \(e\) the chart at \(x\), whose domain lies in the image of \(e\), such that near \(x\), \(y\in N\) if and only if
\(\psi(y)=\Theta(e(y))\in\mathbf C^k\times0\) (as in Lemma~\ref{hir:order-along-flat}); \(\psi^{-1}=e^{-1}\circ\Theta^{-1}\). The construction is that of the
glueing lemma for \(N=M\) (Theorem~\ref{mc:glueing-lemma}) with coordinates that also flatten \(N\).

\begin{definition}[The coordinate change of a flat chart]
\label{mc:flat-change}
\uses{mc:coordinate-change,hir:order-along-flat}
\source{wlodarczyk05-hironaka,wlodarczyk09-analytic-resolution:page-0021}
\lean{Hironaka.HomogenizedGlueingAlong.flatChange,Hironaka.HomogenizedGlueingAlong.flatChange_coordMap,Hironaka.HomogenizedGlueingAlong.continuousAt_flatMap,Hironaka.HomogenizedGlueingAlong.flatMap_comp_symm_eventuallyEq,Hironaka.HomogenizedGlueingAlong.injective_fderiv_toFun,Hironaka.HomogenizedGlueingAlong.exists_flatChart_of_regularAtDim}
\leanok
Let \(\psi=\Theta\circ e\) be a flat chart of \(N\) at \(x\), \(h\) holomorphic at \(x\) with \(h(x)=0\) and \(i\) an index. The \emph{coordinate change of \(\psi\)} is
\(\varphi(p)=\psi^{-1}\bigl(\psi(p)+h(p)\,e_i\bigr)\); it is the coordinate change of Definition~\ref{mc:coordinate-change} with \(W=\psi\) and \(\Psi=\Theta^{-1}\).
The flat chart of a submanifold regular of dimension \(k\) at \(x\) can be taken centred, \(\psi(x)=0\).
\end{definition}
\(W=\psi\) is a coordinate system at \(x\): \(\psi\circ e^{-1}=\Theta\) near \(e(x)\), whose derivative is injective because \(\Theta^{-1}\circ\Theta=\mathrm{id}\); and
\(\Theta^{-1}(\psi(p))=e(p)\) near \(x\).

\begin{lemma}[Tangential coordinate derivatives in a flat chart]
\label{mc:flat-coordinate-iterated}
\uses{mc:flat-change,hir:order-along-flat,mc:derivative-along,mc:homogenized-along,mc:homogenized-basic,mc:coordinate-iterated}
\source{wlodarczyk05-hironaka}
\lean{Hironaka.HomogenizedGlueingAlong.iterDerivIdealAlong_iterDerivIdealAlong,Hironaka.HomogenizedGlueingAlong.tangentIdealAlong_iterDerivIdealAlong,Hironaka.HomogenizedGlueingAlong.homogenizedIdealAlong_one,Hironaka.HomogenizedGlueingAlong.comp_flatMap_flatInv_eventuallyEq,Hironaka.HomogenizedGlueingAlong.fieldDeriv_flatField_comp,Hironaka.HomogenizedGlueingAlong.iterDeriv_comp_flatMap_mem}
\leanok
\begin{enumerate}
\item \(\mathcal D_N^j(\mathcal D_N^i\mathfrak a)=\mathcal D_N^{i+j}\mathfrak a\); for \(i<\mu\) the tangent ideal of \((\mathcal D_N^i\mathfrak a,\mu-i)\) is \(T_N(\mathfrak a)\); and
\(\mathcal H_N(\mathfrak b,1)=\mathfrak b\) for a coherent \(\mathfrak b\).
\item Let \(\psi\) be a flat chart of \(N\) at \(x\) and \(l<k\). For \(G\) with \(G\circ\psi\) holomorphic near a point \(p\) of the domain of \(\psi\), the derivative of
\(G\circ\psi\) along the coordinate field \(\xi_{e_l}\) of \(\psi\) (Lemma~\ref{hir:order-along-flat}) is \((\partial_lG)\circ\psi\) at \(p\). Hence, if \(F\circ\psi\in\mathfrak a_x\),
then \((\partial_l^jF)\circ\psi\in\mathcal D_N^j(\mathfrak a)_x\) for every \(j\).
\end{enumerate}
\end{lemma}
\begin{proof}
\leanok
(1) Induction on \(j\); the tangent ideal of \((\mathcal D_N^i\mathfrak a,\mu-i)\) is \(\mathcal D_N^{\mu-i-1}\mathcal D_N^i\mathfrak a=\mathcal D_N^{\mu-1}\mathfrak a\); for the mark \(1\),
\(\mathcal H_N(\mathfrak b,1)\) has the single summand \(\mathcal D_N^0(\mathfrak b)\cdot\mathcal O=\mathfrak b\) (Lemma~\ref{mc:homogenized-basic}).
(2) By Lemma~\ref{hir:order-along-flat}, \(\xi_{e_l}(G\circ\psi)(p)=D(G\circ\psi\circ\psi^{-1})(\psi(p))\,e_l\), and \(G\circ\psi\circ\psi^{-1}=G\) near \(\psi(p)\). For the second
statement induct on \(j\): \((\partial_l^{j+1}F)\circ\psi=\xi_{e_l}\bigl((\partial_l^jF)\circ\psi\bigr)\) near \(x\), and \(\xi_{e_l}\) is tangent to \(N\) at \(x\) because \(e_l\in\mathbf C^k\times0\)
(Lemma~\ref{hir:order-along-flat}(2)); so it is a generator of \(\mathcal D_N(\mathcal D_N^j\mathfrak a)_x\) (Definition~\ref{mc:derivative-along}). This replaces
Lemma~\ref{mc:coordinate-iterated}, where every coordinate field is admissible.
\end{proof}

\begin{lemma}[The glueing inclusion along \(N\)]
\label{mc:glueing-inclusion-along}
\uses{mc:flat-coordinate-iterated,mc:flat-change,mc:taylor-tail,mc:homogenized-along,mc:glueing-inclusion}
\source{wlodarczyk05-hironaka,wlodarczyk09-analytic-resolution:page-0021}
\lean{Hironaka.HomogenizedGlueingAlong.comp_mem_homogenizedIdealAlong}
\leanok
Let \(\psi\) be a flat chart of \(N\) at \(x\), \(h\in T_N(\mathfrak a)_x\) with \(h(x)=0\), \(i<k\), and \(\varphi(p)=\psi^{-1}(\psi(p)+h(p)e_i)\). Then \(f\circ\varphi\in\mathcal H_N(\mathfrak a)_x\)
for every \(f\in\mathfrak a_x\).
\end{lemma}
\begin{proof}
\leanok
As for Lemma~\ref{mc:glueing-inclusion}: put \(F=f\circ\psi^{-1}\), holomorphic near \(\psi(x)\) with \(F\circ\psi=f\) near \(x\). The convergent Taylor tail
(Lemma~\ref{mc:taylor-tail}) gives \(G\), holomorphic near \((\psi(x),0)\), with
\[
 f\circ\varphi=\sum_{j<\mu}\frac1{j!}\,\bigl((\partial_i^jF)\circ\psi\bigr)\,h^j+G(\psi,h)\cdot h\cdot h^{\mu-1}\qquad\text{near }x.
\]
Since \(i<k\), \((\partial_i^jF)\circ\psi\in\mathcal D_N^j(\mathfrak a)_x\) (Lemma~\ref{mc:flat-coordinate-iterated}(2)), and \(h^j\in T_N(\mathfrak a)^j_x\); so each term is a multiple of a germ
of \(\mathcal D_N^j(\mathfrak a)\cdot T_N(\mathfrak a)^j\), the last one of \(\mathcal D_N^{\mu-1}(\mathfrak a)\cdot T_N(\mathfrak a)^{\mu-1}\).
\end{proof}

\begin{proposition}[The coordinate change of a flat chart preserves \(\mathcal H_N\)]
\label{mc:homogenized-along-invariance}
\uses{mc:glueing-inclusion-along,mc:flat-coordinate-iterated,mc:homogenized-product,mc:homogenized-invariance,mc:homogenized-along-basic,mc:derivative-along-coherent,mc:coordinate-change}
\source{wlodarczyk05-hironaka,wlodarczyk09-analytic-resolution:page-0021}
\lean{Hironaka.HomogenizedGlueingAlong.mul_mem_homogenizedIdealAlong,Hironaka.HomogenizedGlueingAlong.comp_mem_idealPow_tangentIdealAlong,Hironaka.HomogenizedGlueingAlong.comp_mem_homogenizedIdealAlong_of_mem}
\leanok
\(\mathcal H_N(\mathcal D_N^i\mathfrak a,\mu-i)\cdot T_N(\mathfrak a)^i\subset\mathcal H_N(\mathfrak a,\mu)\) for \(i<\mu\). With \(\psi\), \(h\in T_N(\mathfrak a)_x\), \(h(x)=0\), \(i<k\) and \(\varphi\)
as in Lemma~\ref{mc:glueing-inclusion-along}: \(b\circ\varphi\in(T_N(\mathfrak a)^j)_x\) for \(b\in(T_N(\mathfrak a)^j)_x\), and \(f\circ\varphi\in\mathcal H_N(\mathfrak a)_x\) for every
\(f\in\mathcal H_N(\mathfrak a)_x\).
\end{proposition}
\begin{proof}
\leanok
The first statement is the proof of Lemma~\ref{mc:homogenized-product} with \(\mathcal D_N\) in place of \(\mathcal D\), using
Lemma~\ref{mc:flat-coordinate-iterated}(1). The rest is the proof of Proposition~\ref{mc:homogenized-invariance}: the ideals \(\mathcal D_N^i\mathfrak a\) and \(T_N(\mathfrak a)\) are
coherent (Corollary~\ref{mc:derivative-along-coherent}); Lemma~\ref{mc:glueing-inclusion-along} for \((T_N(\mathfrak a),1)\), whose homogenized ideal is \(T_N(\mathfrak a)\), and products
give the statement for \(T_N(\mathfrak a)^j\); for a generator \(ab\) of \(\mathcal H_N(\mathfrak a)_x\), \(a\in\mathcal D_N^i(\mathfrak a)_x\), \(b\in(T_N(\mathfrak a)^i)_x\), the same lemma for
\((\mathcal D_N^i\mathfrak a,\mu-i)\) (tangent ideal \(T_N(\mathfrak a)\ni h\)) gives \(a\circ\varphi\in\mathcal H_N(\mathcal D_N^i\mathfrak a,\mu-i)_x\), and the first statement concludes.
Composition with \(\varphi\) respects holomorphic combinations (Definition~\ref{mc:coordinate-change}).
\end{proof}

\begin{lemma}[Flat charts from two tangent directions along \(N\)]
\label{mc:coordinate-pair-along}
\uses{mc:coordinate-pair,mc:flat-change,mc:contact-hypersurface-along,mc:field-chain-rule,mc:tangent-direction-along,hir:order-along-flat,bcalc:snc-chart,ana:local-holomorphic-inverse}
\source{wlodarczyk05-hironaka,wlodarczyk09-analytic-resolution:page-0021}
\lean{Hironaka.HomogenizedGlueingAlong.alongCoords,Hironaka.HomogenizedGlueingAlong.alongCoords_apply_self,Hironaka.HomogenizedGlueingAlong.alongCoords_apply_of_le,Hironaka.HomogenizedGlueingAlong.alongCoords_mem_plane_iff,Hironaka.HomogenizedGlueingAlong.alongCoords_sub,Hironaka.HomogenizedGlueingAlong.mdifferentiableAt_alongCoords,Hironaka.HomogenizedGlueingAlong.injective_killMap_add,Hironaka.HomogenizedGlueingAlong.injective_fderiv_alongCoords,Hironaka.HomogenizedGlueingAlong.exists_flatChart_alongCoords,Hironaka.HomogenizedGlueingAlong.exists_common_direction,Hironaka.HomogenizedGlueingAlong.eventually_flatChange_pair}
\leanok
Let \(\psi=\Phi\circ e\) be a flat chart of \(N\) at \(x\in N\) with \(\psi(x)=0\).
\begin{enumerate}
\item Let \(u\), \(v\) be holomorphic at \(x\) with \(\xi(u)(x)\ne0\) and \(\xi'(v)(x)\ne0\) for vector fields \(\xi\), \(\xi'\) tangent to \(N\) at \(x\). There are
\(w\in\mathbf C^k\times0\) and \(i<k\) with \(w_i\ne0\), \(\ell_u(w)\ne0\) and \(\ell_v(w)\ne0\), where \(\ell_u=D(u\circ\psi^{-1})(0)\), \(\ell_v=D(v\circ\psi^{-1})(0)\).
\item For such \(w\), \(i\) and \(u\) holomorphic on an open \(U\ni x\) with \(u(x)=0\), the coordinates \(W_u=P_w\circ\psi+u\,e_i\) (\(P_w(\zeta)=\zeta-(\zeta_i/w_i)\,w\))
have \(u\) as \(i\)-th coordinate and the normal coordinates \(\psi_r\), \(r\ge k\), of \(\psi\); they form a flat chart of \(N\) at \(x\), centred at \(0\), on a neighbourhood of \(x\) in \(U\);
and \(W_v=W_u+(v-u)\,e_i\).
\item Let \(\psi_a\), \(\psi_b\) be flat charts of \(N\) at \(x\), centred at \(0\), with \(\psi_b=\psi_a+(b-a)\,e_i\) and \((\psi_a)_i=a\) near \(x\), and let \(\varphi\), \(\varphi'\) be the
coordinate changes of \(\psi_a\) with \(h=b-a\) and of \(\psi_b\) with \(h=a-b\) (Definition~\ref{mc:flat-change}). Then near \(x\): \(a\circ\varphi=b\), \(\varphi'\circ\varphi=\mathrm{id}\), and
\(\varphi(p)\in N\iff p\in N\).
\end{enumerate}
\end{lemma}
\begin{proof}
\leanok
(1) By Lemma~\ref{mc:contact-hypersurface-along}(1), \(\partial u/\partial\psi_j(x)=\ell_u(e_j)\ne0\) for some \(j<k\), so \(\ell_u\circ P\ne0\) for the projection \(P\) onto
\(\mathbf C^k\times0\); likewise \(\ell_v\circ P\ne0\). Two nonzero linear forms are both nonzero on some vector \(w_0\) (proof of Lemma~\ref{mc:coordinate-pair}); put \(w=Pw_0\ne0\)
and choose \(i\) with \(w_i\ne0\); then \(i<k\) since \(w\in\mathbf C^k\times0\).
(2) The \(i\)-th coordinate of \(P_w\zeta\) is \(0\), and for \(r\ge k\), \((P_w\zeta)_r=\zeta_r\) since \(w_r=0\); also \(r\ne i\). Hence \(W_u(p)\in\mathbf C^k\times0\iff\psi(p)\in\mathbf C^k\times0\iff
p\in N\) on the domain of \(\psi\). In the chart, \(W_u\circ e^{-1}=(P_w+(u\circ\psi^{-1})\,e_i)\circ\Phi\) near \(e(x)\), with derivative \((P_w+\ell_u\otimes e_i)\circ D\Phi(e(x))\); as in
Lemma~\ref{mc:coordinate-pair}, \(P_w\xi+\ell_u(\xi)e_i=0\) forces \(\ell_u(\xi)=0\), then \(\xi=c\,w\) and \(c\,\ell_u(w)=0\), so \(\xi=0\). By the inverse function theorem
(Lemma~\ref{ana:local-holomorphic-inverse}) \(W_u-W_u(x)=W_u\) is a local coordinate system centred at \(x\) (Lemma~\ref{bcalc:snc-chart}) and it flattens \(N\). The last identity is the
definition.
(3) Near \(x\), \(\varphi(p)=\psi_a^{-1}(\psi_a(p)+(b-a)(p)e_i)=\psi_a^{-1}(\psi_b(p))\); as \(\psi_b(x)=0=\psi_a(x)\), \(\psi_a(\varphi(p))=\psi_b(p)\) for \(p\) near \(x\), and
\(\varphi(p)\) is near \(x\). So \(a(\varphi(p))=(\psi_a(\varphi(p)))_i=(\psi_b(p))_i=a(p)+(b(p)-a(p))=b(p)\); \(\psi_b(q)+(a-b)(q)e_i=\psi_a(q)\) at \(q=\varphi(p)\) gives
\(\varphi'(\varphi(p))=\psi_b^{-1}(\psi_a(\varphi(p)))=\psi_b^{-1}(\psi_b(p))=p\); and \(\varphi(p)\in N\iff\psi_a(\varphi(p))=\psi_b(p)\in\mathbf C^k\times0\iff p\in N\).
\end{proof}

\begin{theorem}[Glueing Lemma along a submanifold]
\label{mc:glueing-lemma-along}
\uses{mc:coordinate-pair-along,mc:homogenized-along-invariance,mc:homogenized-along-basic,mc:coherent-germs,mc:tangent-direction-along,mc:tangent-direction-along-basic,mc:glueing-lemma,mc:derivative-along-coherent}
\source{wlodarczyk05-hironaka,wlodarczyk09-analytic-resolution:page-0021}
\lean{Hironaka.HomogenizedGlueingAlong.eventually_pullback_iff_of_inverse,Hironaka.HomogenizedGlueingAlong.glueing_lemma_along}
\leanok
Let \(u\) and \(v\) be tangent directions of \((\mathfrak a|_N,\mu)\) along \(N\) (Definition~\ref{mc:tangent-direction-along}) on open sets containing
\(x\in\operatorname{cosupp}(\mathfrak a|_N,\mu)\). There are maps \(\varphi,\psi\colon M\to M\), holomorphic near \(x\), inverse to each other near \(x\), with
\(\varphi(x)=\psi(x)=x\), such that near \(x\):
\begin{enumerate}
\item \(\varphi^*\mathcal H_N(\mathfrak a)=\mathcal H_N(\mathfrak a)\) and \(\psi^*\mathcal H_N(\mathfrak a)=\mathcal H_N(\mathfrak a)\);
\item \(u\circ\varphi=v\) and \(v\circ\psi=u\);
\item \(\varphi\) and \(\psi\) fix every point of \(\operatorname{cosupp}(\mathfrak a|_N,\mu)\);
\item \(\varphi^{-1}(N)=N\) and \(\psi^{-1}(N)=N\).
\end{enumerate}
\end{theorem}
\begin{proof}
\leanok
The germs of \(u\), \(v\) at \(x\) lie in \(T_N(\mathfrak a)_x\), so they vanish at \(x\) (Lemma~\ref{mc:tangent-direction-along-basic}), and fields tangent to \(N\) have nonzero
derivatives of \(u\), \(v\) at \(x\). Take a flat chart \(\psi\) of \(N\) at \(x\) centred at \(0\), and \(w\), \(i\) as in Lemma~\ref{mc:coordinate-pair-along}(1); by (2) of that lemma
\(\psi_u=W_u\) and \(\psi_v=W_v\) are flat charts of \(N\) at \(x\), centred at \(0\), with \(\psi_v=\psi_u+(v-u)e_i\), \((\psi_u)_i=u\), \((\psi_v)_i=v\). Let \(\varphi\) be the coordinate change
of \(\psi_u\) with \(h=v-u\) and \(\psi\) that of \(\psi_v\) with \(h=u-v\) (Definition~\ref{mc:flat-change}). Lemma~\ref{mc:coordinate-pair-along}(3), applied to \((\psi_u,\psi_v)\) and to
\((\psi_v,\psi_u)\), gives (2), (4) and \(\psi\circ\varphi=\mathrm{id}=\varphi\circ\psi\) near \(x\). \(T_N(\mathfrak a)\) is coherent (Corollary~\ref{mc:derivative-along-coherent}), so \(\pm(v-u)\)
are germs of \(T_N(\mathfrak a)\) at \(x\) (Lemma~\ref{mc:coherent-germs}(1)), and Proposition~\ref{mc:homogenized-along-invariance} gives \(\varphi^*\mathcal H_N(\mathfrak a)\subset\mathcal H_N(\mathfrak a)\)
and \(\psi^*\mathcal H_N(\mathfrak a)\subset\mathcal H_N(\mathfrak a)\) at \(x\); as \(\mathcal H_N(\mathfrak a)\) is coherent (Lemma~\ref{mc:homogenized-along-basic}), Lemma~\ref{mc:coherent-germs}(2) extends
both inclusions to the points near \(x\), and the inverse map gives the reverse inclusions, as in the proof of Theorem~\ref{mc:glueing-lemma}. (3): the germs of \(u\), \(v\) at the
points \(y\) near \(x\) lie in \(T_N(\mathfrak a)_y\), so \(u\), \(v\) vanish on \(\operatorname{cosupp}(\mathfrak a|_N,\mu)\) near \(x\) (Lemma~\ref{mc:tangent-direction-along-basic}); there
\(h=0\), and a coordinate change fixes the points where its \(h\) vanishes (Lemma~\ref{mc:coordinate-pair}).
\end{proof}
This is Włodarczyk's Glueing Lemma ``le:homo'' (1)--(3) for the marked ideal \((\mathfrak a|_N,\mu)\) on \(N\), in the ambient encoding: the coordinate change is a
biholomorphism of neighbourhoods of \(x\) in \(M\) which preserves \(N\), so that its restriction to \(N\) is the coordinate change of the source on \(N\). Part (4) of the source
along multiple test blow-ups needs the transforms along \(N\) (Lemma~\ref{hir:along-calculus}(2)).

\subsection{The restricted coefficient ideal along a submanifold}

\begin{lemma}[The support of the restricted coefficient ideal along \(N\)]
\label{mc:coefficient-restriction-along}
\uses{mc:homogenized-along,mc:homogenized-along-basic,mc:coefficient-restriction-support,mc:restrict-order,mc:restrict-marked,hir:giraud-along,mc:tangent-direction-along,mc:tangent-direction-along-basic,hir:order-basic}
\source{wlodarczyk05-hironaka,wlodarczyk09-analytic-resolution:page-0023}
\lean{Hironaka.CoefficientIdealAlongRestriction.cosupport_inter_subset_of_subset,Hironaka.CoefficientIdealAlongRestriction.cosupport_coeffMarkedAlong_subset,Hironaka.CoefficientIdealAlongRestriction.cosupp_coeffMarkedAlong_withN,Hironaka.CoefficientIdealAlongRestriction.cosupp_contact_along}
\leanok
Let \(\underline{\mathcal I}=(M,N,E,\mathfrak a,d)\) with \(N\) regular at each of its points, \(\mathfrak a\) an ideal of holomorphic germs and \(d\ge1\), and let
\(\mathcal C_N(\underline{\mathcal I})|_S\) be the coefficient ideal along \(N\) (Definition~\ref{mc:homogenized-along}) with \(N\) replaced by \(S\)
(Definition~\ref{mc:restrict-marked}).
\begin{enumerate}
\item For \(S\subset N\), \(\operatorname{cosupp}(\mathfrak b|_N,d)\cap S\subset\operatorname{cosupp}(\mathfrak b|_S,d)\) for every ideal \(\mathfrak b\).
\item For \(A\subset N\), \(\operatorname{cosupp}(\mathcal C_N(\underline{\mathcal I})|_A)\subset\operatorname{cosupp}\underline{\mathcal I}\); hence
\(\operatorname{cosupp}(\mathcal C_N(\underline{\mathcal I})|_S)=\operatorname{cosupp}\underline{\mathcal I}\cap S\) for every \(S\subset N\).
\item For a tangent direction \(u\) of \((\mathfrak a|_N,d)\) along \(N\) on \(U\) (Definition~\ref{mc:tangent-direction-along}),
\(\operatorname{cosupp}(\mathcal C_N(\underline{\mathcal I})|_{N\cap V(u)\cap U})=\operatorname{cosupp}\underline{\mathcal I}\cap U\).
\end{enumerate}
\end{lemma}
\begin{proof}
\leanok
(1) A function agreeing on \(N\) near \(x\) with a germ of \(\mathfrak m_x^d\) agrees with it on \(S\). (2) Let \(x\in A\) with
\(\operatorname{ord}_x(\mathcal C_N(\mathfrak a)|_A)\ge\mu_{\mathcal C}\), and \(g\in\mathcal D_N^{d-1}(\mathfrak a)_x\). Then \(g^{\mu_{\mathcal C}}\) is a germ of the last summand of
\(\mathcal C_N(\mathfrak a)\) (Definition~\ref{mc:homogenized-along}, \(\mu_{\mathcal C}=e_{d-1}\ge1\)), so it vanishes at \(x\in A\), and so does \(g\). Hence
\(x\in N\cap V(\mathcal D_N^{d-1}\mathfrak a)=\operatorname{cosupp}(\mathfrak a|_N,d)\) (Lemma~\ref{hir:giraud-along}). The reverse inclusion is
\(\operatorname{cosupp}\underline{\mathcal I}=\operatorname{cosupp}\mathcal C_N(\underline{\mathcal I})\) (Lemma~\ref{mc:homogenized-along-basic}) followed by (1).
(3) \(\operatorname{cosupp}\underline{\mathcal I}\cap U\subset N\cap V(u)\) (Lemma~\ref{mc:tangent-direction-along-basic}); apply (2) with \(S=N\cap V(u)\cap U\).
\end{proof}
This is the first statement of Włodarczyk's Lemma ``le:S'' (Proposition~5.6.4 (1) of the analytic paper) for the marked ideal on \(N\), and the single-stage part of
the reduction to a hypersurface of maximal contact inside \(N\) (Lemmas~\ref{hir:coefficient-restriction} and~\ref{hir:contact-equivalence} along \(N\)); the
statements along multiple test blow-ups need the transforms along \(N\) (Lemma~\ref{hir:along-calculus}(2)).

\subsection{The glueing lemma along a submanifold on open sets and along multiple test blow-ups}

In this subsection \(\mathfrak a\) is coherent, \(N\) a closed submanifold, \(\mu\ge1\), and \(\underline{\mathcal H}_N=(M,N,\emptyset,\mathcal H_N(\mathfrak a),\mu)\) is the
homogenized marked ideal along \(N\); for an open \(U\subset M\), \(\underline{\mathcal H}_N|_U\) is its restriction (Lemma~\ref{hir:open-restriction-marked}).

\begin{lemma}[The glueing biholomorphism along \(N\)]
\label{mc:glueing-biholomorphism-along}
\uses{mc:glueing-lemma-along,mc:local-inverse-biholomorphism,mc:glueing-biholomorphism,mc:homogenized-along-basic,mc:homogenized-along,hir:open-restriction-marked,hir:marked-ideal}
\source{wlodarczyk05-hironaka,wlodarczyk09-analytic-resolution:page-0021}
\lean{Hironaka.HomogenizedGlueingAlongTower.homogenizedMarkedAlong,Hironaka.HomogenizedGlueingAlongTower.cosupp_homogenizedMarkedAlong,Hironaka.HomogenizedGlueingAlongTower.GlueingDataAlong,Hironaka.HomogenizedGlueingAlongTower.exists_glueingDataAlong,Hironaka.HomogenizedGlueingAlongTower.GlueingDataAlong.dom,Hironaka.HomogenizedGlueingAlongTower.GlueingDataAlong.cod,Hironaka.HomogenizedGlueingAlongTower.GlueingDataAlong.equiv,Hironaka.HomogenizedGlueingAlongTower.GlueingDataAlong.equiv_apply,Hironaka.HomogenizedGlueingAlongTower.GlueingDataAlong.mem_dom_and_mem_cod,Hironaka.HomogenizedGlueingAlongTower.GlueingDataAlong.apply_eq_of_mem_cosupport,Hironaka.HomogenizedGlueingAlongTower.GlueingDataAlong.preimage_inter_eq,Hironaka.HomogenizedGlueingAlongTower.GlueingDataAlong.preimage_zeroSet,Hironaka.HomogenizedGlueingAlongTower.GlueingDataAlong.preimage_N,Hironaka.HomogenizedGlueingAlongTower.GlueingDataAlong.pullback_comp_val,Hironaka.HomogenizedGlueingAlongTower.GlueingDataAlong.pullback_restrict,Hironaka.HomogenizedGlueingAlongTower.GlueingDataAlong.preimage_inter_eq_of_isAdmissibleCentre}
\leanok
Let \(u\) and \(v\) be tangent directions of \((\mathfrak a|_N,\mu)\) along \(N\) on open sets containing \(x\in\operatorname{cosupp}(\mathfrak a|_N,\mu)\). There are an open
\(W\ni x\) and maps \(\varphi,\psi\) as in Lemma~\ref{mc:local-inverse-biholomorphism} such that, with \(\Phi\colon U\to U'\) the resulting biholomorphism, \(x\in U\cap U'\),
\(\Phi(x)=x\), and
\begin{enumerate}
\item \(\Phi^*(\underline{\mathcal H}_N|_{U'})=\underline{\mathcal H}_N|_U\); in particular \(\Phi^{-1}(N\cap U')=N\cap U\);
\item \(\Phi^{-1}(V(u)\cap U')=V(v)\cap U\);
\item \(\Phi\) fixes every point of \(\operatorname{cosupp}(\mathfrak a|_N,\mu)\cap U\), and \(\Phi^{-1}(C\cap U')=C\cap U\) for every \(C\subset\operatorname{cosupp}(\mathfrak a|_N,\mu)\), in
particular for every admissible centre \(C\) of \(\underline{\mathcal H}_N\).
\end{enumerate}
\end{lemma}
\begin{proof}
\leanok
As for Lemma~\ref{mc:glueing-biholomorphism}, with Theorem~\ref{mc:glueing-lemma-along} in place of Theorem~\ref{mc:glueing-lemma}: choose \(W\) open on which \(\varphi,\psi\) are
holomorphic and mutually inverse, \(\varphi^*\mathcal H_N(\mathfrak a)=\mathcal H_N(\mathfrak a)\) stalkwise, \(u\circ\varphi=v\), \(\varphi\) fixes \(\operatorname{cosupp}(\mathfrak a|_N,\mu)\), and
\(\varphi(p)\in N\iff p\in N\). (1) The submanifold of \(\Phi^*(\underline{\mathcal H}_N|_{U'})\) is \(\Phi^{-1}(N\cap U')=N\cap U\) by the last property; both have \(E=\emptyset\) and mark
\(\mu\); the ideals agree as in Lemma~\ref{mc:glueing-biholomorphism}, \(\mathcal H_N(\mathfrak a)\) being coherent (Lemma~\ref{mc:homogenized-along-basic}). (2), (3): as there; an
admissible centre lies in \(\operatorname{cosupp}(\mathcal H_N(\mathfrak a)|_N,\mu)=\operatorname{cosupp}(\mathfrak a|_N,\mu)\) (Lemma~\ref{mc:homogenized-along-basic}).
\end{proof}

\begin{theorem}[Glueing Lemma along \(N\) along multiple test blow-ups, pullback form]
\label{mc:glueing-tower-pullback-along}
\uses{mc:glueing-biholomorphism-along,mc:glueing-tower-pullback,mc:strict-along,hir:tower-transport-open,hir:open-restriction-tower,hir:tower-transport-resolution,hir:open-restriction-marked}
\source{wlodarczyk05-hironaka,wlodarczyk09-analytic-resolution:page-0021}
\lean{Hironaka.HomogenizedGlueingAlongTower.isCoherent_homogenizedMarkedAlong,Hironaka.HomogenizedGlueingAlongTower.isCoherent_restrict,Hironaka.HomogenizedGlueingAlongTower.GlueingDataAlong.pullbackTower,Hironaka.HomogenizedGlueingAlongTower.GlueingDataAlong.isGeoTestTower_restrict_dom,Hironaka.HomogenizedGlueingAlongTower.GlueingDataAlong.isGeoTestTower_pullbackTower,Hironaka.HomogenizedGlueingAlongTower.GlueingDataAlong.geoTransformAlong_pullbackTower_step,Hironaka.HomogenizedGlueingAlongTower.GlueingDataAlong.cosupp_pullbackTower,Hironaka.HomogenizedGlueingAlongTower.GlueingDataAlong.isGeoResolution_pullbackTower,Hironaka.HomogenizedGlueingAlongTower.GlueingDataAlong.strictAlong_pullbackTower_zeroSet,Hironaka.HomogenizedGlueingAlongTower.glueing_lemma_along_tower}
\leanok
With the notation of Lemma~\ref{mc:glueing-biholomorphism-along}, let \(T\) be a multiple test blow-up of \(\underline{\mathcal H}_N\), and \(\varphi^*T=\Phi^*(T|_{U'})\). Then
\(\varphi^*T\) and \(T|_U\) are multiple test blow-ups of \(\underline{\mathcal H}_N|_U\), \(\varphi^*T\) is a resolution if \(T\) is one; from the first blowup on the transforms
along \(\varphi^*T\) are the restrictions of those along \(T\), the cosupports correspond under the last-stage biholomorphism; and the strict transform of
\(V(v)\cap U\) along \(\varphi^*T\) corresponds to that of \(V(u)\) along \(T\).
\end{theorem}
\begin{proof}
\leanok
The proof of Theorem~\ref{mc:glueing-tower-pullback} applies verbatim: the transport and restriction statements it uses
(Lemmas~\ref{hir:tower-transport-open}, \ref{hir:open-restriction-tower} and~\ref{hir:tower-transport-resolution}) hold for marked ideals whose submanifold is any closed
submanifold with coherent ideal, here \(N\) and \(\mathcal H_N(\mathfrak a)\), and Lemma~\ref{mc:glueing-biholomorphism-along}(1) identifies \(\Phi^*(\underline{\mathcal H}_N|_{U'})\) with
\(\underline{\mathcal H}_N|_U\).
\end{proof}
This is the part of Włodarczyk's le:homo'' (4) along \(N\) that follows from the functoriality of multiple test blow-ups under étale maps. The stage-by-stage
lifts of \(\Phi\) (the form of Theorem~\ref{mc:glueing-lift} along \(N\)) also need the inclusions of the derivative ideals along \(N\) under blowups
(Lemma~\ref{hir:along-calculus}(2)).

\subsection{Homogenized and coefficient ideals along a submanifold, along multiple test blow-ups}

In this subsection \(N\) is a closed submanifold, \(\pi\colon M'\to M\) the blowup along an admissible centre \(C\subset N\), \(N'\) the strict transform of \(N\), and
\(\mathfrak l=\pi^*\mathcal I_C\); ideals are coherent. Along \(N\) only the restriction of the ideal of a marked ideal to \(N\) matters, so the statements below hold
up to restricted equivalence (Lemma~\ref{hir:restricted-equivalence}), which has the same cosupports, the same admissible centres and the same transforms.

\begin{lemma}[Calculus of the transform along \(N\)]
\label{mc:transform-along-calculus}
\uses{hir:geo-transform,hir:geo-transform-coherent,hir:restricted-calculus,hir:restricted-equivalence,hir:marked-operations-along,mc:finite-sum,mc:pullback-finite-sum,mc:hom-sum}
\source{wlodarczyk05-hironaka}
\lean{Hironaka.HomogenizedIdealAlongTower.restrictLe_finSumIdeal,Hironaka.HomogenizedIdealAlongTower.pullback_finSumIdeal_of_lt,Hironaka.HomogenizedIdealAlongTower.mul_finSumIdeal,Hironaka.HomogenizedIdealAlongTower.mul_mul_mul_comm',Hironaka.HomogenizedIdealAlongTower.restrictEquiv_geoTransform_of_pullback,Hironaka.HomogenizedIdealAlongTower.restrictLe_geoTransform_of_restrictLe,Hironaka.HomogenizedIdealAlongTower.pullback_homSum_of_lt,Hironaka.HomogenizedIdealAlongTower.restrictLe_homSum}
\leanok
\begin{enumerate}
\item Restricted inclusions \(\operatorname{Res}_W\mathfrak a_i\subset\operatorname{Res}_W\mathfrak a_i'\) pass to finite sums and to \(\sum_iA_iB^i\); pullbacks commute with
finite sums and with \(\sum_iA_iB^i\); \(\mathfrak l\sum_i\mathfrak k_i=\sum_i\mathfrak l\mathfrak k_i\).
\item (The transform is characterised by its defining equation.) If \(\underline{\mathcal I}=(M,N,E,\mathfrak a,d)\) and \(\mathfrak q\) on \(M'\) satisfy
\(\operatorname{Res}_{N'}\pi^*\mathfrak a=\operatorname{Res}_{N'}\mathfrak l^d\mathfrak q\), then the transform \(\underline{\mathcal I}'\) is restricted-equivalent to
\(\underline{\mathcal I}'\) with ideal replaced by \(\mathfrak q\).
\item (Monotonicity.) For two marked ideals on \(N\) with the same mark and \(C\) admissible for both, \(\operatorname{Res}_N\mathfrak a_1\subset\operatorname{Res}_N\mathfrak a_2\)
implies \(\operatorname{Res}_{N'}\mathfrak a_1'\subset\operatorname{Res}_{N'}\mathfrak a_2'\).
\end{enumerate}
\end{lemma}
\begin{proof}
\leanok
(1) Combine the restricted span relations summand by summand; for pullbacks and the distributive law, induct on the number of summands
(Lemma~\ref{mc:pullback-finite-sum}). (2) By Lemma~\ref{hir:marked-operations-along} (the defining equation of the transform along \(N'\)),
\(\operatorname{Res}_{N'}\mathfrak l^d\mathfrak a'=\operatorname{Res}_{N'}\pi^*\mathfrak a=\operatorname{Res}_{N'}\mathfrak l^d\mathfrak q\); the ideal \(\mathfrak l^d\) is invertible along
\(N'\) and cancels (Lemma~\ref{hir:restricted-calculus}). (3) Pull back the inclusion to \(N'\subset\pi^{-1}N\) and cancel \(\mathfrak l^d\) in
\(\mathfrak l^d\mathfrak a_1'\equiv\pi^*\mathfrak a_1\subset\pi^*\mathfrak a_2\equiv\mathfrak l^d\mathfrak a_2'\) on \(N'\).
\end{proof}

\begin{lemma}[Włodarczyk's sum of \(m\) marked ideals along \(N\)]
\label{mc:sum-along}
\uses{mc:transform-along-calculus,mc:marked-finite-sum,mc:marked-finite-sum-tower,hir:order-power-along,hir:marked-operations-along,hir:restricted-equivalence}
\source{wlodarczyk05-hironaka}
\lean{Hironaka.HomogenizedIdealAlongTower.IsSubmanifoldFamily,Hironaka.HomogenizedIdealAlongTower.IsSubmanifoldFamily.isCoherent_msum,Hironaka.HomogenizedIdealAlongTower.IsSubmanifoldFamily.isClosedSubmanifold,Hironaka.HomogenizedIdealAlongTower.cosupp_msum_along,Hironaka.HomogenizedIdealAlongTower.isAdmissibleCentre_msum_along_iff,Hironaka.HomogenizedIdealAlongTower.geoTransform_msum_restrictEquiv,Hironaka.HomogenizedIdealAlongTower.IsSubmanifoldFamily.geoTransform,Hironaka.HomogenizedIdealAlongTower.geoTransformAlong_msum_along}
\leanok
Let \(\underline{\mathcal I}_0,\dots,\underline{\mathcal I}_{m-1}\) (\(m\ge1\)) be marked ideals on one closed submanifold \(N\) with a common \(E\), coherent ideals and marks \(\ge1\), and
\(\underline{\mathcal I}_0+\dots+\underline{\mathcal I}_{m-1}\) their sum (Definition~\ref{mc:marked-finite-sum}). Then its cosupport is the intersection of the cosupports, a
centre is admissible for it iff it is admissible for every summand, and its transform by an admissible blowup is restricted-equivalent to the sum of the transforms. Along
towers: the multiple test blow-ups of the sum are exactly the simultaneous multiple test blow-ups of the \(\underline{\mathcal I}_i\), and along them the transform of the
sum is restricted-equivalent to the sum of the transforms (Lemma~\ref{mc:marked-finite-sum-tower} along \(N\)).
\end{lemma}
\begin{proof}
\leanok
Cosupport: \(\operatorname{cosupp}(\mathfrak a^\nu|_N,\mu\nu)=\operatorname{cosupp}(\mathfrak a|_N,\mu)\) (Lemma~\ref{hir:order-power-along}) with \(\nu=e_i\ge1\). One step: with
\(\mu=\prod_jd_j=d_ie_i\), \(\pi^*\sum_i\mathfrak a_i^{e_i}=\sum_i(\pi^*\mathfrak a_i)^{e_i}\equiv\sum_i(\mathfrak l^{d_i}\mathfrak a_i')^{e_i}=\mathfrak l^\mu\sum_i\mathfrak a_i'^{e_i}\) on
\(N'\), and Lemma~\ref{mc:transform-along-calculus}(2). Towers: induction; at each stage replace the transform of the sum by the restricted-equivalent sum of the transforms,
which has the same admissible centres and the same transform (Lemma~\ref{hir:restricted-equivalence}).
\end{proof}

\begin{lemma}[The transform of \(\sum_iA_iB^i\) along \(N\)]
\label{mc:hom-sum-transform-along}
\uses{mc:transform-along-calculus,mc:hom-sum,mc:hom-sum-transform}
\source{wlodarczyk05-hironaka}
\lean{Hironaka.HomogenizedIdealAlongTower.geoTransform_homSum_restrictEquiv}
\leanok
Let \(\mu\ge1\), \(\underline{\mathcal A}_i\) (\(i<\mu\), mark \(\mu-i\)) and \(\underline{\mathcal B}\) (mark \(1\)) marked ideals on \(N\) with coherent ideals, and \(\underline{\mathcal H}\) a
marked ideal on \(N\) with mark \(\mu\) and ideal \(\sum_iA_iB^i\), \(C\) admissible for all of them. Then \(\underline{\mathcal H}'\) is restricted-equivalent to
\(\underline{\mathcal H}'\) with ideal \(\sum_iA_i'B'^i\) (Lemma~\ref{mc:hom-sum-transform} along \(N\)).
\end{lemma}
\begin{proof}
\leanok
On \(N'\): \(\pi^*\sum_iA_iB^i=\sum_i\pi^*A_i(\pi^*B)^i\equiv\sum_i\mathfrak l^{\mu-i}A_i'(\mathfrak lB')^i=\mathfrak l^\mu\sum_iA_i'B'^i\); apply
Lemma~\ref{mc:transform-along-calculus}(2).
\end{proof}

\begin{lemma}[Inclusions along \(N\) along multiple test blow-ups]
\label{mc:inclusions-tower-along}
\uses{hir:along-transform,hir:giraud-along,mc:transform-along-calculus,mc:homogenized-along,mc:derivative-along-coherent,mc:inclusions-tower,hir:geo-transform-coherent}
\source{wlodarczyk05-hironaka}
\lean{Hironaka.HomogenizedIdealAlongTower.derivMarkedAlong_zero,Hironaka.HomogenizedIdealAlongTower.inclusions_tower_along,Hironaka.HomogenizedIdealAlongTower.mem_cosupp_derivMarkedAlong_tower}
\leanok
Let \(\underline{\mathcal I}=(M,N,E,\mathfrak a,d)\), \(j\le d\), and \(T\) a multiple test blow-up of \(\underline{\mathcal I}\). Then \(T\) is a multiple test blow-up of
\(\mathcal D_N^j(\underline{\mathcal I})\), the transforms have the same submanifold \(N_T\) (a closed submanifold) and the same \(E\), both are coherent, and on \(N_T\)
\[[\mathcal D_N^j(\underline{\mathcal I})]_T\subset\mathcal D_{N_T}^j(\underline{\mathcal I}_T);\]
hence \(\operatorname{cosupp}\underline{\mathcal I}_T\subset\operatorname{cosupp}[\mathcal D_N^i(\underline{\mathcal I})]_T\) for \(i\le d\) (Lemma~\ref{mc:inclusions-tower} along \(N\)).
\end{lemma}
\begin{proof}
\leanok
Induction on \(T\). At a stage, \(C\subset\operatorname{cosupp}\underline{\mathcal I}_T\subset\operatorname{cosupp}(\mathcal D_{N_T}^j\mathfrak a_T|_{N_T},d-j)\) (Lemma~\ref{hir:giraud-along})
\(\subset\operatorname{cosupp}[\mathcal D_N^j(\underline{\mathcal I})]_T\) by the induction hypothesis, so \(C\) is admissible for both. By monotonicity
(Lemma~\ref{mc:transform-along-calculus}(3)) the transform of \([\mathcal D_N^j(\underline{\mathcal I})]_T\) lies in that of \(\mathcal D_{N_T}^j(\underline{\mathcal I}_T)\), which lies in
\(\mathcal D_{N_{T'}}^j(\underline{\mathcal I}_{T'})\) by Lemma~\ref{hir:along-transform}.
\end{proof}

\begin{theorem}[\(\mathcal C_N(\underline{\mathcal I})\simeq\underline{\mathcal I}\)]
\label{mc:coefficient-equivalent-along}
\uses{mc:sum-along,mc:inclusions-tower-along,mc:homogenized-along,mc:coefficient-wequiv,hir:marked-equivalence}
\source{wlodarczyk05-hironaka,wlodarczyk09-analytic-resolution:page-0022}
\lean{Hironaka.HomogenizedIdealAlongTower.isSubmanifoldFamily_derivMarkedAlong,Hironaka.HomogenizedIdealAlongTower.coeffAlong_tower,Hironaka.HomogenizedIdealAlongTower.coeffAlong_wEquiv}
\leanok
Let \(\underline{\mathcal I}=(M,N,E,\mathfrak a,d)\) with \(d\ge1\). A tower is a multiple test blow-up of the coefficient ideal along \(N\), \(\mathcal C_N(\underline{\mathcal I})\)
(Definition~\ref{mc:homogenized-along}), iff it is one of \(\underline{\mathcal I}\), and along it the cosupports agree; i.e.
\(\mathcal C_N(\underline{\mathcal I})\simeq\underline{\mathcal I}\) (Definition~\ref{hir:marked-equivalence}; Lemma~\ref{mc:coefficient-wequiv} along \(N\)).
\end{theorem}
\begin{proof}
\leanok
\(\mathcal C_N(\underline{\mathcal I})\) is the sum of the family \(\mathcal D_N^i(\underline{\mathcal I})\), \(i<d\), on \(N\). By Lemma~\ref{mc:sum-along} its multiple test blow-ups
are the simultaneous ones of the \(\mathcal D_N^i(\underline{\mathcal I})\), which are those of \(\underline{\mathcal I}=\mathcal D_N^0(\underline{\mathcal I})\)
(Lemma~\ref{mc:inclusions-tower-along}), and its cosupport along them is \(\bigcap_i\operatorname{cosupp}[\mathcal D_N^i(\underline{\mathcal I})]_T=\operatorname{cosupp}\underline{\mathcal I}_T\)
(Lemma~\ref{mc:inclusions-tower-along}).
\end{proof}

\begin{lemma}[The homogenized ideal along \(N\) along multiple test blow-ups]
\label{mc:homogenized-tower-along}
\uses{mc:hom-sum-transform-along,mc:inclusions-tower-along,mc:hom-sum-order,mc:homogenized-along,mc:homogenized-along-basic,mc:homogenized-tower,hir:restricted-equivalence}
\source{wlodarczyk05-hironaka}
\lean{Hironaka.HomogenizedIdealAlongTower.homTower,Hironaka.HomogenizedIdealAlongTower.cosupp_homTower,Hironaka.HomogenizedIdealAlongTower.hom_tower_along,Hironaka.HomogenizedIdealAlongTower.isGeoTestTower_of_homAlong}
\leanok
Let \(\underline{\mathcal I}=(M,N,E,\mathfrak a,d)\) with \(d\ge1\), and \(\mathcal H_N(\underline{\mathcal I})=(M,N,E,\mathcal H_N(\mathfrak a),d)\). A tower is a multiple test blow-up of
\(\mathcal H_N(\underline{\mathcal I})\) iff it is one of \(\underline{\mathcal I}\); along it \(\mathcal H_N(\underline{\mathcal I})_T\) is coherent and restricted-equivalent to
\(\underline{\mathcal I}_T\) with ideal \(\sum_{i<d}[\mathcal D_N^i\underline{\mathcal I}]_T\,[T_N(\underline{\mathcal I})]_T^i\), \(T_N=\mathcal D_N^{d-1}\), whose cosupport is
\(\operatorname{cosupp}\underline{\mathcal I}_T\) (Lemma~\ref{mc:homogenized-tower} along \(N\)).
\end{lemma}
\begin{proof}
\leanok
Cosupport: the term \(i=0\) is \(\underline{\mathcal I}_T\), and on \(\operatorname{cosupp}\underline{\mathcal I}_T\) the order of \([\mathcal D_N^i\underline{\mathcal I}]_T\) along \(N_T\) is
\(\ge d-i\) (Lemma~\ref{mc:inclusions-tower-along}), so that of the sum is \(\ge d\) (Lemma~\ref{mc:hom-sum-order}). Induction on \(T\): at each stage replace
\(\mathcal H_N(\underline{\mathcal I})_T\) by the restricted-equivalent expected transform (same admissible centres, same transform, Lemma~\ref{hir:restricted-equivalence}) and apply
Lemma~\ref{mc:hom-sum-transform-along} with \(\underline{\mathcal A}_i=[\mathcal D_N^i\underline{\mathcal I}]_T\), \(\underline{\mathcal B}=[\mathcal D_N^{d-1}\underline{\mathcal I}]_T\), for
which the centre is admissible by Lemma~\ref{mc:inclusions-tower-along}.
\end{proof}

\begin{theorem}[\(\mathcal H_N(\underline{\mathcal I})\simeq\underline{\mathcal I}\)]
\label{mc:homogenized-equivalent-along}
\uses{mc:homogenized-tower-along,mc:coefficient-equivalent-along,mc:homogenized-along-basic,mc:homogenized-equivalent,hir:marked-equivalence,hir:marked-equivalence-basic}
\source{wlodarczyk05-hironaka,wlodarczyk09-analytic-resolution:page-0020,wlodarczyk09-analytic-resolution:page-0022}
\lean{Hironaka.HomogenizedIdealAlongTower.homAlong_wEquiv,Hironaka.HomogenizedIdealAlongTower.coeff_homAlong_wEquiv}
\leanok
Let \(\underline{\mathcal I}=(M,N,E,\mathfrak a,d)\) with \(N\) a closed submanifold, \(\mathfrak a\) coherent and \(d\ge1\). Then
\(\mathcal H_N(\underline{\mathcal I})\simeq\underline{\mathcal I}\) and \(\mathcal C_N(\mathcal H_N(\underline{\mathcal I}))\simeq\underline{\mathcal I}\) (Lemma~\ref{mc:homogenized-equivalent} along \(N\)).
\end{theorem}
\begin{proof}
\leanok
The first is Lemma~\ref{mc:homogenized-tower-along}; the second is Theorem~\ref{mc:coefficient-equivalent-along} for \(\mathcal H_N(\underline{\mathcal I})\) (coherent,
Lemma~\ref{mc:homogenized-along-basic}) followed by the first, \(\simeq\) being transitive (Lemma~\ref{hir:marked-equivalence-basic}).
\end{proof}
With these, the statements of Lemma~\ref{hir:homogenized-coefficient} along multiple test blow-ups hold along \(N\) (the transform formula and the two equivalences;
maximal order at one stage is Lemma~\ref{mc:homogenized-along-maximal-order}). The restricted coefficient ideal along towers whose centres lie in the strict
transforms of \(S\) (Lemmas~\ref{mc:coefficient-restriction-tower} and~\ref{mc:coefficient-restriction-full} along \(N\)) is Lemma~\ref{mc:coefficient-restriction-tower-along}
and Theorem~\ref{mc:coefficient-restriction-full-along} below. Still missing in Lemma~\ref{hir:along-calculus}(4) along towers: the stage-by-stage lifts of the
glueing biholomorphism (Theorem~\ref{mc:glueing-lift} along \(N\)); their inputs along \(N\) are Lemma~\ref{mc:inclusions-tower-along},
Lemma~\ref{mc:homogenized-tower-along} and Theorem~\ref{mc:glueing-tower-pullback-along}.

\subsection{The restricted coefficient ideal along a submanifold, along multiple test blow-ups}

\begin{lemma}[Restriction along \(N\) along multiple test blow-ups]
\label{mc:restrict-tower-along}
\uses{mc:restrict-step,mc:restrict-tower,mc:restrict-marked,mc:strict-along,mc:transform-along-calculus,mc:coefficient-restriction-along,hir:restricted-equivalence,hir:geo-transform-coherent}
\source{wlodarczyk05-hironaka}
\lean{Hironaka.CoefficientIdealAlongRestrictionTower.geoStrictTransform_mono',Hironaka.CoefficientIdealAlongRestrictionTower.geoTransform_withN_restrictEquiv_along,Hironaka.CoefficientIdealAlongRestrictionTower.restrict_tower_along}
\leanok
Let \(\underline{\mathcal I}=(M,N,E,\mathfrak a,d)\) with \(N\) a closed submanifold and \(\mathfrak a\) coherent, and \(S\subset N\) a closed submanifold.
\begin{enumerate}
\item For a centre \(C\) admissible for \(\underline{\mathcal I}\) and for \(\underline{\mathcal I}|_S\), the transform of \(\underline{\mathcal I}|_S\) is restricted-equivalent to
\((\underline{\mathcal I}')|_{S'}\), and \(S'\subset N'\).
\item Every multiple test blow-up \(T\) of \(\underline{\mathcal I}\) whose centres lie in the strict transforms of \(S\) is one of \(\underline{\mathcal I}|_S\), and along it
\([\underline{\mathcal I}|_S]_T\) is restricted-equivalent to \((\underline{\mathcal I}_T)|_{S_T}\), \(S_T\subset N_T\) closed submanifolds
(Lemmas~\ref{mc:restrict-step} and~\ref{mc:restrict-tower} along \(N\)).
\end{enumerate}
\end{lemma}
\begin{proof}
\leanok
(1) On \(S'\subset N'\) both transforms satisfy the defining equation \(\operatorname{Res}_{S'}\pi^*\mathfrak a=\operatorname{Res}_{S'}(\pi^*\mathcal I_C)^d\mathfrak a'\)
(Lemma~\ref{mc:transform-along-calculus}(2), the equation on \(N'\) restricted to \(S'\)). (2) Induction on \(T\); the centre is admissible for \((\underline{\mathcal I}_T)|_{S_T}\)
because orders along \(N_T\) do not drop on \(S_T\subset N_T\) (Lemma~\ref{mc:coefficient-restriction-along}(1)), hence for the restricted-equivalent
\([\underline{\mathcal I}|_S]_T\), whose transform equals that of \((\underline{\mathcal I}_T)|_{S_T}\) (Lemma~\ref{hir:restricted-equivalence}); apply (1).
\end{proof}

\begin{lemma}[The restricted coefficient ideal along \(N\) along multiple test blow-ups]
\label{mc:coefficient-restriction-tower-along}
\uses{mc:restrict-tower-along,mc:coefficient-equivalent-along,mc:coefficient-restriction-along,mc:coefficient-restriction-tower,mc:homogenized-along-basic}
\source{wlodarczyk05-hironaka,wlodarczyk09-analytic-resolution:page-0023}
\lean{Hironaka.CoefficientIdealAlongRestrictionTower.coeffAlong_restrict_tower}
\leanok
Let \(\underline{\mathcal I}=(M,N,E,\mathfrak a,d)\) with \(N\) a closed submanifold, \(\mathfrak a\) coherent, \(d\ge1\), \(S\subset N\) a closed submanifold, and \(T\) a multiple
test blow-up of \(\underline{\mathcal I}\) whose centres lie in the strict transforms of \(S\). Then \(T\) is a multiple test blow-up of \(\mathcal C_N(\underline{\mathcal I})|_S\),
\([\mathcal C_N(\underline{\mathcal I})|_S]_T\) is restricted-equivalent to \(\mathcal C_N(\underline{\mathcal I})_T|_{S_T}\), and
\(\operatorname{cosupp}\underline{\mathcal I}_T\cap S_T\subset\operatorname{cosupp}[\mathcal C_N(\underline{\mathcal I})|_S]_T\) (Lemma~\ref{mc:coefficient-restriction-tower} along \(N\):
Włodarczyk's ``le:S'' (1) and the inclusion \(\subset\) of (2)).
\end{lemma}
\begin{proof}
\leanok
\(T\) is a multiple test blow-up of \(\mathcal C_N(\underline{\mathcal I})\) with \(\operatorname{cosupp}\mathcal C_N(\underline{\mathcal I})_T=\operatorname{cosupp}\underline{\mathcal I}_T\)
(Theorem~\ref{mc:coefficient-equivalent-along}); apply Lemma~\ref{mc:restrict-tower-along} to \(\mathcal C_N(\underline{\mathcal I})\) (coherent,
Lemma~\ref{mc:homogenized-along-basic}) and Lemma~\ref{mc:coefficient-restriction-along}(1) on \(S_T\subset N_T\).
\end{proof}
The reverse inclusion \(\operatorname{cosupp}[\mathcal C_N(\underline{\mathcal I})|_S]_T\subset\operatorname{cosupp}\underline{\mathcal I}_T\) for \(T\) nonempty (``le:S'' (2) in full,
Lemma~\ref{mc:coefficient-restriction-full} along \(N\)) is Theorem~\ref{mc:coefficient-restriction-full-along} below: for \(N=M\) it rests on the normal-coefficient invariant of
Lemma~\ref{mc:normal-coefficient-invariant}, which along \(N\) is carried in flat charts of the pair \(S_T\subset N_T\).

\subsection{The restricted coefficient ideal along a submanifold: the full statement}

This subsection proves the remaining inclusion of Włodarczyk's Lemma ``le:S'' along a submanifold \(N\), following Lemmas~\ref{mc:normal-coefficient-invariant}
to~\ref{mc:coefficient-restriction-full}. The derivatives playing the role of the coefficients are those in the directions normal to \(S_T\) and tangent to \(N_T\), read in
flat charts of the pair \(S_T\subset N_T\). Throughout, \(P_k\colon\mathbf C^n\to\mathbf C^n\) is the linear projection killing the coordinates of index \(\ge k\); its image is
the plane \(\mathbf C^k\times0\).

\begin{lemma}[Derivatives and the plane retraction]
\label{mc:plane-retraction-derivatives}
\uses{mc:normal-derivative-monomial}
\lean{Hironaka.NormalDerivativeAlong.plane_mono,Hironaka.NormalDerivativeAlong.planeProj_add_smul_bv,Hironaka.NormalDerivativeAlong.planeProj_add_smul_bv_of_le,Hironaka.NormalDerivativeAlong.iterDeriv_comp_planeProj,Hironaka.NormalDerivativeAlong.iterDeriv_comp_planeProj_eq_zero,Hironaka.NormalDerivativeAlong.iterDeriv_eventuallyEq_on_plane}
\leanok
Let \(F\) be holomorphic on an open \(V\subset\mathbf C^n\) and \(L\) a list of directions \(e_l\). If every \(l<k\), then \(\partial^L(F\circ P_k)=(\partial^LF)\circ P_k\)
on \(P_k^{-1}V\); if some \(l\ge k\), then \(\partial^L(F\circ P_k)=0\) on \(P_k^{-1}V\). Consequently, if \(F\) and \(G\) are holomorphic on \(V\ni0\) and agree on
\(\mathbf C^k\times0\) near \(0\), then \(\partial^LF=\partial^LG\) on \(\mathbf C^k\times0\) near \(0\) for every list \(L\) of directions \(<k\).
\end{lemma}
\begin{proof}
\leanok
\(P_k(v+te_l)=P_kv+te_l\) for \(l<k\) and \(P_k(v+te_l)=P_kv\) for \(l\ge k\). The first gives the first formula (Lemma~\ref{mc:normal-derivative-monomial}, maps
translating along the directions); for the second induct on \(L\): if the tail of \(L\) contains a direction \(\ge k\) its derivative vanishes near the point, and
otherwise it is \((\partial^{L'}F)\circ P_k\), constant along the head direction \(e_l\), \(l\ge k\) (derivatives along lines). For the last statement,
\(F\circ P_k=G\circ P_k\) near \(0\), and \(P_k\) is the identity on the plane.
\end{proof}

\begin{lemma}[The order along \(N\) from the derivatives normal to \(S\)]
\label{mc:normal-taylor-along}
\uses{mc:plane-retraction-derivatives,mc:normal-taylor-order,mc:flat-chart}
\source{wlodarczyk05-hironaka}
\lean{Hironaka.NormalDerivativeAlong.FlatChart.vanishesToOrderOn_of_nd}
\leanok
Let \(c\) be a flat chart for \(S\) at \(p\) (plane \(\mathbf C^{k_S}\times0\)) in which the points of \(N\) near \(p\) lie in \(\mathbf C^{k_N}\times0\), \(k_S\le k_N\), and \(g\)
holomorphic at \(p\). If for every list \(L\) of directions in \([k_S,k_N)\) with \(|L|<\mu\) the derivative \(\partial^Lg\) in \(c\) vanishes to order \(\mu-|L|\) along \(S\)
at \(p\), then \(g|_N\in\mathfrak m_{N,p}^\mu\).
\end{lemma}
\begin{proof}
\leanok
Put \(h=g\circ\psi^{-1}\circ P_{k_N}\circ\psi\) (\(\psi\) the coordinates of \(c\)). By Lemma~\ref{mc:plane-retraction-derivatives} its derivatives normal to
\(\mathbf C^{k_S}\times0\) are either those of \(g\) (directions in \([k_S,k_N)\); on \(S\) the coordinates lie in \(\mathbf C^{k_S}\times0\subset\mathbf C^{k_N}\times0\), where
\(P_{k_N}\) is the identity) or zero (some direction \(\ge k_N\)). So \(h\in\mathfrak m_p^\mu\) (Lemma~\ref{mc:normal-taylor-order}), and \(h=g\) on \(N\) near \(p\).
\end{proof}

\begin{lemma}[Derivatives tangent to \(N\) of an ideal]
\label{mc:normal-derivative-along-ideal}
\uses{mc:flat-coordinate-iterated,mc:derivative-along,mc:flat-chart}
\source{wlodarczyk05-hironaka}
\lean{Hironaka.NormalDerivativeAlong.iterDeriv_comp_flatMap_mem_list,Hironaka.NormalDerivativeAlong.FlatChart.nd_mem_iterDerivIdealAlong}
\leanok
In a flat chart at \(p\) in which \(N\) is \(\mathbf C^{k_N}\times0\), for a coherent \(\mathfrak a\), \(g\in\mathfrak a_p\) and a list \(L\) of directions \(<k_N\),
\(\partial^Lg\in\mathcal D_N^{|L|}(\mathfrak a)_p\).
\end{lemma}
\begin{proof}
\leanok
The coordinate field \(\partial/\partial\psi_l\), \(l<k_N\), is tangent to \(N\) and \(\partial/\partial\psi_l(F\circ\psi)=(\partial_lF)\circ\psi\)
(Lemma~\ref{mc:flat-coordinate-iterated}); induct on \(L\).
\end{proof}

\begin{lemma}[Nested flat charts of a pair]
\label{mc:nested-pair-chart}
\uses{mc:nested-product-chart,bcalc:nested-flattening}
\source{wlodarczyk05-hironaka}
\lean{Hironaka.NormalDerivativeAlong.prodTo_apply_of_le,Hironaka.NormalDerivativeAlong.prodTo_mem_plane_iff,Hironaka.NormalDerivativeAlong.FlatChart.exists_nested_pair}
\leanok
Let \(c\) be a flat chart for \(S\) at \(p\) (plane \(\mathbf C^{k_S}\times0\)) in which \(N\) is the plane \(\mathbf C^{k_N}\times0\), \(k_S\le k_N\), and \(C\subset S\) regular of
dimension \(k_C\) at \(p\). There is a flat chart \(c'\) for \(S\) at \(p\) in which \(C\) is \(\mathbf C^{k_C}\times0\) and \(N\) is still \(\mathbf C^{k_N}\times0\), with the same
derivatives normal to \(S\).
\end{lemma}
\begin{proof}
\leanok
Take the product-form chart of Lemma~\ref{mc:nested-product-chart}: \(T(u)=\Phi_m(P_{k_S}u)+(u-P_{k_S}u)\) with \(\Phi_m(\mathbf C^{k_S}\times0)\subset\mathbf C^{k_S}\times0\).
It does not change the coordinates of index \(\ge k_S\), hence it preserves \(\mathbf C^{k_N}\times0\).
\end{proof}

\begin{lemma}[The strict transforms of a pair in one monomial chart]
\label{mc:strict-pair-monomial}
\uses{bcalc:strict-submanifold,bcalc:monomial-lift}
\source{wlodarczyk05-hironaka}
\lean{Hironaka.NormalDerivativeAlong.mem_geoStrictTransform_iff_of_monomial,Hironaka.NormalDerivativeAlong.exists_local_geoStrictTransform_pair,Hironaka.NormalDerivativeAlong.FlatChart.exists_of_map'}
\leanok
Let \(\pi\) be the blowup along \(C\) and \(\Phi\) coordinates at \(\pi(z)\in C\) in which \(C=\mathbf C^{k_C}\times0\), \(S=\mathbf C^{k_S}\times0\) and \(N=\mathbf C^{k_N}\times0\).
There are \(j\ge k_C\) and the local inverse \(W\) of the \(j\)-th monomial chart on an open \(O\ni z\) such that on \(O\): \(u\in S'\) iff \(j<k_S\) and
\(W(u)\in\mathbf C^{k_S}\times0\), and \(u\in N'\) iff \(j<k_N\) and \(W(u)\in\mathbf C^{k_N}\times0\). Moreover a holomorphic map with injective differential at \(z\), in
which a set is a plane on an open set \(O\), is a flat chart for that set on a domain contained in \(O\).
\end{lemma}
\begin{proof}
\leanok
The monomial chart and \(j\) depend only on \(\Phi\) and \(C\); the description of the strict transform of a plane through the centre is that of
Lemma~\ref{bcalc:strict-submanifold}, applied once to \(S\) and once to \(N\). The last statement is the holomorphic inverse function theorem.
\end{proof}

\begin{lemma}[Derivatives through a monomial chart, along \(N'\)]
\label{mc:normal-derivative-monomial-along}
\uses{mc:normal-derivative-monomial,mc:plane-retraction-derivatives}
\source{wlodarczyk05-hironaka}
\lean{Hironaka.NormalDerivativeAlong.nd_monomial_identity_along}
\leanok
Let \(c''\) be a flat chart at \(\pi(z)\), \(W\) coordinates near \(z\) with \(c''\circ\pi=\sigma_j\circ W\) on an open \(O\ni z\), and \(c'\) a flat chart at \(z\) with
coordinates \(W-W(z)\) near \(z\) in which \(N'\) is \(\mathbf C^{k_N}\times0\). If \(g\circ\pi=W_j^dh'\) on \(N'\) near \(z\), then
\(W_j^d\,\partial^Lh'=W_j^{|L|}(\partial^Lg)\circ\pi\) on \(N'\) near \(z\) for every list \(L\) of directions in \([k_C,k_N)\setminus\{j\}\).
\end{lemma}
\begin{proof}
\leanok
In the coordinates of \(c'\) the functions \(u\mapsto(u+W(z))_j^d\,h'(c'^{-1}u)\) and \(u\mapsto g(c''^{-1}\sigma_j(u+W(z)))\) agree on \(\mathbf C^{k_N}\times0\) near
\(0\), so their derivatives along \(L\) agree there (Lemma~\ref{mc:plane-retraction-derivatives}). Compute both with Lemma~\ref{mc:normal-derivative-monomial}: the factor
\((u+W(z))_j^d\) is constant along the directions of \(L\), and \(\partial^L(G\circ\sigma_j)=w_j^{|L|}(\partial^LG)\circ\sigma_j\).
\end{proof}

\begin{lemma}[The transform along \(N'\) and the pullback]
\label{mc:transform-along-pullback}
\uses{hir:geo-transform,hir:hadamard-bridge-along,bcalc:function-transform,bcalc:isomorphism-off-centre,ana:restricted-colon}
\source{wlodarczyk05-hironaka}
\lean{Hironaka.CoefficientIdealAlongRestrictionFull.mem_geoTransform_of_mem_pullback,Hironaka.CoefficientIdealAlongRestrictionFull.isInSpanOnAt_pullback_of_mem_geoTransform_of_notMem,Hironaka.CoefficientIdealAlongRestrictionFull.isInSpanOnAt_pow_mul_of_mem_geoTransform,Hironaka.CoefficientIdealAlongRestrictionFull.exists_mem_geoTransform_eq_pow_mul,Hironaka.CoefficientIdealAlongRestrictionFull.isInSpanOnAt_comp}
\leanok
Let \(\underline{\mathcal I}=(M,N,E,\mathfrak a,d)\), \(\pi\) the blowup along \(C\), and \(\mathfrak a'\) the transformed ideal, the colon along \(N'\) of \(\pi^*\mathfrak a\) by
\((\pi^*\mathcal I_C)^d\) (Definition~\ref{hir:geo-transform}).
\begin{enumerate}
\item \(\pi^*\mathfrak a\subset\mathfrak a'\).
\item If \(\pi(z)\notin C\), every germ of \(\mathfrak a'_z\) is on \(N'\) near \(z\) a combination of germs of \(\pi^*\mathfrak a\).
\item For a local equation \(l\) of the exceptional divisor near \(z\), \(l^d\mathfrak a'_z\subset\pi^*\mathfrak a\) on \(N'\) near \(z\).
\item If \(N\) is a closed submanifold, \(\mathfrak a\) coherent, \(C\) admissible and \(z\in N'\), then every \(g\in\mathfrak a_{\pi(z)}\) has \(g\circ\pi=l^dh'\) on \(N'\)
near \(z\) for some \(h'\in\mathfrak a'_z\).
\end{enumerate}
Restricted span relations pull back along holomorphic maps.
\end{lemma}
\begin{proof}
\leanok
(1) An ideal lies in its colon. (2) Near \(z\) the ideal of \(\pi^{-1}C\) contains \(1\). (3) \(l^d\in(\pi^*\mathcal I_C)^d\). (4) By the Hadamard bridge along \(N\)
(Lemma~\ref{hir:hadamard-bridge-along}) \(g=\sum a_is_i\) on \(N\) near \(\pi(z)\) with \(s_i\in\mathcal I_C^d\), and \(s_i\circ\pi=l^ds'_i\) (Lemma~\ref{bcalc:function-transform});
\(h'=\sum(a_i\circ\pi)s'_i\) satisfies \(g\circ\pi=l^dh'\) on \(N'\subset\pi^{-1}N\). Every \(t\in(\pi^*\mathcal I_C)^d_z\) is \(bl^d\), so \(th'=b\,(g\circ\pi)\) on \(N'\):
\(h'\) lies in the colon.
\end{proof}

\begin{lemma}[The coefficient invariant along \(N\)]
\label{mc:normal-coefficient-invariant-along}
\uses{mc:normal-coefficient-invariant,mc:normal-taylor-along,mc:normal-derivative-along-ideal,mc:nested-product-chart,mc:derivative-along,mc:restrict-marked}
\source{wlodarczyk05-hironaka}
\lean{Hironaka.CoefficientIdealAlongRestrictionFull.GoodAtAlong,Hironaka.CoefficientIdealAlongRestrictionFull.GoodAtAlong.orderGE,Hironaka.CoefficientIdealAlongRestrictionFull.goodAtAlong_refl}
\leanok
Let \(\underline{\mathcal I}=(X,N,E,\mathfrak a,\mu)\) with \(N\) a closed submanifold and \(\mathfrak a\) coherent, \(S\subset N\) a closed submanifold and \(T\) a tower. A point
\(p\in S_T\) is \emph{good} if there are a flat chart for \(S_T\) at \(p\) (plane \(\mathbf C^{k_S}\times0\)) in which \(N_T\) is the plane \(\mathbf C^{k_N}\times0\), and germs of
\(\mathfrak a_T\) at \(p\) generating \(\mathfrak a_T\) along \(N_T\) at \(p\), whose derivatives \(\partial^L\) along lists \(L\) of directions in \([k_S,k_N)\) with \(|L|<\mu\) restrict
on \(S_T\) near \(p\) into \([\mathcal D_N^{|L|}\underline{\mathcal I}]_T\). At a good point where every \([\mathcal D_N^i\underline{\mathcal I}]_T\) has order \(\ge\mu-i\) along \(S_T\),
\(\mathfrak a_T\) has order \(\ge\mu\) along \(N_T\). Every point of \(S\) is good for the empty tower.
\end{lemma}
\begin{proof}
\leanok
The first statement is Lemma~\ref{mc:normal-taylor-along} for the generators; combinations along \(N_T\) keep the order. For the empty tower take a flat chart of \(N\)
in which \(S\) is a plane (Lemma~\ref{mc:nested-product-chart}), all germs of \(\mathfrak a\) as generators, and Lemma~\ref{mc:normal-derivative-along-ideal}.
\end{proof}

\begin{lemma}[The coefficient invariant along \(N\) along multiple test blow-ups]
\label{mc:normal-coefficient-step-along}
\uses{mc:normal-coefficient-invariant-along,mc:normal-coefficient-step,mc:nested-pair-chart,mc:strict-pair-monomial,mc:normal-derivative-monomial-along,mc:transform-along-pullback,mc:monomial-local-equation,mc:restrict-tower-along,mc:inclusions-tower-along,bcalc:isomorphism-off-centre,ana:restricted-calculus}
\source{wlodarczyk05-hironaka}
\lean{Hironaka.CoefficientIdealAlongRestrictionFull.goodAtAlong_step_of_notMem,Hironaka.CoefficientIdealAlongRestrictionFull.goodAtAlong_step_of_mem,Hironaka.CoefficientIdealAlongRestrictionFull.goodAtAlong_tower}
\leanok
Along a multiple test blow-up of \(\underline{\mathcal I}\) with centres in the strict transforms of \(S\), every point of \(S_T\) is good.
\end{lemma}
\begin{proof}
\leanok
Induct on the tower; let \(\pi\) blow up \(C\subset S_T\) and \(z\in S_{T'}\). If \(\pi(z)\notin C\), \(\pi\) is a local biholomorphism near \(z\) and \(S_{T'}\), \(N_{T'}\) are the
preimages of \(S_T\), \(N_T\) there: pull back the chart and the generators (Lemma~\ref{mc:transform-along-pullback}(1)--(2)). If \(\pi(z)\in C\), refine the chart at \(\pi(z)\)
by Lemma~\ref{mc:nested-pair-chart} and take the monomial chart \(W\) of Lemma~\ref{mc:strict-pair-monomial}; \(z\in S'\) forces \(j<k_S\), so the directions in
\([k_S,k_N)\) are \(\ge k_C\) and \(\ne j\), and \(W_j\) is a local equation of the exceptional divisor (Lemma~\ref{mc:monomial-local-equation}). The new chart is \(W-W(z)\); the new
generators are the \(h'\in\mathfrak a_{T'}\) with \(g\circ\pi=W_j^\mu h'\) on \(N_{T'}\) (Lemma~\ref{mc:transform-along-pullback}(4)). They generate \(\mathfrak a_{T'}\) along
\(N_{T'}\): for \(f\in\mathfrak a_{T'}\), \(W_j^\mu f\) is on \(N_{T'}\) a combination of germs \(a\circ\pi\), \(a\in\mathfrak a_T\) (Lemma~\ref{mc:transform-along-pullback}(3)),
hence of the \(g\circ\pi=W_j^\mu h'\), and \(W_j^\mu\) cancels on \(N_{T'}\), the closure of \(\pi^{-1}(N_T\setminus C)\), where \(W_j\ne0\) (Lemma~\ref{ana:restricted-calculus}). By
Lemma~\ref{mc:normal-derivative-monomial-along}, on \(S_{T'}\subset N_{T'}\): \(W_j^\mu\partial^Lh'=W_j^{|L|}(\partial^Lg)\circ\pi=W_j^{|L|}\sum(a_i\circ\pi)(s_i\circ\pi)\) with
\(s_i\in[\mathcal D_N^{|L|}\underline{\mathcal I}]_T\); the centre is admissible for \([\mathcal D_N^{|L|}\underline{\mathcal I}]_T\) (Lemma~\ref{mc:inclusions-tower-along}), so
\(s_i\circ\pi=W_j^{\mu-|L|}s'_i\) on \(N_{T'}\) with \(s'_i\in[\mathcal D_N^{|L|}\underline{\mathcal I}]_{T'}\) (Lemma~\ref{mc:transform-along-pullback}(4)); cancel \(W_j^\mu\) on \(S_{T'}\).
\end{proof}

\begin{theorem}[The restricted coefficient ideal along \(N\) along multiple test blow-ups: full statement]
\label{mc:coefficient-restriction-full-along}
\uses{mc:normal-coefficient-step-along,mc:coefficient-restriction-tower-along,mc:restrict-tower-along,mc:sum-along,mc:coefficient-equivalent-along,mc:coefficient-restriction-along,mc:finite-sum,hir:order-power-along,hir:restricted-equivalence-basic}
\source{wlodarczyk05-hironaka,wlodarczyk09-analytic-resolution:page-0023}
\lean{Hironaka.CoefficientIdealAlongRestrictionFull.cosupp_coeffAlong_withN_tower_subset,Hironaka.CoefficientIdealAlongRestrictionFull.cosupp_coeffAlong_withN_tower,Hironaka.CoefficientIdealAlongRestrictionFull.coeffAlong_restrict_tower_iff}
\leanok
Let \(\underline{\mathcal I}=(X,N,E,\mathfrak a,\mu)\) with \(N\) a closed submanifold, \(\mathfrak a\) coherent and \(\mu\ge1\), and \(S\subset N\) a closed submanifold. A tower is a
multiple test blow-up of \(\mathcal C_N(\underline{\mathcal I})|_S\) if and only if it is a multiple test blow-up of \(\underline{\mathcal I}\) with centres in the strict
transforms of \(S\), and along it \(\operatorname{cosupp}[\mathcal C_N(\underline{\mathcal I})|_S]_T=\operatorname{cosupp}\underline{\mathcal I}_T\cap S_T\).
\end{theorem}
\begin{proof}
\leanok
The inclusion \(\supset\) is Lemma~\ref{mc:coefficient-restriction-tower-along}. Conversely, at \(x\) in the left side, \([\mathcal C_N(\underline{\mathcal I})|_S]_T\) is
restricted-equivalent to \(\mathcal C_N(\underline{\mathcal I})_T|_{S_T}\), and on \(N_T\supset S_T\) the transform \(\mathcal C_N(\underline{\mathcal I})_T\) is restricted-equivalent
to \(\sum_i[\mathcal D_N^i\underline{\mathcal I}]_T\) (Lemma~\ref{mc:sum-along}); along \(S_T\) the order of the sum is the minimum of the scaled orders
(Lemmas~\ref{mc:finite-sum} and~\ref{hir:order-power-along}), so \([\mathcal D_N^i\underline{\mathcal I}]_T\) has order \(\ge\mu-i\) along \(S_T\) at \(x\), and the good point \(x\)
(Lemma~\ref{mc:normal-coefficient-step-along}) lies in \(\operatorname{cosupp}\underline{\mathcal I}_T\). For the empty tower the identity is
Lemma~\ref{mc:coefficient-restriction-along}; the test blow-ups correspond by induction on the tower, a centre inside \(S_T\) being admissible for one side if and only if
for the other (the cosupports correspond and \(E\) is the same).
\end{proof}
This is Włodarczyk's Lemma ``le:S'' (2)--(3) along \(N\), for the restriction keeping \(E\).

\begin{corollary}[The restricted coefficient ideal along \(N\) with \(E_S\)]
\label{mc:coefficient-restriction-drop-along}
\uses{mc:coefficient-restriction-full-along,hir:drop-containing-members,hir:wellformed-tower}
\source{wlodarczyk05-hironaka}
\lean{Hironaka.CoefficientIdealAlongRestrictionDrop.coeffAlong_restrict_withE_iff}
\leanok
Let \(\underline{\mathcal I}=(X,N,E,\mathfrak a,\mu)\) be well formed with \(\mu\ge1\), \(S\subset N\) a closed submanifold, and \(E_S\) a list of members of \(E\) such that every
member of \(E\) missing from \(E_S\) contains \(S\). A tower is a multiple test blow-up of \((X,S,E_S,\mathcal C_N(\mathfrak a),\mu_{\mathcal C})\) if and only if it is a
multiple test blow-up of \(\underline{\mathcal I}\) with centres in the strict transforms of \(S\), and along it the cosupport is
\(\operatorname{cosupp}\underline{\mathcal I}_T\cap S_T\).
\end{corollary}
\begin{proof}
\leanok
With \(E\) in place of \(E_S\) this is Theorem~\ref{mc:coefficient-restriction-full-along}. Along towers with centres in the strict transforms of \(S\), deleting
members that contain \(S\) changes neither the multiple test blow-ups nor the cosupports (Lemma~\ref{hir:drop-containing-members}); the simple normal crossings of
the transformed divisor that this uses hold along the multiple test blow-ups of the well-formed \(\underline{\mathcal I}\) (Lemma~\ref{hir:wellformed-tower}).
\end{proof}
This is Lemma~\ref{hir:coefficient-restriction} inside \(N\), as Step 1a of Proposition~\ref{hir:maximal-order-step} uses it along \(N\).

\subsection{The glueing lemma along a submanifold: the lifts along multiple test blow-ups}

In this subsection \(\underline{\mathcal H}=(M,N,\emptyset,\mathcal H_N(\mathfrak a),\mu)\) is the homogenized marked ideal along a closed submanifold \(N\) of
Lemma~\ref{mc:glueing-biholomorphism-along}, \(\mathfrak a\) coherent, \(\mu\ge1\), and \(\mathfrak J=\mathcal D_N^{\mu-1}(\mathcal H_N(\mathfrak a))\), an ideal of \(\mathcal O_M\).

\begin{lemma}[Glueing data along \(N\) acting trivially]
\label{mc:glueing-trivial-along}
\uses{mc:glueing-trivial,mc:glueing-lemma-along,mc:glueing-biholomorphism-along,mc:coordinate-pair-along,mc:coordinate-change-trivial,mc:homogenized-along-basic,mc:derivative-along}
\source{wlodarczyk05-hironaka,wlodarczyk09-analytic-resolution:page-0022}
\lean{Hironaka.HomogenizedGlueingLiftAlong.iterDerivIdealAlong_mono,Hironaka.HomogenizedGlueingLiftAlong.glueIdealAlong,Hironaka.HomogenizedGlueingLiftAlong.isCoherent_glueIdealAlong,Hironaka.HomogenizedGlueingLiftAlong.tangentIdealAlong_le_glueIdealAlong,Hironaka.HomogenizedGlueingLiftAlong.exists_glueingDataAlong_trivial}
\leanok
\(\mathfrak J\) is coherent and contains \(T_N(\mathfrak a)\). In Lemma~\ref{mc:glueing-biholomorphism-along} the open set \(W\) can be chosen so that, moreover, \(\varphi\) and
\(\psi\) act trivially modulo \(\mathfrak J\) at every point of \(V(\mathfrak J)\cap W\).
\end{lemma}
\begin{proof}
\leanok
\(\mathcal D_N^{\mu-1}\) is monotone and \(\mathfrak a\subset\mathcal H_N(\mathfrak a)\) (Lemma~\ref{mc:homogenized-along-basic}). The maps of the glueing lemma along \(N\)
(Lemma~\ref{mc:glueing-lemma-along}) are coordinate changes \(p\mapsto\psi^{-1}(\psi(p)+h(p)e_i)\) of flat charts (Lemma~\ref{mc:coordinate-pair-along}) with
\(h=v-u\) and \(u-v\), germs of \(T_N(\mathfrak a)\subset\mathfrak J\) at \(x\); apply Lemma~\ref{mc:coordinate-change-trivial} and shrink \(W\), as in
Lemma~\ref{mc:glueing-trivial}.
\end{proof}

\begin{lemma}[The glue ideal along multiple test blow-ups]
\label{mc:glue-tower-along}
\uses{mc:inclusions-tower-along,mc:restrict-tower,mc:glueing-trivial-along,hir:test-tower,hir:geo-transform-ambient}
\source{wlodarczyk05-hironaka}
\lean{Hironaka.HomogenizedGlueingLiftAlong.glueMarkedAlong,Hironaka.HomogenizedGlueingLiftAlong.withN_glueMarkedAlong,Hironaka.HomogenizedGlueingLiftAlong.cosupp_subset_zeroSet_glue,Hironaka.HomogenizedGlueingLiftAlong.glue_tower}
\leanok
Let \(\underline{\mathcal J}=(M,M,\emptyset,\mathfrak J,1)\), so that \(\underline{\mathcal J}|_N=\mathcal D_N^{\mu-1}(\underline{\mathcal H})\). Every multiple test blow-up \(T\) of
\(\underline{\mathcal H}\) is a multiple test blow-up of \(\underline{\mathcal J}\) with centres in the strict transforms of \(N\); along it
\(\operatorname{cosupp}\underline{\mathcal H}_T\subset V(\mathfrak J_T)\), and \(\mathfrak J_T\) is coherent (\(\underline{\mathcal J}_T\) has \(N\) the whole space, so \(\mathfrak J_T\) is
the controlled transform of \(\mathfrak J\) with mark \(1\)).
\end{lemma}
\begin{proof}
\leanok
Induction on \(T\). If \(T\) is a multiple test blow-up of \(\underline{\mathcal J}\) with centres in the strict transforms of \(N\), then
\([\underline{\mathcal J}|_N]_T\) is restricted-equivalent to \(\underline{\mathcal J}_T|_{N_T}\) (Lemma~\ref{mc:restrict-tower}), and
\(\operatorname{cosupp}\underline{\mathcal H}_T\subset\operatorname{cosupp}[\mathcal D_N^{\mu-1}\underline{\mathcal H}]_T=\operatorname{cosupp}(\mathfrak J_T|_{N_T},1)\subset V(\mathfrak J_T)\)
(Lemma~\ref{mc:inclusions-tower-along}). A centre \(C\) of the next step is admissible for \(\underline{\mathcal H}_T\), so \(C\subset\operatorname{cosupp}\underline{\mathcal H}_T\subset
N_T\cap V(\mathfrak J_T)=N_T\cap\operatorname{cosupp}(\mathfrak J_T,1)\); both marked ideals have the same \(E_T\), so \(C\) is admissible for \(\underline{\mathcal J}_T\).
\end{proof}

\begin{theorem}[Glueing Lemma along \(N\): the lifts along multiple test blow-ups]
\label{mc:glueing-lift-along}
\uses{mc:glueing-lift,mc:glueing-trivial-along,mc:glue-tower-along,mc:lift-step,mc:lift-biholomorphism,mc:trivial-action-lift,mc:glueing-biholomorphism-along,mc:glueing-tower-pullback-along,mc:strict-along}
\source{wlodarczyk05-hironaka,wlodarczyk09-analytic-resolution:page-0021,wlodarczyk09-analytic-resolution:page-0022}
\lean{Hironaka.HomogenizedGlueingLiftAlong.IsGlueingLiftAlong,Hironaka.HomogenizedGlueingLiftAlong.IsGlueingLiftAlong.preimage_inter_eq,Hironaka.HomogenizedGlueingLiftAlong.exists_isGlueingLiftAlong,Hironaka.HomogenizedGlueingLiftAlong.IsGlueingLiftAlong.apply_eq_of_mem_cosupp,Hironaka.HomogenizedGlueingLiftAlong.glueing_lemma_lift_along}
\leanok
Let \(u\), \(v\) be tangent directions along \(N\) of \((\mathfrak a|_N,\mu)\) on open sets containing \(x\in\operatorname{cosupp}(\mathfrak a|_N,\mu)\), and let \(\Phi\colon U\to U'\)
be the glueing biholomorphism of Lemma~\ref{mc:glueing-trivial-along}. For every multiple test blow-up \(T\) of \(\underline{\mathcal H}\) there is a biholomorphism
\(\Phi_T\) from the preimage \(U_T\) of \(U\) onto the preimage \(U'_T\) of \(U'\) such that
\begin{enumerate}
\item \(\Phi_T\) lifts \(\varphi\);
\item \(\Phi_T^*(\underline{\mathcal H}_T|_{U'_T})=\underline{\mathcal H}_T|_{U_T}\) as marked ideals, in particular \(\Phi_T^{-1}(N_T\cap U'_T)=N_T\cap U_T\);
\item \(\Phi_T\) maps the strict transforms of \(V(v)\) and of \(N\cap V(v)\) onto those of \(V(u)\) and of \(N\cap V(u)\);
\item \(\Phi_T\) and \(\Phi_T^{-1}\) act trivially modulo \(\mathfrak J_T\) at its zeros; in particular \(\Phi_T\) fixes \(\operatorname{cosupp}\underline{\mathcal H}_T\cap U_T\) pointwise.
\end{enumerate}
\end{theorem}
\begin{proof}
\leanok
As for Theorem~\ref{mc:glueing-lift}, by induction on \(T\), the empty tower being Lemmas~\ref{mc:glueing-biholomorphism-along} and~\ref{mc:glueing-trivial-along}.
At a step with centre \(C\), \(C\subset V(\mathfrak J_T)\) (Lemma~\ref{mc:glue-tower-along}), which \(\Phi_T\) and \(\Phi_T^{-1}\) fix pointwise, so
\(\Phi_T^{-1}(C\cap U'_T)=C\cap U_T\) and \(\Phi_T\) lifts (Lemma~\ref{mc:lift-biholomorphism}); (2) and (3) are Lemma~\ref{mc:lift-step}, and (4) is
Theorem~\ref{mc:trivial-action-lift}, \(\mathfrak J_{T'}\) being the controlled transform of \(\mathfrak J_T\) with mark \(1\). The fixed points:
\(\operatorname{cosupp}\underline{\mathcal H}_T\subset V(\mathfrak J_T)\) (Lemma~\ref{mc:glue-tower-along}).
\end{proof}
This is Włodarczyk's ``le:homo'' (4)(a)--(c) along \(N\); with Theorem~\ref{mc:glueing-tower-pullback-along} it completes item (4) of Lemma~\ref{hir:along-calculus}
along multiple test blow-ups.

\subsection{Reduction to a hypersurface of maximal contact inside a submanifold}

In this subsection \(\underline{\mathcal I}=(X,N,E,\mathfrak a,\mu)\) with \(N\) a closed submanifold, \(\mathfrak a\) coherent and \(\mu\ge1\).

\begin{lemma}[Tangent directions along \(N\) on open sets]
\label{mc:tangent-direction-along-open}
\uses{mc:tangent-direction-along,mc:derivative-along-etale,mc:field-transfer,mc:tangent-direction-open}
\source{wlodarczyk05-hironaka}
\lean{Hironaka.ContactEquivalenceAlong.isTangentDirectionAlong_pullback_val,Hironaka.ContactEquivalenceAlong.isTangentDirectionAlong_pullback}
\leanok
If \(u\) is a tangent direction of \((\mathfrak a|_N,\mu)\) along \(N\) on an open set \(U\), then \(u|_U\) is a tangent direction of \((\iota^*\mathfrak a|_{\iota^{-1}N},\mu)\) along
\(\iota^{-1}N\) on all of \(U\), \(\iota\colon U\to X\) the inclusion. More generally, for a holomorphic map \(\varphi\colon X'\to X\) with holomorphic local inverses at
every point, \(u\circ\varphi\) is a tangent direction of \((\varphi^*\mathfrak a|_{\varphi^{-1}N},\mu)\) along \(\varphi^{-1}N\) on \(\varphi^{-1}U\).
\end{lemma}
\begin{proof}
\leanok
\(\varphi^*\mathcal D_N^{\mu-1}(\mathfrak a)=\mathcal D_{\varphi^{-1}N}^{\mu-1}(\varphi^*\mathfrak a)\) (Lemma~\ref{mc:derivative-along-etale}); a field tangent to \(N\) with \(\xi(u)\ne0\) at a
zero of \(u\) transfers along \(\varphi\) and its local inverse to a field tangent to \(\varphi^{-1}N\) with the same derivative of \(u\circ\varphi\)
(Lemma~\ref{mc:field-transfer}).
\end{proof}

\begin{lemma}[Reduction to a hypersurface of maximal contact inside \(N\)]
\label{mc:contact-equivalence-along}
\uses{mc:contact-equivalence,mc:coefficient-restriction-full-along,hir:giraud-tower-along,mc:contact-hypersurface-along,hir:marked-equivalence}
\source{wlodarczyk05-hironaka}
\lean{Hironaka.ContactEquivalenceAlong.isClosedSubmanifold_contact_along,Hironaka.ContactEquivalenceAlong.centresIn_contact_along,Hironaka.ContactEquivalenceAlong.contact_wEquiv_along}
\leanok
Let \(u\) be a tangent direction of \((\mathfrak a|_N,\mu)\) along \(N\) on all of \(X\). Then \(N\cap V(u)\) is a closed submanifold, every multiple test blow-up of
\(\underline{\mathcal I}\) has its centres in the strict transforms of \(N\cap V(u)\), and \(\underline{\mathcal I}\simeq\mathcal C_N(\underline{\mathcal I})|_{N\cap V(u)}\)
(Definition~\ref{hir:marked-equivalence}).
\end{lemma}
\begin{proof}
\leanok
\(N\cap V(u)\) is closed and regular of dimension \(\dim N-1\) at its points (Lemma~\ref{mc:contact-hypersurface-along}). By Lemma~\ref{hir:giraud-tower-along} the admissible centres
and the cosupports of the transforms along multiple test blow-ups lie in the strict transforms of \(N\cap V(u)\); Theorem~\ref{mc:coefficient-restriction-full-along}
identifies the multiple test blow-ups and gives \(\operatorname{cosupp}[\mathcal C_N(\underline{\mathcal I})|_{N\cap V(u)}]_T=\operatorname{cosupp}\underline{\mathcal I}_T\cap(N\cap V(u))_T=\operatorname{cosupp}\underline{\mathcal I}_T\).
\end{proof}
This is Lemma~\ref{mc:contact-equivalence} along \(N\): Włodarczyk's ``le:codim2'' inside \(N\), with any \(E\) and without its hypothesis on the codimension of the support.

\begin{lemma}[Reduction to a hypersurface of maximal contact inside \(N\), on an open set]
\label{mc:contact-equivalence-open-along}
\uses{mc:contact-equivalence-along,mc:tangent-direction-along-open,mc:homogenized-along-etale,hir:open-restriction-marked,mc:contact-equivalence-open}
\source{wlodarczyk05-hironaka}
\lean{Hironaka.ContactEquivalenceAlong.contact_wEquiv_open_along}
\leanok
Let \(u\) be a tangent direction of \((\mathfrak a|_N,\mu)\) along \(N\) on an open set \(U\). Then
\(\underline{\mathcal I}|_U\simeq(U,\iota^{-1}(N\cap V(u)),E|_U,\mathcal C_N(\mathfrak a)|_U,\mu_{\mathcal C})\) (Definition~\ref{hir:marked-equivalence} on \(U\)).
\end{lemma}
\begin{proof}
\leanok
By Lemma~\ref{mc:tangent-direction-along-open}, \(u|_U\) is a tangent direction along \(\iota^{-1}N\) on all of \(U\); Lemma~\ref{mc:contact-equivalence-along} on \(U\) gives
\(\underline{\mathcal I}|_U\simeq\mathcal C_{\iota^{-1}N}(\underline{\mathcal I}|_U)|_{\iota^{-1}(N\cap V(u))}\), and \(\mathcal C_{\iota^{-1}N}(\underline{\mathcal I}|_U)=\mathcal C_N(\underline{\mathcal I})|_U\)
(Lemma~\ref{mc:homogenized-along-etale}).
\end{proof}
This is Lemma~\ref{hir:contact-equivalence} inside \(N\), as Step 1b of Proposition~\ref{hir:maximal-order-step} uses it along \(N\).

\begin{lemma}[Restricted tangent directions along \(N\) are locally genuine]
\label{mc:tangent-direction-along-restricted}
\uses{mc:tangent-direction-along,mc:contact-equivalence-open-along,mc:derivative-along-coherent,hir:giraud-tower-along}
\source{wlodarczyk05-hironaka}
\lean{Hironaka.ContactEquivalenceAlong.fieldDeriv_eq_of_eqOn_N,Hironaka.ContactEquivalenceAlong.exists_isTangentDirectionAlong_of_isInSpanOnAt,Hironaka.ContactEquivalenceAlong.contact_wEquiv_open_along_of_isInSpanOnAt}
\leanok
Let \(w\) be holomorphic on an open set \(O\), at every point of \(N\cap O\) equal on \(N\) to a combination of germs of \(\mathcal D_N^{\mu-1}(\mathfrak a)\), and such that
at every zero \(y\in N\cap O\) some field \(\xi\) tangent to \(N\) has \(\xi(w)(y)\ne0\) (the form in which tangent directions along \(N_T\) are produced along
multiple test blow-ups, Lemma~\ref{hir:giraud-tower-along}). Then every \(x\in N\cap O\) has an open neighbourhood \(U\subset O\) carrying a tangent direction
\(\tilde w\) of \((\mathfrak a|_N,\mu)\) along \(N\) with \(\tilde w=w\) on \(N\cap U\); consequently
\(\underline{\mathcal I}|_U\simeq(U,\iota^{-1}(N\cap V(w)),E|_U,\mathcal C_N(\mathfrak a)|_U,\mu_{\mathcal C})\).
\end{lemma}
\begin{proof}
\leanok
At \(x\), \(w=\sum a_is_i\) on \(N\) near \(x\) with \(s_i\in\mathcal D_N^{\mu-1}(\mathfrak a)_x\); put \(\tilde w=\sum a_is_i\). As \(\mathcal D_N^{\mu-1}(\mathfrak a)\) is coherent
(Lemma~\ref{mc:derivative-along-coherent}), \(\tilde w\) is a germ of it at every point near \(x\). Derivatives along fields tangent to \(N\) of \(\tilde w\) and \(w\) agree at
points of \(N\), since \(\tilde w-w\) vanishes on \(N\). At zeros of \(\tilde w\) on \(N\) use the field given for \(w\). If \(w(x)\ne0\), \(\tilde w\) has no zeros near \(x\). If
\(w(x)=0\), the field \(\xi\) at \(x\) has \(\xi(\tilde w)\ne0\) near \(x\), and it is tangent to \(N\) at the points off the closed set \(N\), which settles the zeros off \(N\).
The equivalence is Lemma~\ref{mc:contact-equivalence-open-along} for \(\tilde w\), whose zero set meets \(N\cap U\) in that of \(w\).
\end{proof}

\subsection{Independence of the hypersurface of maximal contact inside a submanifold, along multiple test blow-ups}

In this subsection \(\underline{\mathcal H}=(M,N,\emptyset,\mathcal H_N(\mathfrak a),\mu)\), \(\mathfrak J\) and the glueing biholomorphism \(\Phi\colon U\to U'\) of two tangent
directions \(u\), \(v\) along \(N\) are as in Theorem~\ref{mc:glueing-lift-along}, and \(\overline{\underline{\mathcal J}}=\mathcal C_N(\underline{\mathcal H})\) is the coefficient
ideal along \(N\) of the homogenized marked ideal, Step 1b's marked ideal (with \(E=\emptyset\)). For a stage \(T\) and a marked ideal \(\underline{\mathcal I}\) on the last
stage of \(T\) put \(\underline{\mathcal S}_{u,T}(\underline{\mathcal I})=\mathcal C_{N(\underline{\mathcal I})}(\underline{\mathcal I})|_{(N\cap V(u))_T}\), the coefficient ideal along the
submanifold of \(\underline{\mathcal I}\) restricted to the strict transform of the hypersurface of maximal contact \(N\cap V(u)\).

\begin{lemma}[Transforms of two marked ideals under the lift of a biholomorphism]
\label{mc:lift-step-pair}
\uses{mc:lift-step,mc:lift-biholomorphism,mc:transform-source-naturality,hir:biholomorphism-transform,hir:open-restriction-marked}
\lean{Hironaka.ContactIndependenceAlong.pullback_lift_geoTransform_pair}
\leanok
In Lemma~\ref{mc:lift-step}, let \(\underline{\mathcal I}\) and \(\underline{\mathcal I}'\) be marked ideals with coherent ideals and coherent transforms
\(\underline{\mathcal I}_1\), \(\underline{\mathcal I}'_1\) along \(\pi\). If \(\Phi^*(\underline{\mathcal I}|_B)=\underline{\mathcal I}'|_A\), then
\(\Phi_1^*(\underline{\mathcal I}_1|_{\pi^{-1}B})=\underline{\mathcal I}'_1|_{\pi^{-1}A}\).
\end{lemma}
\begin{proof}
\leanok
The proof of Lemma~\ref{mc:lift-step} with \(\underline{\mathcal I}'\) on the side of \(A\): restriction to \(A\) commutes with the transform of \(\underline{\mathcal I}'\)
(Lemma~\ref{hir:open-restriction-marked}), and the naturality of the transform under \(\Phi\) (Lemma~\ref{hir:biholomorphism-transform}) carries the transform of
\(\underline{\mathcal I}|_B\) along \(\pi_B\circ\Phi_1=\Phi\circ\pi_A\) to that of \(\Phi^*(\underline{\mathcal I}|_B)=\underline{\mathcal I}'|_A\) along \(\pi_A\).
\end{proof}

\begin{lemma}[The lifts along \(N\) transport preserved marked ideals]
\label{mc:lift-transport-along}
\uses{mc:glueing-lift-along,mc:lift-step-pair,mc:glue-tower-along,mc:lift-biholomorphism,mc:trivial-action-lift,mc:lift-step,hir:open-restriction-tower,hir:geo-transform-coherent}
\source{wlodarczyk05-hironaka}
\lean{Hironaka.ContactIndependenceAlong.exists_isGlueingLiftAlong_step,Hironaka.ContactIndependenceAlong.lift_eq_equiv_refl,Hironaka.ContactIndependenceAlong.exists_isGlueingLiftAlong_transport,Hironaka.ContactIndependenceAlong.exists_isGlueingLiftAlong_transport_one}
\leanok
Let \((\underline{\mathcal I}_i,\underline{\mathcal I}'_i)_{i}\) be a family of pairs of marked ideals on \(M\) with closed submanifolds and coherent ideals such that
\(\Phi^*(\underline{\mathcal I}_i|_{U'})=\underline{\mathcal I}'_i|_U\) for every \(i\). Along every tower \(T\) that is a multiple test blow-up of \(\underline{\mathcal H}\) and of every
\(\underline{\mathcal I}_i\) and \(\underline{\mathcal I}'_i\) there is a biholomorphism \(\Phi_T\) with the properties (1)--(4) of Theorem~\ref{mc:glueing-lift-along} and with
\(\Phi_T^*((\underline{\mathcal I}_i)_T|_{U'_T})=(\underline{\mathcal I}'_i)_T|_{U_T}\) for every \(i\). Along the empty tower every such \(\Phi_T\) is \(\Phi\).
\end{lemma}
\begin{proof}
\leanok
Induction on \(T\), carrying the lift of the previous stage. For the empty tower, (1) forces \(\Phi_T=\Phi\). At a step with centre \(C\), the proof of
Theorem~\ref{mc:glueing-lift-along} gives \(\Phi_T^{-1}(C\cap U'_T)=C\cap U_T\) and a lift \(\Phi_{T'}\) of \(\Phi_T\) (\(\pi\circ\Phi_{T'}=\Phi_T\circ\pi\),
Lemma~\ref{mc:lift-biholomorphism}) with (1)--(4). Along \(T\) the transforms of \(\underline{\mathcal I}_i\) and \(\underline{\mathcal I}'_i\) have closed submanifolds and coherent
ideals, and \(C\) is admissible for both, so their transforms along the blowup are coherent (Lemmas~\ref{hir:open-restriction-tower} and~\ref{hir:geo-transform-coherent});
Lemma~\ref{mc:lift-step-pair} carries \(\Phi_T^*((\underline{\mathcal I}_i)_T|_{U'_T})=(\underline{\mathcal I}'_i)_T|_{U_T}\) to the next stage.
\end{proof}

\begin{lemma}[Restricted coefficient ideals under biholomorphisms of open sets]
\label{mc:restricted-coefficient-transport}
\uses{mc:homogenized-along-etale,mc:restrict-marked,hir:open-restriction-marked}
\lean{Hironaka.ContactIndependenceAlong.pullback_withN_of_pullback,Hironaka.ContactIndependenceAlong.pullback_withN_coeffMarkedAlong_of_pullback}
\leanok
Let \(\Phi\colon A\to B\) be a biholomorphism of open sets of a complex manifold \(Y\), \(\underline{\mathcal I}\), \(\underline{\mathcal I}'\) marked ideals on \(Y\) with
\(\Phi^*(\underline{\mathcal I}|_B)=\underline{\mathcal I}'|_A\), and \(S,S'\subset Y\) with \(\Phi^{-1}(S\cap B)=S'\cap A\). Then
\(\Phi^*(\underline{\mathcal I}|_S\,|_B)=\underline{\mathcal I}'|_{S'}\,|_A\). If moreover both have closed submanifolds, coherent ideals and marks \(\ge1\), then
\(\Phi^*(\mathcal C(\underline{\mathcal I})|_S\,|_B)=\mathcal C(\underline{\mathcal I}')|_{S'}\,|_A\), the coefficient ideals being taken along the submanifolds of
\(\underline{\mathcal I}\) and \(\underline{\mathcal I}'\).
\end{lemma}
\begin{proof}
\leanok
Restricting the submanifold to \(S\) only replaces it, and the pullback of the restriction is the restriction of the pullback to the preimage of \(S\). For the
coefficient ideals, \(\mathcal C_N\) commutes with restriction to the open set \(B\) and with the biholomorphism \(\Phi\) (Lemma~\ref{mc:homogenized-along-etale}), so
\(\Phi^*(\mathcal C(\underline{\mathcal I})|_B)=\mathcal C(\Phi^*(\underline{\mathcal I}|_B))=\mathcal C(\underline{\mathcal I}'|_A)=\mathcal C(\underline{\mathcal I}')|_A\).
\end{proof}

\begin{lemma}[The restricted coefficient ideals correspond under the lifts]
\label{mc:restricted-coefficient-lift-along}
\uses{mc:glueing-lift-along,mc:restricted-coefficient-transport,mc:lift-transport-along,mc:strict-along}
\source{wlodarczyk05-hironaka}
\lean{Hironaka.ContactIndependenceAlong.lift_preimage_proj_preimage,Hironaka.ContactIndependenceAlong.lift_preimage_strict_inter_proj,Hironaka.ContactIndependenceAlong.lift_preimage_strict_zeroSet_inter_proj,Hironaka.ContactIndependenceAlong.lift_pullback_restrictedCoeff,Hironaka.ContactIndependenceAlong.lift_pullback_restricted,Hironaka.ContactIndependenceAlong.lift_pullback_restrictedCoeff_inter_proj}
\leanok
Let \(\Phi_T\) be as in Theorem~\ref{mc:glueing-lift-along} along a tower \(T\), \(\pi_T\) the projection of its last stage, and \(\underline{\mathcal I}\), \(\underline{\mathcal I}'\)
marked ideals on the last stage with \(\Phi_T^*(\underline{\mathcal I}|_{U'_T})=\underline{\mathcal I}'|_{U_T}\) (for instance \(\underline{\mathcal I}=\underline{\mathcal I}'\) a transform
as in Lemma~\ref{mc:lift-transport-along}). Then
\(\Phi_T^*(\underline{\mathcal I}|_{(N\cap V(u))_T}\,|_{U'_T})=\underline{\mathcal I}'|_{(N\cap V(v))_T}\,|_{U_T}\), and, with closed submanifolds, coherent ideals and marks
\(\ge1\), \(\Phi_T^*(\underline{\mathcal S}_{u,T}(\underline{\mathcal I})|_{U'_T})=\underline{\mathcal S}_{v,T}(\underline{\mathcal I}')|_{U_T}\). For sets \(A,A'\subset M\) with
\(\varphi^{-1}(A)\cap U=A'\cap U\), \(\Phi_T\) maps \(\pi_T^{-1}(A')\cap U_T\), \((N\cap V(v))_T\cap\pi_T^{-1}(A')\cap U_T\) and \(V(v)_T\cap\pi_T^{-1}(A')\cap U_T\) onto
\(\pi_T^{-1}(A)\cap U'_T\), \((N\cap V(u))_T\cap\pi_T^{-1}(A)\cap U'_T\) and \(V(u)_T\cap\pi_T^{-1}(A)\cap U'_T\), and the restricted coefficient ideals to these sets
correspond in the same way.
\end{lemma}
\begin{proof}
\leanok
\(\Phi_T\) maps \((N\cap V(v))_T\cap U_T\) onto \((N\cap V(u))_T\cap U'_T\) and \(V(v)_T\cap U_T\) onto \(V(u)_T\cap U'_T\) (Theorem~\ref{mc:glueing-lift-along}(3)), and
\(\pi_T\circ\Phi_T=\varphi\circ\pi_T\) (Theorem~\ref{mc:glueing-lift-along}(1)), so \(\Phi_T^{-1}(\pi_T^{-1}(A)\cap U'_T)=\pi_T^{-1}(\varphi^{-1}(A))\cap U_T=\pi_T^{-1}(A')\cap U_T\).
Apply Lemma~\ref{mc:restricted-coefficient-transport} to \(\Phi_T\) and these sets.
\end{proof}

\begin{theorem}[Independence of the hypersurface of maximal contact inside \(N\)]
\label{mc:contact-independence-along}
\uses{mc:lift-transport-along,mc:restricted-coefficient-lift-along,mc:restricted-coefficient-transport,mc:glueing-lift-along,mc:glueing-trivial-along,mc:glueing-biholomorphism-along,mc:homogenized-along-etale,mc:homogenized-along-basic,mc:coefficient-equivalent-along,hir:marked-equivalence,hir:open-restriction-tower}
\source{wlodarczyk05-hironaka}
\lean{Hironaka.ContactIndependenceAlong.one_le_coeffMarkedAlong_d,Hironaka.ContactIndependenceAlong.coeffHom_basic,Hironaka.ContactIndependenceAlong.wEquiv_coeffHom,Hironaka.ContactIndependenceAlong.pullback_restrict_coeffHom,Hironaka.ContactIndependenceAlong.exists_isGlueingLiftAlong_coeffHom,Hironaka.ContactIndependenceAlong.contact_independence_along}
\leanok
\(\overline{\underline{\mathcal J}}\) has submanifold \(N\), coherent ideal and mark \(\mu_{\mathcal C}\ge1\); \(\Phi^*(\overline{\underline{\mathcal J}}|_{U'})=\overline{\underline{\mathcal J}}|_U\), and
\(\overline{\underline{\mathcal J}}\simeq\underline{\mathcal H}\) (Definition~\ref{hir:marked-equivalence}). Consequently, for tangent directions \(u\), \(v\) along \(N\) of
\((\mathfrak a|_N,\mu)\) on open sets containing \(x\in\operatorname{cosupp}(\mathfrak a|_N,\mu)\) there is a glueing biholomorphism \(\Phi\colon U\to U'\) with
\(x\in U\cap U'\) such that along every multiple test blow-up \(T\) of \(\overline{\underline{\mathcal J}}\) some \(\Phi_T\) as in Theorem~\ref{mc:glueing-lift-along} satisfies
\begin{enumerate}
\item \(\Phi_T^*(\overline{\underline{\mathcal J}}_T|_{U'_T})=\overline{\underline{\mathcal J}}_T|_{U_T}\);
\item \(\Phi_T^*(\underline{\mathcal S}_{u,T}(\overline{\underline{\mathcal J}}_T)|_{U'_T})=\underline{\mathcal S}_{v,T}(\overline{\underline{\mathcal J}}_T)|_{U_T}\);
\item \(\Phi_T^*(\overline{\underline{\mathcal J}}_T|_{(N\cap V(u))_T}\,|_{U'_T})=\overline{\underline{\mathcal J}}_T|_{(N\cap V(v))_T}\,|_{U_T}\);
\item \(\Phi_T\) fixes \(\operatorname{cosupp}\overline{\underline{\mathcal J}}_T\cap U_T\) pointwise.
\end{enumerate}
\end{theorem}
\begin{proof}
\leanok
The mark of a coefficient ideal is \(\prod_{j<\mu}(\mu-j)\ge1\), and its ideal is coherent (Lemma~\ref{mc:homogenized-along-basic}). \(\mathcal C_N\) commutes with restriction to
open sets and with biholomorphisms (Lemma~\ref{mc:homogenized-along-etale}) and \(\Phi^*(\underline{\mathcal H}|_{U'})=\underline{\mathcal H}|_U\)
(Lemma~\ref{mc:glueing-biholomorphism-along}), which gives the first claim; \(\overline{\underline{\mathcal J}}\simeq\underline{\mathcal H}\) is
Lemma~\ref{mc:coefficient-equivalent-along}. Take the glueing data of Lemma~\ref{mc:glueing-trivial-along}. A multiple test blow-up of \(\overline{\underline{\mathcal J}}\) is one of
\(\underline{\mathcal H}\), so Lemma~\ref{mc:lift-transport-along} with the single pair \((\overline{\underline{\mathcal J}},\overline{\underline{\mathcal J}})\) gives \(\Phi_T\) with (1); along
\(T\) the transform \(\overline{\underline{\mathcal J}}_T\) has closed submanifold, coherent ideal and the same mark (Lemma~\ref{hir:open-restriction-tower}), and
Lemma~\ref{mc:restricted-coefficient-lift-along} gives (2) and (3). The cosupports of \(\overline{\underline{\mathcal J}}_T\) and \(\underline{\mathcal H}_T\) agree, and \(\Phi_T\) fixes the
latter (Theorem~\ref{mc:glueing-lift-along}(4)).
\end{proof}
This is the independence of the invariants of Step 1bb of the source from the tangent direction, inside \(N\): there \(\varphi_{pu}^{-1}(V(u)_p)=\varphi_{pv}^{-1}(V(v)_p)\) and
\(\varphi_{pu}^*(U_p,\overline{\mathcal J}_p,E_p,\overline\mu)=\varphi_{pv}^*(U_p,\overline{\mathcal J}_p,E_p,\overline\mu)\), the restrictions to the strict transforms of the
hypersurfaces of maximal contact being compared through functoriality; here \(\varphi_v\) is the inclusion of \(U\) and \(\varphi_u=\Phi\), and (3) is the source's
\(\overline{\mathcal J}_{p|V(u)_p}\) while (2) is the form \(\mathcal C_{N_T}(\overline{\underline{\mathcal J}}_T)|_{(N\cap V(u))_T}\) read by Step 1b of
Proposition~\ref{hir:maximal-order-step} inside \(N\), through Lemmas~\ref{mc:contact-equivalence-open-along} and~\ref{mc:tangent-direction-along-restricted}.

\subsection{Canonicity of the coefficient ideals along a submanifold and of their restrictions, along multiple test blow-ups}

In this subsection \(\underline{\mathcal I}=(X,N,E,\mathfrak a,d)\) has \(N\) a closed submanifold, \(\mathfrak a\) coherent and \(d\ge1\), and \(T\) is a multiple test blow-up
of \(\underline{\mathcal I}\) with transform \(\underline{\mathcal I}_T\) and submanifold \(N_T\). For \(A\subset N\) let \(A_T\) be its strict transform along \(T\)
(Definition~\ref{mc:strict-along}) and
\[\underline{\mathcal S}_A(\underline{\mathcal I}_T)=\mathcal C_{N_T}(\underline{\mathcal I}_T)|_{A_T},\]
the coefficient ideal along \(N_T\) restricted to \(A_T\) (Definition~\ref{mc:restrict-marked}); Step 1a of Proposition~\ref{hir:maximal-order-step} inside \(N\) uses
\(A=N\cap H_S\), Step 1b uses \(A=N\cap V(u)\). Restricted equivalence \(\simeq\) is that of Definition~\ref{hir:restricted-equivalence}.

\begin{lemma}[Representatives of the coefficient ideal along \(N\)]
\label{mc:coefficient-representatives-along}
\uses{mc:homogenized-along,mc:derivative-along,hir:along-transform,hir:restricted-calculus,mc:transform-along-calculus,mc:restrict-marked,hir:restricted-equivalence}
\lean{Hironaka.ContactCanonicalAlong.holomorphicAt_of_mem_coeffMarkedAlong,Hironaka.ContactCanonicalAlong.restrictLe_iterDerivIdealAlong,Hironaka.ContactCanonicalAlong.restrictLe_coeffMarkedAlong,Hironaka.ContactCanonicalAlong.restrictEquiv_coeffMarkedAlong,Hironaka.ContactCanonicalAlong.restrictEquiv_withN,Hironaka.ContactCanonicalAlong.restrictEquiv_restrictedCoeff}
\leanok
Let \(\mathfrak a\), \(\mathfrak b\) be ideals of holomorphic germs and \(W\) a set with \(\mathfrak a|_W\subset\mathfrak b|_W\). Then
\(\mathcal D^i_W(\mathfrak a)|_W\subset\mathcal D^i_W(\mathfrak b)|_W\) for every \(i\). If two marked ideals \(\underline{\mathcal I}\), \(\underline{\mathcal I}'\) have the same submanifold
\(N\) and the same mark and \(\mathfrak a|_N\subset\mathfrak a'|_N\), then \(\mathcal C_N(\underline{\mathcal I})|_N\subset\mathcal C_N(\underline{\mathcal I}')|_N\). Hence
\(\underline{\mathcal I}\simeq\underline{\mathcal I}'\) (ideals of holomorphic germs) gives \(\mathcal C_N(\underline{\mathcal I})\simeq\mathcal C_N(\underline{\mathcal I}')\) and
\(\mathcal C_N(\underline{\mathcal I})|_A\simeq\mathcal C_N(\underline{\mathcal I}')|_A\) for every \(A\subset N\). The germs of \(\mathcal C_N(\underline{\mathcal I})\) are holomorphic.
\end{lemma}
\begin{proof}
\leanok
The first claim is induction on \(i\), one derivative at a time by Lemma~\ref{hir:along-transform} (for \(f=\sum a_kg_k+t\) on \(W\) with \(g_k\in\mathfrak b\) and
\(t\in\mathcal I_W\), a field tangent to \(W\) gives \(\xi(f)=\sum\xi(a_k)g_k+a_k\xi(g_k)+\xi(t)\) with \(\xi(t)\in\mathcal I_W\)). The coefficient ideal is the finite sum
\(\sum_{i<d}\mathcal D^i_N(\mathfrak a)^{e_i}\) with exponents depending only on \(d\) (Definition~\ref{mc:homogenized-along}); restricted inclusions pass to powers and to
finite sums (Lemmas~\ref{hir:restricted-calculus} and~\ref{mc:transform-along-calculus}). The mark of \(\mathcal C_N\) depends only on \(d\), and \(N\) and \(E\) are kept, so
the two inclusions give the restricted equivalence. A restricted inclusion on \(N\) holds on every subset \(A\subset N\) (Lemma~\ref{hir:restricted-calculus}), and
restricting to \(A\) changes only the submanifold. Germs of a finite sum of ideals are holomorphic combinations of holomorphic germs.
\end{proof}

\begin{lemma}[Representatives along towers]
\label{mc:coefficient-representatives-tower}
\uses{mc:coefficient-representatives-along,hir:marked-equivalence-basic,mc:strict-along,bcalc:member-intersection-tower}
\source{wlodarczyk05-hironaka}
\lean{Hironaka.ContactCanonicalAlong.restrictEquiv_self,Hironaka.ContactCanonicalAlong.restrictEquiv_coeffMarkedAlong_geoTransformAlong,Hironaka.ContactCanonicalAlong.restrictEquiv_restrictedCoeff_geoTransformAlong}
\leanok
Let \(\underline{\mathcal I}\simeq\underline{\mathcal I}'\) with \(N\) closed and coherent ideals. Then for every tower \(T\) (not necessarily a multiple test blow-up)
\(\mathcal C_{N_T}(\underline{\mathcal I}_T)\simeq\mathcal C_{N_T}(\underline{\mathcal I}'_T)\), and \(\underline{\mathcal S}_A(\underline{\mathcal I}_T)\simeq\underline{\mathcal S}_A(\underline{\mathcal I}'_T)\) for every
\(A\subset N\).
\end{lemma}
\begin{proof}
\leanok
For the empty tower this is Lemma~\ref{mc:coefficient-representatives-along}. From the first blowup on, restricted-equivalent marked ideals have equal transforms
(Lemma~\ref{hir:marked-equivalence-basic}), and a marked ideal of holomorphic germs is restricted-equivalent to itself. Strict transforms are monotone
(Lemma~\ref{bcalc:member-intersection-tower}) and \(N_T\) is the strict transform of \(N\) (Definition~\ref{mc:strict-along}), so \(A_T\subset N_T\).
\end{proof}

\begin{lemma}[Restriction, dropping members and pullbacks]
\label{mc:restriction-operations-along}
\uses{mc:restrict-marked,hir:biholomorphism-transform,hir:geo-transform,mc:homogenized-along-etale,hir:restricted-equivalence,bcalc:cancellation}
\lean{Hironaka.ContactCanonicalAlong.pullback_withN,Hironaka.ContactCanonicalAlong.pullback_dropE,Hironaka.ContactCanonicalAlong.pullback_id,Hironaka.ContactCanonicalAlong.pullback_restrictedCoeff,Hironaka.ContactCanonicalAlong.pullback_val_restrictedCoeff,Hironaka.ContactCanonicalAlong.restrictEquiv_dropE,Hironaka.ContactCanonicalAlong.geoTransform_dropE}
\leanok
For every map \(\varphi\), \(\varphi^*(\underline{\mathcal I}|_S)=(\varphi^*\underline{\mathcal I})|_{\varphi^{-1}S}\), and \(\varphi^*\) commutes with deleting the first \(k\) members of \(E\);
the pullback along the identity of a marked ideal with coherent ideal is the marked ideal. For a biholomorphism \(\Phi\colon X'\to X\), and for the inclusion of an
open set \(U\subset X\), \(\Phi^*(\mathcal C_N(\underline{\mathcal I})|_A)=\mathcal C_{\Phi^{-1}N}(\Phi^*\underline{\mathcal I})|_{\Phi^{-1}A}\) and
\((\mathcal C_N(\underline{\mathcal I})|_A)|_U=\mathcal C_{N\cap U}(\underline{\mathcal I}|_U)|_{A\cap U}\). Deleting the first \(k\le\#E\) members commutes with restricted
equivalence and with the transform along any blowup.
\end{lemma}
\begin{proof}
\leanok
The pullback of \((X,N,E,\mathfrak a,d)\) is \((\varphi^{-1}N,\varphi^{-1}E,\varphi^*\mathfrak a,d)\) (Lemma~\ref{hir:biholomorphism-transform}), and restriction to \(S\) only replaces the
submanifold (Definition~\ref{mc:restrict-marked}); the pullback along the identity of a coherent ideal is the ideal (Lemma~\ref{bcalc:cancellation}). \(\mathcal C_N\) commutes
with biholomorphisms and with restriction to open sets (Lemma~\ref{mc:homogenized-along-etale}). Restricted equivalence does not involve \(E\) beyond its equality, and the
transform (Definition~\ref{hir:geo-transform}) replaces the members by their strict transforms in order and appends the exceptional divisor, so deleting the first
\(k\le\#E\) members before or after the transform gives the same list; submanifold, ideal and mark of the transform do not depend on \(E\).
\end{proof}

\begin{lemma}[Coefficient ideals and their restrictions along transported towers]
\label{mc:coefficient-transported-tower}
\uses{hir:tower-transport,hir:tower-transport-resolution,hir:open-restriction-tower,hir:biholomorphism-transform,mc:restriction-operations-along,mc:homogenized-along-etale,mc:strict-along}
\source{wlodarczyk05-hironaka}
\lean{Hironaka.ContactCanonicalAlong.geoTransformAlong_transport,Hironaka.ContactCanonicalAlong.coeffMarkedAlong_geoTransformAlong_transport,Hironaka.ContactCanonicalAlong.restrictedCoeff_geoTransformAlong_transport,Hironaka.ContactCanonicalAlong.restricted_geoTransformAlong_transport,Hironaka.ContactCanonicalAlong.cosupp_restrictedCoeff_geoTransformAlong_transport}
\leanok
Let \(\Phi\colon X'\to X\) be a biholomorphism, \(\Phi^*T\) the transported tower and \(\Phi_T\) its last-stage biholomorphism (Definition~\ref{hir:tower-transport}). Then
\((\Phi^*\underline{\mathcal I})_{\Phi^*T}=\Phi_T^*(\underline{\mathcal I}_T)\), \(\mathcal C((\Phi^*\underline{\mathcal I})_{\Phi^*T})=\Phi_T^*\mathcal C_{N_T}(\underline{\mathcal I}_T)\), and for every
\(A\subset X\)
\[\underline{\mathcal S}_{\Phi^{-1}A}\bigl((\Phi^*\underline{\mathcal I})_{\Phi^*T}\bigr)=\Phi_T^*\,\underline{\mathcal S}_A(\underline{\mathcal I}_T),\qquad
(\Phi^*\underline{\mathcal I})_{\Phi^*T}|_{(\Phi^{-1}A)_{\Phi^*T}}=\Phi_T^*(\underline{\mathcal I}_T|_{A_T}),\]
so the cosupport of the left-hand restricted coefficient ideal is \(\Phi_T^{-1}\) of that of \(\underline{\mathcal S}_A(\underline{\mathcal I}_T)\).
\end{lemma}
\begin{proof}
\leanok
For the empty tower \(\Phi_T=\Phi\). From the first blowup on, the transform of \(\Phi^*\underline{\mathcal I}\) along \(\Phi^*T\) equals that of \(\underline{\mathcal I}\) along \(T\)
(Lemma~\ref{hir:tower-transport-resolution}) and \(\Phi_T\) is the identity; the transform has coherent ideal (Lemma~\ref{hir:open-restriction-tower}), so its pullback along
the identity is itself (Lemma~\ref{mc:restriction-operations-along}). \(\underline{\mathcal I}_T\) has closed submanifold \(N_T\), coherent ideal and mark \(d\)
(Lemma~\ref{hir:open-restriction-tower}), so \(\mathcal C_{N_T}\) commutes with \(\Phi_T\) (Lemma~\ref{mc:homogenized-along-etale}). Strict transforms commute with
transport, \((\Phi^{-1}A)_{\Phi^*T}=\Phi_T^{-1}(A_T)\) (Definition~\ref{mc:strict-along}), and restriction commutes with pullback
(Lemma~\ref{mc:restriction-operations-along}). Cosupports of pullbacks by biholomorphisms are preimages (Lemma~\ref{hir:biholomorphism-transform}).
\end{proof}

\begin{lemma}[Coefficient ideals and their restrictions along restricted towers]
\label{mc:coefficient-restricted-tower}
\uses{hir:open-restriction-tower,hir:open-restriction-marked,mc:restriction-operations-along,mc:homogenized-along-etale,mc:strict-along}
\source{wlodarczyk05-hironaka}
\lean{Hironaka.ContactCanonicalAlong.coeffMarkedAlong_geoTransformAlong_restrict,Hironaka.ContactCanonicalAlong.restrictedCoeff_geoTransformAlong_restrict,Hironaka.ContactCanonicalAlong.restricted_geoTransformAlong_restrict,Hironaka.ContactCanonicalAlong.cosupp_restrictedCoeff_geoTransformAlong_restrict}
\leanok
Let \(U\subset X\) be open, \(T|_U\) the restricted tower and \(T^{-1}U\) the open preimage of \(U\) in the last stage. Then
\(\mathcal C((\underline{\mathcal I}|_U)_{T|_U})=\mathcal C_{N_T}(\underline{\mathcal I}_T)|_{T^{-1}U}\), and for every \(A\subset X\)
\[\underline{\mathcal S}_{A\cap U}\bigl((\underline{\mathcal I}|_U)_{T|_U}\bigr)=\underline{\mathcal S}_A(\underline{\mathcal I}_T)|_{T^{-1}U},\qquad
(\underline{\mathcal I}|_U)_{T|_U}|_{(A\cap U)_{T|_U}}=(\underline{\mathcal I}_T|_{A_T})|_{T^{-1}U},\]
so the cosupport of the left-hand restricted coefficient ideal is the trace on \(T^{-1}U\) of that of \(\underline{\mathcal S}_A(\underline{\mathcal I}_T)\).
\end{lemma}
\begin{proof}
\leanok
The transform of \(\underline{\mathcal I}|_U\) along \(T|_U\) is the restriction of \(\underline{\mathcal I}_T\) to \(T^{-1}U\), which has closed submanifold, coherent ideal and mark
\(d\) (Lemma~\ref{hir:open-restriction-tower}); \(\mathcal C_{N_T}\) commutes with restriction to open sets (Lemma~\ref{mc:homogenized-along-etale}); strict transforms commute
with restriction, \((A\cap U)_{T|_U}=A_T\cap T^{-1}U\) (Definition~\ref{mc:strict-along}); restriction to sets commutes with pullback
(Lemma~\ref{mc:restriction-operations-along}); cosupports restrict (Lemma~\ref{hir:open-restriction-marked}).
\end{proof}

\begin{lemma}[Blowups along the empty set]
\label{mc:empty-blowup-inverse}
\uses{bcalc:isomorphism-off-centre,mc:field-transfer,hir:geo-transform,mc:transform-along-pullback}
\lean{Hironaka.ContactCanonicalAlong.exists_isLocalInverseAt_of_isBlowupAlong_empty,Hironaka.ContactCanonicalAlong.geoStrictTransform_empty,Hironaka.ContactCanonicalAlong.restrictLe_geoTransform_empty_le_pullback,Hironaka.ContactCanonicalAlong.restrictLe_pullback_le_geoTransform}
\leanok
Let \(\pi\colon Z\to Y\) be a blowup along \(\emptyset\). Then \(\pi\) has a holomorphic local inverse at every point, in the sense of
Lemma~\ref{mc:field-transfer}; the strict transform of a closed set \(A\) along \(\emptyset\) is \(\pi^{-1}A\). For a marked ideal \(\underline{\mathcal I}\) on \(Y\) with
transform \(\underline{\mathcal I}'=(Z,N',E',\mathfrak a',d)\) along any centre, \(\pi^*\mathfrak a|_{N'}\subset\mathfrak a'|_{N'}\); along \(\emptyset\) also
\(\mathfrak a'|_{N'}\subset\pi^*\mathfrak a|_{N'}\).
\end{lemma}
\begin{proof}
\leanok
Over the complement of the centre, here all of \(Y\), the blowup has a holomorphic section \(s\) with \(\pi\circ s=\mathrm{id}\) and \(s(\pi(z))=z\)
(Lemma~\ref{bcalc:isomorphism-off-centre}); \(s\) is the local inverse at every point. The strict transform is \(\overline{\pi^{-1}(A\setminus\emptyset)}=\pi^{-1}A\) for closed \(A\). The first inclusion holds because the
pullback of \(\mathfrak a\) lies in its colon by \((\pi^*\mathcal I_C)^d\) along \(N'\) (Definition~\ref{hir:geo-transform}); off the centre \((\pi^*\mathcal I_C)^d\) contains the
constant \(1\), so every germ of \(\mathfrak a'\) is on \(N'\) a combination of germs of \(\pi^*\mathfrak a\) (Lemma~\ref{mc:transform-along-pullback}).
\end{proof}

\begin{lemma}[Coefficient ideals and their restrictions across steps with empty centre]
\label{mc:coefficient-empty-step}
\uses{mc:empty-blowup-inverse,mc:coefficient-representatives-along,mc:homogenized-along-etale,hir:restricted-calculus,hir:geo-transform-coherent,bcalc:transform-off-centre,hir:open-restriction-tower,bcalc:member-intersection-tower,mc:strict-along}
\source{wlodarczyk05-hironaka}
\lean{Hironaka.ContactCanonicalAlong.restrictEquiv_coeffMarkedAlong_geoTransform_empty,Hironaka.ContactCanonicalAlong.restrictEquiv_restrictedCoeff_geoTransform_empty,Hironaka.ContactCanonicalAlong.cosupp_restrictedCoeff_geoTransform_empty,Hironaka.ContactCanonicalAlong.restrictEquiv_restrictedCoeff_geoTransformAlong_empty}
\leanok
Let \(\pi\) be a blowup along \(\emptyset\) and \(\underline{\mathcal I}'\) the transform of \(\underline{\mathcal I}\). Then
\(\mathcal C_{N'}(\underline{\mathcal I}')\simeq\mathcal C_N(\underline{\mathcal I})'\), and for \(A\subset N\) with strict transform \(A'\)
\[\mathcal C_{N'}(\underline{\mathcal I}')|_{A'}\simeq\bigl(\mathcal C_N(\underline{\mathcal I})|_A\bigr)';\]
if \(A\) is closed, the cosupport of the left-hand side is \(\pi^{-1}\) of that of \(\mathcal C_N(\underline{\mathcal I})|_A\). Along a multiple test blow-up \(T\) followed by
\(\pi\): \(\underline{\mathcal S}_A(\underline{\mathcal I}_{T\cdot\pi})\simeq\underline{\mathcal S}_A(\underline{\mathcal I}_T)'\).
\end{lemma}
\begin{proof}
\leanok
\(N'=\pi^{-1}N\) and on \(N'\) the ideal \(\mathfrak a'\) has the same restriction as \(\pi^*\mathfrak a\) (Lemma~\ref{mc:empty-blowup-inverse}); its germs are holomorphic, the
transform having coherent ideal (Lemma~\ref{hir:geo-transform-coherent}, \(\emptyset\) being admissible). By Lemma~\ref{mc:coefficient-representatives-along},
\(\mathcal C_{N'}(\mathfrak a')\) and \(\mathcal C_{\pi^{-1}N}(\pi^*\mathfrak a)\) have the same restriction to \(N'\); the latter is \(\pi^*\mathcal C_N(\mathfrak a)\), since \(\pi\) has
holomorphic local inverses (Lemmas~\ref{mc:empty-blowup-inverse} and~\ref{mc:homogenized-along-etale}); and \(\pi^*\mathcal C_N(\mathfrak a)\) has the same restriction to
\(N'\) as the ideal of \(\mathcal C_N(\underline{\mathcal I})'\) (Lemma~\ref{mc:empty-blowup-inverse} for \(\mathcal C_N(\underline{\mathcal I})\), whose transform has the same
submanifold \(N'\)). Submanifolds, members and marks agree by definition, which gives the first equivalence. For \(A\subset N\), \(A'\subset N'\); restricting the chain to
\(A'\) (Lemma~\ref{hir:restricted-calculus}) and comparing with \((\mathcal C_N(\underline{\mathcal I})|_A)'\), whose ideal has on \(A'\) the restriction of
\(\pi^*\mathcal C_N(\mathfrak a)\) (Lemma~\ref{mc:empty-blowup-inverse} once more), gives the second. Restricted-equivalent marked ideals have the same cosupport, and away from
the centre, here everywhere, the cosupport of a transform is the preimage of the cosupport (Lemma~\ref{bcalc:transform-off-centre}). The tower form is the one-step form
for \(\underline{\mathcal I}_T\), which has closed submanifold \(N_T\), coherent ideal and mark \(d\) (Lemma~\ref{hir:open-restriction-tower}), and \(A_T\subset N_T\)
(Lemma~\ref{bcalc:member-intersection-tower}, Definition~\ref{mc:strict-along}).
\end{proof}

\begin{lemma}[Tangent directions along transported and restricted towers and across empty steps]
\label{mc:tangent-direction-canonical-along}
\uses{mc:tangent-direction-along,mc:tangent-direction-along-open,mc:coefficient-transported-tower,mc:coefficient-restricted-tower,mc:empty-blowup-inverse,mc:derivative-along-basic,mc:derivative-along-etale,mc:homogenized-along-etale}
\source{wlodarczyk05-hironaka}
\lean{Hironaka.ContactCanonicalAlong.iterDerivIdealAlong_le,Hironaka.ContactCanonicalAlong.isTangentDirectionAlong_geoTransformAlong_transport,Hironaka.ContactCanonicalAlong.N_inter_zeroSet_geoTransformAlong_transport,Hironaka.ContactCanonicalAlong.isTangentDirectionAlong_geoTransformAlong_restrict,Hironaka.ContactCanonicalAlong.N_inter_zeroSet_geoTransformAlong_restrict,Hironaka.ContactCanonicalAlong.isTangentDirectionAlong_geoTransform_empty}
\leanok
Let \(u\) be a tangent direction of \((\mathfrak a_T|_{N_T},\mu)\) on an open set \(V\) of the last stage (Definition~\ref{mc:tangent-direction-along}). Then \(u\circ\Phi_T\) is a
tangent direction of the transform of \(\Phi^*\underline{\mathcal I}\) along \(\Phi^*T\) on \(\Phi_T^{-1}V\), \(u|_{T^{-1}U}\) is one of the transform of \(\underline{\mathcal I}|_U\)
along \(T|_U\) on \(V\cap T^{-1}U\), and the hypersurfaces correspond: \(N'_{\Phi^*T}\cap V(u\circ\Phi_T)=\Phi_T^{-1}(N_T\cap V(u))\) and
\(N_{T|_U}\cap V(u|_{T^{-1}U})=(N_T\cap V(u))\cap T^{-1}U\). If \(\pi\) is a blowup along \(\emptyset\), \(N\) is closed and \(u\) is a tangent direction of
\((\mathfrak a|_N,\mu)\) on \(V\), then \(u\circ\pi\) is a tangent direction of \((\mathfrak a'|_{N'},\mu)\) on \(\pi^{-1}V\), and \(N'\cap V(u\circ\pi)=\pi^{-1}(N\cap V(u))\).
\end{lemma}
\begin{proof}
\leanok
Along the transported and restricted towers the transforms are the pullbacks of \(\underline{\mathcal I}_T\) by \(\Phi_T\) and by the inclusion of \(T^{-1}U\)
(Lemmas~\ref{mc:coefficient-transported-tower} and~\ref{mc:coefficient-restricted-tower}), which have holomorphic local inverses
(Lemmas~\ref{mc:homogenized-along-etale} and~\ref{mc:derivative-along-etale}), and tangent directions pull back by such maps (Lemma~\ref{mc:tangent-direction-along-open}).
Across an empty step, \(u\circ\pi\) is a tangent direction of \((\pi^*\mathfrak a|_{\pi^{-1}N},\mu)\) by the same lemma and Lemma~\ref{mc:empty-blowup-inverse}; since
\(\pi^*\mathfrak a\subset\mathfrak a'\) (Lemma~\ref{mc:empty-blowup-inverse}) and the derivative ideals along \(N'=\pi^{-1}N\) are monotone
(Lemma~\ref{mc:derivative-along-basic}), its germs lie in \(\mathcal D^{\mu-1}_{N'}(\mathfrak a')\), and the multiplicity-one condition does not involve the ideal.
\end{proof}
These lemmas are the canonicity of the ingredients that the maximal contact provides to Step 1 of Proposition~\ref{hir:maximal-order-step} inside \(N\): property (6) of
Definition~\ref{hir:centre-function} for the restricted coefficient ideals \(\underline{\mathcal S}_A\) and the tangent directions is
Lemmas~\ref{mc:coefficient-transported-tower} and~\ref{mc:tangent-direction-canonical-along}, property (5) is Lemmas~\ref{mc:coefficient-restricted-tower}
and~\ref{mc:tangent-direction-canonical-along}, property (4) is Lemmas~\ref{mc:coefficient-empty-step} and~\ref{mc:tangent-direction-canonical-along} (restricting at the stage
after an empty step or transforming the restriction made before it gives restricted-equivalent marked ideals), and property (7) is
Lemma~\ref{mc:coefficient-representatives-tower}. The codimension-one part and the boundary count are the other half (Lemmas~\ref{bcalc:codim-one-part-canonical}
and~\ref{bcalc:boundary-count-canonical}). In the source, the lemma ``Commutativity of the resolution procedure in Step 1 with étale morphisms'' in the proof of
Proposition~pr:resol reads these statements for an étale \(\varphi\): \(\varphi^*(\overline{\mathcal J}_{i|H^s_\alpha})=\overline{\mathcal J}'_{j|(H')^s_\alpha}\) and
\(\overline{\mathcal J}'_{|V(u)}=\varphi^*(\overline{\mathcal J}_{|V(u)})\), the inverse image of a hypersurface of maximal contact being one, while blow-ups whose centres
miss the image pull back to isomorphisms; here étale maps are open embeddings and biholomorphisms, and those isomorphisms are steps with empty centre.

\subsection{Isomorphisms of towers and the glueing lifts}

In this subsection \(\Phi\colon X'\to X\) is a biholomorphism, \(T\) a tower over \(X\), \(T'\) a tower over \(X'\), and \(\chi\colon T'\to T\) an isomorphism over \(\Phi\)
(Definition~\ref{hir:tower-isomorphism}) with last-stage biholomorphism \(\chi_r\colon X'_r\to X_r\).

\begin{lemma}[Isomorphisms of towers]
\label{mc:tower-isomorphism-basic}
\uses{hir:tower-isomorphism,hir:tower-transport,mc:strict-along,mc:transform-source-naturality}
\lean{Hironaka.BlowupTowerIsomorphism.IsTowerIso.length_eq,Hironaka.BlowupTowerIsomorphism.IsTowerIso.proj_comm,Hironaka.BlowupTowerIsomorphism.IsTowerIso.strictAlong_eq,Hironaka.BlowupTowerIsomorphism.transport_isTowerIso}
\leanok
Isomorphic towers have the same length, \(\Phi\circ p_{T'}=p_T\circ\chi_r\) for the projections, and for every \(A\subset X\) the strict transforms correspond:
\((\Phi^{-1}A)_{T'}=\chi_r^{-1}(A_T)\). The transport \(\Phi^*T\) (Definition~\ref{hir:tower-transport}) is isomorphic to \(T\) over \(\Phi\), with its last-stage
biholomorphism.
\end{lemma}
\begin{proof}
\leanok
Induction on the isomorphism. At a step, \(p_{T'}\circ\sigma'=p_{T'}\) of the previous stage composed with \(\sigma'\), and \(\sigma\circ\chi_{i}=\chi_{i-1}\circ\sigma'\). For the strict
transforms, \(\chi_{i-1}^{-1}C=C'\) and \(\sigma'^{-1}(\chi_{i-1}^{-1}(A_{i-1}\setminus C))=\chi_i^{-1}(\sigma^{-1}(A_{i-1}\setminus C))\), and closures commute with the preimage
under the homeomorphism \(\chi_i\) (Lemma~\ref{mc:transform-source-naturality}). For the transport, \(\chi_i=\mathrm{id}\) for \(i\ge1\) and the first blowup is \(\Phi^{-1}\circ\sigma_1\) along
\(\Phi^{-1}C_0\), so the conditions hold with \(\chi_0=\Phi\); after the first step every blowup is the composite of the old one with \(\chi_{i-1}^{-1}\) and
\(\chi_{i-1}\circ\chi_{i-1}^{-1}=\mathrm{id}\).
\end{proof}

\begin{lemma}[Transforms along isomorphic towers]
\label{mc:tower-isomorphism-transform}
\uses{mc:tower-isomorphism-basic,hir:tower-isomorphism,hir:biholomorphism-transform,hir:biholomorphism-invariance,hir:geo-transform-coherent,hir:open-restriction-tower,mc:transform-source-naturality}
\source{wlodarczyk05-hironaka}
\lean{Hironaka.BlowupTowerIsomorphism.IsTowerIso.geoTransformAlong_eq,Hironaka.BlowupTowerIsomorphism.IsTowerIso.isGeoTestTower_iff,Hironaka.BlowupTowerIsomorphism.IsTowerIso.cosupp_geoTransformAlong,Hironaka.BlowupTowerIsomorphism.IsTowerIso.isGeoResolution_iff}
\leanok
Let \(\underline{\mathcal I}\) be a marked ideal on \(X\) with closed submanifold \(N\) and coherent ideal. Then \(T'\) is a multiple test blow-up of
\(\Phi^*\underline{\mathcal I}\) if and only if \(T\) is one of \(\underline{\mathcal I}\), and then \((\Phi^*\underline{\mathcal I})_{T'}=\chi_r^*(\underline{\mathcal I}_T)\); in particular
the cosupports correspond under \(\chi_r\), and \(T'\) is a resolution of \(\Phi^*\underline{\mathcal I}\) if and only if \(T\) is one of \(\underline{\mathcal I}\).
\end{lemma}
\begin{proof}
\leanok
Induction on the isomorphism, the case of the empty towers being trivial. At a step, write the new blowup of \(T'\) as \(\sigma'=\chi_{i-1}^{-1}\circ(\sigma\circ\chi_i)\) and its centre
as \(\chi_{i-1}^{-1}C\). By induction the transform along the previous stages is \(\chi_{i-1}^*\) of that of \(\underline{\mathcal I}\), which has closed submanifold and
coherent ideal (Lemma~\ref{hir:open-restriction-tower}). The transform of \(\chi_{i-1}^*\underline{\mathcal J}\) along \(\chi_{i-1}^{-1}\circ(\sigma\circ\chi_i)\) with centre
\(\chi_{i-1}^{-1}C\) is the transform of \(\underline{\mathcal J}\) along \(\sigma\circ\chi_i\) (Lemma~\ref{hir:biholomorphism-transform}), and that is \(\chi_i^*\) of the
transform along \(\sigma\) (Lemma~\ref{mc:transform-source-naturality}; the transform along \(\sigma\) is coherent, Lemma~\ref{hir:geo-transform-coherent}). Admissible centres correspond
under biholomorphisms (Lemma~\ref{hir:biholomorphism-invariance}), which gives the multiple test blow-ups step by step; cosupports of pullbacks by
biholomorphisms are preimages (the same lemma), and \(\chi_r^{-1}S=\emptyset\) if and only if \(S=\emptyset\).
\end{proof}

\begin{lemma}[Coefficient ideals, their restrictions and tangent directions along isomorphic towers]
\label{mc:tower-isomorphism-coefficient}
\uses{mc:tower-isomorphism-transform,mc:tower-isomorphism-basic,mc:homogenized-along-etale,mc:restriction-operations-along,mc:tangent-direction-along-open,mc:coefficient-representatives-along}
\source{wlodarczyk05-hironaka}
\lean{Hironaka.BlowupTowerIsomorphism.IsTowerIso.coeffMarkedAlong_eq,Hironaka.BlowupTowerIsomorphism.IsTowerIso.restrictedCoeff_eq,Hironaka.BlowupTowerIsomorphism.IsTowerIso.restricted_eq,Hironaka.BlowupTowerIsomorphism.IsTowerIso.cosupp_restrictedCoeff_eq,Hironaka.BlowupTowerIsomorphism.IsTowerIso.isTangentDirectionAlong}
\leanok
In Lemma~\ref{mc:tower-isomorphism-transform}, let \(d\ge1\) and \(T\) be a multiple test blow-up of \(\underline{\mathcal I}\). Then
\(\mathcal C((\Phi^*\underline{\mathcal I})_{T'})=\chi_r^*\mathcal C_{N_T}(\underline{\mathcal I}_T)\) and, for every \(A\subset X\),
\[\underline{\mathcal S}_{\Phi^{-1}A}\bigl((\Phi^*\underline{\mathcal I})_{T'}\bigr)=\chi_r^*\,\underline{\mathcal S}_A(\underline{\mathcal I}_T),\qquad
(\Phi^*\underline{\mathcal I})_{T'}|_{(\Phi^{-1}A)_{T'}}=\chi_r^*(\underline{\mathcal I}_T|_{A_T}),\]
with \(\underline{\mathcal S}_A(\underline{\mathcal I}_T)=\mathcal C_{N_T}(\underline{\mathcal I}_T)|_{A_T}\); the cosupports of the restricted coefficient ideals correspond under
\(\chi_r\). If \(u\) is a tangent direction of \((\mathfrak a_T|_{N_T},\mu)\) on \(V\), then \(u\circ\chi_r\) is one of the transform of \(\Phi^*\underline{\mathcal I}\) along \(T'\) on
\(\chi_r^{-1}V\), and \(N'_{T'}\cap V(u\circ\chi_r)=\chi_r^{-1}(N_T\cap V(u))\).
\end{lemma}
\begin{proof}
\leanok
By Lemma~\ref{mc:tower-isomorphism-transform} the transform along \(T'\) is \(\chi_r^*(\underline{\mathcal I}_T)\), with closed submanifold, coherent ideal and mark \(d\);
\(\mathcal C_{N_T}\) commutes with the biholomorphism \(\chi_r\) (Lemma~\ref{mc:homogenized-along-etale}), the strict transforms correspond (Lemma~\ref{mc:tower-isomorphism-basic}),
restriction to sets commutes with pullback (Lemma~\ref{mc:restriction-operations-along}), and tangent directions pull back by biholomorphisms
(Lemma~\ref{mc:tangent-direction-along-open}). Germs of coefficient ideals are holomorphic (Lemma~\ref{mc:coefficient-representatives-along}), so cosupports of their
pullbacks are preimages.
\end{proof}

\begin{theorem}[The glueing lifts are an isomorphism of the restricted segments]
\label{mc:glueing-lift-tower-isomorphism}
\uses{mc:glueing-lift-along,mc:lift-transport-along,mc:contact-independence-along,mc:restricted-coefficient-lift-along,hir:tower-isomorphism,hir:open-restriction-tower}
\source{wlodarczyk05-hironaka}
\lean{Hironaka.BlowupTowerIsomorphism.exists_isGlueingLiftAlong_isTowerIso,Hironaka.BlowupTowerIsomorphism.exists_isGlueingLiftAlong_isTowerIso_hom,Hironaka.BlowupTowerIsomorphism.exists_isTowerIso_coeffHom,Hironaka.BlowupTowerIsomorphism.contact_independence_along_isTowerIso}
\leanok
In the setting of Lemma~\ref{mc:lift-transport-along}, along every tower \(T\) that is a multiple test blow-up of \(\underline{\mathcal H}\) and of every \(\underline{\mathcal I}_i\),
\(\underline{\mathcal I}'_i\), there is a lift \(\Phi_T\) as in Theorem~\ref{mc:glueing-lift-along} transporting every \((\underline{\mathcal I}_i)_T\) to \((\underline{\mathcal I}'_i)_T\) which
is the last-stage biholomorphism of an isomorphism \(T|_U\to T|_{U'}\) of the restricted towers over the glueing biholomorphism \(\Phi\colon U\to U'\)
(Definition~\ref{hir:tower-isomorphism}). For \(\overline{\underline{\mathcal J}}\) the lift of Theorem~\ref{mc:contact-independence-along} can be chosen to be such an isomorphism,
with its four properties; in particular \(\Phi_T^*(\underline{\mathcal S}_{u,T}(\overline{\underline{\mathcal J}}_T)|_{U'_T})=\underline{\mathcal S}_{v,T}(\overline{\underline{\mathcal J}}_T)|_{U_T}\)
and \(\Phi_T\) fixes \(\operatorname{cosupp}\overline{\underline{\mathcal J}}_T\cap U_T\) pointwise.
\end{theorem}
\begin{proof}
\leanok
Induction on \(T\), as in Lemma~\ref{mc:lift-transport-along}, carrying the isomorphism. Over the empty tower the lift is \(\Phi\), the isomorphism of the empty towers.
At a step with centre \(C\), the restricted towers are extended by the blowups of the open preimages along \(C\cap U_T\) and \(C\cap U'_T\)
(Lemma~\ref{hir:open-restriction-tower}); the new lift \(\Phi_{T'}\) satisfies \(\pi\circ\Phi_{T'}=\Phi_T\circ\pi\) and \(\Phi_T^{-1}(C\cap U'_T)=C\cap U_T\) (Lemma~\ref{mc:lift-transport-along}), which are the conditions of Definition~\ref{hir:tower-isomorphism} for the new stage. For
\(\overline{\underline{\mathcal J}}\), the proof of Theorem~\ref{mc:contact-independence-along} applies to this lift (Lemma~\ref{mc:restricted-coefficient-lift-along}).
\end{proof}
This is the form in which Step 1b of Proposition~\ref{hir:maximal-order-step} reads the independence of \(w\): the restricted segments over \(U\) and \(U'\) are isomorphic over the
glueing biholomorphism, and the restricted coefficient ideals on the strict transforms of the two hypersurfaces of maximal contact correspond under the isomorphism;
property (6) of Definition~\ref{hir:centre-function} is stated for such isomorphisms. In the source the induced towers \(\varphi^*(X_i)=X'\times_XX_i\) of Proposition~pr:funct
are defined up to isomorphism, with lifts \(\varphi_i\) commuting with the blow-ups, and Step 1bb compares \(\varphi_{pu}^*\) and \(\varphi_{pv}^*\) of the same data.

\subsection{Isomorphisms of towers along concatenations, and the glueing lifts along a tail}

In this subsection towers and their isomorphisms are as in the previous subsection. For a tower \(T_0\) over \(X\) and a tower \(T_1\) over the last stage of
\(T_0\), \(T_0T_1\) denotes the concatenation (first the blowups of \(T_0\), then those of \(T_1\)); its projection is \(p_{T_0T_1}=p_{T_0}\circ p_{T_1}\)
(Lemma~\ref{hir:tower-sequence}).

\begin{lemma}[Identities, inverses, composites and concatenations of isomorphisms of towers]
\label{mc:tower-isomorphism-append}
\uses{hir:tower-isomorphism,mc:tower-isomorphism-basic}
\lean{Hironaka.BlowupTowerIsomorphism.IsTowerIso.refl_self,Hironaka.BlowupTowerIsomorphism.IsTowerIso.symm,Hironaka.BlowupTowerIsomorphism.IsTowerIso.step_inv,Hironaka.BlowupTowerIsomorphism.IsTowerIso.refl_inv,Hironaka.BlowupTowerIsomorphism.IsTowerIso.trans,Hironaka.BlowupTowerIsomorphism.IsTowerIso.append}
\leanok
\begin{enumerate}
\item The identities \(\chi_i=\mathrm{id}\) form an isomorphism of every tower with itself over \(\mathrm{id}_X\).
\item If \(\chi\colon T'\to T\) is an isomorphism over \(\Phi\), then \((\chi_i^{-1})_i\colon T\to T'\) is an isomorphism over \(\Phi^{-1}\).
\item An isomorphism over \(\Phi\) of two towers of length \(r\ge1\) with last blowups \(\sigma'_r\), \(\sigma_r\) restricts to an isomorphism of the towers of length
\(r-1\) over \(\Phi\), with \(\sigma_r\circ\chi_r=\chi_{r-1}\circ\sigma'_r\) and \(\chi_{r-1}^{-1}(C_{r-1})=C'_{r-1}\); an isomorphism over \(\Phi\) of the empty towers has
last-stage map \(\Phi\).
\item If \(\chi'\colon T''\to T'\) is an isomorphism over \(\Phi\colon X''\to X'\) and \(\chi\colon T'\to T\) one over \(\Psi\colon X'\to X\), then
\((\chi_i\circ\chi'_i)_i\colon T''\to T\) is an isomorphism over \(\Psi\circ\Phi\).
\item If \(\chi^0\colon T'_0\to T_0\) is an isomorphism over \(\Phi\) with last-stage map \(\chi^0_r\) and \(\chi^1\colon T'_1\to T_1\) is an isomorphism over \(\chi^0_r\), then
\(\chi^0\) followed by \(\chi^1\) is an isomorphism \(T'_0T'_1\to T_0T_1\) over \(\Phi\) whose last-stage map is that of \(\chi^1\).
\end{enumerate}
\end{lemma}
\begin{proof}
\leanok
(1) \(\sigma\circ\mathrm{id}=\mathrm{id}\circ\sigma\) and \(\mathrm{id}^{-1}(C)=C\) at every step. (2) Induction on \(\chi\): from \(\sigma_i\circ\chi_i=\chi_{i-1}\circ\sigma'_i\), applied
at \(\chi_i^{-1}(q)\) and composed with \(\chi_{i-1}^{-1}\), \(\sigma'_i\circ\chi_i^{-1}=\chi_{i-1}^{-1}\circ\sigma_i\); and \((\chi_{i-1}^{-1})^{-1}(C'_{i-1})=\chi_{i-1}(\chi_{i-1}^{-1}(C_{i-1}))=C_{i-1}\).
(3) An isomorphism of two towers ending with blowups is, by definition, obtained by one step from an isomorphism of the shorter towers; an isomorphism of the
empty towers is the one with \(\chi_0=\Phi\). (4) Induction on \(\chi'\). If \(T''\) and \(T'\) are empty, \(T\) has the length of \(T'\) (Lemma~\ref{mc:tower-isomorphism-basic}), so
it is empty and \(\chi_0=\Psi\) by (3); the composite is the isomorphism of the empty towers over \(\Psi\circ\Phi\). At a step, \(T'\) and hence \(T\) have positive length, (3)
gives the previous stage \(\chi_{i-1}\) of \(\chi\) with \(\sigma_i\circ\chi_i=\chi_{i-1}\circ\sigma'_i\), and
\(\sigma_i\circ\chi_i\circ\chi'_i=\chi_{i-1}\circ\sigma'_i\circ\chi'_i=\chi_{i-1}\circ\chi'_{i-1}\circ\sigma''_i\), \((\chi_{i-1}\circ\chi'_{i-1})^{-1}(C_{i-1})=\chi'^{-1}_{i-1}(C'_{i-1})=C''_{i-1}\).
(5) Induction on \(\chi^1\): the stages of \(T_0T_1\) after those of \(T_0\) are those of \(T_1\), whose base is the last stage of \(T_0\), where \(\chi^1_0=\chi^0_r\); the
conditions at these stages are those of \(\chi^1\).
\end{proof}

\begin{lemma}[Restriction of a concatenation]
\label{mc:tower-isomorphism-restrict-append}
\uses{bcalc:tower-restrict,hir:tower-sequence,hir:tower-isomorphism,mc:tower-isomorphism-append}
\lean{Hironaka.BlowupTowerIsoTail.opensCongr,Hironaka.BlowupTowerIsoTail.coe_opensCongr_apply,Hironaka.BlowupTowerIsoTail.coe_opensCongr_symm_apply,Hironaka.BlowupTowerIsoTail.openPreimage_append,Hironaka.BlowupTowerIsoTail.isTowerIso_restrict_append}
\leanok
Let \(U\subset X\) be open, \(T_0\) a tower over \(X\) and \(T_1\) a tower over the last stage of \(T_0\). Then \(p_{T_1}^{-1}(p_{T_0}^{-1}U)=p_{T_0T_1}^{-1}U\), and the restriction
\((T_0T_1)|_U\) (Lemma~\ref{bcalc:tower-restrict}) is isomorphic over \(\mathrm{id}_U\) to the concatenation of \(T_0|_U\) and \(T_1|_{p_{T_0}^{-1}U}\), every stage map being the identity
of the common open subset of the corresponding stage of \(T_0T_1\).
\end{lemma}
\begin{proof}
\leanok
The equality of open sets is \(p_{T_0T_1}=p_{T_0}\circ p_{T_1}\) (Lemma~\ref{hir:tower-sequence}); equal open subsets of a complex manifold are biholomorphic by the identity.
Induction on \(T_1\). If \(T_1\) is empty, both towers are \(T_0|_U\) and the identity is an isomorphism (Lemma~\ref{mc:tower-isomorphism-append}(1)). At a step with centre \(C\)
and blowup \(\sigma\), both towers are extended by the restriction of \(\sigma\) to the open preimage, along the trace of \(C\) (Lemma~\ref{bcalc:tower-restrict}); by the equality
just proved the two open preimages are the same open subset of the new stage, the identity commutes with the two restrictions of \(\sigma\), and it maps the trace of \(C\) onto
the trace of \(C\).
\end{proof}

\begin{theorem}[The glueing lifts along the tail of a segment]
\label{mc:glueing-lift-tail}
\uses{mc:glueing-lift-tower-isomorphism,mc:lift-transport-along,mc:lift-step-pair,mc:glueing-lift-along,mc:contact-independence-along,mc:tower-isomorphism-restrict-append,mc:tower-isomorphism-append,hir:tower-isomorphism,hir:open-restriction-tower,hir:geo-transform-basic}
\source{wlodarczyk05-hironaka}
\lean{Hironaka.BlowupTowerIsoTail.exists_isGlueingLiftAlong_tail,Hironaka.BlowupTowerIsoTail.exists_isTowerIso_tail_coeffHom,Hironaka.BlowupTowerIsoTail.contact_independence_along_tail}
\leanok
In the setting of Lemma~\ref{mc:lift-transport-along}, let \(T_0\) be a tower over \(M\), \(U_{T_0}=p_{T_0}^{-1}U\), \(U'_{T_0}=p_{T_0}^{-1}U'\), and \(\Phi_{T_0}\) a lift along
\(T_0\) as in Theorem~\ref{mc:glueing-lift-along} which is the last-stage map of an isomorphism \(T_0|_U\to T_0|_{U'}\) over \(\Phi\) and transports every
\((\underline{\mathcal I}_i)_{T_0}\) to \((\underline{\mathcal I}'_i)_{T_0}\). For every tower \(T_1\) over the last stage of \(T_0\) such that \(T_0T_1\) is a multiple test blow-up of
\(\underline{\mathcal H}\) and of every \(\underline{\mathcal I}_i\), \(\underline{\mathcal I}'_i\), there are
\begin{enumerate}
\item a lift \(\Phi_{T_0T_1}\) along \(T_0T_1\) as in Theorem~\ref{mc:glueing-lift-along}, the last-stage map of an isomorphism \((T_0T_1)|_U\to(T_0T_1)|_{U'}\) over \(\Phi\),
transporting every \((\underline{\mathcal I}_i)_{T_0T_1}\) to \((\underline{\mathcal I}'_i)_{T_0T_1}\), and
\item an isomorphism \(\chi\colon T_1|_{U_{T_0}}\to T_1|_{U'_{T_0}}\) over \(\Phi_{T_0}\) whose last-stage map agrees with \(\Phi_{T_0T_1}\) at every point.
\end{enumerate}
For \(\overline{\underline{\mathcal J}}\) and towers \(T_0\), \(T_1\) with \(T_0T_1\) a multiple test blow-up of \(\overline{\underline{\mathcal J}}\), there are a lift \(\Phi_{T_0}\) along \(T_0\)
with the properties of Theorem~\ref{mc:glueing-lift-tower-isomorphism} at \(T_0\) (in particular
\(\Phi_{T_0}^*(\underline{\mathcal S}_{u,T_0}(\overline{\underline{\mathcal J}}_{T_0})|_{U'_{T_0}})=\underline{\mathcal S}_{v,T_0}(\overline{\underline{\mathcal J}}_{T_0})|_{U_{T_0}}\), with and
without \(\mathcal C\), and \(\Phi_{T_0}\) fixes \(\operatorname{cosupp}\overline{\underline{\mathcal J}}_{T_0}\cap U_{T_0}\) pointwise) and an isomorphism
\(\chi\colon T_1|_{U_{T_0}}\to T_1|_{U'_{T_0}}\) over \(\Phi_{T_0}\) whose last-stage map fixes \(\operatorname{cosupp}\overline{\underline{\mathcal J}}_{T_0T_1}\cap U_{T_0T_1}\) pointwise.
For two tangent directions at a point \(x\) of the cosupport, glueing data with \(x\in U\cap U'\) can be chosen with this property for all \(T_0\), \(T_1\).
\end{theorem}
\begin{proof}
\leanok
Induction on \(T_1\), carrying \(\Phi_{T_0T_1}\) and \(\chi\). For the empty \(T_1\) take \(\Phi_{T_0}\) and the isomorphism of the empty towers over \(\Phi_{T_0}\). At a step with centre \(C\)
and blowup \(\sigma\), Lemma~\ref{mc:lift-transport-along}, applied along \(T_0T_1\) to the lift \(\Phi_{T_0T_1}\), gives \(\Phi_{T_0T_1}^{-1}(C\cap U'_{T_0T_1})=C\cap U_{T_0T_1}\) and a lift
\(\Phi'\) with \(\sigma\circ\Phi'=\Phi_{T_0T_1}\circ\sigma\) and the properties of Theorem~\ref{mc:glueing-lift-along}; as in Theorem~\ref{mc:glueing-lift-tower-isomorphism} these
are the conditions extending the isomorphism of the whole restricted towers, and Lemma~\ref{mc:lift-step-pair} carries the family to the new stage (the transforms have closed
submanifolds and coherent ideals, Lemma~\ref{hir:open-restriction-tower}). By Lemma~\ref{mc:tower-isomorphism-restrict-append} the new stages of \(T_1|_{U_{T_0}}\) and
\(T_1|_{U'_{T_0}}\) are \(\sigma^{-1}(U_{T_0T_1})\) and \(\sigma^{-1}(U'_{T_0T_1})\), so \(\Phi'\), read between them, is a biholomorphism \(\chi'\) agreeing with \(\Phi'\) at every point.
Since the previous \(\chi\) agrees with \(\Phi_{T_0T_1}\) pointwise, \(\sigma\circ\chi'=\chi\circ\sigma\) on the restricted stages, and
\(\chi^{-1}(C\cap U'_{T_0T_1})=\Phi_{T_0T_1}^{-1}(C\cap U'_{T_0T_1})=C\cap U_{T_0T_1}\): these are the conditions of Definition~\ref{hir:tower-isomorphism} for the new stage of the tails.
For \(\overline{\underline{\mathcal J}}\): \(T_0\) is a multiple test blow-up of \(\overline{\underline{\mathcal J}}\) because \(T_0T_1\) is (Lemma~\ref{hir:geo-transform-basic}), and
Theorem~\ref{mc:glueing-lift-tower-isomorphism} at \(T_0\) gives \(\Phi_{T_0}\) with its properties. As \(\overline{\underline{\mathcal J}}\) and \(\underline{\mathcal H}\) have the same multiple
test blow-ups and cosupports along them (Lemma~\ref{mc:contact-independence-along}), the first part applies with the empty family; the last-stage map of \(\chi\) agrees with
\(\Phi_{T_0T_1}\), which fixes \(\operatorname{cosupp}\underline{\mathcal H}_{T_0T_1}=\operatorname{cosupp}\overline{\underline{\mathcal J}}_{T_0T_1}\) pointwise
(Theorem~\ref{mc:glueing-lift-along}). The last statement takes the glueing data of Lemma~\ref{mc:contact-independence-along}.
\end{proof}
This is the form in which Step 1b of Proposition~\ref{hir:maximal-order-step} uses the independence of \(w\): the restricted marked ideal \(\underline{\mathcal S}_{w}\) is formed at
the stage \(T_0\) where Step 1b starts, and its invariant is computed along the tail \(T_1\) of the segment. Property (6) of Definition~\ref{hir:centre-function}, over
\(\Phi_{T_0}\) for \(\underline{\mathcal S}_{u,T_0}|_{U'_{T_0}}\) and the isomorphism \(\chi\) of the restricted tails, gives
\(\operatorname{inv}_{\underline{\mathcal S}_{v,T_0}|_{U_{T_0}}}(T_1|_{U_{T_0}})(x)=\operatorname{inv}_{\underline{\mathcal S}_{u,T_0}|_{U'_{T_0}}}(T_1|_{U'_{T_0}})(\chi_r(x))\) on the cosupport, and
\(\chi_r(x)=x\) there; property (5) passes between the tails and their restrictions. The restriction of \(T_0T_1\) to \(U\) is not literally the concatenation of the restrictions
(Lemma~\ref{mc:tower-isomorphism-restrict-append}), which is why the tail isomorphism is constructed directly rather than cut out of the isomorphism of
Theorem~\ref{mc:glueing-lift-tower-isomorphism}. In the source (Step 1bb) the lifts \(\varphi_{pu}\), \(\varphi_{pv}\) at the stage \(p\) where the tangent directions are chosen
compare the restrictions \(\overline{\mathcal J}_{p|V(u)_p}\), \(\overline{\mathcal J}_{p|V(v)_p}\), whose canonical resolutions start at that stage.

\subsection{Reduction to a hypersurface of maximal contact at a later stage of a tower}

In this subsection \(\underline{\mathcal I}=(X,N,E,\mathfrak a,d)\) has \(N\) a closed submanifold, \(\mathfrak a\) coherent and \(d\ge1\). For a tower \(T\) over \(X\) and
\(A\subset X\), \(A_T\) is the strict transform of \(A\) along \(T\) (Definition~\ref{mc:strict-along}), and for an open set \(U\subset X\), \(U_T=p_T^{-1}U\).

\begin{lemma}[Strict transforms along towers are local over the base]
\label{mc:strict-along-local}
\uses{mc:strict-along,mc:strict-transform-local}
\lean{Hironaka.ContactEquivalenceStage.strictAlong_inter_preimage_eq}
\leanok
If \(A\cap O=B\cap O\) for an open set \(O\subset X\), then \(A_T\cap O_T=B_T\cap O_T\) for every tower \(T\) over \(X\).
\end{lemma}
\begin{proof}
\leanok
Induction on \(T\). At a step with blowup \(\sigma\) along \(C\) after the tower \(T'\), the sets \(A_{T'}\) and \(B_{T'}\) agree on the open set \(O_{T'}\), so their strict
transforms agree on \(\sigma^{-1}(O_{T'})=O_T\) (Lemma~\ref{mc:strict-transform-local}).
\end{proof}

\begin{lemma}[Giraud from local hypersurfaces of maximal contact]
\label{mc:giraud-local}
\uses{hir:giraud-tower-along,mc:tangent-direction-along,mc:coefficient-restriction-full-along,mc:restrict-marked,bcalc:member-intersection-tower,mc:strict-along,hir:marked-equivalence}
\source{wlodarczyk05-hironaka}
\lean{Hironaka.ContactEquivalenceStage.centresIn_of_cosupp_subset,Hironaka.ContactEquivalenceStage.cosupp_subset_strictAlong_of_forall_exists,Hironaka.ContactEquivalenceStage.wEquiv_withN_coeffMarkedAlong_of_cosupp_subset,Hironaka.ContactEquivalenceStage.cosupp_subset_of_wEquiv}
\leanok
Let \(S\subset X\).
\begin{enumerate}
\item If every point of \(N\) has an open neighbourhood \(O\) carrying a tangent direction \(w\) of \((\mathfrak a|_N,d)\) along \(N\) on \(O\)
(Definition~\ref{mc:tangent-direction-along}) with \(N\cap V(w)\cap O\subset S\), then \(\operatorname{cosupp}\underline{\mathcal I}_T\subset S_T\) along every multiple test blow-up
\(T\) of \(\underline{\mathcal I}\).
\item If \(\operatorname{cosupp}\underline{\mathcal I}_T\subset S_T\) along every multiple test blow-up \(T\), then the centres of every multiple test blow-up lie in the strict
transforms of \(S\), and the same inclusions hold for every marked ideal equivalent to \(\underline{\mathcal I}\) (Definition~\ref{hir:marked-equivalence}).
\item If moreover \(S\) is a closed submanifold contained in \(N\), then \(\underline{\mathcal I}\simeq\mathcal C_N(\underline{\mathcal I})|_S\).
\end{enumerate}
\end{lemma}
\begin{proof}
\leanok
(1) A point \(z\) of \(\operatorname{cosupp}\underline{\mathcal I}_T\) lies on \(N_T\), the strict transform of the closed set \(N\), which lies over \(N\)
(Lemma~\ref{bcalc:member-intersection-tower}); take \(O\) and \(w\) at \(p_T(z)\). By Theorem~\ref{hir:giraud-tower-along},
\(\operatorname{cosupp}\underline{\mathcal I}_T\cap O_T\subset(N\cap V(w)\cap O)_T\), which lies in \(S_T\) since strict transforms are monotone
(Lemma~\ref{bcalc:member-intersection-tower}). (2) The centre of the step after a part \(T'\) of a multiple test blow-up is an admissible centre of
\(\underline{\mathcal I}_{T'}\), hence lies in \(\operatorname{cosupp}\underline{\mathcal I}_{T'}\subset S_{T'}\). Equivalent marked ideals have the same multiple test blow-ups
and the same cosupports along them. (3) By Theorem~\ref{mc:coefficient-restriction-full-along}, a tower is a multiple test blow-up of \(\mathcal C_N(\underline{\mathcal I})|_S\)
if and only if it is one of \(\underline{\mathcal I}\) with centres in the strict transforms of \(S\) (Definition~\ref{mc:restrict-marked}), which by (2) is every multiple
test blow-up of \(\underline{\mathcal I}\), and along it \(\operatorname{cosupp}[\mathcal C_N(\underline{\mathcal I})|_S]_T=\operatorname{cosupp}\underline{\mathcal I}_T\cap S_T=\operatorname{cosupp}\underline{\mathcal I}_T\).
Both marked ideals have the divisor \(E\).
\end{proof}
This is Lemma~\ref{mc:contact-equivalence-along} with the hypersurface of maximal contact replaced by any closed submanifold in whose strict transforms the
cosupports lie, which is the form in which it applies at a later stage of a tower, where the hypersurfaces of maximal contact are only given locally.

\begin{lemma}[Tangent directions at a later stage]
\label{mc:tangent-direction-stage}
\uses{hir:giraud-tower-along,mc:tangent-direction-along-restricted,mc:tangent-direction-along-open,mc:strict-along-local,mc:contact-hypersurface-along,hir:open-restriction-tower,lag:regular-point-basics,rem:open-subset-analytic}
\source{wlodarczyk05-hironaka}
\lean{Hironaka.ContactEquivalenceStage.exists_isTangentDirectionAlong_stage,Hironaka.ContactEquivalenceStage.isClosedSubmanifold_strictAlong_stage,Hironaka.ContactEquivalenceStage.forall_exists_isTangentDirectionAlong_stage_val}
\leanok
Let \(u\) be a tangent direction of \((\mathfrak a|_N,d)\) along \(N\) on an open set \(U\), \(T_0\) a multiple test blow-up of \(\underline{\mathcal I}\), and
\(S=(N\cap V(u))_{T_0}\). Every point \(z\in N_{T_0}\cap U_{T_0}\) has an open neighbourhood \(W\subset U_{T_0}\) carrying a tangent direction \(w\) of
\((\mathfrak a_{T_0}|_{N_{T_0}},d)\) along \(N_{T_0}\) on \(W\) with \(N_{T_0}\cap V(w)\cap W=S\cap W\); restricted to \(U_{T_0}\), these are tangent directions of
\(\underline{\mathcal I}_{T_0}|_{U_{T_0}}\). The set \(S\cap U_{T_0}\) is a closed submanifold of \(U_{T_0}\).
\end{lemma}
\begin{proof}
\leanok
Theorem~\ref{hir:giraud-tower-along} gives an open set \(O\ni z\) over \(U\) and a function \(w_0\), holomorphic on \(O\), at every point of \(N_{T_0}\cap O\) equal on \(N_{T_0}\) to a
combination of germs of \(\mathcal D_{N_{T_0}}^{d-1}(\mathfrak a_{T_0})\), with a field tangent to \(N_{T_0}\) not annihilating \(w_0\) at each of its zeros on \(N_{T_0}\), and with
\(N_{T_0}\cap V(w_0)\cap O=(N\cap V(u)\cap U)_{T_0}\cap O\). Lemma~\ref{mc:tangent-direction-along-restricted} replaces \(w_0\) near \(z\) by a tangent direction \(w\) on an open
\(W\subset O\) equal to \(w_0\) on \(N_{T_0}\cap W\). Since \((N\cap V(u)\cap U)\cap U=(N\cap V(u))\cap U\), Lemma~\ref{mc:strict-along-local} gives
\((N\cap V(u)\cap U)_{T_0}\cap U_{T_0}=S\cap U_{T_0}\), hence \(N_{T_0}\cap V(w)\cap W=S\cap W\). Tangent directions restrict to open sets
(Lemma~\ref{mc:tangent-direction-along-open}). If \(y\in U_{T_0}\) is a limit of points of \(S\cap U_{T_0}\), then \(y\in N_{T_0}\), which is closed
(Lemma~\ref{hir:open-restriction-tower}); for \(W\), \(w\) at \(y\), \(w\) vanishes at the points of \(S\cap W\) accumulating at \(y\), so \(w(y)=0\) by continuity and
\(y\in N_{T_0}\cap V(w)\cap W=S\cap W\). At a point of \(S\cap U_{T_0}\), the set \(N_{T_0}\cap V(w)\cap W\) is regular (Lemma~\ref{mc:contact-hypersurface-along}) and agrees with
\(S\) on the open set \(W\), so \(S\) is regular there (Lemma~\ref{lag:regular-point-basics}), also as a subset of \(U_{T_0}\) (Lemma~\ref{rem:open-subset-analytic}).
\end{proof}

\begin{theorem}[Reduction to a hypersurface of maximal contact at a later stage]
\label{mc:contact-equivalence-stage}
\uses{mc:giraud-local,mc:tangent-direction-stage,hir:open-restriction-tower,hir:open-restriction-marked,bcalc:restrict-open,mc:homogenized-along-etale,hir:marked-equivalence-basic,hir:geo-transform-basic,mc:strict-along,bcalc:member-intersection-tower}
\source{wlodarczyk05-hironaka}
\lean{Hironaka.ContactEquivalenceStage.wEquiv_stage_of_wEquiv,Hironaka.ContactEquivalenceStage.wEquiv_stage,Hironaka.ContactEquivalenceStage.wEquiv_pullback_val_geoTransformAlong,Hironaka.ContactEquivalenceStage.isGeoTestTower_tail_of_wEquiv}
\leanok
In the setting of Lemma~\ref{mc:tangent-direction-stage} write \(V=U_{T_0}\).
\begin{enumerate}
\item \(\underline{\mathcal I}_{T_0}|_V\simeq(\mathcal C_{N_{T_0}}(\underline{\mathcal I}_{T_0})|_S)|_V\).
\item Let \(\underline{\mathcal I}'\) have the same \(N\), a coherent ideal and a mark \(\ge1\), let \(T_0\) be a multiple test blow-up of \(\underline{\mathcal I}'\), and assume
\(\underline{\mathcal I}'_{T_0}|_V\simeq\underline{\mathcal I}_{T_0}|_V\). Then \(\underline{\mathcal I}'_{T_0}|_V\simeq(\mathcal C_{N_{T_0}}(\underline{\mathcal I}'_{T_0})|_S)|_V\). The assumption
holds when \(\underline{\mathcal I}'|_U\simeq\underline{\mathcal I}|_U\) and \(T_0\) is a multiple test blow-up of both.
\item If a marked ideal \(\underline{\mathcal S}\) on \(V\) is equivalent to \(\underline{\mathcal I}'_{T_0}|_V\) and \(T_0T_1\) is a multiple test blow-up of \(\underline{\mathcal I}'\), then
\(T_1|_V\) is a multiple test blow-up of \(\underline{\mathcal S}\), and \(\operatorname{cosupp}\underline{\mathcal S}_{T_1|_V}=\operatorname{cosupp}\underline{\mathcal I}'_{T_0T_1}\cap V_{T_1}\).
\end{enumerate}
\end{theorem}
\begin{proof}
\leanok
(2) The marked ideals \(\underline{\mathcal I}_{T_0}|_V\) and \(\underline{\mathcal I}'_{T_0}|_V\) have the closed submanifold \(N_{T_0}\cap V\), coherent ideals and marks \(\ge1\)
(Lemmas~\ref{hir:open-restriction-tower} and~\ref{bcalc:restrict-open}). By Lemma~\ref{mc:tangent-direction-stage} and Lemma~\ref{mc:giraud-local}(1), along every multiple test
blow-up of \(\underline{\mathcal I}_{T_0}|_V\) the cosupports lie in the strict transforms of \(S\cap V\); by Lemma~\ref{mc:giraud-local}(2) the same holds for
\(\underline{\mathcal I}'_{T_0}|_V\). The set \(S\cap V\) is a closed submanifold of \(V\) (Lemma~\ref{mc:tangent-direction-stage}) contained in \(N_{T_0}\), the strict transform of
\(N\supset N\cap V(u)\) (strict transforms are monotone, Lemma~\ref{bcalc:member-intersection-tower}). Lemma~\ref{mc:giraud-local}(3) gives
\(\underline{\mathcal I}'_{T_0}|_V\simeq\mathcal C_{N_{T_0}\cap V}(\underline{\mathcal I}'_{T_0}|_V)|_{S\cap V}\), and \(\mathcal C\) commutes with restriction to open sets
(Lemma~\ref{mc:homogenized-along-etale}). For the last statement, \(T_0|_U\) is a multiple test blow-up of \(\underline{\mathcal I}'|_U\), the equivalence passes to the
transforms along it (Lemma~\ref{hir:marked-equivalence-basic}), and the transforms along \(T_0|_U\) are the restrictions to \(V\) of the transforms along \(T_0\)
(Lemma~\ref{hir:open-restriction-tower}). (1) is (2) with \(\underline{\mathcal I}'=\underline{\mathcal I}\). (3) \(T_1\) is a multiple test blow-up of
\(\underline{\mathcal I}'_{T_0}\) (Lemma~\ref{hir:geo-transform-basic}); its restriction to \(V\) is one of \(\underline{\mathcal I}'_{T_0}|_V\), whose transform along it is the restriction
of \(\underline{\mathcal I}'_{T_0T_1}\) (Lemma~\ref{hir:open-restriction-tower}), with the trace of \(\operatorname{cosupp}\underline{\mathcal I}'_{T_0T_1}\) as cosupport
(Lemma~\ref{hir:open-restriction-marked}); equivalent marked ideals have the same multiple test blow-ups and the same cosupports along them.
\end{proof}

\begin{corollary}[The restricted coefficient ideal of Step 1b at a later stage]
\label{mc:contact-equivalence-stage-coeff-hom}
\uses{mc:contact-equivalence-stage,mc:homogenized-equivalent-along,mc:homogenized-along-etale,mc:contact-independence-along,mc:glueing-lift-tail,hir:open-restriction-marked,bcalc:restrict-open}
\source{wlodarczyk05-hironaka}
\lean{Hironaka.ContactEquivalenceStage.homMarkedAlong_mk,Hironaka.ContactEquivalenceStage.wEquiv_pullback_val_coeffHom,Hironaka.ContactEquivalenceStage.wEquiv_stage_coeffHom,Hironaka.ContactEquivalenceStage.exists_wEquiv_stage_coeffHom,Hironaka.ContactEquivalenceStage.wEquiv_stage_coeffHom_of_isTangentDirectionAlong_coeffHom,Hironaka.ContactEquivalenceStage.isGeoTestTower_tail_coeffHom}
\leanok
Let \(\mathfrak a\) be coherent on \(M\), \(N\) a closed submanifold, \(\mu\ge1\), \(\overline{\underline{\mathcal J}}=\mathcal C_N(\mathcal H_N(M,N,\emptyset,\mathfrak a,\mu))\), \(T_0\) a
multiple test blow-up of \(\overline{\underline{\mathcal J}}\), \(U\subset M\) open, \(V=U_{T_0}\), and for a function \(u\) let
\(\underline{\mathcal S}_{u,T_0}=\mathcal C_{N_{T_0}}(\overline{\underline{\mathcal J}}_{T_0})|_{(N\cap V(u))_{T_0}}\). If \(u\) is a tangent direction along \(N\) on \(U\) of
\((\mathfrak a|_N,\mu)\), or of \(\overline{\underline{\mathcal J}}\), then \(\overline{\underline{\mathcal J}}_{T_0}|_V\simeq\underline{\mathcal S}_{u,T_0}|_V\). In the first case, for every
tower \(T_1\) such that \(T_0T_1\) is a multiple test blow-up of \(\overline{\underline{\mathcal J}}\), \(T_1|_V\) is a multiple test blow-up of \(\underline{\mathcal S}_{u,T_0}|_V\) and
\(\operatorname{cosupp}(\underline{\mathcal S}_{u,T_0}|_V)_{T_1|_V}=\operatorname{cosupp}\overline{\underline{\mathcal J}}_{T_0T_1}\cap V_{T_1}\).
\end{corollary}
\begin{proof}
\leanok
\(\mathcal H_N\) of \((M,N,\emptyset,\mathfrak a,\mu)\) is the homogenized marked ideal of Theorem~\ref{mc:contact-independence-along}, and \(\overline{\underline{\mathcal J}}\) has a closed
submanifold, a coherent ideal and a mark \(\ge1\) there. As \(\mathcal C_N\) and \(\mathcal H_N\) commute with restriction to open sets (Lemma~\ref{mc:homogenized-along-etale}),
\(\overline{\underline{\mathcal J}}|_U=\mathcal C_{N\cap U}(\mathcal H_{N\cap U}((\mathfrak a,\mu)|_U))\), which is equivalent to \((\mathfrak a,\mu)|_U\) (Theorem~\ref{mc:homogenized-equivalent-along}
on \(U\)); on \(M\), \(\overline{\underline{\mathcal J}}\simeq(\mathfrak a,\mu)\), so \(T_0\) is a multiple test blow-up of \((M,N,\emptyset,\mathfrak a,\mu)\) as well.
Theorem~\ref{mc:contact-equivalence-stage}(2), with \(\underline{\mathcal I}=(M,N,\emptyset,\mathfrak a,\mu)\) and \(\underline{\mathcal I}'=\overline{\underline{\mathcal J}}\), gives the first case,
and Theorem~\ref{mc:contact-equivalence-stage}(3) the tails. The second case is Theorem~\ref{mc:contact-equivalence-stage}(1) with \(\underline{\mathcal I}=\overline{\underline{\mathcal J}}\).
\end{proof}
In Step 1b of Proposition~\ref{hir:maximal-order-step} the tangent directions are those of \((\mathfrak a|_N,\mu)\) at the start of Step 1, the ones glued by the glueing
data of Theorem~\ref{mc:glueing-lift-tail}, which homogenize the ideal whose tangent directions they glue; the source (Step 1bb) takes a tangent direction
\(u\in\mathcal D^{\overline\mu-1}(\overline{\mathcal J})\), the second case. The corollary supplies the hypotheses on \(\underline{\mathcal S}_{u,T_0}\) under which property (6) of
Definition~\ref{hir:centre-function} applies along the tails: by Theorem~\ref{mc:glueing-lift-tail}, the last-stage map of the tail isomorphism fixes
\(\operatorname{cosupp}\overline{\underline{\mathcal J}}_{T_0T_1}\) pointwise, so it maps points of that cosupport into the cosupport of \((\underline{\mathcal S}_{u,T_0}|_V)_{T_1|_V}\). In the
source this is the statement that \((V(u)_i)_{0\le i\le p}\) is a multiple test blow-up of \((V(u),\overline{\mathcal J}|_{V(u)},\emptyset,\overline\mu)\), so that the induced marked
ideal at the stage \(p\) is \(\overline{\mathcal J}_{p|V(u)_p}\) (by Lemma ``le:codim2'').

\begin{corollary}[The restricted coefficient ideal of Step 1b at a later stage, with a divisor]
\label{mc:contact-equivalence-stage-coeff-hom-divisor}
\uses{mc:contact-equivalence-stage,mc:homogenized-equivalent-along,mc:homogenized-along-etale,mc:contact-equivalence-stage-coeff-hom}
\source{wlodarczyk05-hironaka}
\lean{Hironaka.ContactEquivalenceStage.wEquiv_pullback_val_coeffHomAlong,Hironaka.ContactEquivalenceStage.wEquiv_stage_coeffHomAlong,Hironaka.ContactEquivalenceStage.exists_wEquiv_stage_coeffHomAlong,Hironaka.ContactEquivalenceStage.isGeoTestTower_tail_coeffHomAlong}
\leanok
Corollary~\ref{mc:contact-equivalence-stage-coeff-hom} holds, in its first case, for \(\overline{\underline{\mathcal J}}=\mathcal C_N(\mathcal H_N(\underline{\mathcal I}))\) and any marked ideal
\(\underline{\mathcal I}=(X,N,E,\mathfrak a,d)\) with \(N\) a closed submanifold, \(\mathfrak a\) coherent and \(d\ge1\), whatever its divisor \(E\): for a tangent direction \(u\) of
\((\mathfrak a|_N,d)\) along \(N\) on an open set \(U\) and a multiple test blow-up \(T_0\) of \(\overline{\underline{\mathcal J}}\),
\(\overline{\underline{\mathcal J}}_{T_0}|_V\simeq\underline{\mathcal S}_{u,T_0}|_V\) on \(V=U_{T_0}\), both with the divisor \(E_{T_0}\) of \(\overline{\underline{\mathcal J}}_{T_0}\), and the tails
\(T_1|_V\) of the multiple test blow-ups \(T_0T_1\) of \(\overline{\underline{\mathcal J}}\) are multiple test blow-ups of \(\underline{\mathcal S}_{u,T_0}|_V\) with
\(\operatorname{cosupp}(\underline{\mathcal S}_{u,T_0}|_V)_{T_1|_V}=\operatorname{cosupp}\overline{\underline{\mathcal J}}_{T_0T_1}\cap V_{T_1}\).
\end{corollary}
\begin{proof}
\leanok
The proof of Corollary~\ref{mc:contact-equivalence-stage-coeff-hom} does not use \(E=\emptyset\): \(\mathcal C_N\) and \(\mathcal H_N\) commute with restriction to \(U\)
(Lemma~\ref{mc:homogenized-along-etale}) and \(\mathcal C_N(\mathcal H_N(\underline{\mathcal I}))\simeq\underline{\mathcal I}\) for any \(E\) (Theorem~\ref{mc:homogenized-equivalent-along}), on
\(U\) and on \(X\); Theorem~\ref{mc:contact-equivalence-stage}(2)--(3) applies with \(\underline{\mathcal I}'=\overline{\underline{\mathcal J}}\).
\end{proof}
In Step 1 the divisor of the marked ideal of maximal order may be kept, its original members being masked by restricting to open sets that miss them once
Step 1a has separated them from the cosupport; this corollary is the equivalence at the stage \(T_0\) in that form. The glueing lifts of
Theorem~\ref{mc:glueing-lift-tail} are for \(E=\emptyset\): the glueing data do not preserve the members of a divisor.

\subsection{The restricted marked ideal of Step 1b in the domain of the inductive invariant}

In this subsection \(\overline{\underline{\mathcal J}}=\mathcal C_N(\mathcal H_N(M,N,\emptyset,\mathfrak a,\mu))\) with \(\mathfrak a\) coherent, \(N\) a closed submanifold and \(\mu\ge1\),
\(u\) is a tangent direction of \((\mathfrak a|_N,\mu)\) along \(N\) on an open set \(U\), \(T_0\) is a multiple test blow-up of \(\overline{\underline{\mathcal J}}\),
\(V=U_{T_0}\), \(S=(N\cap V(u))_{T_0}\) and \(\underline{\mathcal S}_{u,T_0}=\mathcal C_{N_{T_0}}(\overline{\underline{\mathcal J}}_{T_0})|_S\), as in
Corollary~\ref{mc:contact-equivalence-stage-coeff-hom}.

\begin{lemma}[Dimension of the hypersurface of maximal contact at a later stage]
\label{mc:contact-stage-dimension}
\uses{mc:tangent-direction-stage,bcalc:contact-restriction-nonzero,bcalc:inv-domain-stable,lag:regular-point-basics,rem:open-subset-analytic,hir:centre-function}
\source{wlodarczyk05-hironaka}
\lean{Hironaka.ContactStageInvDomain.exists_regularAtDim_strictAlong_stage,Hironaka.ContactStageInvDomain.dimLE_strictAlong_stage}
\leanok
Let \(\underline{\mathcal I}=(X,N,E,\mathfrak a,d)\) have \(N\) a closed submanifold, \(\mathfrak a\) coherent and \(d\ge1\), let \(u\) be a tangent direction of
\((\mathfrak a|_N,d)\) along \(N\) on an open set \(U\) and \(T_0\) a multiple test blow-up of \(\underline{\mathcal I}\). At every point of \((N\cap V(u))_{T_0}\) over \(U\),
\(N_{T_0}\) is regular of some dimension \(j+1\) and \((N\cap V(u))_{T_0}\) of dimension \(j\). Hence, if \(\dim N\le k\), then \((N\cap V(u))_{T_0}\cap U_{T_0}\) has dimension
\(\le k-1\) at each of its points.
\end{lemma}
\begin{proof}
\leanok
By Lemma~\ref{mc:tangent-direction-stage}, near such a point \(y\), \((N\cap V(u))_{T_0}\) agrees with \(N_{T_0}\cap V(w)\cap W\) for a tangent direction \(w\) of
\((\mathfrak a_{T_0}|_{N_{T_0}},d)\) on an open \(W\ni y\); a hypersurface of maximal contact has dimension one less than \(N_{T_0}\)
(Proposition~\ref{bcalc:contact-restriction-nonzero}), and regularity depends only on the germ of a set (Lemma~\ref{lag:regular-point-basics}). If \(\dim N\le k\), then
\(\dim N_{T_0}\le k\) along the multiple test blow-up (Proposition~\ref{bcalc:inv-domain-stable}); a point is regular of one dimension only
(Lemma~\ref{lag:regular-point-basics}), so \(j+1\le k\). Regularity in \(U_{T_0}\) is regularity in the last stage (Lemma~\ref{rem:open-subset-analytic}).
\end{proof}

\begin{theorem}[The restricted marked ideal of Step 1b is in the domain of the inductive invariant]
\label{mc:contact-stage-domain}
\uses{mc:contact-stage-dimension,mc:contact-equivalence-stage-coeff-hom,mc:contact-equivalence-stage,mc:tangent-direction-stage,mc:contact-equivalence-along,mc:tangent-direction-along-open,bcalc:contact-restriction-nonzero,bcalc:codim-one-part-canonical,bcalc:exceptional-transverse-open,hir:giraud-tower-along,hir:snc-basic,hir:snc-invariance,hir:wellformed-tower,hir:open-restriction-marked,hir:open-restriction-tower,mc:homogenized-along-basic,mc:homogenized-equivalent-along,mc:contact-independence-along,hir:centre-function}
\source{wlodarczyk05-hironaka}
\lean{Hironaka.ContactStageInvDomain.isNonzeroAlongAt_stage_of_wEquiv,Hironaka.ContactStageInvDomain.isWellFormed_coeffHom,Hironaka.ContactStageInvDomain.isTransverse_strictAlong_stage_coeffHom,Hironaka.ContactStageInvDomain.isInvDomain_stage_coeffHom}
\leanok
Assume \(\dim N\le k\) and that no point of \(S\cap V\) lies in the codimension-one part \(Z_1(\overline{\underline{\mathcal J}}_{T_0})\) of the cosupport along \(N_{T_0}\) (for
instance after Step 1ba, when \(Z_1\) is empty). Then \(\underline{\mathcal S}_{u,T_0}|_V\) lies in the domain of a resolution invariant in dimension \(\le k-1\)
(Definition~\ref{hir:centre-function}): it is well formed, \(\dim(S\cap V)\le k-1\), its mark is \(\ge1\), and its ideal vanishes identically near no point of \(S\cap V\).
More generally, for \(\underline{\mathcal I}\), \(u\), \(T_0\) as in Lemma~\ref{mc:contact-stage-dimension} and a marked ideal \(\underline{\mathcal I}'\) with the same \(N\), a coherent ideal and
\(T_0\) as a multiple test blow-up, if \(\underline{\mathcal I}'_{T_0}|_V\simeq(\mathcal C_{N_{T_0}}(\underline{\mathcal I}'_{T_0})|_S)|_V\), then the coefficient ideal of \(\underline{\mathcal I}'_{T_0}\) is
nonzero along \(S\) at every point of \(S\cap V\) outside \(Z_1(\underline{\mathcal I}'_{T_0})\).
\end{theorem}
\begin{proof}
\leanok
Nonvanishing: at \(y\in S\cap V\), \(N_{T_0}\) is regular of dimension \(j+1\) and \(S\) of dimension \(j\) (Lemma~\ref{mc:contact-stage-dimension}); the cosupport of
\(\underline{\mathcal I}'_{T_0}|_V\) lies in \(S\) and equals that of \((\mathcal C_{N_{T_0}}(\underline{\mathcal I}'_{T_0})|_S)|_V\) (the equivalence for the empty tower). If the coefficient ideal
vanished on \(S\) near \(y\), every point of \(S\) near \(y\) would lie in the cosupport, which would then agree with \(S\) near \(y\) and be regular of dimension
\(j\) there, i.e.\ \(y\in Z_1\) (Proposition~\ref{bcalc:contact-restriction-nonzero}); \(Z_1\) commutes with restriction to open sets (Proposition~\ref{bcalc:codim-one-part-canonical}).
For \(\overline{\underline{\mathcal J}}\) the equivalence is Corollary~\ref{mc:contact-equivalence-stage-coeff-hom}.
Well formedness: \(S\cap V\) is a closed submanifold of \(V\) (Lemma~\ref{mc:tangent-direction-stage}); the coefficient ideal of the coherent
\(\overline{\underline{\mathcal J}}_{T_0}\) is coherent (Lemma~\ref{mc:homogenized-along-basic}) and restricts to a coherent ideal. \(\overline{\underline{\mathcal J}}\) is well formed
(\(E=\emptyset\), Lemma~\ref{hir:snc-basic}), hence so are \(\overline{\underline{\mathcal J}}_{T_0}\) (Lemma~\ref{hir:wellformed-tower}) and its restriction to \(V\)
(Lemma~\ref{hir:open-restriction-marked}); its divisor, which is the divisor of \(\underline{\mathcal S}_{u,T_0}|_V\), has simple normal crossings. All its members are
exceptional; the centres of the multiple test blow-ups of \(\overline{\underline{\mathcal J}}\) lie in its cosupports, which equal those of \((M,N,\emptyset,\mathfrak a,\mu)\)
(Theorem~\ref{mc:homogenized-equivalent-along}) and lie over \(U\) in the strict transforms of \(N\cap V(u)\) (Theorem~\ref{hir:giraud-tower-along}); \(N\cap V(u)\cap U\) is a closed
submanifold of \(U\) (Lemmas~\ref{mc:tangent-direction-along-open} and~\ref{mc:contact-equivalence-along}). So the exceptional members are transverse to \(S\cap V\)
(Lemma~\ref{bcalc:exceptional-transverse-open}), also in \(V\) (Lemma~\ref{hir:snc-invariance}). The dimension is Lemma~\ref{mc:contact-stage-dimension}, and the mark of a
coefficient ideal is \(\ge1\) (Theorem~\ref{mc:contact-independence-along}).
\end{proof}

\begin{lemma}[Glueing data inside the sets of the two tangent directions]
\label{mc:glueing-data-shrink}
\uses{mc:tangent-direction-along,mc:glueing-trivial-along,mc:glueing-biholomorphism-along,mc:glueing-lift-tail}
\source{wlodarczyk05-hironaka}
\lean{Hironaka.ContactStageInvDomain.isTangentDirectionAlong_mono,Hironaka.ContactStageInvDomain.exists_glueingDataAlong_trivial_subset,Hironaka.ContactStageInvDomain.contact_independence_along_tail_subset}
\leanok
A tangent direction along \(N\) on an open set \(U\) is one on every open subset of \(U\). In Lemma~\ref{mc:glueing-trivial-along} and in the last statement of
Theorem~\ref{mc:glueing-lift-tail}, for tangent directions \(u\), \(v\) on open sets \(U_1\), \(U_2\) containing \(x\), the glueing data can be chosen with their open
set inside \(U_1\cap U_2\); then the domain \(U'\) and the codomain \(U\) of the glueing biholomorphism lie in \(U_1\cap U_2\), and \(u\) and \(v\) are tangent directions on
\(U\) and \(U'\).
\end{lemma}
\begin{proof}
\leanok
The conditions defining tangent directions and glueing data (Definition~\ref{mc:tangent-direction-along} and Lemma~\ref{mc:glueing-biholomorphism-along}) are pointwise
conditions on the open set, so they hold on its intersection with the open set \(U_1\cap U_2\ni x\); so does the triviality of the action modulo the glue ideal. The
domain and the codomain of the glueing biholomorphism lie in the open set of the data, and Theorem~\ref{mc:glueing-lift-tail} applies to the shrunk data.
\end{proof}
With Corollary~\ref{mc:contact-equivalence-stage-coeff-hom} for \(U\) the codomain of the glueing biholomorphism, these are the three hypotheses on
\(\underline{\mathcal S}_{u,T_0}\) under which the value of Step 1b of Proposition~\ref{hir:maximal-order-step} does not depend on the tangent direction: membership in the
domain of the inductive invariant, the tails as multiple test blow-ups (or canonical segments) of \(\underline{\mathcal S}_{u,T_0}|_V\), and the cosupport along them.
In the source the codimension-one components are blown up in Step 1ba before the tangent directions are chosen in Step 1bb, and the induced marked ideal on
\(V(u)_p\) has dimension one less.

\begin{theorem}[The restricted marked ideal of Step 1b on smaller open sets, with a masked divisor]
\label{mc:contact-stage-domain-masked}
\uses{mc:contact-stage-domain,mc:contact-stage-dimension,mc:tangent-direction-stage,mc:coefficient-restriction-along,hir:giraud-tower-along,mc:homogenized-equivalent-along,bcalc:exceptional-transverse-open,bcalc:contact-restriction-nonzero,bcalc:codim-one-part-canonical,hir:wellformed-tower,hir:open-restriction-marked,hir:snc-basic,hir:snc-invariance,hir:centre-function}
\source{wlodarczyk05-hironaka}
\lean{Hironaka.ContactStageInvDomain.isTransverse_append_of_forall_notMem,Hironaka.ContactStageInvDomain.isClosedSubmanifold_strictAlong_stage_of_le,Hironaka.ContactStageInvDomain.isInvDomain_stage_coeffHomAlong_of_le,Hironaka.ContactStageInvDomain.isInvDomain_stage_coeffHom_of_le}
\leanok
Let \(\underline{\mathcal I}=(X,N,E,\mathfrak a,d)\) be well formed with \(d\ge1\) and \(\dim N\le k\), \(\overline{\underline{\mathcal J}}=\mathcal C_N(\mathcal H_N(\underline{\mathcal I}))\), \(u\) a tangent
direction of \((\mathfrak a|_N,d)\) along \(N\) on an open set \(U\), \(T_0\) a multiple test blow-up of \(\overline{\underline{\mathcal J}}\), \(S=(N\cap V(u))_{T_0}\), and \(W\subset U_{T_0}\)
open. Assume that no point of \(S\cap W\) lies in \(Z_1(\overline{\underline{\mathcal J}}_{T_0})\) and that the strict transforms of the members of \(E\) (the first \(\#E\) members of
the divisor of \(\overline{\underline{\mathcal J}}_{T_0}\)) miss \(S\cap W\). Then \((\mathcal C_{N_{T_0}}(\overline{\underline{\mathcal J}}_{T_0})|_S)|_W\) lies in the domain of a resolution invariant in
dimension \(\le k-1\) (Definition~\ref{hir:centre-function}). For \(E=\emptyset\) the second assumption is empty, and this is Theorem~\ref{mc:contact-stage-domain} with
any open \(W\subset U_{T_0}\) in place of \(U_{T_0}\). A divisor is transverse to a set if one part of it is and the members of the other part miss the set, and \(S\cap W\) is a
closed submanifold of \(W\).
\end{theorem}
\begin{proof}
\leanok
The proof of Theorem~\ref{mc:contact-stage-domain} applies on \(W\), with two changes. Instead of the equivalence over \(U_{T_0}\), the two cosupport inclusions used for
the nonvanishing come from the last stage: \(\operatorname{cosupp}\overline{\underline{\mathcal J}}_{T_0}\) equals \(\operatorname{cosupp}\underline{\mathcal I}_{T_0}\)
(Theorem~\ref{mc:homogenized-equivalent-along}), which lies over \(U\) in \(S\) (Theorem~\ref{hir:giraud-tower-along}), and
\(\operatorname{cosupp}(\mathcal C_{N_{T_0}}(\overline{\underline{\mathcal J}}_{T_0})|_S)=\operatorname{cosupp}\overline{\underline{\mathcal J}}_{T_0}\cap S\) (Lemma~\ref{mc:coefficient-restriction-along});
both pass to \(W\) (Lemma~\ref{hir:open-restriction-marked}). For the transversality, the exceptional members of the divisor of \(\overline{\underline{\mathcal J}}_{T_0}\) are
transverse to \(S\) over \(U\) (Lemma~\ref{bcalc:exceptional-transverse-open}, with \(\overline{\underline{\mathcal J}}\) well formed because \(\underline{\mathcal I}\) is), hence to
\(S\cap W\), which agrees with \(S\cap U_{T_0}\) near each of its points. At a point \(x\in S\cap W\) take the coordinates of that transversality; the strict transforms of the
original members do not contain \(x\), so they impose no condition at \(x\), and the same coordinates exhibit the transversality to the whole divisor. Closedness and
regularity of \(S\cap W\) in \(W\) are those of Lemma~\ref{mc:tangent-direction-stage}, at the points of \(W\), and the dimension is Lemma~\ref{mc:contact-stage-dimension}.
\end{proof}
In Step 1 with the divisor of the marked ideal kept, the original members are masked by passing to an open set \(W\) on which their strict transforms do not meet
the hypersurface of maximal contact; after Step 1a they no longer meet the cosupport, so near each point of the cosupport such a \(W\) exists.

\subsection{The tails of Step 1b on smaller open sets}

In this subsection \(\underline{\mathcal I}=(X,N,E,\mathfrak a,d)\) has \(N\) a closed submanifold, \(\mathfrak a\) coherent and \(d\ge1\),
\(\overline{\underline{\mathcal J}}=\mathcal C_N(\mathcal H_N(\underline{\mathcal I}))\), \(u\) is a tangent direction of \((\mathfrak a|_N,d)\) along \(N\) on an open set \(U\),
and for a multiple test blow-up \(T_0\) of \(\overline{\underline{\mathcal J}}\), \(S=(N\cap V(u))_{T_0}\) and
\(\underline{\mathcal S}_{u,T_0}=\mathcal C_{N_{T_0}}(\overline{\underline{\mathcal J}}_{T_0})|_S\).

\begin{theorem}[The tails of Step 1b on smaller open sets]
\label{mc:contact-stage-tail}
\uses{mc:coefficient-restriction-full-along,bcalc:centres-localise,bcalc:member-locus-level,bcalc:member-intersection-tower,hir:open-restriction-tower,hir:open-restriction-marked,hir:giraud-tower-along,hir:geo-transform-basic,mc:homogenized-equivalent-along,mc:homogenized-along-basic,mc:homogenized-along-etale,mc:restriction-operations-along,mc:strict-along,mc:contact-stage-domain-masked,mc:contact-equivalence-stage-coeff-hom-divisor,hir:step-one-b-independence}
\source{wlodarczyk05-hironaka}
\lean{Hironaka.ContactStageTail.isGeoTestTower_restrict_withN_coeff_of_cosupp_subset,Hironaka.ContactStageTail.cosupp_tail_subset_strictAlong,Hironaka.ContactStageTail.isGeoTestTower_tail_coeffHomAlong_of_le,Hironaka.ContactStageTail.isGeoTestTower_tail_contactMarked_of_le}
\leanok
\begin{enumerate}
\item Let \(\underline{\mathcal I}'=(Y,N',E',\mathfrak b,d')\) have \(N'\) a closed submanifold, \(\mathfrak b\) coherent and \(d'\ge1\), let \(W\subset Y\) be open and \(A\subset N'\)
with \(A\cap W\) a closed submanifold of \(W\). If along every multiple test blow-up \(T\) of \(\underline{\mathcal I}'\) the cosupport of \(\underline{\mathcal I}'_T\) over \(W\) lies in
\(A_T\), then for every multiple test blow-up \(T_1\) of \(\underline{\mathcal I}'\) the restriction \(T_1|_W\) is a multiple test blow-up of
\((\mathcal C_{N'}(\underline{\mathcal I}')|_A)|_W\), and the cosupport of its transform along \(T_1|_W\) is \(\operatorname{cosupp}\underline{\mathcal I}'_{T_1}\cap W_{T_1}\).
\item For every multiple test blow-up \(T_0T\) of \(\overline{\underline{\mathcal J}}\), the cosupport of \(\overline{\underline{\mathcal J}}_{T_0T}\) over \(U_{T_0}\) lies in \(S_T\).
\item For every multiple test blow-up \(T_0T_1\) of \(\overline{\underline{\mathcal J}}\) and every open \(W\subset U_{T_0}\) of the last stage of \(T_0\), \(T_1|_W\) is a multiple
test blow-up of \(\underline{\mathcal S}_{u,T_0}|_W\), and the cosupport of its transform along \(T_1|_W\) is \(\operatorname{cosupp}\overline{\underline{\mathcal J}}_{T_0T_1}\cap W_{T_1}\).
For \(E=\emptyset\) this is the marked ideal of Lemma~\ref{hir:step-one-b-independence} on \(W\).
\end{enumerate}
\end{theorem}
\begin{proof}
\leanok
(1) The restriction \(T_1|_W\) is a multiple test blow-up of \(\underline{\mathcal I}'|_W\), whose transforms are the restrictions of those of \(\underline{\mathcal I}'\)
(Lemma~\ref{hir:open-restriction-tower}). Its centres lie in the strict transforms of \(A\cap W\): each is the trace over \(W\) of an admissible centre of a transform
\(\underline{\mathcal I}'_T\) along a prefix \(T\) of \(T_1\), which lies in \(\operatorname{cosupp}\underline{\mathcal I}'_T\), hence over \(W\) in \(A_T\)
(Lemma~\ref{bcalc:centres-localise}), and strict transforms commute with restriction to open sets (Definition~\ref{mc:strict-along}). The coefficient ideal and the
restriction to \(A\) commute with restriction to \(W\) (Lemmas~\ref{mc:homogenized-along-etale} and~\ref{mc:restriction-operations-along}), so
Theorem~\ref{mc:coefficient-restriction-full-along} on \(W\), with \(A\cap W\subset N'\cap W\), shows that \(T_1|_W\) is a multiple test blow-up of
\((\mathcal C_{N'}(\underline{\mathcal I}')|_A)|_W\) whose cosupport is \(\operatorname{cosupp}(\underline{\mathcal I}'|_W)_{T_1|_W}\cap(A\cap W)_{T_1|_W}
=\operatorname{cosupp}\underline{\mathcal I}'_{T_1}\cap A_{T_1}\cap W_{T_1}\) (Lemma~\ref{hir:open-restriction-marked}); by the hypothesis for \(T_1\) the factor \(A_{T_1}\) can be dropped.

(2) \(T_0T\) is a multiple test blow-up of \(\overline{\underline{\mathcal J}}\) (Lemma~\ref{hir:geo-transform-basic}), hence of \(\underline{\mathcal I}\), with the same cosupports
(Theorem~\ref{mc:homogenized-equivalent-along}), and the transform of \(\overline{\underline{\mathcal J}}\) along \(T_0T\) is that of \(\overline{\underline{\mathcal J}}_{T_0}\) along \(T\).
By Theorem~\ref{hir:giraud-tower-along} the cosupport of \(\underline{\mathcal I}_{T_0T}\) over \(U\) lies in \((N\cap V(u)\cap U)_{T_0T}\subset S_T\): strict transforms are monotone
and compose along concatenations (Lemma~\ref{bcalc:member-intersection-tower} and Theorem~\ref{bcalc:member-locus-level}). The points of the last stage of \(T\) over \(U_{T_0}\) are
those over \(U\).

(3) Apply (1) to \(\underline{\mathcal I}'=\overline{\underline{\mathcal J}}_{T_0}\), whose submanifold \(N_{T_0}\) is closed, whose ideal is coherent and whose mark is \(\ge1\)
(Lemmas~\ref{hir:open-restriction-tower} and~\ref{mc:homogenized-along-basic}), to \(A=S\subset N_{T_0}\) (Theorem~\ref{hir:giraud-tower-along}) and to the tower \(T_1\): \(S\cap W\) is a closed
submanifold of \(W\) (Theorem~\ref{mc:contact-stage-domain-masked}), the hypothesis of (1) is (2) because \(W\subset U_{T_0}\), and
\(\overline{\underline{\mathcal J}}_{T_0T_1}=(\overline{\underline{\mathcal J}}_{T_0})_{T_1}\).
\end{proof}
On \(U_{T_0}\) itself the tails are those of Corollary~\ref{mc:contact-equivalence-stage-coeff-hom-divisor}, which comes from an equivalence of marked ideals over \(U_{T_0}\).
An equivalence is a statement about all multiple test blow-ups of its base and does not pass to a smaller open set \(W\), whose towers need not extend; the theorem
is therefore stated for the segment at hand, which is what the value of Step 1b in Proposition~\ref{hir:maximal-order-step} uses on a neighbourhood \(W\) of a point
that meets no original member.

\subsection{Translations along a submanifold}

\begin{lemma}[Translations along a submanifold]
\label{mc:smooth-centre-translation}
\uses{mc:flat-chart,fpgeom:complex-analytic-subvariety,hir:open-biholomorphism,bcalc:restrict-open,hir:biholomorphism-blowup,hir:biholomorphism-transform}
\source{wlodarczyk05-hironaka}
\lean{Hironaka.SmoothCentreHomogeneity.add_mem_complexCoordinatePlane_iff,Hironaka.SmoothCentreHomogeneity.neg_mem_complexCoordinatePlane,Hironaka.SmoothCentreHomogeneity.exists_diffeomorph_opens_of_mapsTo,Hironaka.SmoothCentreHomogeneity.exists_translation_of_regularAtDim,Hironaka.SmoothCentreHomogeneity.pullback_val_vanishingIdeal,Hironaka.SmoothCentreHomogeneity.pullback_pullback_val_vanishingIdeal,Hironaka.SmoothCentreHomogeneity.pullback_pullback_val_smoothCentreMarked,Hironaka.SmoothCentreHomogeneity.exists_smoothCentre_homogeneity,Hironaka.SmoothCentreHomogeneity.exists_chart_transfer_of_regularAtDim,Hironaka.SmoothCentreHomogeneity.exists_smoothCentre_iso}
\leanok
Let \(Z\subset M\) be regular of dimension \(k\) at \(x\). There is an open set \(W\ni x\) such that for every \(x'\in Z\cap W\) there are open sets \(U_1\ni x\),
\(U_2\ni x'\) and a biholomorphism \(\Phi\colon U_1\to U_2\) with \(\Phi(x)=x'\) and \(\Phi^{-1}(Z\cap U_2)=Z\cap U_1\). For every such \(\Phi\), every \(N\subset M\) with
\(\Phi^{-1}(N\cap U_2)=N\cap U_1\) and every \(d\),
\[\Phi^*\bigl((M,N,\emptyset,\mathcal I_Z,d)|_{U_2}\bigr)=(M,N,\emptyset,\mathcal I_Z,d)|_{U_1};\]
in particular this holds for \(N=M\). (Homogeneous local model.) The same holds for any two points \(x,x'\), not necessarily close, at which \(Z\) is
regular of the same dimension \(k\): there are open \(U_1\ni x\), \(U_2\ni x'\) and a biholomorphism \(\Phi\colon U_1\to U_2\) with \(\Phi(x)=x'\) and
\(\Phi^{-1}(Z\cap U_2)=Z\cap U_1\).
\end{lemma}
\begin{proof}
\leanok
Take a flat chart \(\psi\) of \(Z\) at \(x\) (Definition~\ref{mc:flat-chart}, from the regularity of Definition~\ref{fpgeom:complex-analytic-subvariety}): \(\psi(x)=0\),
\(\psi\) has a holomorphic inverse \(\psi^{-1}\) near \(0\), and near \(x\) a point lies in \(Z\) exactly when its image lies in the plane \(P=\mathbf C^k\times0\). Let \(A\ni x\)
be open with \(\psi\) holomorphic on \(A\), \(\psi^{-1}\psi=\mathrm{id}\) on \(A\) and \(Z\cap A=A\cap\psi^{-1}(P)\), and let \(\rho>0\) be such that on the ball \(B(0,\rho)\) the map
\(\psi^{-1}\) is holomorphic, \(\psi\psi^{-1}=\mathrm{id}\) and \(\psi^{-1}(B(0,\rho))\subset A\). Put \(r=\rho/2\), \(U(v)=A\cap\psi^{-1}(B(v,r))\) and \(W=U(0)\ni x\).
For \(x'\in Z\cap W\) the vector \(v=\psi(x')\) lies in \(P\) and \(|v|<r\). The maps \(\varphi(y)=\psi^{-1}(\psi(y)+v)\) and \(\varphi'(q)=\psi^{-1}(\psi(q)-v)\) are defined and
holomorphic on \(U(0)\) and \(U(v)\), since \(|\psi(y)\pm v|<2r=\rho\) there, map \(U(0)\) to \(U(v)\) and back, because \(\psi(\varphi(y))=\psi(y)+v\) and
\(\psi(\varphi'(q))=\psi(q)-v\), and are inverse to each other. Inverse holomorphic maps between open sets form a biholomorphism of the open submanifolds, as in
Lemma~\ref{hir:open-biholomorphism}. Then \(\varphi(x)=\psi^{-1}(v)=x'\), and since \(P\) is a subgroup containing \(v\), \(\psi(\varphi(y))\in P\) if and only if
\(\psi(y)\in P\), that is \(\varphi(y)\in Z\) if and only if \(y\in Z\).
For the marked ideals: \(\mathcal I_Z|_U=\mathcal I_{Z\cap U}\) on an open set \(U\) (Lemma~\ref{bcalc:restrict-open}), and the pullback of a vanishing ideal under a
biholomorphism is the vanishing ideal of the preimage (Lemma~\ref{hir:biholomorphism-blowup}); so \(\Phi^*(\mathcal I_Z|_{U_2})=\mathcal I_{\Phi^{-1}(Z\cap U_2)}=\mathcal I_Z|_{U_1}\).
The submanifold and the (empty) divisor of the pulled-back marked ideal are the preimages (Lemma~\ref{hir:biholomorphism-transform}), and the mark is unchanged.
For two points \(x,x'\) of the same dimension take flat charts \(\psi_1\) at \(x\) and \(\psi_2\) at \(x'\), both onto neighbourhoods of \(0\) with \(Z\) corresponding to the
same plane \(P\), good neighbourhoods \(A_1\ni x\), \(A_2\ni x'\) as above, and \(r>0\) such that on \(B(0,r)\) both inverses are holomorphic, right inverse, and map into
\(A_1\), resp.\ \(A_2\). On \(U_i=A_i\cap\psi_i^{-1}(B(0,r))\) the maps \(\psi_2^{-1}\circ\psi_1\) and \(\psi_1^{-1}\circ\psi_2\) are holomorphic, inverse to each other and map
\(U_1\) to \(U_2\) and back; the first sends \(x\) to \(\psi_2^{-1}(0)=x'\), and it preserves \(Z\) since both charts send \(Z\) to \(P\).
\end{proof}
This is the homogeneity of a smooth centre used in Lemma~\ref{hir:smooth-centre}: with properties (5) and (6) of Definition~\ref{hir:centre-function} for the empty towers,
it makes the invariant of \((M,M,\emptyset,\mathcal I_Z,1)\) locally constant on \(Z\).

\begin{corollary}[A resolution invariant is constant along a smooth centre]
\label{mc:smooth-centre-invariant}
\uses{mc:smooth-centre-translation,hir:centre-function,hir:tower-isomorphism,hir:open-restriction-marked,bcalc:smooth-centre-locus,bcalc:inv-domain-stable,bcalc:transform-off-centre,hir:canonical-segments,hir:wellformed-tower}
\source{wlodarczyk05-hironaka}
\lean{Hironaka.SmoothCentreInvariant.inv_refl_locallyConstant,Hironaka.SmoothCentreInvariant.inv_refl_eq_of_isPreconnected,Hironaka.SmoothCentreInvariant.mem_cosupp_geoTransformAlong_of_isEmptyTower,Hironaka.SmoothCentreInvariant.inv_eq_inv_refl_of_isEmptyTower,Hironaka.SmoothCentreInvariant.inv_eq_of_isEmptyTower_of_isPreconnected,Hironaka.SmoothCentreInvariant.maxLocus_eq_cosupp_of_isEmptyTower,Hironaka.SmoothCentreInvariant.inv_refl_eq_of_regularAtDim,Hironaka.SmoothCentreInvariant.inv_refl_eq_of_pureDim,Hironaka.SmoothCentreInvariant.inv_eq_of_isEmptyTower_of_pureDim,Hironaka.SmoothCentreInvariant.maxLocus_eq_cosupp_of_isEmptyTower_of_pureDim}
\leanok
Let \(\operatorname{inv}\) be a resolution invariant in dimension \(\le k\) (Definition~\ref{hir:centre-function}), \(Z\subset X\) a closed submanifold and
\(\underline{\mathcal I}=(X,X,\emptyset,\mathcal I_Z,1)\) in its domain. Then \(\operatorname{inv}_{\underline{\mathcal I}}\) of the empty tower is locally constant on
\(Z=\operatorname{cosupp}\underline{\mathcal I}\), and constant on \(Z\) if \(Z\) is connected. Along a tower \(T\) all of whose centres are empty,
\(\operatorname{cosupp}\underline{\mathcal I}_T\) is the preimage of \(Z\) and \(\operatorname{inv}_{\underline{\mathcal I}}(T)\) is \(\operatorname{inv}_{\underline{\mathcal I}}\) of
the empty tower composed with the projection; so, \(Z\) being connected, \(\operatorname{inv}_{\underline{\mathcal I}}(T)\) is constant on
\(\operatorname{cosupp}\underline{\mathcal I}_T\), and its maximum locus there is the whole cosupport. (The first statement along such a tower holds for every
marked ideal with a closed submanifold and a coherent ideal, the second for every marked ideal in the domain.) Connectedness is not needed
when \(Z\) has pure dimension: then the values at any two points of \(Z\) agree, and so do the conclusions along \(T\).
\end{corollary}
\begin{proof}
\leanok
The cosupport is \(Z\) (Proposition~\ref{bcalc:smooth-centre-locus}). Let \(x\in Z\) and take \(W\) as in Lemma~\ref{mc:smooth-centre-translation}; for \(x'\in Z\cap W\) let
\(\Phi\colon U_1\to U_2\) be the biholomorphism given there, so \(\Phi(x)=x'\) and \(\Phi^*(\underline{\mathcal I}|_{U_2})=\underline{\mathcal I}|_{U_1}\). The restrictions of
\(\underline{\mathcal I}\) to open sets lie in the domain (Proposition~\ref{bcalc:inv-domain-stable}) and have the traces of \(Z\) as cosupports
(Lemma~\ref{hir:open-restriction-marked}); the empty towers over \(U_1\) and \(U_2\) are isomorphic over \(\Phi\) with last-stage map \(\Phi\)
(Definition~\ref{hir:tower-isomorphism}) and are canonical segments. By property (5) at \(x\), property (6) for these empty towers, and property (5) at \(x'\),
\[\operatorname{inv}_{\underline{\mathcal I}}(x)=\operatorname{inv}_{\underline{\mathcal I}|_{U_1}}(x)=\operatorname{inv}_{\Phi^*(\underline{\mathcal I}|_{U_2})}(x)
=\operatorname{inv}_{\underline{\mathcal I}|_{U_2}}(\Phi(x))=\operatorname{inv}_{\underline{\mathcal I}}(x').\]
A locally constant function on a connected space is constant.
Along \(T\) with empty centres, at a point \(y\) of the last stage, the image of \(y\) under the last blowup lies off its (empty) centre, so \(y\) lies in the
cosupport of the transform iff its image lies in the cosupport of the preceding stage (Theorem~\ref{bcalc:transform-off-centre}; the stages have closed
submanifolds and coherent ideals, Lemma~\ref{hir:wellformed-tower}); by induction, \(\operatorname{cosupp}\underline{\mathcal I}_T\) is the preimage of
\(\operatorname{cosupp}\underline{\mathcal I}\). Such a tower is a canonical segment (Lemma~\ref{hir:canonical-segments}), so property (4) of
Definition~\ref{hir:centre-function}, one step at a time, gives \(\operatorname{inv}_{\underline{\mathcal I}}(T)(y)=\operatorname{inv}_{\underline{\mathcal I}}(p_T(y))\),
\(p_T\) the projection of \(T\), with the empty tower on the right. For \(\underline{\mathcal I}=(X,X,\emptyset,\mathcal I_Z,1)\) with \(Z\) connected the right-hand side is constant on \(Z\) by the
first part, so the values on \(\operatorname{cosupp}\underline{\mathcal I}_T\) all agree and every point of it lies in the maximum locus.
If \(Z\) has pure dimension, for any \(x,x'\in Z\) the biholomorphism \(\Phi\colon U_1\to U_2\) of the homogeneous local model in
Lemma~\ref{mc:smooth-centre-translation} pulls \(\underline{\mathcal I}|_{U_2}\) back to \(\underline{\mathcal I}|_{U_1}\), and the chain of equalities above, with
\(\Phi(x)=x'\), compares the values at \(x\) and \(x'\) directly; the conclusions along \(T\) follow as before.
\end{proof}

\begin{lemma}[The invariant on the smooth part of a subvariety]
\label{mc:smooth-part-invariant}
\uses{mc:smooth-centre-invariant,hir:centre-function,hir:canonical-segments,hir:open-restriction-marked,hir:open-restriction-tower,bcalc:restrict-open,bcalc:inv-domain-stable,bcalc:restrict-empty-tower}
\source{wlodarczyk05-hironaka}
\lean{Hironaka.SmoothCentreInvariant.pullback_val_smoothCentreMarked,Hironaka.SmoothCentreInvariant.inv_eq_over_open_of_isEmptyTower,Hironaka.SmoothCentreInvariant.inv_eq_over_open_of_isEmptyTower_of_pureDim,Hironaka.SmoothCentreInvariant.mem_maxLocus_over_open_of_isEmptyTower}
\leanok
Let \(\operatorname{inv}\) be a resolution invariant in dimension \(\le k\) (Definition~\ref{hir:centre-function}), \(V\subset X\) a closed set with
\(\underline{\mathcal I}=(X,X,\emptyset,\mathcal I_V,1)\) in its domain, and \(U\subset X\) open such that \(V\cap U\) is a closed submanifold of \(U\) which is
connected, or of pure dimension. Then \(\underline{\mathcal I}|_U=(U,U,\emptyset,\mathcal I_{V\cap U},1)\). Let \(T\) be a canonical segment of \(\underline{\mathcal I}\)
whose restriction to \(U\) has only empty centres (for instance, all its centres miss the preimages of \(U\)). Then
\(\operatorname{inv}_{\underline{\mathcal I}}(T)\) is constant on the points of \(\operatorname{cosupp}\underline{\mathcal I}_T\) over \(U\); in particular, if a point over
\(U\) lies in the maximum locus of \(\operatorname{inv}_{\underline{\mathcal I}}(T)\) on the cosupport, all points of the cosupport over \(U\) do.
\end{lemma}
\begin{proof}
\leanok
The restriction of a vanishing ideal to an open set is the vanishing ideal of the trace (Lemma~\ref{bcalc:restrict-open}), and the submanifold \(X\) restricts to \(U\).
The restriction \(\underline{\mathcal I}|_U\) lies in the domain (Proposition~\ref{bcalc:inv-domain-stable}), and the cosupport along \(T|_U\) is the trace of the cosupport
along \(T\) (Lemmas~\ref{hir:open-restriction-tower} and~\ref{hir:canonical-segments}). By property (5) of Definition~\ref{hir:centre-function}, at a point \(p\) of the
cosupport over \(U\) the value of \(\underline{\mathcal I}\) along \(T\) is the value of \(\underline{\mathcal I}|_U\) along \(T|_U\), which has only empty centres; by
Corollary~\ref{mc:smooth-centre-invariant} for \((U,U,\emptyset,\mathcal I_{V\cap U},1)\) these values agree at all such points. If \(c\) over \(U\) is in the maximum
locus, every \(y\) of the cosupport over \(U\) has the value of \(c\), which bounds all values on the cosupport.
\end{proof}
This is the middle step of Lemma~\ref{hir:embedded-from-resolution} (with \(U=X\setminus V_{\mathrm{sing}}\) and the segment before the first centre meeting the preimage
of \(U\); that its restriction to \(U\) has only empty centres is Lemma~\ref{bcalc:restrict-empty-tower}).

\subsection{The first centre meeting the smooth part of a subvariety}

In this subsection \(X\) is a complex manifold of dimension \(n\), \(T\) a tower of blowups over \(X\) with projection \(p_T\colon X_T\to X\), and for a closed
irreducible analytic subvariety \(V\subset X\) we write \(V_{\mathrm{sing}}\) for the set of points of \(V\) at which \(V\) is not regular of dimension \(\dim V\)
(Definition~\ref{ana:nonkahler-locus}), \(V_{\mathrm{reg}}=V\setminus V_{\mathrm{sing}}\) and \(\underline{\mathcal I}_V=(X,X,\emptyset,\mathcal I_V,1)\).

\begin{lemma}[Towers whose centres miss an open set]
\label{mc:centres-avoid-open}
\uses{ana:blowups,bcalc:isomorphism-off-centre,ana:strict-transform-composition,bcalc:restrict-empty-tower,bcalc:restrict-open,hir:open-restriction-tower,hir:canonical-segments,mc:smooth-centre-invariant,bcalc:smooth-centre-locus,bcalc:inv-domain-stable,hir:centre-function}
\source{wlodarczyk05-hironaka}
\lean{Hironaka.EmbeddedFromResolution.notMem_exceptionalLocus_of_centresAvoid,Hironaka.EmbeddedFromResolution.image_exceptionalLocus_subset_compl,Hironaka.EmbeddedFromResolution.strictTransform_subset_preimage,Hironaka.EmbeddedFromResolution.mem_strictTransform_iff_of_centresAvoid,Hironaka.EmbeddedFromResolution.mem_cosupp_iff_of_centresAvoid}
\leanok
Let \(U\subset X\) be open and suppose that every centre of \(T\) misses the preimage of \(U\) at its stage. Then:
\begin{enumerate}
\item \(p_T\) has bijective differential at every point of \(p_T^{-1}(U)\); hence \(p_T(\operatorname{Exc}p_T)\subset X\setminus U\);
\item for a closed \(A\subset X\) the strict transform \(\overline{p_T^{-1}(A)\setminus\operatorname{Exc}p_T}\) lies in \(p_T^{-1}(A)\) and agrees with it over \(U\);
\item if \(V\subset X\) is closed, \((X,X,\emptyset,\mathcal I_V,1)\) lies in the domain of the invariants in dimension \(\le k\)
(Definition~\ref{hir:centre-function}) and \(T\) is a multiple test blow-up of it, then the points of its cosupport along \(T\) lying over \(U\) are exactly
the points of \(p_T^{-1}(V)\) over \(U\).
\end{enumerate}
\end{lemma}
\begin{proof}
\leanok
(1) Induction on the tower. For \(T\) followed by the blowup \(\pi\) along \(C\), \(\operatorname{Exc}(p_T\circ\pi)=\operatorname{Exc}\pi\cup\pi^{-1}(\operatorname{Exc}p_T)\) by the
chain rule, and \(\operatorname{Exc}\pi\subset\pi^{-1}(C)\) (Lemma~\ref{bcalc:isomorphism-off-centre}); a point \(y\) over \(U\) has \(\pi(y)\notin C\), and
\(\pi(y)\notin\operatorname{Exc}p_T\) by induction. (2) \(p_T^{-1}(A)\) is closed and contains \(p_T^{-1}(A)\setminus\operatorname{Exc}p_T\); over \(U\) no point is
exceptional by (1). (3) The restriction of \(T\) to \(U\) has only empty centres (Lemma~\ref{bcalc:restrict-empty-tower}); the cosupport along \(T|_U\) of the
restriction \((U,U,\emptyset,\mathcal I_{V\cap U},1)\) (Lemma~\ref{bcalc:restrict-open}), which lies in the domain (Proposition~\ref{bcalc:inv-domain-stable}),
is the trace of the cosupport along \(T\) (Lemmas~\ref{hir:open-restriction-tower} and~\ref{hir:canonical-segments}); along a tower with only empty centres
the cosupport is the preimage of the cosupport (Corollary~\ref{mc:smooth-centre-invariant}), and the cosupport of \((U,U,\emptyset,\mathcal I_{V\cap U},1)\)
is \(V\cap U\) (Proposition~\ref{bcalc:smooth-centre-locus}).
\end{proof}

\begin{definition}[The properties used for embedded resolution]
\label{mc:embedding-properties}
\uses{hir:centre-function}
\lean{Hironaka.EmbeddedFromResolution.EmbeddingFields,Hironaka.EmbeddedFromResolution.embeddingFields_of_resolutionInvariant}
\leanok
A function \(\operatorname{inv}\) assigning to every marked ideal \(\underline{\mathcal I}\) on a complex manifold of dimension \(n\), every tower \(T\) over its base
and every point of the last stage of \(T\) a value in a totally ordered set has the \emph{embedding properties in dimension \(\le k\)} if properties (1)
(admissibility), (4) (empty steps), (5) (open sets) and (6) (isomorphisms of towers) of Definition~\ref{hir:centre-function} hold for it, in the same form.
A resolution invariant in dimension \(\le k\) has them.
\end{definition}

\begin{corollary}[Smooth centres for the embedding properties]
\label{mc:embedding-properties-smooth-centre}
\uses{mc:embedding-properties,hir:invariant-driver,mc:smooth-centre-invariant,mc:smooth-part-invariant}
\source{wlodarczyk05-hironaka}
\lean{Hironaka.EmbeddedFromResolution.EmbeddingFields.isAdmissibleCentre_maxLocus,Hironaka.EmbeddedFromResolution.EmbeddingFields.isGeoTestTower_of_isCanonical,Hironaka.EmbeddedFromResolution.EmbeddingFields.inv_eq_inv_refl_of_isEmptyTower,Hironaka.EmbeddedFromResolution.EmbeddingFields.inv_refl_eq_of_regularAtDim,Hironaka.EmbeddedFromResolution.EmbeddingFields.inv_eq_of_isEmptyTower_of_pureDim,Hironaka.EmbeddedFromResolution.EmbeddingFields.inv_eq_over_open_of_isEmptyTower_of_pureDim}
\leanok
Let \(\operatorname{inv}\) have the embedding properties in dimension \(\le k\). Then every canonical segment of a marked ideal of the domain is a multiple test
blow-up, along which the maximum locus is an admissible centre (Theorem~\ref{hir:invariant-driver}(3)); and the conclusions of
Corollary~\ref{mc:smooth-centre-invariant} for a smooth centre of pure dimension and of Lemma~\ref{mc:smooth-part-invariant} for \(V\cap U\) of pure dimension
hold for \(\operatorname{inv}\).
\end{corollary}
\begin{proof}
\leanok
The cited proofs use only properties (1), (4), (5) and (6) of Definition~\ref{hir:centre-function}.
\end{proof}

\begin{lemma}[The first centre meeting the smooth part]
\label{mc:first-centre-smooth-part}
\uses{mc:centres-avoid-open,mc:embedding-properties,mc:embedding-properties-smooth-centre,ana:nonkahler-locus,lag:regular-point-basics,ana:regular-transport,ana:strict-transform-sequence,ana:strict-transform-blowup,ana:regular-component,ana:analytic-preimage,rem:open-subset-analytic,hir:resolution-driver}
\source{wlodarczyk05-hironaka}
\lean{Hironaka.EmbeddedFromResolution.regularOpen,Hironaka.EmbeddedFromResolution.isClosed_singularLocus,Hironaka.EmbeddedFromResolution.mem_regularOpen_iff,Hironaka.EmbeddedFromResolution.regularAtDim_of_mem_regularOpen,Hironaka.EmbeddedFromResolution.mem_regularOpen_of_regularAtDim,Hironaka.EmbeddedFromResolution.strictTransform_eq_connectedComponentIn,Hironaka.EmbeddedFromResolution.mem_maxLocus_iff_of_centresAvoid,Hironaka.EmbeddedFromResolution.strictTransform_eq_connectedComponentIn_maxLocus}
\leanok
Let \(X\) be compact, \(V\subset X\) a closed irreducible analytic subvariety and \(U=X\setminus V_{\mathrm{sing}}\), and suppose that the centres of \(T\) miss the
preimages of \(U\).
\begin{enumerate}
\item \(V_{\mathrm{sing}}\) is closed. If \(C\) is a closed submanifold of \(X_T\) which over \(U\) is the preimage of \(V\), and \(c\in C\) lies over \(U\), then
\(p_T(\operatorname{Exc}p_T)\subset V_{\mathrm{sing}}\) and the strict transform \(V_T\) of \(V\) is the connected component of \(C\) through \(c\); in particular
\(V_T\) is regular of dimension \(\dim V\) at each of its points.
\item Let \(\operatorname{inv}\) have the embedding properties in dimension \(\le k\), with \(\underline{\mathcal I}_V\) in its domain, and let \(T\) be a canonical
segment of \(\underline{\mathcal I}_V\). If the maximum locus \(C\) of \(\operatorname{inv}_{\underline{\mathcal I}_V}(T)\) on \(\operatorname{cosupp}(\underline{\mathcal I}_V)_T\)
contains a point \(c\) over \(U\), then over \(U\) it is the preimage of \(V\), and (1) applies to it.
\end{enumerate}
\end{lemma}
\begin{proof}
\leanok
\(V_{\mathrm{sing}}\) is closed because regularity of a fixed dimension is an open condition on \(V\) (Lemma~\ref{lag:regular-point-basics}).
(1) By Lemma~\ref{mc:centres-avoid-open}, \(p_T(\operatorname{Exc}p_T)\subset X\setminus U=V_{\mathrm{sing}}\), and over the open set \(O=p_T^{-1}(U)\) the strict transform \(S=V_T\)
is \(p_T^{-1}(V)\). As \(p_T\) has bijective differential at \(c\) and \(V\) is regular of dimension \(\dim V\) at \(p_T(c)\in V_{\mathrm{reg}}\), \(S\) is regular of dimension
\(\dim V\) at \(c\) (Lemma~\ref{ana:regular-transport}), and so is \(C\), which agrees with \(S\) on \(O\) (Lemma~\ref{lag:regular-point-basics}). Since
\(V\not\subset p_T(\operatorname{Exc}p_T)\), \(S\) is a closed irreducible analytic subvariety (Lemma~\ref{ana:strict-transform-sequence}), hence the closure of the
connected component \(S_c\) of its regular locus through \(c\) (Lemma~\ref{ana:strict-transform-blowup}). \(C\) is analytic (Lemma~\ref{ana:analytic-preimage}) and all its
points are regular; let \(K\) be its connected component through \(c\). By the identity theorem along a component of a regular locus (an analytic set containing
a nonempty relatively open part of the component contains the component; proof of Proposition~\ref{ana:regular-component}), applied on \(O\ni c\) where \(S\) and
\(C\) agree, \(S_c\subset C\) and \(K\subset S\). The connected set \(S_c\) of regular points of \(C\) meets \(K\), so \(S_c\subset K\), and \(S=\overline{S_c}\subset C\).
Let \(x\in S\); then \(x\in C\) has a neighbourhood \(W\) with \(C\cap W\) inside the component \(K_x\) of \(C\) through \(x\)
(Proposition~\ref{ana:regular-component}(1)), and \(W\) meets \(S_c\subset K\), so \(K_x=K\) and \(x\in K\). Hence \(S=K\). All points of \(K\) are regular points of \(C\)
of dimension \(\dim V\) (Proposition~\ref{ana:regular-component}(1)), and near each of them \(C\) lies in \(K=S\subset C\), so \(S\) is regular of dimension \(\dim V\) there.
(2) \(T\) is a multiple test blow-up and \(C\) an admissible centre, in particular a closed submanifold (Corollary~\ref{mc:embedding-properties-smooth-centre}).
Over \(U\) the cosupport is the preimage of \(V\) (Lemma~\ref{mc:centres-avoid-open}(3)) and contains \(C\). \(V\cap U=V_{\mathrm{reg}}\) is a closed submanifold of \(U\) of pure
dimension \(\dim V\) (Lemma~\ref{rem:open-subset-analytic} for the traces), so by Corollary~\ref{mc:embedding-properties-smooth-centre} (the form of
Lemma~\ref{mc:smooth-part-invariant}) the invariant is constant on the cosupport over \(U\); as \(c\) lies in the maximum locus, so does every point of the cosupport
over \(U\).
\end{proof}

\begin{lemma}[The first centre meeting an open set]
\label{mc:first-centre-exists}
\uses{hir:canonical-segments,hir:resolution-driver,mc:centres-avoid-open,ana:blowup-surjective,hir:blowup-tower}
\source{wlodarczyk05-hironaka}
\lean{Hironaka.EmbeddedFromResolution.exists_first_centre,Hironaka.EmbeddedFromResolution.not_centresAvoid_of_cosupp_eq_empty}
\leanok
\begin{enumerate}
\item Let \(\operatorname{inv}\) be any value function, \(U\subset X\) open and \(T\) a canonical segment of a marked ideal on \(X\) (each centre empty or the
maximum locus of \(\operatorname{inv}\) of the preceding part on the cosupport) whose centres do not all miss the preimages of \(U\). Then \(T=T'S\) where \(T'\) is a
canonical segment whose centres miss the preimages of \(U\) and whose maximum locus, the next centre of \(T\), meets the preimage of \(U\).
\item Let \(X\) be compact, \(V\subset X\) a closed irreducible analytic subvariety with \(\underline{\mathcal I}_V\) in the domain of the invariants in
dimension \(\le k\), and \(T\) a multiple test blow-up of \(\underline{\mathcal I}_V\) with empty final cosupport. Then not all centres of \(T\) miss the preimages of
\(X\setminus V_{\mathrm{sing}}\).
\end{enumerate}
\end{lemma}
\begin{proof}
\leanok
(1) Induction on \(T\). Let \(T\) be \(T_0\) followed by the blowup along \(C\). If not all centres of \(T_0\) miss the preimages of \(U\), the induction hypothesis
applies to \(T_0\). Otherwise \(C\) meets the preimage of \(U\); so \(C\neq\emptyset\), \(C\) is the maximum locus at \(T_0\), and \(T'=T_0\).
(2) Otherwise, take \(x\in V_{\mathrm{reg}}\) (nonempty by definition of the subvariety); \(p_T\) is onto (Lemma~\ref{ana:blowup-surjective}), and a point over \(x\) lies in the
cosupport by Lemma~\ref{mc:centres-avoid-open}(3), against its emptiness.
\end{proof}

\begin{theorem}[Embedded resolution from the first centre meeting the smooth part]
\label{mc:embedded-from-first-centre}
\uses{mc:first-centre-smooth-part,mc:first-centre-exists,mc:embedding-properties,hir:invariant-driver,hir:blowup-tower}
\source{wlodarczyk05-hironaka,bierstone-milman08-functoriality}
\lean{Hironaka.EmbeddedFromResolution.exists_embedded_of_isCanonical,Hironaka.EmbeddedFromResolution.exists_embeddedResolution_data_of_isCanonical,Hironaka.EmbeddedFromResolution.exists_embeddedResolution_data_of_resolutionInvariant}
\leanok
Let \(X\) be compact, \(V\subset X\) a closed irreducible analytic subvariety, \(\operatorname{inv}\) a function with the embedding properties in dimension
\(\le k\) with \(\underline{\mathcal I}_V\) in its domain, and \(T\) a canonical segment of \(\underline{\mathcal I}_V\) with empty final cosupport. Then the part \(T'\) of
\(T\) before the first centre meeting the preimage of \(X\setminus V_{\mathrm{sing}}\) is a canonical segment, and its projection, a finite composition of blowups with
smooth centres, has its exceptional image in \(V_{\mathrm{sing}}\) and a strict transform of \(V\) regular of dimension \(\dim V\) at each of its points. In
particular this holds for a resolution invariant in dimension \(\le k\) with \(\underline{\mathcal I}_V\) in its domain, with \(T\) the canonical resolution of
Theorem~\ref{hir:invariant-driver}(4).
\end{theorem}
\begin{proof}
\leanok
Lemma~\ref{mc:first-centre-exists}(2) and (1) give \(T'\) and a point of its maximum locus over \(X\setminus V_{\mathrm{sing}}\); Lemma~\ref{mc:first-centre-smooth-part}(2)
and (1) conclude. The projection of a tower is a finite composition of blowups with smooth centres (Definition~\ref{hir:blowup-tower}). For a resolution
invariant, Theorem~\ref{hir:invariant-driver}(4) gives the canonical segment \(T\).
\end{proof}
With Lemma~\ref{mc:smooth-part-invariant} this is the proof of Lemma~\ref{hir:embedded-from-resolution} for \(\dim V<n\), given a function with the embedding properties
whose domain contains \(\underline{\mathcal I}_V\) and a canonical segment of \(\underline{\mathcal I}_V\) with empty final cosupport.

\begin{lemma}[Subvarieties of dimension less than \(n\)]
\label{mc:subvariety-thin}
\uses{fpgeom:complex-analytic-subvariety,ana:ideal-sheaves,lag:regular-point-basics,ana:regular-component,ana:strict-transform-blowup,hir:centre-function}
\source{wlodarczyk05-hironaka}
\lean{Hironaka.EmbeddedFromResolution.regularAtDim_of_mem_nhds,Hironaka.EmbeddedFromResolution.mem_nhds_of_regularAtDim_self,Hironaka.EmbeddedFromResolution.interior_eq_empty_of_dim_lt,Hironaka.EmbeddedFromResolution.isNonzeroAlongAt_vanishingIdeal_of_dim_lt,Hironaka.EmbeddedFromResolution.isInvDomain_vanishingIdeal_of_dim_lt}
\leanok
Let \(M\) be a complex manifold of dimension \(n\). A subset of \(M\) is regular of dimension \(n\) at a point if and only if it is a neighbourhood of that point.
Let \(V\subset M\) be a closed irreducible analytic subvariety with \(\dim V<n\). Then \(V\) has empty interior, and at every point of \(M\) the ideal
\(\mathcal I_V\) contains a germ which does not vanish identically near the point. Hence, if \(\mathcal I_V\) is coherent and \(n\le k\), the marked ideal
\((M,M,\emptyset,\mathcal I_V,1)\) lies in the domain of the invariants in dimension \(\le k\) (Definition~\ref{hir:centre-function}).
\end{lemma}
\begin{proof}
\leanok
A regular point of dimension \(n\) has a flattening onto all of \(\mathbb C^n\), so a neighbourhood of it lies in the set; conversely a set agreeing near a point
with \(M\), which is regular of dimension \(n\) everywhere, is regular of dimension \(n\) there (Lemma~\ref{lag:regular-point-basics}). If \(x\) were an interior point of
\(V\), it would be a regular point of dimension \(n\); so \(V\) is the closure of the connected component \(C_x\) of its regular locus through \(x\)
(Lemma~\ref{ana:strict-transform-blowup}), and all points of \(C_x\) are regular of dimension \(n\) (Proposition~\ref{ana:regular-component}(1)). A regular point \(c\)
of dimension \(\dim V\) has a neighbourhood \(W\) with \(V\cap W\) inside the component \(C_c\) through \(c\) (Proposition~\ref{ana:regular-component}(1)); as
\(c\in\overline{C_x}\), \(W\) meets \(C_x\), so \(C_c=C_x\) and \(c\) is regular of dimension both \(\dim V\) and \(n\), whence \(\dim V=n\)
(Lemma~\ref{lag:regular-point-basics}), a contradiction. Off \(V\) the constant \(1\) lies in \(\mathcal I_V\). At \(x\in V\) take local equations
\(g_1,\dots,g_m\) of \(V\) in a chart at \(x\) (Definition~\ref{fpgeom:complex-analytic-subvariety}); they are holomorphic near \(x\) and vanish on \(V\) near \(x\), so their
germs lie in \(\mathcal I_{V,x}\) (Definition~\ref{ana:ideal-sheaves}), and if they all vanished identically near \(x\), \(V\) would contain a neighbourhood of \(x\). The
conditions of the domain: \(M\) is a closed submanifold of dimension \(n\le k\), \(E=\emptyset\), \(d=1\), \(\mathcal I_V\) is coherent and nonzero along \(M\) at
every point.
\end{proof}

\begin{lemma}[Subvarieties of dimension \(n\)]
\label{mc:subvariety-full-dimension}
\uses{fpgeom:complex-analytic-subvariety,mc:subvariety-thin,ana:regular-component,lag:regular-point-basics,ana:blowups}
\source{wlodarczyk05-hironaka}
\lean{Hironaka.EmbeddedFromResolution.isClosedConnectedSubmanifold_of_dim_eq,Hironaka.EmbeddedFromResolution.exists_embeddedResolution_data_of_dim_eq}
\leanok
A closed irreducible analytic subvariety \(V\) of a complex manifold \(M\) of dimension \(n\) with \(\dim V=n\) is a connected component of \(M\); in particular it is a
closed connected submanifold of dimension \(n\). The identity of \(M\) has empty exceptional locus and strict transform \(V\), so it is an embedded resolution
of \(V\) in the sense of Lemma~\ref{hir:embedded-from-resolution}.
\end{lemma}
\begin{proof}
\leanok
At a regular point \(c\) of dimension \(n\), \(V\) is a neighbourhood of \(c\) (Lemma~\ref{mc:subvariety-thin}). All points of \(M\) are regular of dimension \(n\), and the
component of the regular locus of \(M\) through \(c\) is the connected component \(K\) of \(c\) in \(M\). By the identity theorem along a component of a regular locus
(proof of Proposition~\ref{ana:regular-component}) the analytic set \(V\), which contains a neighbourhood of \(c\), contains \(K\). As \(M\) is locally connected, \(K\)
is open and closed, so \(K\) and \(V\setminus K\) are analytic and \(V=K\cup(V\setminus K)\); by irreducibility \(V=K\), because \(c\in K\) does not lie in
\(V\setminus K\). Every point of the open set \(K\) is a regular point of \(K\) of dimension \(n\) (Lemma~\ref{lag:regular-point-basics}). The identity has bijective
differential, and the strict transform of the closed set \(V\) under it is \(V\).
\end{proof}

\begin{theorem}[Embedded resolution from a resolution invariant]
\label{mc:embedded-from-invariant}
\uses{mc:embedded-from-first-centre,mc:embedding-properties,mc:subvariety-thin,mc:subvariety-full-dimension,hir:invariant-driver,hir:centre-function}
\source{wlodarczyk05-hironaka,bierstone-milman08-functoriality}
\lean{Hironaka.EmbeddedFromResolution.exists_embeddedResolution_data_of_resolutionInvariant_of_isCoherent}
\leanok
Let \(X\) be a compact complex manifold of dimension \(n\), \(\operatorname{inv}\) a resolution invariant in dimension \(\le k\) with \(n\le k\)
(Definition~\ref{hir:centre-function}), and \(V\subset X\) a closed irreducible analytic subvariety whose ideal sheaf \(\mathcal I_V\) is coherent. There is a finite
composition of blowups with smooth centres whose exceptional image lies in \(V_{\mathrm{sing}}\) and whose strict transform of \(V\) is regular of dimension \(\dim V\)
at each of its points.
\end{theorem}
\begin{proof}
\leanok
If \(\dim V=n\), the identity (Lemma~\ref{mc:subvariety-full-dimension}). Otherwise \(\underline{\mathcal I}_V\) lies in the domain (Lemma~\ref{mc:subvariety-thin}), and
Theorem~\ref{mc:embedded-from-first-centre} applies to the canonical resolution of Theorem~\ref{hir:invariant-driver}(4).
\end{proof}
By Cartan's theorem (Theorem~\ref{ana:cartan-coherence}) the coherence hypothesis holds for every closed analytic subvariety; with it,
Theorem~\ref{mc:embedded-from-invariant} is Lemma~\ref{hir:embedded-from-resolution} for a resolution invariant in dimension \(\le n\).

\subsection{Re-basing the restricted marked ideal of Step 1b}

In this subsection \(\underline{\mathcal I}=(X,N,E,\mathfrak a,d)\), \(\overline{\underline{\mathcal J}}\), \(u\), \(U\) and \(\underline{\mathcal S}_{u,T_0}\) are as in
Theorem~\ref{mc:contact-stage-tail}; \(T_0T\) is a multiple test blow-up of \(\overline{\underline{\mathcal J}}\) and \(W\subset U_{T_0}\) an open set of the last stage of
\(T_0\) which the centres of \(T\) miss.

\begin{theorem}[Re-basing the restricted marked ideal of Step 1b]
\label{mc:contact-stage-rebase}
\uses{mc:coefficient-empty-step,mc:coefficient-restricted-tower,hir:restricted-equivalence-basic,bcalc:restrict-empty-tower,bcalc:member-intersection-tower,bcalc:member-locus-level,mc:contact-stage-domain-masked,mc:homogenized-equivalent-along,hir:giraud-tower-along,hir:geo-transform-basic,hir:canonical-segments,hir:centre-function,hir:invariant-driver,bcalc:inv-rebase,hir:step-one-b-independence}
\source{wlodarczyk05-hironaka}
\lean{Hironaka.ContactStageRebase.restrictEquiv_restrictedCoeff_geoTransformAlong_of_isEmptyTower,Hironaka.ContactStageRebase.restrictEquiv_restrictedCoeffHomAlong_rebase,Hironaka.ContactStageRebase.restrictEquiv_contactMarked_rebase,Hironaka.ContactStageRebase.contactMarked_rebase}
\leanok
\begin{enumerate}
\item Let \(\underline{\mathcal I}'=(Y,N',E',\mathfrak b,d')\) have \(N'\) a closed submanifold, \(\mathfrak b\) coherent and \(d'\ge1\), \(A\subset N'\) a closed submanifold and \(T'\) a
tower over \(Y\) all of whose centres are empty. Then \(\mathcal C_{N'_{T'}}(\underline{\mathcal I}'_{T'})|_{A_{T'}}\simeq(\mathcal C_{N'}(\underline{\mathcal I}')|_A)_{T'}\) (restricted
equivalence).
\item \(\underline{\mathcal S}_{u,T_0T}|_{W_T}\simeq(\underline{\mathcal S}_{u,T_0}|_W)_{T|_W}\): the restricted marked ideal formed at the later stage is
restricted-equivalent to the transform along \(T|_W\) of the one formed at \(T_0\). For \(E=\emptyset\) these are the marked ideals of
Lemma~\ref{hir:step-one-b-independence} on \(W_T\) and on \(W\).
\item Let \(\operatorname{inv}\) be a resolution invariant in dimension \(\le k\) (Definition~\ref{hir:centre-function}) and \(E=\emptyset\). If
\(\underline{\mathcal S}_{u,T_0}|_W\) lies in its domain and \(T|_W\,T_1\) is a canonical segment of it, then \(\underline{\mathcal S}_{u,T_0T}|_{W_T}\) lies in the domain,
\(T_1\) is a canonical segment of it, its transforms along \(T_1\) lie in the domain, and
\(\operatorname{inv}_{\underline{\mathcal S}_{u,T_0}|_W}(T|_W\,T_1)=\operatorname{inv}_{\underline{\mathcal S}_{u,T_0T}|_{W_T}}(T_1)\) on the cosupport.
\end{enumerate}
\end{theorem}
\begin{proof}
\leanok
(1) Induction on \(T'\). A step with empty centre after \(T''\) is a multiple test blow-up step, and by Lemma~\ref{mc:coefficient-empty-step}
\(\mathcal C(\underline{\mathcal I}'_{T''\sigma})|_{A_{T''\sigma}}\simeq(\mathcal C(\underline{\mathcal I}'_{T''})|_{A_{T''}})'\). By induction
\(\mathcal C(\underline{\mathcal I}'_{T''})|_{A_{T''}}\simeq(\mathcal C(\underline{\mathcal I}')|_A)_{T''}\); restricted-equivalent marked ideals with closed submanifolds
(\(A_{T''}\) is closed, Lemma~\ref{bcalc:member-intersection-tower}) and coherent ideals have equal transforms (Lemma~\ref{hir:restricted-equivalence-basic}), so the
right-hand side transforms to \((\mathcal C(\underline{\mathcal I}')|_A)_{T''\sigma}\).

(2) The restriction \(T|_W\) has only empty centres (Lemma~\ref{bcalc:restrict-empty-tower}). Apply (1) on \(W\) to
\(\underline{\mathcal I}'=\overline{\underline{\mathcal J}}_{T_0}|_W\) and \(A=S\cap W\), \(S=(N\cap V(u))_{T_0}\), a closed submanifold of \(W\) contained in \(N_{T_0}\)
(Theorems~\ref{mc:contact-stage-domain-masked} and~\ref{hir:giraud-tower-along}); \(T_0\) is a multiple test blow-up of \(\underline{\mathcal I}\) as well
(Theorem~\ref{mc:homogenized-equivalent-along}). The left-hand side is, by Lemma~\ref{mc:coefficient-restricted-tower}, the restriction to \(W_T\) of
\(\mathcal C(\overline{\underline{\mathcal J}}_{T_0T})|_{S_T}\), and \(S_T=(N\cap V(u))_{T_0T}\), \(\overline{\underline{\mathcal J}}_{T_0T}=(\overline{\underline{\mathcal J}}_{T_0})_T\)
(Theorem~\ref{bcalc:member-locus-level} and Lemma~\ref{hir:geo-transform-basic}); the coefficient ideal and the restriction to \(S\) commute with restriction to \(W\), so the
right-hand side is the transform of \(\underline{\mathcal S}_{u,T_0}|_W\).

(3) Re-basing after the empty prefix \(T|_W\) (Lemma~\ref{bcalc:inv-rebase}, from property (4') and property (7) of the invariant) with the restricted
equivalence of (2), whose left-hand side has a coherent ideal; the transforms of a marked ideal of the domain along its canonical segments lie in the domain
(Theorem~\ref{hir:invariant-driver}).
\end{proof}
This is the comparison of nearby points in the value of Step 1b of Proposition~\ref{hir:maximal-order-step}: a point whose regime started at \(T_0\) and a
nearby point whose regime started later, at \(T_0T\), with the steps of \(T\) empty near them (Lemma~\ref{hir:canonical-segments}); Step 1a has the
same comparison for its restricted marked ideals (Lemma~\ref{bcalc:inv-rebase}).

\subsection{Independence of the representative in Step 1}

\begin{theorem}[Step 1 does not depend on the representative]
\label{mc:contact-representative}
\uses{mc:homogenized-along,mc:transform-along-calculus,mc:coefficient-representatives-along,mc:coefficient-representatives-tower,mc:homogenized-along-basic,ana:restricted-calculus,mc:tangent-direction-along,mc:tangent-direction-along-restricted,mc:strict-along-local,mc:contact-independence-along,hir:marked-equivalence-basic,hir:centre-function,bcalc:inv-rebase,hir:step-one-b-independence}
\source{wlodarczyk05-hironaka}
\lean{Hironaka.ContactRepresentative.restrictLe_homogenizedIdealAlong,Hironaka.ContactRepresentative.restrictEquiv_homMarkedAlong,Hironaka.ContactRepresentative.restrictEquiv_barMarked,Hironaka.ContactRepresentative.restrictEquiv_pullback_val,Hironaka.ContactRepresentative.exists_isTangentDirectionAlong_of_restrictLe,Hironaka.ContactRepresentative.restrictEquiv_contactMarked,Hironaka.ContactRepresentative.contactMarked_representative}
\leanok
Let \(N\) be a closed submanifold.
\begin{enumerate}
\item If \(\mathfrak a|_N\subset\mathfrak b|_N\), then \(\mathcal H_N(\mathfrak a,\mu)|_N\subset\mathcal H_N(\mathfrak b,\mu)|_N\). Hence restricted-equivalent marked ideals
\(\underline{\mathcal I}\simeq\underline{\mathcal I}'\) (same \(N\), \(E\), mark, and \(\mathfrak a|_N=\mathfrak a'|_N\)) with coherent ideals have restricted-equivalent homogenized
marked ideals and restricted-equivalent marked ideals of Step 1, \(\mathcal C_N(\mathcal H_N(\underline{\mathcal I}))\simeq\mathcal C_N(\mathcal H_N(\underline{\mathcal I}'))\);
restricted equivalence passes to open sets.
\item If \(\mathfrak a|_N\subset\mathfrak b|_N\), \(\mathfrak b\) coherent, and \(u\) is a tangent direction of \((\mathfrak a|_N,\mu)\) along \(N\) on \(U\), then every point of
\(N\cap U\) has an open neighbourhood \(O\subset U\) carrying a tangent direction \(w\) of \((\mathfrak b|_N,\mu)\) along \(N\) with \(w=u\) on \(N\cap O\).
\item Let \((M,N,\emptyset,\mathfrak a,\mu)\simeq(M,N,\emptyset,\mathfrak a',\mu)\) with coherent ideals, \(u\), \(w\) functions equal on \(N\cap O\) for an open \(O\), \(T_0\) a
multiple test blow-up of \(\overline{\underline{\mathcal J}}=\mathcal C_N(\mathcal H_N(\mathfrak a,\mu))\) and \(W\subset O_{T_0}\) open. Then the marked ideals of
Lemma~\ref{hir:step-one-b-independence} formed from \((\mathfrak a,u)\) and from \((\mathfrak a',w)\) on \(W\) are restricted-equivalent. Consequently, for a resolution
invariant (Definition~\ref{hir:centre-function}), if the first lies in its domain so does the second, they have the same canonical segments, and along them the
invariant takes the same values on the cosupport.
\end{enumerate}
\end{theorem}
\begin{proof}
\leanok
(1) \(\mathcal H_N(\mathfrak a,\mu)=\sum_{i<\mu}\mathcal D_N^i(\mathfrak a)\,\mathcal D_N^{\mu-1}(\mathfrak a)^i\) (Definition~\ref{mc:homogenized-along}); the derivative ideals along \(N\) respect
restricted inclusions (Lemma~\ref{mc:coefficient-representatives-along}), and so do products, powers and finite sums (Lemma~\ref{mc:transform-along-calculus}). The
homogenized ideals are coherent (Lemma~\ref{mc:homogenized-along-basic}), so the coefficient ideals of the restricted-equivalent homogenized marked ideals are
restricted-equivalent (Lemma~\ref{mc:coefficient-representatives-along}). A relation on \(N\) near \(f(y)\) composes with the inclusion of an open set
(Lemma~\ref{ana:restricted-calculus}).
(2) On \(N\) near a point of \(U\), the germ of \(u\in\mathcal D_N^{\mu-1}(\mathfrak a)\) is a combination of germs of \(\mathcal D_N^{\mu-1}(\mathfrak b)\) (Lemma~\ref{mc:coefficient-representatives-along}),
it is holomorphic, and it has multiplicity one along \(N\) at its zeros (Definition~\ref{mc:tangent-direction-along}); Lemma~\ref{mc:tangent-direction-along-restricted}
gives \(w\).
(3) By (1) the two marked ideals of Step 1 are restricted-equivalent, so they have the same multiple test blow-ups (Lemma~\ref{hir:marked-equivalence-basic}) and, at
every stage, restricted-equivalent restrictions of their coefficient ideals to the strict transform of any subset of \(N\) (Lemma~\ref{mc:coefficient-representatives-tower},
with the mark and coherence of Theorem~\ref{mc:contact-independence-along}); restrict to \(W\) by (1). The sets \(N\cap V(u)\) and \(N\cap V(w)\) agree on \(O\), so their
strict transforms agree over \(O\) (Lemma~\ref{mc:strict-along-local}), in particular on \(W\). Domain and canonical segments pass along restricted equivalence
(Lemma~\ref{bcalc:inv-rebase}), and the values agree by property (7) of Definition~\ref{hir:centre-function}.
\end{proof}
With (2), a tangent direction chosen for one representative of the marked ideal of maximal order gives, near each point, one for any other representative with
the same hypersurface of maximal contact inside \(N\); with (3) and Lemma~\ref{hir:step-one-b-independence} (independence of the tangent direction), the value of
Step 1b in Proposition~\ref{hir:maximal-order-step} does not depend on the representative, which is property (7) of Definition~\ref{hir:centre-function} for it.

\subsection{Step 1b along isomorphisms of towers and on restricted towers}

\begin{theorem}[The restricted marked ideal of Step 1b along isomorphisms and restrictions]
\label{mc:contact-marked-canonical}
\uses{hir:step-one-b-independence,mc:homogenized-along-etale,mc:homogenized-along-basic,mc:contact-independence-along,mc:tangent-direction-along-open,hir:tower-isomorphism,mc:tower-isomorphism-transform,mc:tower-isomorphism-coefficient,mc:coefficient-restricted-tower,hir:open-restriction-tower,hir:marked-operations-transform,bcalc:cancellation,bcalc:centre-local,bcalc:restrict-nested,bcalc:inv-domain-stable,hir:canonical-segments,hir:invariant-driver,hir:centre-function}
\source{wlodarczyk05-hironaka}
\lean{Hironaka.ContactMarkedCanonical.pullback_pullback_of_comp_eq,Hironaka.ContactMarkedCanonical.stepOneMarked_pullback,Hironaka.ContactMarkedCanonical.isTangentDirectionAlong_comp,Hironaka.ContactMarkedCanonical.contactMarked_isTowerIso,Hironaka.ContactMarkedCanonical.contactMarked_isTowerIso_refl,Hironaka.ContactMarkedCanonical.inv_contactMarked_isTowerIso,Hironaka.ContactMarkedCanonical.stepOneMarked_pullback_val,Hironaka.ContactMarkedCanonical.isTangentDirectionAlong_comp_val,Hironaka.ContactMarkedCanonical.contactMarked_restrict,Hironaka.ContactMarkedCanonical.inv_contactMarked_restrict}
\leanok
Let \(N\) be a closed submanifold of \(X\), \(\mathfrak a\) coherent, \(\mu\ge1\), \(\overline{\underline{\mathcal J}}=\mathcal C_N(\mathcal H_N(\mathfrak a,\mu))\) and \(u\colon X\to\mathbf C\).
For a multiple test blow-up \(T_0\) of \(\overline{\underline{\mathcal J}}\) and an open set \(W\) of its last stage, write \(\underline{\mathcal S}_{u,T_0}|_W\) for the marked ideal
of Lemma~\ref{hir:step-one-b-independence}, the coefficient ideal of \(\overline{\underline{\mathcal J}}_{T_0}\) restricted to the strict transform of \(N\cap V(u)\) and to \(W\).
\begin{enumerate}
\item (Biholomorphisms.) Let \(\Phi\colon X'\to X\) be a biholomorphism. The marked ideal of Step 1 formed from \((\Phi^{-1}N,\Phi^*\mathfrak a,\mu)\) is
\(\Phi^*\overline{\underline{\mathcal J}}\), and if \(u\) is a tangent direction of \((\mathfrak a|_N,\mu)\) along \(N\) on \(U\), then \(u\circ\Phi\) is one of
\((\Phi^*\mathfrak a|_{\Phi^{-1}N},\mu)\) along \(\Phi^{-1}N\) on \(\Phi^{-1}U\). Let \(T_0'\) be a tower isomorphic to \(T_0\) over \(\Phi\)
(Definition~\ref{hir:tower-isomorphism}) with last-stage biholomorphism \(\chi_0\), let \(W'\), \(W\) be open sets of the two last stages and \(\psi\colon W'\to W\) a
biholomorphism with \(\psi=\chi_0\) on \(W'\). Then \(\underline{\mathcal S}'_{u\circ\Phi,T_0'}|_{W'}=\psi^*\bigl(\underline{\mathcal S}_{u,T_0}|_W\bigr)\), where
\(\underline{\mathcal S}'\) is formed from \((\Phi^{-1}N,\Phi^*\mathfrak a,\mu)\). For \(\Phi=\mathrm{id}\) this reads
\(\underline{\mathcal S}_{u,T_0'}|_{W'}=\psi^*\bigl(\underline{\mathcal S}_{u,T_0}|_W\bigr)\).
\item In (1), let \(\operatorname{inv}\) be a resolution invariant (Definition~\ref{hir:centre-function}), \(\underline{\mathcal S}_{u,T_0}|_W\) in its domain, \(T_1\) a canonical
segment of it and \(T_1'\) a tower isomorphic to \(T_1\) over \(\psi\), with last-stage biholomorphism \(\xi\). Then \(\underline{\mathcal S}'_{u\circ\Phi,T_0'}|_{W'}\) lies
in the domain, \(T_1'\) is a canonical segment of it, and its value along \(T_1'\) at \(p\) is the value of \(\underline{\mathcal S}_{u,T_0}|_W\) along \(T_1\) at
\(\xi(p)\) whenever \(\xi(p)\) lies in the cosupport.
\item (Open sets.) Let \(U\subset X\) be open with inclusion \(\iota\). The marked ideal of Step 1 formed from \((\iota^{-1}N,\iota^*\mathfrak a,\mu)\) is
\(\iota^*\overline{\underline{\mathcal J}}\), and \(u\circ\iota\) is a tangent direction on \(\iota^{-1}U_0\) if \(u\) is one on \(U_0\). For \(W\subset T_0^{-1}U\) open, the
marked ideal formed from \((\iota^{-1}N,\iota^*\mathfrak a,\mu)\) along \(T_0|_U\) with \(u\circ\iota\), on \(W\) seen as an open subset of \(T_0^{-1}U\), is the transport of
\(\underline{\mathcal S}_{u,T_0}|_W\) along the biholomorphism of Lemma~\ref{bcalc:centre-local}(1), the identity on points.
\item Let \(T_1\) be a tower over the last stage of \(T_0\). If \(\underline{\mathcal S}_{u,T_0}|_W\) lies in the domain of \(\operatorname{inv}\) and \(T_1|_W\) is a canonical
segment of it, then the marked ideal of (3) lies in the domain, the restriction to \(W\) of \(T_1|_{T_0^{-1}U}\) is a canonical segment of it, and at two points of the
last stages with the same image, one lies in the cosupport iff the other does, and then the values agree.
\end{enumerate}
\end{theorem}
\begin{proof}
\leanok
(1) Pullback by a biholomorphism commutes with \(\mathcal H_N\) and \(\mathcal C_N\) (Lemma~\ref{mc:homogenized-along-etale}; the homogenized ideal is coherent,
Lemma~\ref{mc:homogenized-along-basic}), and tangent directions pull back by local biholomorphisms (Lemma~\ref{mc:tangent-direction-along-open}). Since
\(\Phi^{-1}N\cap V(u\circ\Phi)=\Phi^{-1}(N\cap V(u))\), Lemma~\ref{mc:tower-isomorphism-coefficient} for \(\overline{\underline{\mathcal J}}\) (closed submanifold, coherent ideal,
mark \(\ge1\); Theorem~\ref{mc:contact-independence-along}) gives
\(\underline{\mathcal S}_{\Phi^{-1}(N\cap V(u))}((\Phi^*\overline{\underline{\mathcal J}})_{T_0'})=\chi_0^*\,\underline{\mathcal S}_{N\cap V(u)}(\overline{\underline{\mathcal J}}_{T_0})\).
The marked ideal on the right has a coherent ideal (Lemmas~\ref{hir:open-restriction-tower} and~\ref{mc:homogenized-along-basic}), and pullback of coherent ideals is
functorial; since \(\chi_0\circ\iota_{W'}=\iota_W\circ\psi\), pulling it back along the two sides of this square gives the same marked ideal
(Lemma~\ref{hir:marked-operations-transform} for the comparison of the components). For \(\Phi=\mathrm{id}\), \(\mathrm{id}^*\mathfrak a=\mathfrak a\)
(Lemma~\ref{bcalc:cancellation}).
(2) By (1) the first marked ideal is \(\psi^*(\underline{\mathcal S}_{u,T_0}|_W)\). The domain is invariant under biholomorphisms (Proposition~\ref{bcalc:inv-domain-stable}),
canonical segments correspond along isomorphisms of towers (Lemma~\ref{hir:canonical-segments}), and the values by property (6) of
Definition~\ref{hir:centre-function}.
(3) As in (1), with Lemma~\ref{mc:homogenized-along-etale} for the inclusion of an open set, Lemma~\ref{mc:coefficient-restricted-tower},
\(\underline{\mathcal S}_{\iota^{-1}A}((\iota^*\overline{\underline{\mathcal J}})_{T_0|_U})=\underline{\mathcal S}_A(\overline{\underline{\mathcal J}}_{T_0})|_{T_0^{-1}U}\), and
Lemma~\ref{bcalc:centre-local}(1) for restricting first to \(T_0^{-1}U\) and then to \(W\).
(4) By Lemma~\ref{bcalc:restrict-nested}(1) the restriction to \(W\) of \(T_1|_{T_0^{-1}U}\) is isomorphic to \(T_1|_W\) over the biholomorphism of (3), with the identity on
points at the last stage. By (3) the marked ideal is the transport of \(\underline{\mathcal S}_{u,T_0}|_W\) along it, so (2) applies (with \(\Phi=\psi\) the biholomorphism
of (3)); \(T_1|_W\) is a multiple test blow-up (Theorem~\ref{hir:invariant-driver}), so the cosupports correspond under the last-stage map
(Lemma~\ref{mc:tower-isomorphism-transform}).
\end{proof}
In the value function of Step 1 (Proposition~\ref{hir:maximal-order-step}), the marked ideal of Step 1b at a point is formed with a tangent direction chosen for the
marked ideal at hand. For \(\Phi^*\underline{\mathcal J}\), resp.\ \(\underline{\mathcal J}|_U\), that choice differs from \(u\circ\Phi\), resp.\ \(u\circ\iota\); after the
change of tangent direction of Lemma~\ref{hir:step-one-b-independence} on small open sets, (1)--(4) are the Step 1b part of properties (5) and (6) of
Definition~\ref{hir:centre-function} for that value function.

\subsection{The change of tangent direction on arbitrary open sets}

\begin{lemma}[Isomorphisms of towers on open sets]
\label{mc:tower-iso-opens}
\uses{hir:tower-isomorphism,mc:tower-isomorphism-basic,hir:open-restriction-tower,bcalc:restrict-open,bcalc:tower-restrict,hir:tower-split-iso,hir:marked-operations-transform}
\source{wlodarczyk05-hironaka}
\lean{Hironaka.StepOneTangentChange.diffeoOpens,Hironaka.StepOneTangentChange.coe_diffeoOpens_apply,Hironaka.StepOneTangentChange.isTowerIso_mem_openPreimage_iff,Hironaka.StepOneTangentChange.isTowerIso_restrictOpens,Hironaka.StepOneTangentChange.pullback_pullback_of_comp,Hironaka.StepOneTangentChange.pullback_replicate_empty}
\leanok
\begin{enumerate}
\item Let \(T'\) be a tower isomorphic to \(T\) over a biholomorphism \(\Phi\colon X'\to X\) (Definition~\ref{hir:tower-isomorphism}) with last-stage biholomorphism
\(\chi\), and let \(V'\subset X'\), \(V\subset X\) be open with \(\Phi^{-1}(V)=V'\). Then \(\Phi\) restricts to a biholomorphism \(V'\to V\), \(\chi\) maps
\(T'^{-1}V'\) onto \(T^{-1}V\), and \(T'|_{V'}\) is isomorphic to \(T|_V\) over \(\Phi|_{V'}\), with last-stage biholomorphism \(\chi|_{T'^{-1}V'}\).
\item For holomorphic maps \(a\), \(b\), \(d\) with \(a\circ b=d\) and a marked ideal \(\underline{\mathcal I}\) with a coherent ideal, \(b^*a^*\underline{\mathcal I}=d^*\underline{\mathcal I}\);
and the pullback of a marked ideal whose divisor begins with \(k\) empty members is the pullback with \(k\) empty members in front of its divisor.
\end{enumerate}
\end{lemma}
\begin{proof}
\leanok
(1) The stage biholomorphisms commute with the projections to the bases (Lemma~\ref{mc:tower-isomorphism-basic}), so a point of a stage of \(T'\) lies over \(V'\) iff
its image lies over \(V\). Induction along the towers: for the empty towers the last-stage map is \(\Phi\) itself, restricted (Lemma~\ref{hir:tower-split-iso} for the
identification of the stages); at a blowup, the restricted towers have as next step the blowup along the trace of the centre on the open preimage
(Lemmas~\ref{hir:open-restriction-tower}, \ref{bcalc:restrict-open} and~\ref{bcalc:tower-restrict}), the restriction of the next stage map commutes with the restricted
blowups because the unrestricted maps do, and it carries the trace of one centre onto the trace of the other.
(2) The submanifold and the members are preimages, which compose, and the preimage of the empty set is empty; the ideals of coherent pullbacks compose; the
components are compared as in Lemma~\ref{hir:marked-operations-transform}.
\end{proof}

\begin{theorem}[The change of tangent direction on arbitrary open sets]
\label{mc:tangent-change-opens}
\uses{mc:tower-iso-opens,hir:step-one-b-independence,mc:glueing-trivial-along,mc:glueing-biholomorphism-along,mc:glueing-lift-tail,mc:contact-stage-tail,bcalc:restrict-nested,hir:drop-original-members,mc:homogenized-along-basic,hir:open-restriction-tower,hir:centre-function}
\source{wlodarczyk05-hironaka}
\lean{Hironaka.StepOneTangentChange.exists_tail_iso_opens,Hironaka.StepOneTangentChange.inv_eq_of_tail_iso,Hironaka.StepOneTangentChange.inv_contactMarked_eq_of_opens,Hironaka.StepOneTangentChange.pullback_val_restrictedCoeffHom_eq,Hironaka.StepOneTangentChange.inv_restrictedCoeffHom_eq_of_opens,Hironaka.StepOneTangentChange.inv_restrictedCoeffHom_eq_of_opens_disjoint}
\leanok
\begin{enumerate}
\item Let \(\mathfrak a\) be coherent, \(N\) a closed submanifold, \(\mu\ge1\), \(\overline{\underline{\mathcal J}}=\mathcal C_N(\mathcal H_N(\mathfrak a,\mu))\), \(x\) a point of
\(\operatorname{cosupp}(\mathfrak a|_N,\mu)\) and \(u\), \(v\) tangent directions of \((\mathfrak a|_N,\mu)\) along \(N\) on open sets \(U_1\ni x\), \(U_2\ni x\). There is an open
\(O\ni x\) with the following property. Let \(T_0T_1\) be a multiple test blow-up of \(\overline{\underline{\mathcal J}}\), \(W\subset T_0^{-1}U_1\) and \(W'\subset T_0^{-1}U_2\)
open subsets of the last stage of \(T_0\), and \(\operatorname{inv}\) a resolution invariant (Definition~\ref{hir:centre-function}). If the marked ideals
\(\underline{\mathcal S}_{u,T_0}|_W\) and \(\underline{\mathcal S}_{v,T_0}|_{W'}\) of Lemma~\ref{hir:step-one-b-independence} lie in the domain of \(\operatorname{inv}\) and \(T_1|_W\),
\(T_1|_{W'}\) are canonical segments of them, then their values along \(T_1|_W\) and \(T_1|_{W'}\) agree at every point of
\(\operatorname{cosupp}\overline{\underline{\mathcal J}}_{T_0T_1}\) over \(W\cap W'\cap T_0^{-1}O\).
\item Let \(\underline{\mathcal J}=(X,N,E,\mathfrak a,d)\) with \(\mathfrak a\) coherent, \(N\) a closed submanifold, \(d\ge1\), \(\overline{\underline{\mathcal J}}=\mathcal C_N(\mathcal H_N(\underline{\mathcal J}))\)
with the divisor \(E\), and write \(\underline{\mathcal S}^E_{w,T_0}|_W\) for the coefficient ideal of \(\overline{\underline{\mathcal J}}_{T_0}\) restricted to
\((N\cap V(w))_{T_0}\) and to \(W\), with the divisor of \(\overline{\underline{\mathcal J}}_{T_0}\). If \(W\) meets none of the first \(\#E\) members of the divisor of
\(\overline{\underline{\mathcal J}}_{T_0}\) (the strict transforms of the members of \(E\)), then \(\underline{\mathcal S}^E_{w,T_0}|_W\) is the marked ideal
\(\underline{\mathcal S}_{w,T_0}|_W\) of (1) for \((\mathfrak a,d)\) with \(\#E\) empty members in front of its divisor. The statement of (1) holds for
\(\underline{\mathcal S}^E_{u,T_0}|_W\) and \(\underline{\mathcal S}^E_{v,T_0}|_{W'}\), for multiple test blow-ups of \(\overline{\underline{\mathcal J}}\) and open sets \(W\), \(W'\)
meeting none of these members.
\end{enumerate}
\end{theorem}
\begin{proof}
\leanok
(1) Take glueing data along \(N\) for \(u\), \(v\) at \(x\) acting trivially modulo the glue ideal (Lemma~\ref{mc:glueing-trivial-along}), with \(x\) in its domain \(U'\) and
codomain \(U\) (Lemma~\ref{mc:glueing-biholomorphism-along}), and \(O=U'\cap U\). For \(T_0T_1\), Theorem~\ref{mc:glueing-lift-tail} gives the lift
\(\Phi_0\colon T_0^{-1}U'\to T_0^{-1}U\) with \(\Phi_0^*(\underline{\mathcal S}_{u,T_0}|_{T_0^{-1}U})=\underline{\mathcal S}_{v,T_0}|_{T_0^{-1}U'}\) and an isomorphism
\(\chi\colon T_1|_{T_0^{-1}U'}\to T_1|_{T_0^{-1}U}\) of towers over \(\Phi_0\) fixing \(\operatorname{cosupp}\overline{\underline{\mathcal J}}_{T_0T_1}\) pointwise. Let \(z\) be a
point of this cosupport over \(W\cap W'\cap T_0^{-1}O\) and \(z_0\) its image in the last stage of \(T_0\); since \(\chi(z)=z\) and \(\chi\) commutes with the projections,
\(\Phi_0(z_0)=z_0\). Let \(V_v\) be the set of points of \(W'\cap T_0^{-1}U'\) mapped into \(W\) by \(\Phi_0\), and \(V_u=\Phi_0(V_v)\subset W\cap T_0^{-1}U\); both are open and
contain \(z_0\). Lemma~\ref{mc:tower-iso-opens}(1) restricts \(\chi\) to an isomorphism of the restrictions of \(T_1|_{T_0^{-1}U'}\) and \(T_1|_{T_0^{-1}U}\) to \(V_v\) and
\(V_u\); with Lemma~\ref{bcalc:restrict-nested}(1) on both sides, \(T_1|_{V_v}\) is isomorphic to \(T_1|_{V_u}\) over \(\psi=\Phi_0|_{V_v}\), fixing \(z\), and
\(\psi^*(\underline{\mathcal S}_{u,T_0}|_{V_u})=\underline{\mathcal S}_{v,T_0}|_{V_v}\) by Lemma~\ref{mc:tower-iso-opens}(2) along the commuting squares of inclusions (the ideals
are coherent, Lemmas~\ref{mc:homogenized-along-basic} and~\ref{hir:open-restriction-tower}). By Lemma~\ref{bcalc:restrict-nested}(2), \(\underline{\mathcal S}_{u,T_0}|_{V_u}\) lies in
the domain, \(T_1|_{V_u}\) is a canonical segment of it, and its value at \(z\) is the value of \(\underline{\mathcal S}_{u,T_0}|_W\) along \(T_1|_W\); likewise for \(v\) on
\(V_v\subset W'\). The point \(z\) lies in both cosupports, which are the traces of \(\operatorname{cosupp}\overline{\underline{\mathcal J}}_{T_0T_1}\)
(Theorem~\ref{mc:contact-stage-tail}). Property (6) of Definition~\ref{hir:centre-function} along \(T_1|_{V_v}\cong T_1|_{V_u}\) identifies the value of
\(\underline{\mathcal S}_{v,T_0}|_{V_v}=\psi^*(\underline{\mathcal S}_{u,T_0}|_{V_u})\) at \(z\) with that of \(\underline{\mathcal S}_{u,T_0}|_{V_u}\) at \(\chi(z)=z\).
(2) Along every tower the divisor of \(\overline{\underline{\mathcal J}}_{T_0}\) is that of the transform of \(\mathcal C_N(\mathcal H_N(\mathfrak a,d))\) with the strict transforms of
the members of \(E\) in front, and the submanifold and ideal are the same (Lemma~\ref{hir:drop-original-members}); the coefficient ideal and the restriction to a
subset of \(N\) do not change the divisor, and the traces of those members on \(W\) are empty. A multiple test blow-up of \(\overline{\underline{\mathcal J}}\) is one of
\(\mathcal C_N(\mathcal H_N(\mathfrak a,d))\) with the same cosupports (Lemma~\ref{hir:drop-original-members}), so the construction of (1) applies; by
Lemma~\ref{mc:tower-iso-opens}(2), \(\psi^*(\underline{\mathcal S}^E_{u,T_0}|_{V_u})=\underline{\mathcal S}^E_{v,T_0}|_{V_v}\) (\(V_u\subset W\) and \(V_v\subset W'\) meet none of
the members), and the values are compared on \(\underline{\mathcal S}^E\) exactly as in (1), with the cosupports of Theorem~\ref{mc:contact-stage-tail} (any divisor).
\end{proof}
In the value function of Step 1 (Proposition~\ref{hir:maximal-order-step}) the marked ideal of Step 1b at a point is formed with the tangent direction chosen at its
base point and restricted to its mask, which misses the original members; (2) compares the values of two points computed with different tangent directions and
masks, with only the hypotheses on the two masks.

\subsection{The restricted marked ideal of Step 1b with the divisor kept}

\begin{theorem}[Step 1b with the divisor kept: transport, restriction, re-basing, representatives]
\label{mc:contact-stage-kept}
\uses{mc:tangent-change-opens,mc:tower-iso-opens,mc:contact-marked-canonical,mc:contact-stage-rebase,mc:contact-representative,mc:homogenized-along-etale,mc:homogenized-along-basic,mc:tower-isomorphism-coefficient,mc:tower-isomorphism-transform,mc:tower-isomorphism-restrict-append,mc:coefficient-restricted-tower,mc:coefficient-representatives-tower,mc:strict-along-local,mc:restriction-operations-along,hir:open-restriction-tower,hir:marked-equivalence-basic,hir:canonical-segments,hir:invariant-driver,hir:centre-function,bcalc:centre-local,bcalc:restrict-nested,bcalc:restrict-empty-tower,bcalc:inv-rebase,bcalc:inv-domain-stable}
\source{wlodarczyk05-hironaka}
\lean{Hironaka.ContactStageKept.contactMarkedKept,Hironaka.ContactStageKept.isCoherent_restrictedCoeff,Hironaka.ContactStageKept.barMarked_pullback,Hironaka.ContactStageKept.contactMarkedKept_isTowerIso,Hironaka.ContactStageKept.inv_contactMarkedKept_isTowerIso,Hironaka.ContactStageKept.barMarked_pullback_val,Hironaka.ContactStageKept.contactMarkedKept_restrict,Hironaka.ContactStageKept.inv_contactMarkedKept_restrict,Hironaka.ContactStageKept.mem_cosupp_of_nested,Hironaka.ContactStageKept.contactMarkedKept_rebase,Hironaka.ContactStageKept.contactMarkedKept_rebase_nested,Hironaka.ContactStageKept.restrictEquiv_contactMarkedKept,Hironaka.ContactStageKept.contactMarkedKept_representative}
\leanok
Let \(\underline{\mathcal J}=(X,N,E,\mathfrak a,d)\) with \(\mathfrak a\) coherent, \(N\) a closed submanifold and \(d\ge1\), \(\overline{\underline{\mathcal J}}=\mathcal C_N(\mathcal H_N(\underline{\mathcal J}))\)
with the divisor \(E\), and write \(\underline{\mathcal S}^E_{w,T_0}|_W\) for the coefficient ideal of \(\overline{\underline{\mathcal J}}_{T_0}\) along \(N_{T_0}\), restricted to the
strict transform of \(N\cap V(w)\) and to an open set \(W\) of the last stage of a tower \(T_0\), with the divisor of \(\overline{\underline{\mathcal J}}_{T_0}\)
(Theorem~\ref{mc:tangent-change-opens}(2)). Let \(\operatorname{inv}\) be a resolution invariant (Definition~\ref{hir:centre-function}) and \(T_0\) a multiple test blow-up
of \(\overline{\underline{\mathcal J}}\).
\begin{enumerate}
\item (Biholomorphisms.) \(\mathcal C(\mathcal H(\Phi^*\underline{\mathcal J}))=\Phi^*\overline{\underline{\mathcal J}}\) for a biholomorphism \(\Phi\colon X'\to X\). If \(T_0'\) is isomorphic to
\(T_0\) over \(\Phi\) with last-stage biholomorphism \(\chi_0\) and \(\psi\colon W'\to W\) is a biholomorphism of open sets of the last stages with \(\psi=\chi_0\) on \(W'\),
then \(\underline{\mathcal S}^{E}_{w\circ\Phi,T_0'}|_{W'}\), formed from \(\Phi^*\underline{\mathcal J}\), is \(\psi^*(\underline{\mathcal S}^E_{w,T_0}|_W)\). If moreover
\(\chi_0^{-1}(W)=W'\) and \(T_1'\) is isomorphic to \(T_1\) over \(\chi_0\), and \(\underline{\mathcal S}^E_{w,T_0}|_W\) lies in the domain of \(\operatorname{inv}\) with \(T_1|_W\) a
canonical segment of it, then the first marked ideal lies in the domain, \(T_1'|_{W'}\) is a canonical segment of it, and the values at corresponding points of the
cosupports agree.
\item (Open sets.) \(\mathcal C(\mathcal H(\underline{\mathcal J}|_U))=\overline{\underline{\mathcal J}}|_U\) for \(U\subset X\) open. For \(W\subset T_0^{-1}U\) open, the marked ideal formed
from \(\underline{\mathcal J}|_U\) along \(T_0|_U\) with \(w|_U\), on \(W\) seen in \(T_0^{-1}U\), is the transport of \(\underline{\mathcal S}^E_{w,T_0}|_W\) along the biholomorphism of
Lemma~\ref{bcalc:centre-local}(1); if \(\underline{\mathcal S}^E_{w,T_0}|_W\) lies in the domain and \(T_1|_W\) is a canonical segment of it, then the marked ideal formed from
\(\underline{\mathcal J}|_U\) lies in the domain, the restriction to \(W\) of \(T_1|_{T_0^{-1}U}\) is a canonical segment of it, and at points with the same image the cosupports
correspond and the values agree.
\item (Re-basing.) Let \(u\) be a tangent direction of \((\mathfrak a|_N,d)\) along \(N\) on \(U\), \(T_0T\) a multiple test blow-up of \(\overline{\underline{\mathcal J}}\) and
\(W\subset T_0^{-1}U\) open with the centres of \(T\) missing \(W\). If \(\underline{\mathcal S}^E_{u,T_0}|_W\) lies in the domain and \(T|_WT_1\) is a canonical segment of it, then
\(\underline{\mathcal S}^E_{u,T_0T}|_{T^{-1}W}\) lies in the domain, \(T_1\) is a canonical segment of it, its transforms along \(T_1\) lie in the domain, and the values agree on the
cosupport. For a tower \(T_1\) over the last stage of \(T_0T\) with \((TT_1)|_W\) a canonical segment of \(\underline{\mathcal S}^E_{u,T_0}|_W\) and an open \(W'\subset T^{-1}W\),
\(\underline{\mathcal S}^E_{u,T_0T}|_{W'}\) lies in the domain, \(T_1|_{W'}\) is a canonical segment of it, and at a point of its cosupport its value is the value of
\(\underline{\mathcal S}^E_{u,T_0}|_W\) along \((TT_1)|_W\) at the same point.
\item (Representatives.) If \(\underline{\mathcal J}\simeq\underline{\mathcal J}'\) are restricted-equivalent with coherent ideals and \(u=w\) on \(N\cap O\) for an open \(O\), then for
\(W\subset T_0^{-1}O\) open \(\underline{\mathcal S}^E_{u,T_0}|_W\simeq\underline{\mathcal S}'^E_{w,T_0}|_W\); if the first lies in the domain, so does the second, they have the same
canonical segments, and along them the values agree on the cosupport.
\end{enumerate}
\end{theorem}
\begin{proof}
\leanok
(1) and (2) are proved as Theorem~\ref{mc:contact-marked-canonical}, with \(\overline{\underline{\mathcal J}}\) carrying \(E\): pullback by a biholomorphism or the inclusion of
an open set commutes with \(\mathcal H_N\) and \(\mathcal C_N\) for any divisor (Lemma~\ref{mc:homogenized-along-etale}; Lemma~\ref{mc:homogenized-along-basic} for coherence),
Lemmas~\ref{mc:tower-isomorphism-coefficient} and~\ref{mc:coefficient-restricted-tower} give the restricted coefficient ideals along isomorphic and restricted towers, and
pullbacks along commuting squares agree (Lemma~\ref{mc:tower-iso-opens}(2)). For the values in (1), Lemma~\ref{mc:tower-iso-opens}(1) restricts the isomorphism of the
tails to \(W'\) and \(W\); the domain is invariant (Proposition~\ref{bcalc:inv-domain-stable}), canonical segments correspond (Lemma~\ref{hir:canonical-segments}), and property (6)
compares the values. In (2), Lemma~\ref{bcalc:restrict-nested}(1) is the isomorphism and the cosupports correspond under it (Lemma~\ref{mc:tower-isomorphism-transform}; canonical
segments are multiple test blow-ups, Theorem~\ref{hir:invariant-driver}).
(3) Theorem~\ref{mc:contact-stage-rebase} gives \(\underline{\mathcal S}^E_{u,T_0T}|_{T^{-1}W}\simeq(\underline{\mathcal S}^E_{u,T_0}|_W)_{T|_W}\) for any divisor; \(T|_W\) has only empty
centres (Lemma~\ref{bcalc:restrict-empty-tower}), and Lemma~\ref{bcalc:inv-rebase} gives the domain, the canonical segments and the values. For the later mask, \((TT_1)|_W\) is
isomorphic over the identity to \(T|_W\) followed by \(T_1|_{T^{-1}W}\) (Lemma~\ref{mc:tower-isomorphism-restrict-append}, \(\mathrm{id}^*=\mathrm{id}\) by
Lemma~\ref{mc:restriction-operations-along}), so the first part applies to \(T_1|_{T^{-1}W}\); Lemma~\ref{bcalc:restrict-nested}(2) passes to \(W'\subset T^{-1}W\). A point of the
cosupport over \(W'\) lies in the cosupport over \(T^{-1}W\) (the restriction to \(T^{-1}W\) and then to \(W'\) is isomorphic to the restriction to \(W'\),
Lemma~\ref{bcalc:restrict-nested}(1), and cosupports along restrictions are traces, Lemma~\ref{hir:canonical-segments}), hence in the cosupport of
\(\underline{\mathcal S}^E_{u,T_0}|_W\) along \(T|_WT_1|_{T^{-1}W}\) (restricted-equivalent marked ideals have the same cosupports along their common multiple test blow-ups,
Lemma~\ref{hir:marked-equivalence-basic}), where the values are compared as above and transported back along the isomorphism by property (6).
(4) As in Theorem~\ref{mc:contact-representative}: the marked ideals of Step 1 are restricted-equivalent for any divisor, so are their restricted coefficient ideals at every
stage (Lemmas~\ref{mc:coefficient-representatives-tower} and~\ref{hir:marked-equivalence-basic}), the strict transforms of \(N\cap V(u)\) and \(N\cap V(w)\) agree over \(O\)
(Lemma~\ref{mc:strict-along-local}), restricted equivalence passes to \(W\) (Lemma~\ref{hir:open-restriction-tower} for coherence), and Lemma~\ref{bcalc:inv-rebase} with property (7)
gives the domain, the canonical segments and the values.
\end{proof}
With the Step 1bb lower marked ideal formed with the divisor kept (its original members have empty traces on the mask), (1)--(4) and
Theorem~\ref{mc:tangent-change-opens}(2) are the Step 1b inputs of properties (4'), (5), (6) and (7) of Definition~\ref{hir:centre-function} for the value function of
Proposition~\ref{hir:maximal-order-step}.

\subsection{Step 1b across an empty prefix of the base}

\begin{theorem}[Step 1b re-based across an empty prefix]
\label{mc:contact-empty-prefix}
\uses{mc:contact-stage-kept,mc:empty-blowup-inverse,mc:coefficient-empty-step,mc:tangent-direction-canonical-along,mc:homogenized-along-etale,mc:homogenized-along-basic,mc:contact-representative,mc:coefficient-representatives-along,mc:coefficient-representatives-tower,mc:strict-along-local,hir:marked-equivalence-basic,hir:geo-transform-wellformed,hir:geo-transform-coherent,hir:geo-transform-basic,hir:open-restriction-tower,hir:canonical-segments,hir:centre-function,bcalc:inv-rebase}
\source{wlodarczyk05-hironaka}
\lean{Hironaka.ContactEmptyPrefix.restrictEquiv_trans,Hironaka.ContactEmptyPrefix.restrictEquiv_homMarkedAlong_geoTransform_empty,Hironaka.ContactEmptyPrefix.restrictEquiv_barMarked_geoTransform_empty,Hironaka.ContactEmptyPrefix.strictAlong_of_isEmptyTower,Hironaka.ContactEmptyPrefix.restrictEquiv_barMarked_geoTransformAlong_of_isEmptyTower,Hironaka.ContactEmptyPrefix.isGeoTestTower_barMarked_geoTransformAlong_iff,Hironaka.ContactEmptyPrefix.isTangentDirectionAlong_geoTransformAlong_of_isEmptyTower,Hironaka.ContactEmptyPrefix.strictAlong_inter_zero_of_isEmptyTower,Hironaka.ContactEmptyPrefix.restrictEquiv_contactMarkedKept_emptyPrefix,Hironaka.ContactEmptyPrefix.contactMarkedKept_emptyPrefix}
\leanok
Let \(\underline{\mathcal J}=(X,N,E,\mathfrak a,d)\) with \(\mathfrak a\) coherent, \(N\) a closed submanifold and \(d\ge1\), and let \(P\) be a tower over \(X\) all of whose centres are
empty, \(\underline{\mathcal J}_P\) the transform and \(p_P\) the projection of its last stage.
\begin{enumerate}
\item For one blowup \(\pi\) along \(\emptyset\): \(\mathcal H_{N'}(\underline{\mathcal J}')\simeq\mathcal H_N(\underline{\mathcal J})'\) and
\(\mathcal C(\mathcal H(\underline{\mathcal J}'))\simeq\mathcal C(\mathcal H(\underline{\mathcal J}))'\) (restricted equivalence, Lemma~\ref{hir:marked-equivalence-basic}).
\item \(\mathcal C(\mathcal H(\underline{\mathcal J}_P))\simeq\mathcal C(\mathcal H(\underline{\mathcal J}))_P\); a tower \(T_0'\) over the last stage of \(P\) is a multiple test blow-up of
\(\mathcal C(\mathcal H(\underline{\mathcal J}_P))\) iff \(PT_0'\) is one of \(\mathcal C(\mathcal H(\underline{\mathcal J}))\). The strict transform along \(P\) of a closed set is its preimage.
\item If \(u\) is a tangent direction of \((\mathfrak a|_N,d)\) along \(N\) on an open \(U\), then \(u\circ p_P\) is one of \((\mathfrak a_P|_{N_P},d)\) along \(N_P\) on \(p_P^{-1}U\),
with \(N_P\cap V(u\circ p_P)=p_P^{-1}(N\cap V(u))\), and the strict transform along \(P\) of \(N\cap V(u)\) agrees with \(p_P^{-1}(N\cap V(u))\) over \(U\).
\item For \(T_0'\) with \(PT_0'\) a multiple test blow-up of \(\mathcal C(\mathcal H(\underline{\mathcal J}))\) and an open \(W\) of the last stage over \(U\), the marked ideal of
Theorem~\ref{mc:contact-stage-kept} formed from \(\underline{\mathcal J}_P\) with \(u\circ p_P\) along \(T_0'\) on \(W\) is restricted-equivalent to the one formed from
\(\underline{\mathcal J}\) with \(u\) along \(PT_0'\) on \(W\). If the latter lies in the domain of a resolution invariant (Definition~\ref{hir:centre-function}), so does the former,
they have the same canonical segments, and along them the invariant takes the same values on the cosupport.
\end{enumerate}
\end{theorem}
\begin{proof}
\leanok
(1) A blowup along \(\emptyset\) is a biholomorphism, the strict transform of the closed \(N\) is \(\pi^{-1}N\), and along \(\emptyset\) the transform agrees with the pullback on
\(\pi^{-1}N\) up to restricted inclusions both ways (Lemma~\ref{mc:empty-blowup-inverse}). The homogenized ideal respects restricted inclusions
(Theorem~\ref{mc:contact-representative}(1)) and commutes with pullback by local biholomorphisms (Lemma~\ref{mc:homogenized-along-etale}); so \(\mathcal H_{N'}(\underline{\mathcal J}')\)
and \(\mathcal H_N(\underline{\mathcal J})'\) are both restricted-equivalent on \(N'\) to \(\pi^*\mathcal H_N(\underline{\mathcal J})\), with the same divisor and mark. For \(\mathcal C\)
compose with Lemma~\ref{mc:coefficient-representatives-along} (restricted equivalence passes to \(\mathcal C\)) and Lemma~\ref{mc:coefficient-empty-step}; the transforms are coherent with
closed submanifolds (Corollary~\ref{hir:geo-transform-coherent}, Lemmas~\ref{hir:geo-transform-wellformed} and~\ref{mc:homogenized-along-basic}).
(2) Induction along \(P\): restricted-equivalent marked ideals have equal transforms from the first blowup on (Lemma~\ref{hir:marked-equivalence-basic}), so (1) at the last step and the
inductive hypothesis compose. Equivalent marked ideals have the same multiple test blow-ups, and \(P\) is one of every marked ideal (empty centres,
Lemma~\ref{hir:canonical-segments}); concatenation (Lemma~\ref{hir:geo-transform-basic}). Strict transforms along \(\emptyset\) of closed sets are preimages
(Lemma~\ref{mc:empty-blowup-inverse}).
(3) Induction with Lemma~\ref{mc:tangent-direction-canonical-along} for each empty step. On \(U\), \(u\) is holomorphic, so \(N\cap V(u)\) agrees over \(U\) with the closed set
\(N\cap\overline{V(u)\cap U}\), whose strict transform is its preimage by (2); strict transforms of sets agreeing over an open set agree over its preimage
(Lemma~\ref{mc:strict-along-local}).
(4) By (2) and Lemma~\ref{mc:coefficient-representatives-tower}, the restricted coefficient ideals of \(\mathcal C(\mathcal H(\underline{\mathcal J}_P))\) and of
\(\mathcal C(\mathcal H(\underline{\mathcal J}))_P\) along \(T_0'\) to the strict transform of \(N_P\cap V(u\circ p_P)\) are restricted-equivalent; on \(W\) (Lemma~\ref{hir:open-restriction-tower}
for coherence) this strict transform agrees with the strict transform along \(PT_0'\) of \(N\cap V(u)\) by (3) and Lemma~\ref{mc:strict-along-local}. The values: Lemma~\ref{bcalc:inv-rebase}
and property (7).
\end{proof}
With the change of tangent direction of Theorem~\ref{mc:tangent-change-opens}(2) from \(u\circ p_P\) to the tangent direction chosen for \(\underline{\mathcal J}_P\), (4) is the Step 1bb input of
property (4') of Definition~\ref{hir:centre-function} for the value function of Proposition~\ref{hir:maximal-order-step} with the divisor kept.

\subsection{The lower value of Step 1bb under biholomorphisms, restrictions, representatives and empty prefixes}

\begin{theorem}[The lower value of Step 1bb: any tangent direction, biholomorphisms, open sets, representatives, empty prefixes]
\label{mc:contact-lower-value}
\uses{mc:tangent-change-opens,mc:contact-stage-kept,mc:contact-marked-canonical,hir:tower-isomorphism,mc:tower-isomorphism-basic,mc:tower-isomorphism-transform,mc:tower-isomorphism-append,hir:step-one-value,hir:step-one-value-masked,hir:step-one-value-masked-basic,hir:step-one-lower-domain,hir:step-one-compare,hir:biholomorphism-invariance,hir:open-restriction-tower,hir:open-restriction-marked,hir:canonical-segments,bcalc:restrict-nested,bcalc:tower-restrict,mc:restriction-operations-along,mc:contact-stage-tail,mc:contact-representative,hir:marked-equivalence-basic,hir:restricted-equivalence,mc:contact-empty-prefix,bcalc:restrict-empty-tower,hir:wellformed-tower,hir:tower-split,hir:geo-transform-basic,hir:tower-sequence,hir:centre-function}
\source{wlodarczyk05-hironaka}
\lean{Hironaka.ContactLowerValue.lowerValueL_zero_eq,Hironaka.ContactLowerValue.origMembers_isTowerIso,Hironaka.ContactLowerValue.lowerValueL_zero_isTowerIso,Hironaka.ContactLowerValue.origMembers_restrict,Hironaka.ContactLowerValue.lowerValueL_zero_restrict,Hironaka.ContactLowerValue.origMembers_restrictEquiv,Hironaka.ContactLowerValue.lowerValueL_zero_representative,Hironaka.ContactLowerValue.lowerValueL_zero_emptyPrefix,Hironaka.ContactLowerValue.centresAvoid_of_isEmptyTower,Hironaka.ContactLowerValue.lowerValueL_zero_emptyPrefix_rebase,Hironaka.ContactLowerValue.lowerValueL_zero_emptyPrefix_refl}
\leanok
Let \(\underline{\mathcal J}=(X,N,E,\mathfrak a,d)\) be well formed with \(d\ge1\), \(\overline{\underline{\mathcal J}}=\mathcal C_N(\mathcal H_N(\underline{\mathcal J}))\) with the divisor
\(E\), \(\operatorname{inv}\) a resolution invariant (Definition~\ref{hir:centre-function}), and \(\underline{\mathcal S}^E_{w,T_0}|_W\) the marked ideal of
Theorem~\ref{mc:tangent-change-opens}(2). In the masked value function of Step 1 (Definition~\ref{hir:step-one-value-masked}) the lower value at level \(0\) of a point
\(p\) along a split \(T_0T_1\) is the value of \(\operatorname{inv}\) for \(\underline{\mathcal S}^E_{w,T_0}|_W\) along \(T_1|_W\) at \(p\), where \(p_0\) is the image of \(p\) at
the last stage of \(T_0\), \(W\) the mask of \(p_0\) and \(w\) the tangent direction chosen at the base point of \(p\); \(p\) has \emph{lower data at level \(0\)} if
\(p_0\in W\), this marked ideal lies in the domain of \(\operatorname{inv}\), \(T_1|_W\) is a canonical segment of it and \(p\) lies in its cosupport along \(T_1|_W\).
\begin{enumerate}
\item (Any tangent direction.) Let \(T_0T_1\) be a multiple test blow-up of \(\overline{\underline{\mathcal J}}\) and \(p\in\operatorname{cosupp}\overline{\underline{\mathcal J}}_{T_0T_1}\)
with lower data at level \(0\). Let \(u\) be a tangent direction of \((\mathfrak a|_N,d)\) along \(N\) on an open set \(U\), and \(V\subset T_0^{-1}U\) an open set containing
\(p_0\) and meeting none of the first \(\#E\) members of the divisor of \(\overline{\underline{\mathcal J}}_{T_0}\), such that \(\underline{\mathcal S}^E_{u,T_0}|_V\) lies in the
domain of \(\operatorname{inv}\) and \(T_1|_V\) is a canonical segment of it. Then the lower value of \(p\) is the value of \(\underline{\mathcal S}^E_{u,T_0}|_V\) along
\(T_1|_V\) at \(p\).
\item (Biholomorphisms.) Let \(\Phi\colon X'\to X\) be a biholomorphism, \(T_0'\) a tower isomorphic to \(T_0\) over \(\Phi\) with last-stage biholomorphism \(\chi_0\)
(Definition~\ref{hir:tower-isomorphism}), and \(T_1'\) isomorphic to \(T_1\) over \(\chi_0\) with last-stage biholomorphism \(\chi\). If \(T_0\) is a multiple test blow-up of
\(\overline{\underline{\mathcal J}}\), the original members of \(\Phi^*\underline{\mathcal J}\) at the last stage of \(T_0'\) are the \(\chi_0\)-preimages of those of
\(\underline{\mathcal J}\) at the last stage of \(T_0\). If \(T_0T_1\) is a multiple test blow-up of \(\overline{\underline{\mathcal J}}\), \(\chi(p')\) lies in
\(\operatorname{cosupp}\overline{\underline{\mathcal J}}_{T_0T_1}\), \(p'\) has lower data at level \(0\) for \(\Phi^*\underline{\mathcal J}\) along \(T_0'T_1'\) and \(\chi(p')\)
has lower data at level \(0\) for \(\underline{\mathcal J}\) along \(T_0T_1\), then the lower values of \(p'\) and of \(\chi(p')\) agree.
\item (Open sets.) Let \(U\subset X\) be open, \(T_0'\) a tower over \(U\) isomorphic to \(T_0|_U\) over the identity of \(U\) with last-stage biholomorphism
\(\chi_0\), and \(T_1'\) isomorphic to \(T_1|_{T_0^{-1}U}\) over \(\chi_0\) with last-stage biholomorphism \(\chi\). If \(T_0\) is a multiple test blow-up of
\(\overline{\underline{\mathcal J}}\), the original members of \(\underline{\mathcal J}|_U\) at the last stage of \(T_0|_U\) are the traces on \(T_0^{-1}U\) of those of
\(\underline{\mathcal J}\) at the last stage of \(T_0\). If \(T_0T_1\) is a multiple test blow-up of \(\overline{\underline{\mathcal J}}\), \(\chi(p')\), a point of the last stage of
\(T_0T_1\) over \(U\), lies in \(\operatorname{cosupp}\overline{\underline{\mathcal J}}_{T_0T_1}\), \(p'\) has lower data at level \(0\) for \(\underline{\mathcal J}|_U\) along
\(T_0'T_1'\) and \(\chi(p')\) has lower data at level \(0\) for \(\underline{\mathcal J}\) along \(T_0T_1\), then the lower values of \(p'\) and of \(\chi(p')\) agree.
\item (Representatives.) Restricted-equivalent marked ideals \(\underline{\mathcal J}\simeq\underline{\mathcal J}'\) (Definition~\ref{hir:restricted-equivalence}) have the same original
members along every tower. If \(\underline{\mathcal J}'\) has a coherent ideal, \(T_0T_1\) is a multiple test blow-up of \(\overline{\underline{\mathcal J}}\) and
\(p\in\operatorname{cosupp}\overline{\underline{\mathcal J}}_{T_0T_1}\) has lower data at level \(0\) along \(T_0T_1\) both for \(\underline{\mathcal J}\) and for
\(\underline{\mathcal J}'\), then its two lower values agree.
\item (Empty prefixes.) Let \(P\) be a tower over \(X\) all of whose centres are empty and \(\underline{\mathcal J}_P\) the transform. If \(PH_0T_1\) is a multiple test
blow-up of \(\overline{\underline{\mathcal J}}\) and \(z\) a point of its cosupport with lower data at level \(0\) for \(\underline{\mathcal J}\) along the split \((PH_0)T_1\)
and for \(\underline{\mathcal J}_P\) along the split \(H_0T_1\), the two lower values of \(z\) agree. If \(P=P_1P_2\) and \(z\) has lower data at level \(0\) for
\(\underline{\mathcal J}\) along the split \(P_1(P_2T_1)\) and for \(\underline{\mathcal J}_P\) along the split of \(T_1\) with the empty tower as head, the two lower values
of \(z\) agree as well; in particular, for \(P_1\) the empty tower, lower data along the splits of \(PT_1\) and of \(T_1\) with the empty tower as
head give equal lower values. (Every tower with only empty centres has its centres missing every open set.)
\end{enumerate}
\end{theorem}
\begin{proof}
\leanok
(1) Let \(x\) be the base point of \(p\) and \(w\) the tangent direction chosen at \(x\). Since \(p_0\) lies in its mask, \(x\) lies in the open set of \(w\), so \(w\) is a
tangent direction there (Lemma~\ref{hir:step-one-compare}); \(x\) lies in \(\operatorname{cosupp}\underline{\mathcal J}\), because the image of a point of the cosupport along a
multiple test blow-up lies in the cosupport and \(\overline{\underline{\mathcal J}}\) has the cosupport of \(\underline{\mathcal J}\) (Lemma~\ref{hir:step-one-lower-domain}). The mask lies
over the open set of \(w\) and misses the original members, and \(V\) lies over \(U\) and misses them. Theorem~\ref{mc:tangent-change-opens}(2) for \(w\) and \(u\) at \(x\), with the
open sets \(W\) and \(V\), compares the two values at \(p\): \(p\) lies over \(W\cap V\), over the neighbourhood of \(x\) given there, and in the cosupport.
(2) By Theorem~\ref{mc:contact-stage-kept}(1), \(\mathcal C(\mathcal H(\Phi^*\underline{\mathcal J}))=\Phi^*\overline{\underline{\mathcal J}}\), and by
Lemma~\ref{mc:tower-isomorphism-transform} its transform along \(T_0'\) is the \(\chi_0\)-pullback of the transform of \(\overline{\underline{\mathcal J}}\) along \(T_0\), whose divisor
consists of the preimages of the members; this gives the original members. The concatenations \(T_0'T_1'\) and \(T_0T_1\) are isomorphic over \(\Phi\) with last-stage
biholomorphism \(\chi\) (Lemma~\ref{mc:tower-isomorphism-append}), so \(T_0'T_1'\) is a multiple test blow-up of \(\Phi^*\overline{\underline{\mathcal J}}\) and \(p'\) lies in its
cosupport (Lemma~\ref{mc:tower-isomorphism-transform}). The image of \(\chi(p')\) at the last stage of \(T_0\) is \(\chi_0(p_0')\), \(p_0'\) the image of \(p'\) at the last stage of
\(T_0'\), and its base point is \(\Phi(x')\), \(x'\) the base point of \(p'\) (Lemma~\ref{mc:tower-isomorphism-basic}). Let \(W\) be the mask of \(\chi_0(p_0')\), \(w\) the tangent
direction chosen at \(\Phi(x')\) and \(W'=\chi_0^{-1}(W)\). By Theorem~\ref{mc:contact-stage-kept}(1), \(\underline{\mathcal S}^E_{w\circ\Phi,T_0'}|_{W'}\), formed from
\(\Phi^*\underline{\mathcal J}\), lies in the domain, \(T_1'|_{W'}\) is a canonical segment of it, and its value at \(p'\) is the value of \(\underline{\mathcal S}^E_{w,T_0}|_W\) at
\(\chi(p')\), the lower value of \(\chi(p')\). The function \(w\circ\Phi\) is a tangent direction of \(\Phi^*\underline{\mathcal J}\) on the preimage of the open set of \(w\)
(Theorem~\ref{mc:contact-marked-canonical}), over which \(W'\) lies; \(W'\) contains \(p_0'\) and misses the original members of \(\Phi^*\underline{\mathcal J}\), the preimages of those
of \(\underline{\mathcal J}\), which \(W\) misses. Since \(\Phi^*\underline{\mathcal J}\) is well formed (Lemma~\ref{hir:biholomorphism-invariance}), (1) for \(\Phi^*\underline{\mathcal J}\),
\(w\circ\Phi\) and \(W'\) gives the lower value of \(p'\) as this value.
(3) By Theorem~\ref{mc:contact-stage-kept}(2), \(\mathcal C(\mathcal H(\underline{\mathcal J}|_U))=\overline{\underline{\mathcal J}}|_U\), and its transform along \(T_0|_U\) is the
restriction of the transform of \(\overline{\underline{\mathcal J}}\) along \(T_0\) to \(T_0^{-1}U\) (Lemma~\ref{hir:open-restriction-tower}); this gives the original members.
The tower \(T_0|_U\) followed by \(T_1|_{T_0^{-1}U}\) is a multiple test blow-up of \(\overline{\underline{\mathcal J}}|_U\) whose cosupport is the trace of
\(\operatorname{cosupp}\overline{\underline{\mathcal J}}_{T_0T_1}\) (Lemmas~\ref{hir:open-restriction-tower} and~\ref{hir:canonical-segments}); along the isomorphism over the
identity, under which \(\overline{\underline{\mathcal J}}|_U\) pulls back to itself (Lemma~\ref{mc:restriction-operations-along}), so is \(T_0'T_1'\), and \(p'\) lies in its
cosupport (Lemmas~\ref{mc:tower-isomorphism-append} and~\ref{mc:tower-isomorphism-transform}). The images of \(\chi(p')\) in the last stages of \(T_0|_U\) and \(T_0\) agree
(Lemma~\ref{bcalc:tower-restrict}). Let \(W\) be the mask of the image \(p_0\) of \(\chi(p')\) at the last stage of \(T_0\), \(w\) the tangent direction chosen at its base
point, and \(W_1=W\cap T_0^{-1}U\), which contains \(p_0\). By Lemma~\ref{bcalc:restrict-nested}(2), \(\underline{\mathcal S}^E_{w,T_0}|_{W_1}\) lies in the domain, \(T_1|_{W_1}\)
is a canonical segment of it, and at \(\chi(p')\), which lies in its cosupport (the trace of \(\operatorname{cosupp}\overline{\underline{\mathcal J}}_{T_0T_1}\),
Theorem~\ref{mc:contact-stage-tail}), its value is the lower value of \(\chi(p')\). Theorem~\ref{mc:contact-stage-kept}(2) passes to the marked ideal formed from
\(\underline{\mathcal J}|_U\) along \(T_0|_U\) with \(w|_U\) on \(W_1\), seen in \(T_0^{-1}U\), and Theorem~\ref{mc:contact-stage-kept}(1), along the isomorphisms over the identity
(with \(\mathrm{id}^*=\mathrm{id}\)), to the one formed along \(T_0'\) on \(W'=\chi_0^{-1}(W_1)\), with the same value at \(p'\). The function \(w|_U\) is a tangent direction of
\(\underline{\mathcal J}|_U\) on the trace of the open set of \(w\) (Theorem~\ref{mc:contact-marked-canonical}), over which \(W'\) lies; \(W'\) contains the image of \(p'\) and
misses the original members of \(\underline{\mathcal J}|_U\) at the last stage of \(T_0'\): by (2) for the identity of \(U\) these are the \(\chi_0\)-preimages of those at
the last stage of \(T_0|_U\), which are the traces of those of \(\underline{\mathcal J}\), missed by \(W\). Since \(\underline{\mathcal J}|_U\) is well formed
(Lemma~\ref{hir:open-restriction-marked}), (1) for \(\underline{\mathcal J}|_U\), \(w|_U\) and \(W'\) gives the lower value of \(p'\) as this value.
(4) The divisor of a transform depends only on the divisor and the tower (Lemma~\ref{hir:marked-equivalence-basic}), and \(\overline{\underline{\mathcal J}}\),
\(\overline{\underline{\mathcal J}'}\) carry the common divisor \(E\); this gives the original members. \(\underline{\mathcal J}'\) is well formed (same submanifold and divisor,
coherent ideal), and \(\overline{\underline{\mathcal J}}\simeq\overline{\underline{\mathcal J}'}\) (Theorem~\ref{mc:contact-representative}), so the two have the same multiple test blow-ups
and cosupports (Lemma~\ref{hir:marked-equivalence-basic}). Let \(x\) be the base point of \(p\), \(w\) the tangent direction of \(\underline{\mathcal J}\) chosen at \(x\) and \(W\)
the mask of the image \(p_0\) of \(p\) for \(\underline{\mathcal J}\). By Theorem~\ref{mc:contact-representative} there are an open \(O\ni x\) and a tangent direction \(v\) of
\((\mathfrak a'|_N,d)\) along \(N\) on \(O\) with \(v=w\) on \(N\cap O\). On \(W_1=W\cap T_0^{-1}O\), which contains \(p_0\), \(\underline{\mathcal S}^E_{w,T_0}|_{W_1}\) lies in the
domain with \(T_1|_{W_1}\) a canonical segment (Lemma~\ref{bcalc:restrict-nested}(2)); by Theorem~\ref{mc:contact-stage-kept}(4) so does the marked ideal formed from
\(\underline{\mathcal J}'\) with \(v\) on \(W_1\), with the same values on the cosupport, where \(p\) lies (Theorem~\ref{mc:contact-stage-tail}). \(W_1\) lies over the open sets of
\(w\) and of \(v\) and misses the common original members; (1) for \(\underline{\mathcal J}\) with \(w\) and for \(\underline{\mathcal J}'\) with \(v\), on \(W_1\), gives the two lower values
as these two equal values.
(5) \(\underline{\mathcal J}_P\) is well formed with mark \(d\) (Lemmas~\ref{hir:wellformed-tower} and~\ref{hir:geo-transform-basic}; \(P\) is a multiple test blow-up,
Lemma~\ref{hir:canonical-segments}), \(H_0T_1\) is a multiple test blow-up of \(\mathcal C(\mathcal H(\underline{\mathcal J}_P))\) and \(z\) lies in its cosupport
(Theorem~\ref{mc:contact-empty-prefix}(2), restricted-equivalent marked ideals having the same multiple test blow-ups and cosupports, Lemma~\ref{hir:marked-equivalence-basic}, and
concatenation is associative, Lemma~\ref{hir:tower-sequence}). Let \(W\) be the mask of the image \(z_0\) of \(z\) for \(\underline{\mathcal J}\) at the last stage of \(PH_0\),
\(w\) the tangent direction chosen at the base point, \(W'\) the mask of \(z_0\) for \(\underline{\mathcal J}_P\) at the last stage of \(H_0\), and \(W_1=W\cap W'\), which
contains \(z_0\). By Lemma~\ref{bcalc:restrict-nested}(2), \(\underline{\mathcal S}^E_{w,PH_0}|_{W_1}\) lies in the domain with \(T_1|_{W_1}\) a canonical segment; by
Theorem~\ref{mc:contact-empty-prefix}(4) so does the marked ideal formed from \(\underline{\mathcal J}_P\) along \(H_0\) with \(w\circ p_P\) on \(W_1\), with the same values on the
cosupport, where \(z\) lies (Theorem~\ref{mc:contact-stage-tail}). \(W_1\) lies over the open set of \(w\), and over the preimage of that open set, on which \(w\circ p_P\) is a
tangent direction of \(\underline{\mathcal J}_P\) (Theorem~\ref{mc:contact-empty-prefix}(3)); it misses the original members of \(\underline{\mathcal J}\) (inside \(W\)) and those of
\(\underline{\mathcal J}_P\) (inside \(W'\)). (1) on both sides gives the two lower values as these two equal values. In the second case, the value for
\(\underline{\mathcal J}\) along \(P_1(P_2T_1)\) is first re-based across \(P_2\), whose centres are empty (Lemma~\ref{hir:tower-split}) and so miss the mask
(Lemma~\ref{bcalc:restrict-empty-tower}): by Lemma~\ref{hir:step-one-compare}(3) it is the value of the marked ideal formed along \(P_1P_2\) with the tangent direction chosen at the
base point, on the preimage \(W_2\) of the mask, along \(T_1|_{W_2}\), which lies in the domain with a canonical restricted tail; then the argument above applies with \(W_2\) in
place of \(W\) and the empty tower in place of \(H_0\). The particular case is the second one with \(P_1\) the empty tower and \(P_2=P\), the empty tower followed
by \(P\) being \(P\) (Lemma~\ref{hir:tower-sequence}).
\end{proof}
Along an isomorphism of towers split at the start of the run of a point (Lemma~\ref{hir:tower-split-iso}), with lower data on both sides along canonical segments of the
value function (Theorem~\ref{hir:step-one-canonical-tail}), (2) is the part at level \(0\) of property (6) of Definition~\ref{hir:centre-function} for the masked value function
of Proposition~\ref{hir:maximal-order-step}; likewise (3), for the split of a restricted tower, which is isomorphic over the identity to the restriction of the split, is the part
at level \(0\) of property (5), (4) the part at level \(0\) of property (7), and (5), in its two cases (the run of the point within \(T_1\), or reaching into \(P\)), the part at
level \(0\) of property (4').

\begin{lemma}[Admissible centres of the marked ideal of Step 1b are admissible for \(\overline{\underline{\mathcal J}}\), locally]
\label{mc:contact-admissible}
\uses{mc:contact-stage-tail,mc:contact-stage-kept,mc:contact-stage-domain-masked,mc:homogenized-equivalent-along,mc:tangent-change-opens,hir:open-restriction-tower,hir:open-restriction-marked,hir:marked-equivalence-basic,hir:geo-transform-basic,hir:transform,hir:step-one-value,hir:step-one-value-masked,hir:step-one-compare}
\source{wlodarczyk05-hironaka}
\lean{Hironaka.ContactLowerValue.isAdmissibleCentre_of_contactMarkedKept,Hironaka.ContactLowerValue.isAdmissibleCentre_of_lowerMarkedM_zero}
\leanok
Let \(\underline{\mathcal J}=(X,N,E,\mathfrak a,d)\) be well formed with \(d\ge1\), \(u\) a tangent direction of \((\mathfrak a|_N,d)\) along \(N\) on an open set \(U\),
\(T_0T_1\) a multiple test blow-up of \(\overline{\underline{\mathcal J}}=\mathcal C_N(\mathcal H_N(\underline{\mathcal J}))\), \(W\subset T_0^{-1}U\) an open set of the last stage of
\(T_0\) and \(P=T_1^{-1}W\). The transform of \(\underline{\mathcal S}^E_{u,T_0}|_W\) (Theorem~\ref{mc:tangent-change-opens}(2)) along \(T_1|_W\) and the restriction of
\(\overline{\underline{\mathcal J}}_{T_0T_1}\) to \(P\) have the same cosupport and the same divisor. Hence, for every open set \(W''\subset P\), an admissible centre
(Definition~\ref{hir:transform}) of the restriction of the first to \(W''\) is an admissible centre of the restriction of the second to \(W''\). In particular this holds for the
masked lower marked ideal of level \(0\) at a point \(y_0\) of the last stage of \(T_0\) lying in its mask (Definition~\ref{hir:step-one-value-masked}), on open sets of
\(T_1^{-1}\) of the mask.
\end{lemma}
\begin{proof}
\leanok
The submanifold of \(\underline{\mathcal S}^E_{u,T_0}|_W\) is a closed submanifold of \(W\) (Theorem~\ref{mc:contact-stage-domain-masked}; \(T_0\) is a multiple test blow-up of
\(\underline{\mathcal J}\), Theorem~\ref{mc:homogenized-equivalent-along}) and its ideal is coherent (Theorem~\ref{mc:contact-stage-kept}), so its transform along the multiple test
blow-up \(T_1|_W\) (Theorem~\ref{mc:contact-stage-tail}) has a coherent ideal (Lemma~\ref{hir:open-restriction-tower}). By Theorem~\ref{mc:contact-stage-tail} the cosupport of that
transform is the trace on \(P\) of \(\operatorname{cosupp}\overline{\underline{\mathcal J}}_{T_0T_1}\), which is the cosupport of the restriction of
\(\overline{\underline{\mathcal J}}_{T_0T_1}\) to \(P\) (Lemma~\ref{hir:open-restriction-marked}). The divisor of a transform depends only on the divisor and the tower
(Lemma~\ref{hir:marked-equivalence-basic}); \(\underline{\mathcal S}^E_{u,T_0}|_W\) has the divisor of \(\overline{\underline{\mathcal J}}_{T_0}|_W\), restriction to \(W\)
commutes with transforms (Lemma~\ref{hir:open-restriction-tower}) and transforms compose along concatenations (Lemma~\ref{hir:geo-transform-basic}); so both divisors are the traces
on \(P\) of the divisor of \(\overline{\underline{\mathcal J}}_{T_0T_1}\). Restricting further to \(W''\) keeps both equalities (Lemma~\ref{hir:open-restriction-marked}), and an
admissible centre is a closed submanifold contained in the cosupport having normal crossings with the divisor (Definition~\ref{hir:transform}). For the masked lower marked ideal
of level \(0\) at \(y_0\): the base point of \(y_0\) lies in the open set of the chosen tangent direction, which is then a tangent direction there
(Lemma~\ref{hir:step-one-compare}), and the mask lies over that open set (Definition~\ref{hir:step-one-value}).
\end{proof}
On an open set of \(P\) on which a resolution invariant of the lower marked ideal is bounded by \(\lambda\), property (1) of Definition~\ref{hir:centre-function} makes the points of
value \(\lambda\) an admissible centre of its restriction; by the lemma it is one of the restriction of \(\overline{\underline{\mathcal J}}_{T_0T_1}\). This is the Step 1bb input of
property (1) for the masked value function of Proposition~\ref{hir:maximal-order-step} at the values \((0,\lambda)\).

\subsection{Relative principalization from given resolutions}

In this subsection \(Y\) is a complex manifold of dimension \(n\), \(W\subset Y\) a closed connected submanifold of dimension \(k\) and \(\mathfrak a\subset\mathcal O_Y\) a
coherent ideal with \(\mathfrak a\not\subset\mathcal I_W\). These are the data of Lemma~\ref{hir:factorization-from-resolution}, whose proof applies the resolution of
marked ideals twice, to \((Y,W,\emptyset,\mathfrak a,1)\) and to \((Y_t,Y_t,\emptyset,\mathfrak a_t,1)\).

\begin{lemma}[The domain condition of relative principalization]
\label{mc:factorization-domain}
\uses{ana:ideal-sheaves,bcalc:contact-restriction-nonzero,bcalc:cosupport-analytic,ana:regular-component,hir:hadamard-bridge,hir:centre-function,hir:marked-ideal,hir:snc-basic,ana:manifold-irreducible}
\source{transcendental-nakamaye-copy4-20260911:lem:relative-principalization,wlodarczyk05-hironaka}
\lean{Hironaka.FactorizationDomain.component_eq_connectedComponentIn,Hironaka.FactorizationDomain.isNonzeroAlongAt_of_not_le,Hironaka.FactorizationDomain.isClosedSubmanifold_of_isClosedConnectedSubmanifold,Hironaka.FactorizationDomain.isInvDomain_of_not_le,Hironaka.FactorizationDomain.isClosedConnectedSubmanifold_univ,Hironaka.FactorizationDomain.isInvDomain_univ_of_not_le}
\leanok
Let \(N\) be a closed connected submanifold of dimension \(k\) of a complex manifold \(M\) and \(\mathfrak a\subset\mathcal O_M\) a coherent ideal with
\(\mathfrak a\not\subset\mathcal I_N\). Then at every point \(y\in N\) some germ of \(\mathfrak a_y\) does not vanish identically on \(N\) near \(y\). Hence the marked ideal
\((M,N,\emptyset,\mathfrak a,1)\) lies in the domain of the invariants in dimension \(\le j\) for every \(j\ge k\) (Definition~\ref{hir:centre-function}). In particular, if \(M\)
is connected and \(\mathfrak a\) is not the zero ideal, \((M,M,\emptyset,\mathfrak a,1)\) lies in the domain in dimension \(\le j\) for every \(j\ge n\).
\end{lemma}
\begin{proof}
\leanok
Every point of \(N\) is a regular point of \(N\) of dimension \(k\), so the component of the regular locus of \(N\) through a point is the connected component of \(N\)
through it, which is \(N\). Suppose that at \(y_0\in N\) every germ of \(\mathfrak a_{y_0}\) vanishes identically on \(N\) near \(y_0\). As \(\mathfrak a\) is coherent, the
same holds at every point of an open neighbourhood \(O\) of \(y_0\) (Lemma~\ref{bcalc:contact-restriction-nonzero}); in particular every germ of \(\mathfrak a\) at a point of
\(O\cap N\) vanishes at that point, so \(O\cap N\) lies in the zero locus \(Z(\mathfrak a)\), which is analytic (Lemma~\ref{bcalc:cosupport-analytic}). By the identity
theorem along a component of a regular locus (an analytic set containing a nonempty relatively open part of the component contains the component; proof of
Proposition~\ref{ana:regular-component}), \(N\subset Z(\mathfrak a)\). Now let \(g\in\mathfrak a_x\) at any \(x\in M\). It is holomorphic near \(x\), and by coherence its germ lies in
\(\mathfrak a_z\) for every \(z\) near \(x\) (Lemma~\ref{hir:hadamard-bridge}); for such \(z\in N\) this gives \(g(z)=0\). So \(g\in\mathcal I_{N,x}\), that is
\(\mathfrak a\subset\mathcal I_N\), a contradiction. For the domain: \(N\) is a closed submanifold of dimension \(k\le j\), the empty divisor has simple normal crossings and
is transverse to \(N\) (Lemma~\ref{hir:snc-basic}), \(\mathfrak a\) is coherent and \(d=1\) (Definition~\ref{hir:marked-ideal}). A connected \(M\) is a closed connected
submanifold of itself of dimension \(n\) (Lemma~\ref{ana:manifold-irreducible}), and \(\mathfrak a\not\subset\mathcal I_M\) says that \(\mathfrak a\) is not the zero ideal.
\end{proof}

\begin{lemma}[The transform along \(N\) and the ambient controlled transform]
\label{mc:factorization-ambient}
\uses{hir:transform,hir:geo-transform,hir:test-tower,hir:geo-transform-ambient,hir:restricted-equivalence,hir:restricted-equivalence-basic,mc:transform-along-calculus,hir:bridge-consumers,hir:cosupport-basic,hir:open-restriction-tower,hir:marked-equivalence-basic}
\source{transcendental-nakamaye-copy4-20260911:lem:relative-principalization,bierstone-milman08-functoriality}
\lean{Hironaka.FactorizationFromResolution.ambient,Hironaka.FactorizationFromResolution.ambient_N,Hironaka.FactorizationFromResolution.ambient_E,Hironaka.FactorizationFromResolution.ambient_ideal,Hironaka.FactorizationFromResolution.ambient_d,Hironaka.FactorizationFromResolution.cosupp_eq_inter_zeroLocus,Hironaka.FactorizationFromResolution.isGeoTestTower_and_restrictEquiv,Hironaka.FactorizationFromResolution.cosupp_geoTransformAlong_eq}
\leanok
Let \(\underline{\mathcal I}=(X,N,E,\mathfrak a,1)\) be a marked ideal with mark one, \(N\) a closed submanifold and \(\mathfrak a\) coherent, and
\(\underline{\mathcal I}'=(X,X,E,\mathfrak a,1)\). Let \(T\) be a multiple test blow-up of \(\underline{\mathcal I}\). Then \(T\) is a multiple test blow-up of
\(\underline{\mathcal I}'\), whose transform \(\underline{\mathcal I}'_T=(X_T,X_T,E_T,\mathfrak a'_T,1)\) carries the ambient controlled transform \(\mathfrak a'_T\), and the ideal of
\(\underline{\mathcal I}_T\) has the same restriction to \(N_T\) as \(\mathfrak a'_T\). In particular \(\operatorname{cosupp}\underline{\mathcal I}_T=N_T\cap Z(\mathfrak a'_T)\).
\end{lemma}
\begin{proof}
\leanok
With mark one and a coherent ideal the cosupport of a marked ideal is \(N\cap Z(\mathfrak a)\) (Lemma~\ref{hir:cosupport-basic}). Induction on \(T\); for the empty tower
there is nothing to prove. Let \(T\) be followed by the blowup \(\pi\) along an admissible centre \(C\) of \(\underline{\mathcal I}_T\), and let \(\underline{\mathcal J}\) be
\(\underline{\mathcal I}_T\) with its ideal replaced by \(\mathfrak a'_T\); by induction \(\underline{\mathcal J}\) is restricted-equivalent to \(\underline{\mathcal I}_T\), so it has
the same admissible centres and the same transform by \(\pi\) (Lemma~\ref{hir:restricted-equivalence-basic}); \(N_T\) is a closed submanifold and both ideals are coherent
(Lemma~\ref{hir:open-restriction-tower}, Lemma~\ref{hir:geo-transform-ambient}). As \(C\subset\operatorname{cosupp}\underline{\mathcal J}=N_T\cap Z(\mathfrak a'_T)\), \(C\) lies in
\(Z(\mathfrak a'_T)=\operatorname{cosupp}\underline{\mathcal I}'_T\); the divisors of \(\underline{\mathcal I}'_T\) and \(\underline{\mathcal I}_T\) agree (they depend only on \(E\)
and the tower, Lemma~\ref{hir:marked-equivalence-basic}), so \(C\) is admissible for \(\underline{\mathcal I}'_T\), and \(T\) followed by \(\pi\) is a multiple test blow-up of
\(\underline{\mathcal I}'\). Since \(C\) lies in the cosupport of \((X_T,X_T,E_T,\mathfrak a'_T,1)\), \(\pi^*\mathfrak a'_T=\mathcal I_D\cdot\sigma^c(\mathfrak a'_T,1)\) with \(D\) the
exceptional divisor (Lemma~\ref{hir:bridge-consumers}), and \(\sigma^c(\mathfrak a'_T,1)\) is the ideal of the next ambient transform (Lemma~\ref{hir:geo-transform-ambient}).
By the characterisation of the transform along \(N\) by its defining equation (Lemma~\ref{mc:transform-along-calculus}), the transform of \(\underline{\mathcal J}\) is
restricted-equivalent to the same marked ideal with its ideal replaced by \(\sigma^c(\mathfrak a'_T,1)\); the transform of \(\underline{\mathcal J}\) is that of
\(\underline{\mathcal I}_T\).
\end{proof}
This is the third identity \(\mathfrak a_i|_{W_i}=\mathfrak b_i\) of the proof of Lemma~\ref{hir:factorization-from-resolution} (Copy-4, \texttt{eq:rp-induction}), with
\(\mathfrak a_i=\mathfrak a'_T\) and \(\mathfrak b_i\) the ideal of the transform along \(W\).

\begin{lemma}[One blowup along a thin centre inside a connected submanifold]
\label{mc:factorization-thin-centre}
\uses{bcalc:analytic-fill,bcalc:exceptional-divisor,bcalc:isomorphism-off-centre,ana:component-complement-connected,bcalc:strict-geo,bcalc:strict-ideal,ana:regular-transport,ana:regular-component,ana:analytic-preimage,mc:factorization-domain}
\source{transcendental-nakamaye-copy4-20260911:lem:relative-principalization}
\lean{Hironaka.FactorizationFromResolution.lt_of_regularAtDim_of_frequently,Hironaka.FactorizationFromResolution.mem_exceptionalLocus_of_regularAtDim,Hironaka.FactorizationFromResolution.isPreconnected_preimage_of_disjoint,Hironaka.FactorizationFromResolution.mem_geoStrictTransform_iff_of_notMem,Hironaka.FactorizationFromResolution.geoStrictTransform_data,Hironaka.FactorizationFromResolution.pullback_vanishingIdeal_eq_of_disjoint}
\leanok
Let \(\pi\colon X'\to X\) be the blowup along a closed submanifold \(C\) of a complex manifold \(X\) of dimension \(n\).
\begin{enumerate}
\item If \(C\subset N\) are regular at \(c\) of dimensions \(j\) and \(k\), and points of \(N\setminus C\) accumulate at \(c\), then \(j<k\). If \(C\) is regular of dimension
\(j\le n-2\) at a point \(\pi(z)\in C\), then \(z\) lies in the exceptional locus of \(\pi\).
\item If \(A\subset X\setminus C\) is connected, so is \(\pi^{-1}(A)\); if \(A\) is closed and disjoint from \(C\), then \(\pi^*\mathcal I_A=\mathcal I_{\pi^{-1}(A)}\).
\item Let \(N\) be a closed connected submanifold of dimension \(k\) with \(C\subset N\), points of \(N\setminus C\) accumulating at every point of \(C\). Then
\(N'=\overline{\pi^{-1}(N\setminus C)}\) is the strict transform \(\overline{\pi^{-1}(N)\setminus\operatorname{Exc}\pi}\), and it is a closed connected submanifold of dimension
\(k\); off \(\pi^{-1}(C)\) it is \(\pi^{-1}(N)\).
\end{enumerate}
\end{lemma}
\begin{proof}
\leanok
(1) If \(j\ge k\), then \(C\), an analytic set (Lemma~\ref{ana:analytic-preimage}) contained in \(N\) and of local dimension \(\ge k\) at \(c\), contains \(N\) near \(c\)
(Lemma~\ref{bcalc:analytic-fill}), against the accumulation. In flattening coordinates in which \(C\) is a coordinate subspace of codimension at least two, every point
over \(C\) is exceptional (the monomial-chart computation behind Lemma~\ref{bcalc:exceptional-divisor}).
(2) \(\pi\) has a holomorphic section \(\sigma\) over \(X\setminus C\) with \(\sigma(\pi(z))=z\) (Lemma~\ref{bcalc:isomorphism-off-centre}); so \(\pi^{-1}(A)=\sigma(A)\) is the continuous image
of a connected set. At a point \(z\) with \(\pi(z)\notin C\) the identity of ideals is Lemma~\ref{bcalc:strict-ideal}; at a point over \(C\), hence off the closed set \(A\), both
ideals contain every holomorphic germ (the constant \(1\) lies in \(\mathcal I_A\) at \(\pi(z)\), and \(\pi^{-1}(A)\) misses a neighbourhood of \(z\)).
(3) The component of the regular locus of \(N\) through any point is \(N\), and it is not contained in \(C\) (\(N\) is nonempty and the accumulation fails at points of
\(C\)); so \(N\setminus C\) is connected and dense in \(N\) (Lemma~\ref{ana:component-complement-connected}). Density gives the comparison of strict transforms
(Lemma~\ref{bcalc:strict-geo}), and \(N'\) is a closed submanifold (Lemma~\ref{bcalc:strict-ideal}). By (2) \(\pi^{-1}(N\setminus C)\) is connected, so its closure \(N'\) is
connected, and nonempty. At a point \(y_0=\sigma(x)\), \(x\in N\setminus C\), \(\pi\) has bijective differential, and near \(y_0\) the set \(N'\) is \(\pi^{-1}(N)\); so \(N'\) is
regular of dimension \(k\) at \(y_0\) (Lemma~\ref{ana:regular-transport}). All points of \(N'\) are regular, and \(N'\) is connected, so the component of its regular locus
through \(y_0\) is \(N'\), along which the dimension is constant (Proposition~\ref{ana:regular-component}(1)).
\end{proof}

\begin{theorem}[Along a multiple test blow-up of \((X,W,\emptyset,\mathfrak a,1)\)]
\label{mc:factorization-tower}
\uses{mc:factorization-thin-centre,mc:factorization-ambient,bcalc:inv-domain-stable,bcalc:tower-exceptional,bcalc:exceptional-divisor,bcalc:strict-ideal,ana:strict-transform-composition,hir:tower-sequence,hir:cosupport-basic,hir:centre-function,ana:blowups}
\source{transcendental-nakamaye-copy4-20260911:lem:relative-principalization}
\lean{Hironaka.FactorizationFromResolution.support_add_subset,Hironaka.FactorizationFromResolution.frequently_mem_N_notMem_of_isInvDomain,Hironaka.FactorizationFromResolution.tower_data}
\leanok
Let \(\underline{\mathcal I}\) be a marked ideal of the domain in dimension \(\le k\) (Definition~\ref{hir:centre-function}) with mark one whose submanifold \(N\) is closed,
connected and of dimension \(k<n\), and \(T\) a multiple test blow-up of \(\underline{\mathcal I}\) with exceptional divisor \(E_T\) (Lemma~\ref{bcalc:tower-exceptional}). Then
\(N_T\) is a closed connected submanifold of dimension \(k\) and the strict transform \(\overline{p_T^{-1}(N)\setminus\operatorname{Exc}p_T}\) of \(N\); the composite \(p_T\) is
a sequence of blowups adapted to \(N\) (Definition~\ref{ana:blowups}: every centre is a proper analytic subset of the current strict transform); \(p_T^*\mathcal I_N=\mathcal
I_{E_T}\cdot\mathcal I_{N_T}\); and the support of \(E_T\) lies in the exceptional locus of \(p_T\).
\end{theorem}
\begin{proof}
\leanok
Induction on \(T\); for the empty tower \(E_T=0\). Let \(T\) be followed by the blowup \(\pi\) along an admissible centre \(C\) of \(\underline{\mathcal I}_T\). The transform
\(\underline{\mathcal I}_T\) lies in the domain (Proposition~\ref{bcalc:inv-domain-stable}), so at every \(c\in C\subset\operatorname{cosupp}\underline{\mathcal I}_T\subset N_T\) some
germ \(g\) of its ideal does not vanish identically on \(N_T\) near \(c\); by coherence \(g\) lies in the ideal at nearby points, and those points of \(N_T\) where \(g\neq0\)
are off the cosupport \(N_T\cap Z\) (Lemma~\ref{hir:cosupport-basic}), hence off \(C\). So points of \(N_T\setminus C\) accumulate at every point of \(C\), and
Lemma~\ref{mc:factorization-thin-centre}(3) applies: \(N_{T\pi}=\overline{\pi^{-1}(N_T\setminus C)}\) is a closed connected submanifold of dimension \(k\) and the strict transform of
\(N_T\), hence of \(N\) (strict transforms compose, Lemma~\ref{ana:strict-transform-composition}). \(C\) is analytic, contained in \(N_T\) and different from it, so the
sequence stays adapted. \(\pi^*\mathcal I_{N_T}=\mathcal I_D\cdot\mathcal I_{N_{T\pi}}\) (Lemma~\ref{bcalc:strict-ideal}), and pulling back the identity for \(T\) gives the identity
for \(T\pi\) with \(E_{T\pi}=\pi^*E_T+D\). By Lemma~\ref{mc:factorization-thin-centre}(1) \(C\) has dimension \(<k\le n-1\) at each of its points, so \(\pi^{-1}(C)=\operatorname{supp}D\)
lies in \(\operatorname{Exc}\pi\) (Lemma~\ref{bcalc:exceptional-divisor}); and \(\operatorname{supp}\pi^*E_T=\pi^{-1}(\operatorname{supp}E_T)\subset\pi^{-1}(\operatorname{Exc}p_T)\). The
support of a sum lies in the union of the supports, and \(\operatorname{Exc}(p_T\circ\pi)=\operatorname{Exc}\pi\cup\pi^{-1}(\operatorname{Exc}p_T)\).
\end{proof}
These are the first identity of \texttt{eq:rp-induction} and the adaptedness of the first stage in the proof of Lemma~\ref{hir:factorization-from-resolution}, for
\(\underline{\mathcal I}=(Y,W,\emptyset,\mathfrak a,1)\), which lies in the domain by Lemma~\ref{mc:factorization-domain}.

\begin{lemma}[The second stage of relative principalization]
\label{mc:factorization-second}
\uses{mc:factorization-thin-centre,mc:factorization-tower,hir:geo-transform-ambient,hir:cosupport-basic,hir:test-tower,bcalc:exceptional-divisor,bcalc:isomorphism-off-centre,ana:regular-transport,ana:strict-transform-composition,bcalc:tower-exceptional}
\source{transcendental-nakamaye-copy4-20260911:lem:relative-principalization}
\lean{Hironaka.FactorizationFromResolution.disjoint_preimage_of_isAdmissibleCentre,Hironaka.FactorizationFromResolution.second_data}
\leanok
Let \(\underline{\mathcal I}=(X,X,E,\mathfrak b,1)\) with \(\mathfrak b\) coherent, and \(A\subset X\) a closed connected submanifold of dimension \(k\) disjoint from the zero
locus \(Z(\mathfrak b)\). Let \(S\) be a multiple test blow-up of \(\underline{\mathcal I}\). Then every admissible centre of \(\underline{\mathcal I}_S\) is disjoint from
\(p_S^{-1}(A)\); \(p_S^{-1}(A)\) is a closed connected submanifold of dimension \(k\) and the strict transform of \(A\); the exceptional divisor \(E_S\) is disjoint from it; and
\(p_S^*\mathcal I_A=\mathcal I_{p_S^{-1}(A)}\).
\end{lemma}
\begin{proof}
\leanok
With \(N=X\) and mark one, \(p_S^*\mathfrak b=\mathcal I_{E_S}\cdot\mathfrak b_S\) with \(\mathfrak b_S\) the ideal of \(\underline{\mathcal I}_S\), coherent (Lemma~\ref{hir:geo-transform-ambient}); so
\(Z(\mathfrak b_S)\subset Z(p_S^*\mathfrak b)=p_S^{-1}(Z(\mathfrak b))\), which misses \(p_S^{-1}(A)\). An admissible centre lies in the cosupport \(Z(\mathfrak b_S)\)
(Lemma~\ref{hir:cosupport-basic}). The rest by induction on \(S\); for the empty tower there is nothing to prove. Let \(S\) be followed by the blowup \(\pi\) along an
admissible centre \(C\), so \(C\cap p_S^{-1}(A)=\emptyset\) and \(\pi\) has bijective differential over \(p_S^{-1}(A)\) (Lemma~\ref{bcalc:isomorphism-off-centre}). By
Lemma~\ref{mc:factorization-thin-centre}(2), \(\pi^{-1}(p_S^{-1}(A))\) is connected, nonempty (the section of \(\pi\) over \(X_S\setminus C\)) and closed, and it is regular of
dimension \(k\) at each of its points (Lemma~\ref{ana:regular-transport}); \(\pi^*\mathcal I_{p_S^{-1}(A)}=\mathcal I_{\pi^{-1}p_S^{-1}(A)}\). No point of it is exceptional for \(\pi\), so
its strict transform under \(\pi\) is itself, and strict transforms compose (Lemma~\ref{ana:strict-transform-composition}). The support of \(E_{S\pi}=\pi^*E_S+D\) lies in
\(\pi^{-1}(\operatorname{supp}E_S)\cup\pi^{-1}(C)\) (Lemma~\ref{bcalc:exceptional-divisor}), which misses \(\pi^{-1}p_S^{-1}(A)\).
\end{proof}
This is the second application in the proof of Lemma~\ref{hir:factorization-from-resolution}, with \(A=W_t\) and \(\mathfrak b=\mathfrak a_t\): all its centres lie in the cosupport
of \(\mathfrak a_t\), hence are disjoint from the strict transforms of \(W_t\).

\begin{theorem}[Relative principalization from two given resolutions]
\label{mc:factorization-assembly}
\uses{mc:factorization-domain,mc:factorization-ambient,mc:factorization-tower,mc:factorization-second,hir:geo-transform-ambient,hir:test-tower,hir:resolution-driver,hir:centre-function,mc:subvariety-thin,ana:strict-transform-composition,bcalc:exceptional-divisor,bcalc:tower-exceptional,ana:blowups,ana:component-complement-connected,bcalc:isomorphism-off-centre}
\source{transcendental-nakamaye-copy4-20260911:lem:relative-principalization}
\lean{Hironaka.FactorizationFromResolution.pullDivisor,Hironaka.FactorizationFromResolution.pullDivisor_ideal,Hironaka.FactorizationFromResolution.support_pullDivisor,Hironaka.FactorizationFromResolution.idealPow_one_of_isCoherent,Hironaka.FactorizationFromResolution.isAdaptedBlowupSequence_comp,Hironaka.FactorizationFromResolution.lt_of_ne_univ,Hironaka.FactorizationFromResolution.factorizationData_of_isGeoResolution,Hironaka.FactorizationFromResolution.factorizationData_of_exists_resolution}
\leanok
Let \(Y\) be a connected complex manifold of dimension \(n\), \(W\subsetneq Y\) a closed connected submanifold of dimension \(k\) and \(\mathfrak a\) a coherent ideal with
\(\mathfrak a\not\subset\mathcal I_W\). Let \(T_1\) be a multiple test blow-up of \((Y,W,\emptyset,\mathfrak a,1)\) with empty final cosupport, \(\mathfrak a_t\) the ambient controlled
transform of \(\mathfrak a\) along \(T_1\) (the ideal of the transform of \((Y,Y,\emptyset,\mathfrak a,1)\) along \(T_1\), Lemma~\ref{mc:factorization-ambient}), and \(T_2\) a multiple
test blow-up of \((Y_t,Y_t,\emptyset,\mathfrak a_t,1)\) with empty final cosupport. Then \(f=p_{T_1}\circ p_{T_2}\colon Y'\to Y\), \(H=p_{T_2}^*E_{T_1}\) and \(H'=E_{T_2}\) are
factorization data in the sense of Lemma~\ref{ana:relative-principalization-from-core}: \(f\) is a sequence of blowups adapted to \(W\), the strict transform \(W'\) of \(W\) is connected
and regular of dimension \(k\) at each of its points, \(\operatorname{supp}H\subset\operatorname{Exc}f\), \(H'\cap W'=\emptyset\), \(f^*\mathcal I_W=\mathcal I_H\cdot\mathcal I_{W'}\) and
\(f^*\mathfrak a=\mathcal I_H\cdot\mathcal I_{H'}\). Consequently, if \(Y\) is compact and every marked ideal of the domain in dimension \(\le j\) (Definition~\ref{hir:centre-function})
on a compact complex manifold of dimension \(n\le j\) has a multiple test blow-up with empty final cosupport, such factorization data exist.
\end{theorem}
\begin{proof}
\leanok
\(W\ne Y\) and \(Y\) connected give \(k<n\): if \(k=n\), \(W\) would be open (Lemma~\ref{mc:subvariety-thin}) and closed. \((Y,W,\emptyset,\mathfrak a,1)\) lies in the domain in dimension
\(\le k\) (Lemma~\ref{mc:factorization-domain}), so Theorem~\ref{mc:factorization-tower} applies to \(T_1\): \(N_1=W_t\) is a closed connected submanifold of dimension \(k\) and the
strict transform of \(W\), \(p_{T_1}\) is adapted to \(W\), \(p_{T_1}^*\mathcal I_W=\mathcal I_{E_{T_1}}\cdot\mathcal I_{W_t}\) and \(\operatorname{supp}E_{T_1}\subset\operatorname{Exc}p_{T_1}\).
By Lemma~\ref{mc:factorization-ambient}, \(\emptyset=\operatorname{cosupp}=W_t\cap Z(\mathfrak a_t)\), and \(p_{T_1}^*\mathfrak a=\mathcal I_{E_{T_1}}\cdot\mathfrak a_t\) with \(\mathfrak a_t\)
coherent (Lemma~\ref{hir:geo-transform-ambient}). Lemma~\ref{mc:factorization-second} with \(A=W_t\), \(\mathfrak b=\mathfrak a_t\) applies to \(T_2\): \(W'=p_{T_2}^{-1}(W_t)\) (strict transforms
compose, Lemma~\ref{ana:strict-transform-composition}) is a closed connected submanifold of dimension \(k\), every centre of \(T_2\) misses the current strict transform, so \(f\) is
adapted to \(W\), \(E_{T_2}\) misses \(W'\), and \(p_{T_2}^*\mathcal I_{W_t}=\mathcal I_{W'}\); and as \(T_2\) resolves a marked ideal with \(N=Y_t\) and mark one,
\(p_{T_2}^*\mathfrak a_t=\mathcal I_{E_{T_2}}\) (Lemma~\ref{hir:geo-transform-ambient}). The total transform of a divisor along a tower is a divisor with ideal \(p_{T_2}^*\mathcal I_{E_{T_1}}\)
and support \(p_{T_2}^{-1}(\operatorname{supp}E_{T_1})\subset p_{T_2}^{-1}(\operatorname{Exc}p_{T_1})\subset\operatorname{Exc}f\) (Lemma~\ref{bcalc:exceptional-divisor}). Pulling back:
\(f^*\mathcal I_W=p_{T_2}^*\mathcal I_{E_{T_1}}\cdot p_{T_2}^*\mathcal I_{W_t}=\mathcal I_H\cdot\mathcal I_{W'}\) and \(f^*\mathfrak a=p_{T_2}^*\mathcal I_{E_{T_1}}\cdot p_{T_2}^*\mathfrak a_t=\mathcal
I_H\cdot\mathcal I_{H'}\). For the last claim: \((Y,W,\emptyset,\mathfrak a,1)\) lies in the domain in dimension \(\le j\) (\(k<n\le j\)), giving \(T_1\); \(Y_t\) is compact
(Theorem~\ref{hir:resolution-driver}) and connected (at each step the complement of the centre is connected, Lemma~\ref{ana:component-complement-connected}, its preimage is the image of a section, Lemma~\ref{bcalc:isomorphism-off-centre}, and it is dense, the exceptional divisor having empty interior, Lemma~\ref{bcalc:exceptional-divisor}); \(\mathfrak a_t\)
is not the zero ideal, since it does not vanish at the points of \(W_t\neq\emptyset\); so \((Y_t,Y_t,\emptyset,\mathfrak a_t,1)\) lies in the domain in dimension \(\le j\)
(Lemma~\ref{mc:factorization-domain}), giving \(T_2\).
\end{proof}
With a resolution invariant in dimension \(\le n\), the driver of Theorem~\ref{hir:invariant-driver-seq} supplies both resolutions, and the theorem gives the geometric core of
Lemma~\ref{hir:factorization-from-resolution} on the nose.

\subsection{The lower marked ideals of Step 1 along a phase}

In this subsection \(\underline{\mathcal J}=(X,N,E,\mathcal J,d)\) is well formed with \(d\ge1\), of maximal order along \(N\) with \(\dim N\le k\) (\(n\ge1\)),
\(\overline{\underline{\mathcal J}}=\mathcal C(\mathcal H(\underline{\mathcal J}))\), \(\operatorname{inv}\) is a resolution invariant in dimension \(\le k-1\) and
\(\operatorname{inv}^M\) is the masked Step 1 value of \(\underline{\mathcal J}\) (Definition~\ref{hir:step-one-value-masked}), with values \((s(y),\lambda)\): \(s(y)\) is the level
(the number of original members through \(y\)) and \(\lambda\) the lower value, \(\infty\) at the special points. For a stage \(T_0\) and a point \(y_0\) of
\(\operatorname{cosupp}\overline{\underline{\mathcal J}}_{T_0}\) in the regime of level \(s_0\) (level \(s_0\), not special), \(W(y_0)\) is its mask and \(\underline{\mathcal L}(y_0)\) the
lower marked ideal formed at \(T_0\) on \(W(y_0)\) (Definition~\ref{hir:step-one-value-masked}). Along a canonical segment the values of the nonempty centres strictly
decrease (Theorem~\ref{hir:step-one-phase}), so the centres of level \(s_0\) come in one block, the \emph{phase} of level \(s_0\): we say that a phase of level \(s_0\)
starts at the stage \(T_0\) of a canonical segment \(T_0S\) if every point of every centre of \(T_0\) has level \(>s_0\), or level \(s_0\) and is special, and it lasts
along \(S\) if every point of every centre of \(S\) has level \(\ge s_0\). At the first nonempty centre of value \((s_0,\lambda)\) with \(\lambda\) finite, a phase of level
\(s_0\) starts, since the earlier values are larger than \((s_0,\lambda)\).

\begin{theorem}[One lower marked ideal for a whole phase]
\label{mc:step-one-phase-tail}
\uses{hir:step-one-value-masked,hir:step-one-value-masked-basic,hir:step-one-canonical-tail,hir:step-one-mask-tower,hir:step-one-compare,mc:tangent-change-opens,hir:step-one-canonical-test,hir:step-one-bridge,hir:step-one-phase,hir:regime-run,hir:tower-split,hir:canonical-segments,hir:centre-function,hir:invariant-driver,hir:centre-function-seq,bcalc:max-locus-local,bcalc:growing-seq-tail,bcalc:general-companion-phase}
\source{wlodarczyk05-hironaka,bierstone-milman08-functoriality}
\lean{Hironaka.StepOnePhaseTail.CentresLevelGt,Hironaka.StepOnePhaseTail.CentresLevelGe,Hironaka.StepOnePhaseTail.centresAvoid_of_centresLevelGt,Hironaka.StepOnePhaseTail.stepOneInvM_eq_of_mem_cosupp_phase,Hironaka.StepOnePhaseTail.isCanonical_restrict_phase,Hironaka.StepOnePhaseTail.isCanonical_tail_restrict,Hironaka.StepOnePhaseTail.isCanonical_tail_restrict_of_isInvDomain}
\leanok
Let a phase of level \(s_0\) start at the stage \(T_0\) of a canonical segment \(T_0S\) of \(\operatorname{inv}^M\) and last along \(S\), and let \(y_0\) be a point of
\(\operatorname{cosupp}\overline{\underline{\mathcal J}}_{T_0}\) in the regime of level \(s_0\), with mask \(W=W(y_0)\) and lower marked ideal \(\underline{\mathcal L}=\underline{\mathcal L}(y_0)\).
\begin{enumerate}
\item For every point \(p\) of level \(s_0\) at a stage \(T'\) and every continuation \(T'M\) whose centres have level \(>s_0\) or are special of level \(s_0\), the centres
of \(M\) miss the preimages of the mask of \(p\).
\item If \(S|_W\) is a canonical segment of \(\underline{\mathcal L}\), then at every point \(q\) of \(\operatorname{cosupp}\underline{\mathcal L}_{S|_W}\) the value is
\(\operatorname{inv}^M(T_0S)(q)=(s_0,\operatorname{inv}_{\underline{\mathcal L}}(S|_W)(q))\).
\item \(S|_W\) is a canonical segment of \(\underline{\mathcal L}\).
\end{enumerate}
Consequently, for a growing sequence \((T_i)\) of canonical segments of \(\operatorname{inv}^M\) (Definition~\ref{hir:centre-function-seq}) and a stage \(a\) at which a
phase of level \(s_0\) starts and which lasts from \(a\) to \(a+i\), the restriction to \(W\) of the part of \(T_{a+i}\) after \(T_a\) is a canonical segment of the one marked
ideal \(\underline{\mathcal L}(y_0)\) formed at \(T_a\); its nonempty centres meeting a set \(L\subset W\) are those of the tail meeting \(L\) (Lemma~\ref{bcalc:growing-seq-tail}).
The same holds for the canonical segments of \(\underline{\mathcal J}\) itself when \(\underline{\mathcal J}\) lies in the domain (Corollary~\ref{hir:step-one-bridge}).
\end{theorem}
\begin{proof}
\leanok
(1) By induction along \(M\). A point \(c\) of a centre of \(M\) lying over the mask of \(p\) has level at most \(s_0\) (over the mask the level is at most that of
\(p\), Lemma~\ref{hir:step-one-mask-tower}); so it has level \(s_0\) and is special, while no point of the level of \(p\) over its mask is special
(Lemma~\ref{hir:step-one-mask-tower}). \(T'M\) is a multiple test blow-up of \(\overline{\underline{\mathcal J}}\) as a prefix of a canonical segment (Corollary~\ref{hir:step-one-canonical-test}).

(2) By Lemma~\ref{hir:step-one-compare}(4) the point \(q\) lies in \(\operatorname{cosupp}\overline{\underline{\mathcal J}}_{T_0S}\) and has level \(s_0\); it is in the regime of
level \(s_0\), and its run started at the stage \(T_0\) or before (Lemma~\ref{hir:step-one-mask-tower}). By the canonical-tail theorem (Theorem~\ref{hir:step-one-canonical-tail})
\(q\) has lower data at the start \(T_0'\) of its run; write \(T_0=T_0'M\) and let \(T_1'=MS\) be the tail of the split at the start of \(q\) (Lemma~\ref{hir:tower-split}). The
centres of \(M\) lie before the start of the phase, so they miss the mask of the image of \(q\) at \(T_0'\) by (1). Lemma~\ref{hir:step-one-compare}(2), resp.\ (3) followed by the change of
tangent direction on the masks (Theorem~\ref{mc:tangent-change-opens}(2)), as in the proof of Theorem~\ref{hir:step-one-canonical-tail}, moves the lower value of \(q\), computed from \(T_0'\), to the value of
\(\operatorname{inv}\) of \(\underline{\mathcal L}\) along \(S|_W\) at \(q\); here one uses that \(S|_W\) is a canonical segment of \(\underline{\mathcal L}\), which lies in the domain
in dimension \(\le k-1\) (Lemma~\ref{hir:step-one-value-masked-basic}).

(3) By induction along \(S\); the empty tower is a canonical segment. Let \(S\) be followed by the blowup along \(C\), empty or the maximum locus of \(\operatorname{inv}^M\) on
\(\operatorname{cosupp}\overline{\underline{\mathcal J}}_{T_0S}\), with \(S|_W\) a canonical segment by induction. If the trace of \(C\) on \(S^{-1}W\) is empty there is nothing to show.
Otherwise it contains a point \(c\), of level \(\ge s_0\) since the phase lasts and of level \(\le s_0\) over the mask, so \(c\in\operatorname{cosupp}\underline{\mathcal L}_{S|_W}\)
(Lemma~\ref{hir:step-one-compare}(4)). On \(\operatorname{cosupp}\underline{\mathcal L}_{S|_W}\subset S^{-1}W\cap\operatorname{cosupp}\overline{\underline{\mathcal J}}_{T_0S}\) the value is
\((s_0,\operatorname{inv}_{\underline{\mathcal L}}(S|_W))\) by (2), a strictly monotone function of \(\operatorname{inv}_{\underline{\mathcal L}}(S|_W)\); every other point of
\(S^{-1}W\cap\operatorname{cosupp}\overline{\underline{\mathcal J}}_{T_0S}\) has level \(<s_0\), hence a smaller value than \(c\). So the trace of the maximum locus is the maximum locus of
\(\operatorname{inv}_{\underline{\mathcal L}}(S|_W)\) on \(\operatorname{cosupp}\underline{\mathcal L}_{S|_W}\) (Lemma~\ref{bcalc:max-locus-local}).

For the growing sequence, the tower \(T_{a+i}\) is \(T_a\) followed by the \(i\)-th tower of the tail (Lemma~\ref{bcalc:growing-seq-tail}), and the \(i\)-th tower of the restricted
tail is its restriction to \(W\). For \(\underline{\mathcal J}\) in the domain, its canonical segments are those of \(\overline{\underline{\mathcal J}}\), with the same cosupports
(Corollary~\ref{hir:step-one-bridge}).
\end{proof}
Without the condition on the centres before \(T_0\) the statement fails in general: a point of level \(s_0\) over \(W\) whose run started before \(T_0\) is valued along its own,
longer, restricted tail, and Definition~\ref{hir:centre-function}(4') re-bases only across prefixes with empty centres.

\begin{lemma}[Finitely many masks cover the part of level at least \(s_0\)]
\label{mc:step-one-phase-cover}
\uses{bcalc:tower-preimage-compact,bcalc:boundary-count,hir:step-one-lower-domain,hir:step-one-value-masked-basic,hir:open-restriction-tower,hir:giraud-along}
\source{wlodarczyk05-hironaka,bierstone-milman08-functoriality}
\lean{Hironaka.StepOnePhaseTail.isClosed_setOf_le_level,Hironaka.StepOnePhaseTail.exists_cover_masks}
\leanok
Let \(T_0\) be a multiple test blow-up of \(\overline{\underline{\mathcal J}}\) and \(K\subset X\) compact. The points of level \(\ge s_0\) form a closed set. If every point of
\(p_{T_0}^{-1}(K)\cap\operatorname{cosupp}\overline{\underline{\mathcal J}}_{T_0}\) of level \(\ge s_0\) is in the regime of level \(s_0\), then there are finitely many such points
\(y_1,\dots,y_r\) and compact sets \(L_j\subset W(y_j)\) covering \(\{y\in p_{T_0}^{-1}(K)\cap\operatorname{cosupp}\overline{\underline{\mathcal J}}_{T_0}\mid s(y)\ge s_0\}\).
\end{lemma}
\begin{proof}
\leanok
The original members are closed (Lemma~\ref{hir:step-one-lower-domain}), so near each point the level is at most its value there (Lemma~\ref{bcalc:boundary-count}), and
\(\{s<s_0\}\) is open. The set \(A=\{y\in p_{T_0}^{-1}(K)\cap\operatorname{cosupp}\overline{\underline{\mathcal J}}_{T_0}\mid s(y)\ge s_0\}\) is compact: \(p_{T_0}^{-1}(K)\) is compact
(Lemma~\ref{bcalc:tower-preimage-compact}) and the cosupport is closed (the transform is well formed with coherent ideal, Lemmas~\ref{hir:open-restriction-tower}
and~\ref{hir:giraud-along}). Each \(y\in A\) is in the regime of level \(s_0\), so it lies in its own mask (Lemma~\ref{hir:step-one-value-masked-basic}), which is open; the stage is
a manifold, hence locally compact, so \(y\) has a compact neighbourhood \(L_y\subset W(y)\). Finitely many of the interiors of the \(L_y\) cover \(A\).
\end{proof}

\begin{theorem}[Finitely many centres in an infinite phase]
\label{mc:step-one-phase-bound}
\uses{mc:step-one-phase-tail,mc:step-one-phase-cover,bcalc:growing-seq-tail,hir:centre-function-seq,hir:step-one-bridge,hir:step-one-mask-tower,hir:step-one-value-masked-basic}
\source{wlodarczyk05-hironaka,bierstone-milman08-functoriality}
\lean{Hironaka.StepOnePhaseTail.countCentresMeeting_setOf_level_lt_eq_zero,Hironaka.StepOnePhaseTail.exists_bound_of_phase}
\leanok
Let \(\underline{\mathcal J}\) lie in the domain in dimension \(\le k\) (Definition~\ref{hir:centre-function}) and suppose \(\operatorname{inv}\) has property (8\('\)) of
Definition~\ref{hir:centre-function-seq} (this follows from its property (8)). Let \((T_i)\) be a growing sequence of canonical segments of \(\operatorname{inv}^M\) for
\(\underline{\mathcal J}\), \(K\subset X\) compact, and \(a\) a stage at which a phase of level \(s_0\) starts and lasts forever, such that every point of
\(p_{T_a}^{-1}(K)\cap\operatorname{cosupp}\underline{\mathcal J}_{T_a}\) of level \(\ge s_0\) is in the regime of level \(s_0\). Then the numbers of centres of \(T_i\) meeting
\(p_{T_i}^{-1}(K)\) are bounded. In a lasting phase no centre after \(T_a\) lies over a point of level \(<s_0\) at \(T_a\).
\end{theorem}
\begin{proof}
\leanok
The level does not increase along a tower (Lemma~\ref{hir:step-one-mask-tower}), so a point of a centre after \(T_a\), of level \(\ge s_0\), lies over a point of level \(\ge s_0\)
at \(T_a\); this is the last claim. The masks \(W(y_j)\) and compact sets \(L_j\) of Lemma~\ref{mc:step-one-phase-cover} (for \(\overline{\underline{\mathcal J}}\), whose cosupports along the
sequence are those of \(\underline{\mathcal J}\), Corollary~\ref{hir:step-one-bridge}) together with \(L_0=\{s<s_0\}\subset X_a\) cover
\(p_{T_a}^{-1}(K)\cap\operatorname{cosupp}\underline{\mathcal J}_{T_a}\). Along the tail restricted to \(W(y_j)\), a growing sequence of canonical segments of the one lower marked ideal
\(\underline{\mathcal L}(y_j)\) (Theorem~\ref{mc:step-one-phase-tail}), which lies in the domain in dimension \(\le k-1\) (Lemma~\ref{hir:step-one-value-masked-basic}), property
(8\('\)) of \(\operatorname{inv}\) bounds the numbers of centres meeting the compact set \(L_j\); no centre after \(T_a\) meets the preimage of \(L_0\). Lemma~\ref{bcalc:growing-seq-tail}
turns bounded counts along the restricted tails over a finite cover into bounded counts along the sequence.
\end{proof}
The remaining bookkeeping for property (8\('\)) of the Step 1 value -- the phases which end, the at most \(\#E+1\) special centres and the sum over the levels
\(s_0\le\#E\) -- is that of Theorem~\ref{hir:step-one-phase}.

\subsection{The two Hironaka conclusions from a resolution invariant with local finiteness along sequences}

\begin{theorem}[Relative principalization and embedded resolution from a resolution invariant]
\label{mc:hironaka-from-seq-invariant}
\uses{mc:factorization-assembly,mc:embedded-from-first-centre,mc:embedding-properties,mc:subvariety-thin,mc:subvariety-full-dimension,mc:embedded-from-invariant,hir:centre-function-seq,hir:invariant-driver-seq}
\source{transcendental-nakamaye-copy4-20260911:lem:relative-principalization,wlodarczyk05-hironaka,bierstone-milman08-functoriality}
\lean{Hironaka.HironakaFromSeqInvariant.embeddingFields_of_seq,Hironaka.HironakaFromSeqInvariant.exists_isGeoResolution_of_seq,Hironaka.HironakaFromSeqInvariant.factorizationData_of_seq,Hironaka.HironakaFromSeqInvariant.exists_embeddedResolution_data_of_seq,Hironaka.HironakaFromSeqInvariant.exists_embeddedResolution_data_of_seq_of_isCoherent}
\leanok
Let \(\operatorname{inv}\) be a resolution invariant in the sense of Definition~\ref{hir:centre-function-seq} (property (8) along growing sequences).
\begin{enumerate}
\item If \(\operatorname{inv}\) is in dimension \(\le j\), then over a compact base every marked ideal of the domain in dimension \(\le j\) has a multiple test blow-up with
empty final cosupport.
\item If \(\operatorname{inv}\) is in dimension \(\le j\) with \(n\le j\), \(Y\) is a compact connected complex manifold of dimension \(n\), \(W\subsetneq Y\) a closed connected
submanifold of dimension \(k\) and \(\mathfrak a\) a coherent ideal with \(\mathfrak a\not\subset\mathcal I_W\), then factorization data as in
Theorem~\ref{mc:factorization-assembly} exist.
\item If \(\operatorname{inv}\) is in dimension \(\le k\) with \(n\le k\), \(X\) is a compact complex manifold of dimension \(n\) and \(V\subset X\) a closed irreducible analytic
subvariety whose ideal sheaf \(\mathcal I_V\) is coherent, then there is a finite composition of blowups with smooth centres whose exceptional image lies in
\(V_{\mathrm{sing}}\) and whose strict transform of \(V\) is regular of dimension \(\dim V\) at each of its points.
\end{enumerate}
\end{theorem}
\begin{proof}
\leanok
(1) The driver of Theorem~\ref{hir:invariant-driver-seq} resolves the marked ideal by a canonical segment. (2) Theorem~\ref{mc:factorization-assembly} with the
existence of resolutions from (1). (3) Properties (1), (4), (5) and (6) of Definition~\ref{hir:centre-function-seq} are those of Definition~\ref{mc:embedding-properties}; the
proof of Theorem~\ref{mc:embedded-from-invariant} applies verbatim with the canonical resolution of (1): for \(\dim V=n\) the identity
(Lemma~\ref{mc:subvariety-full-dimension}), otherwise \((X,X,\emptyset,\mathcal I_V,1)\) lies in the domain (Lemma~\ref{mc:subvariety-thin}) and its canonical resolution gives the
modification (Theorem~\ref{mc:embedded-from-first-centre}).
\end{proof}
This is Theorem~\ref{mc:embedded-from-invariant} and the last claim of Theorem~\ref{mc:factorization-assembly} for the form of resolution invariant that the
induction of Theorem~\ref{hir:canonical-centres} produces; with a resolution invariant in dimension \(\le n\), both conclusions are those of
Lemmas~\ref{hir:factorization-from-resolution} and~\ref{hir:embedded-from-resolution}, the coherence of \(\mathcal I_V\) being Cartan's theorem.
